\input{preamble}

% OK, start here.
%
\begin{document}

\title{Groupes fondamentaux des schémas}


\maketitle

\phantomsection
\label{section-phantom}

\tableofcontents

\section{Introduction}
\label{section-introduction}

\noindent
Dans ce chapitre, nous étudions le groupe fondamental d'un schéma au sens de
Grothendieck ainsi que ses applications. Une référence fondamentale est
\cite{SGA1}. On trouvera une bonne introduction dans \cite{Lenstra}.
Voir aussi \cite{Murre-lectures} et \cite{Grothendieck-Murre}.










\section{Schémas étales sur un point}
\label{section-schemes-etale-point}

\noindent
Dans cette section, nous décrivons les schémas étales sur le spectre d'un
corps. Avant d'énoncer le résultat, introduisons la catégorie des
$G$-ensembles pour un groupe topologique $G$.

\begin{definition}
\label{definition-G-set-continuous}
Soit $G$ un groupe topologique.
Un {\it $G$-ensemble}, parfois appelé {\it $G$-ensemble discret},
est un ensemble $X$ muni d'une action à gauche $a : G \times X \to X$
telle que $a$ soit continue lorsque $X$ est muni de la topologie discrète et
$G \times X$ de la topologie produit.
Un {\it morphisme de $G$-ensembles} $f : X \to Y$ est simplement une
application $G$-équivariante de $X$ dans $Y$.
La catégorie des $G$-ensembles est notée {\it $G\textit{-Ensembles}$}.
\end{definition}

\noindent
La continuité de $a : G \times X \to X$ signifie simplement que le
stabilisateur de tout $x \in X$ est ouvert dans $G$.
Si $G$ est un groupe abstrait (c'est-à-dire un groupe non muni d'une
topologie), ceci concorde avec notre définition précédente (voir par exemple
Sites, exemple \ref{sites-example-site-on-group}), pourvu que l'on munisse
$G$ de la topologie discrète.

\medskip\noindent
Rappelons que, si $L/K$ est une extension galoisienne infinie, le groupe de
Galois $G = \text{Gal}(L/K)$ est muni d'une topologie canonique ; voir
Corps, section \ref{fields-section-infinite-galois}.

\begin{lemma}
\label{lemma-sheaves-point}
Soit $K$ un corps. Soit $K^{sep}$ une clôture séparable de $K$.
Considérons le groupe profini $G = \text{Gal}(K^{sep}/K)$.
Le foncteur
$$
\begin{matrix}
\text{schémas étales sur }K &
\longrightarrow &
G\textit{-Ensembles} \\
X/K & \longmapsto &
\Mor_{\Spec(K)}(\Spec(K^{sep}), X)
\end{matrix}
$$
est une équivalence de catégories.
\end{lemma}

\begin{proof}
Un schéma $X$ sur $K$ est étale sur $K$ si et seulement si
$X \cong \coprod_{i\in I} \Spec(K_i)$, où chaque $K_i$ est une extension
séparable finie de $K$
(Morphismes, lemme \ref{morphisms-lemma-etale-over-field}).
Le foncteur du lemme associe à $X$ le $G$-ensemble
$$
\coprod\nolimits_i \Hom_K(K_i, K^{sep})
$$
muni de son action à gauche naturelle de $G$. Par définition de la topologie
sur $G$, chaque élément a un stabilisateur ouvert. Réciproquement, tout
$G$-ensemble $S$ est réunion disjointe de ses orbites. Écrivons
$S = \coprod S_i$. Choisissons $s_i \in S_i$ et notons $G_i \subset G$ son
stabilisateur ouvert. D'après la théorie de Galois
(Corps, théorème \ref{fields-theorem-inifinite-galois-theory})
les corps $(K^{sep})^{G_i}$ sont des extensions séparables finies de $K$ ;
par conséquent, le schéma
$$
\coprod\nolimits_i \Spec((K^{sep})^{G_i})
$$
est étale sur $K$. On obtient ainsi un quasi-inverse du foncteur du lemme.
Nous omettons certains détails.
\end{proof}

\begin{remark}
\label{remark-covering-surjective}
Sous la correspondance du lemme \ref{lemma-sheaves-point}, les recouvrements
du petit site étale $\Spec(K)_\etale$ de $K$ correspondent aux familles
surjectives de morphismes dans $G\textit{-Ensembles}$.
\end{remark}









\section{Catégories galoisiennes}
\label{section-galois}

\noindent
Dans cette section, nous présentons une partie du contenu que le lecteur
trouvera dans \cite[Exposé V, sections 4, 5 et 6]{SGA1}.

\medskip\noindent
Soit $F : \mathcal{C} \to \textit{Ensembles}$ un foncteur.
Rappelons que, suivant nos conventions, les catégories ont un ensemble
d'objets et, pour toute paire d'objets, un ensemble de morphismes. Il existe
une application injective canonique
\begin{equation}
\label{equation-embedding-product}
\text{Aut}(F)
\longrightarrow
\prod\nolimits_{X \in \Ob(\mathcal{C})} \text{Aut}(F(X))
\end{equation}
Pour un ensemble $E$, nous munissons $\text{Aut}(E)$ de la topologie
compacte-ouverte ; voir Topologie, exemple
\ref{topology-example-automorphisms-of-a-set}. Il s'agit bien entendu de la
topologie discrète lorsque $E$ est fini, ce qui est le cas qui nous intéresse
dans cette section\footnote{Lorsque nous étudierons le groupe fondamental
pro-étale, le cas général nous intéressera.}. Nous munissons
$\text{Aut}(F)$ de la topologie induite par la topologie produit du membre de
droite de (\ref{equation-embedding-product}). En particulier, les applications
d'action
$$
\text{Aut}(F) \times F(X) \longrightarrow F(X)
$$
sont continues lorsque $F(X)$ est muni de la topologie discrète, puisque
c'est le cas des applications d'action $\text{Aut}(E) \times E \to E$ pour
tout ensemble $E$. La propriété universelle de notre topologie sur
$\text{Aut}(F)$ est la suivante : supposons que $G$ soit un groupe topologique
et que $G \to \text{Aut}(F)$ soit un homomorphisme de groupes tel que les
actions induites $G \times F(X) \to F(X)$ soient continues pour tout
$X \in \Ob(\mathcal{C})$, où $F(X)$ est muni de la topologie discrète.
Alors $G \to \text{Aut}(F)$ est continue.

\medskip\noindent
Le lemme suivant affirme que le groupe des automorphismes d'un foncteur à
valeurs dans la catégorie des ensembles finis est automatiquement profini.

\begin{lemma}
\label{lemma-aut-inverse-limit}
Soit $\mathcal{C}$ une catégorie et soit
$F : \mathcal{C} \to \textit{Ensembles}$ un foncteur. L'application
(\ref{equation-embedding-product}) identifie $\text{Aut}(F)$ à un sous-groupe
fermé de
$\prod_{X \in \Ob(\mathcal{C})} \text{Aut}(F(X))$.
En particulier, si $F(X)$ est fini pour tout $X$, alors
$\text{Aut}(F)$ est un groupe profini.
\end{lemma}

\begin{proof}
Soit $\xi = (\gamma_X) \in \prod \text{Aut}(F(X))$ un élément qui
n'appartient pas à $\text{Aut}(F)$. Il existe alors un morphisme
$f : X \to X'$ de $\mathcal{C}$ et un élément $x \in F(X)$ tels que
$F(f)(\gamma_X(x)) \not = \gamma_{X'}(F(f)(x))$.
Considérons le voisinage ouvert
$U = \{\gamma \in \text{Aut}(F(X)) \mid \gamma(x) = \gamma_X(x)\}$
de $\gamma_X$ et le voisinage ouvert
$U' = \{\gamma' \in \text{Aut}(F(X')) \mid \gamma'(F(f)(x)) =
\gamma_{X'}(F(f)(x))\}$.
Alors
$U \times U' \times \prod_{X'' \not = X, X'} \text{Aut}(F(X''))$
est un voisinage ouvert de $\xi$ qui ne rencontre pas $\text{Aut}(F)$.
La dernière assertion résulte du fait que
$\prod \text{Aut}(F(X))$ est un espace profini si chaque $F(X)$ est fini.
\end{proof}

\begin{example}
\label{example-galois-category-G-sets}
Soit $G$ un groupe topologique. Un exemple important sera le foncteur d'oubli
\begin{equation}
\label{equation-forgetful}
\textit{Finite-}G\textit{-Ensembles} \longrightarrow \textit{Ensembles}
\end{equation}
où $\textit{Finite-}G\textit{-Ensembles}$ est la sous-catégorie pleine de
$G\textit{-Ensembles}$ dont les objets sont les $G$-ensembles finis.
La catégorie $G\textit{-Ensembles}$ des $G$-ensembles est définie dans la
définition \ref{definition-G-set-continuous}.
\end{example}

\noindent
Soit $G$ un groupe topologique. La {\it complétion profinie} de $G$ est le
groupe profini
$$
G^\wedge =
\lim_{U \subset G\text{ ouvert, distingué, d'indice fini}} G/U
$$
muni de sa topologie profinie. Remarquons que la limite est cofiltrante,
car toute intersection finie de sous-groupes ouverts distingués d'indice fini
est encore de cette forme. La propriété universelle de la complétion profinie
affirme que toute application continue $G \to H$ vers un groupe profini $H$
se factorise canoniquement sous la forme $G \to G^\wedge \to H$.

\begin{lemma}
\label{lemma-single-out-profinite}
Soit $G$ un groupe topologique. Le groupe des automorphismes du foncteur
(\ref{equation-forgetful}), muni de la topologie profinie du lemme
\ref{lemma-aut-inverse-limit}, est la complétion profinie de $G$.
\end{lemma}

\begin{proof}
Notons $F_G$ le foncteur (\ref{equation-forgetful}). Tout morphisme
$X \to Y$ dans $\textit{Finite-}G\textit{-Ensembles}$ commute à l'action de $G$.
Ainsi, tout $g \in G$ définit un automorphisme de $F_G$, et l'on obtient un
homomorphisme canonique de groupes $G \to \text{Aut}(F_G)$.
Remarquons que tout $G$-ensemble fini $X$ est réunion disjointe finie de
$G$-ensembles de la forme $G/H_i$, munis de l'action canonique de $G$, où
$H_i \subset G$ est un sous-groupe ouvert d'indice fini. Alors
$U_i = \bigcap gH_ig^{-1}$ est ouvert, distingué et d'indice fini.
De plus, $U_i$ agit trivialement sur $G/H_i$ ; ainsi,
$U = \bigcap U_i$ agit trivialement sur $F_G(X)$. L'action
$G \times F_G(X) \to F_G(X)$ est donc continue. Par la propriété universelle
de la topologie sur $\text{Aut}(F_G)$, l'application
$G \to \text{Aut}(F_G)$ est continue.
D'après le lemme \ref{lemma-aut-inverse-limit} et la propriété universelle
de la complétion profinie, on en déduit un homomorphisme continu de groupes
$$
G^\wedge \longrightarrow \text{Aut}(F_G)
$$
De plus, puisque $G/U$ agit fidèlement sur $G/U$, cette application est
injective. Si son image est dense, l'application est surjective, donc c'est
un homéomorphisme d'après Topologie, lemme
\ref{topology-lemma-bijective-map}.

\medskip\noindent
Soient $\gamma \in \text{Aut}(F_G)$ et $X \in \Ob(\mathcal{C})$.
Montrons qu'il existe $g \in G$ tel que $\gamma$ et $g$ induisent la même
action sur $F_G(X)$, ce qui achèvera la démonstration. Comme ci-dessus,
$X$ est une réunion disjointe finie de $G/H_i$. Avec les notations $U_i$ et
$U$ précédentes, le $G$-ensemble fini $Y = G/U$ se surjecte sur $G/H_i$ pour
tout $i$ ; il suffit donc de trouver $g \in G$ tel que $\gamma$ et $g$
induisent la même action sur $F_G(G/U) = G/U$. Soit $e \in G$ l'élément
neutre, et écrivons $\gamma(eU) = g_0U$ pour un certain $g_0 \in G$.
Pour tout $g_1 \in G$, le morphisme
$$
R_{g_1} : G/U \longrightarrow G/U,\quad gU \longmapsto gg_1U
$$
de $\textit{Finite-}G\textit{-Ensembles}$ commute à l'action de $\gamma$. Ainsi
$$
\gamma(g_1U) = \gamma(R_{g_1}(eU)) = R_{g_1}(\gamma(eU)) =
R_{g_1}(g_0U) = g_0g_1U
$$
On voit donc que $g = g_0$ convient.
\end{proof}

\noindent
Rappelons qu'un foncteur exact est un foncteur qui commute à toutes les
limites finies et à toutes les colimites finies. En particulier, un tel
foncteur commute aux égalisateurs, aux coégalisateurs, aux produits fibrés,
aux sommes amalgamées, etc.

\begin{lemma}
\label{lemma-second-fundamental-functor}
Soit $G$ un groupe topologique. Soit
$F : \textit{Finite-}G\textit{-Ensembles} \to \textit{Ensembles}$
un foncteur exact tel que $F(X)$ soit fini pour tout $X$.
Alors $F$ est isomorphe au foncteur (\ref{equation-forgetful}).
\end{lemma}

\begin{proof}
Soit $X$ un objet non vide de $\textit{Finite-}G\textit{-Ensembles}$. Le diagramme
$$
\xymatrix{
X \ar[r] \ar[d] & \{*\} \ar[d] \\
\{*\} \ar[r] & \{*\}
}
$$
est cocartésien. On en déduit que $F(X)$ est non vide.
Soit $U \subset G$ un sous-groupe ouvert distingué d'indice fini. Remarquons que
$$
G/U \times G/U = \coprod\nolimits_{gU \in G/U} G/U
$$
où le facteur correspondant à $gU$ correspond, dans le membre de gauche, à
l'orbite de $(eU, gU)$. On a alors
$$
F(G/U) \times F(G/U) = F(G/U \times G/U) = \coprod\nolimits_{gU \in G/U} F(G/U)
$$
Puisque $F(G/U)$ est non vide, il s'ensuit que
$|F(G/U)| = |G/U|$. Par conséquent,
$$
\lim_{U \subset G\text{ ouvert, distingué, d'indice fini}} F(G/U)
$$
est non vide (Catégories, lemme \ref{categories-lemma-nonempty-limit}).
Choisissons dans cette limite un élément $\gamma = (\gamma_U)$.
Notons $F_G$ le foncteur (\ref{equation-forgetful}). On peut identifier
$F_G$ au foncteur
$$
X \longmapsto \colim_U \Mor(G/U, X)
$$
où $f : G/U \to X$ correspond à $f(eU) \in X = F_G(X)$
(nous omettons les détails). L'élément $\gamma$ détermine donc une
application bien définie
$$
t : F_G \longrightarrow F
$$
En effet, étant donné $x \in X$, choisissons $U$ et
$f : G/U \to X$ envoyant $eU$ sur $x$, puis posons
$t_X(x) = F(f)(\gamma_U)$. Montrons que $t$ induit une application bijective
$t_{G/U} : F_G(G/U) \to F(G/U)$ pour tout $U$. Il en résulte immédiatement
que $t$ est un isomorphisme (nous omettons les détails).
Comme $|F_G(G/U)| = |F(G/U)|$, il suffit de montrer que $t_{G/U}$ est
surjective. Son image contient au moins un élément, à savoir
$t_{G/U}(eU) = F(\text{id}_{G/U})(\gamma_U) = \gamma_U$.
Pour $g \in G$, notons $R_g : G/U \to G/U$ la multiplication à droite.
Si $g \not \in U$, l'ensemble des points fixes de
$F(R_g) : F(G/U) \to F(G/U)$ est égal à
$F(\emptyset) = \emptyset$, car $F$ commute aux égalisateurs. Il s'ensuit que,
si $g_1, \ldots, g_{|G/U|}$ est un système de représentants de $G/U$, les
éléments $F(R_{g_i})(\gamma_U)$ sont deux à deux distincts et constituent
donc tout $F(G/U)$. Enfin,
$$
t_{G/U}(g_iU) = F(R_{g_i})(\gamma_U)
$$
ce qui achève la démonstration.
\end{proof}

\begin{example}
\label{example-from-C-F-to-G-sets}
Soit $\mathcal{C}$ une catégorie et soit
$F : \mathcal{C} \to \textit{Ensembles}$ un foncteur tel que $F(X)$ soit fini
pour tout $X \in \Ob(\mathcal{C})$. Le lemme
\ref{lemma-aut-inverse-limit} montre que $G = \text{Aut}(F)$ est muni de
manière canonique d'une structure de groupe topologique profini. On obtient
un foncteur
\begin{equation}
\label{equation-remember}
\mathcal{C} \longrightarrow \textit{Finite-}G\textit{-Ensembles},\quad
X \longmapsto F(X)
\end{equation}
où $F(X)$ est muni de l'action induite de $G$. Cette action est continue par
construction de la topologie sur $\text{Aut}(F)$.
\end{example}

\noindent
Le but de la définition des catégories galoisiennes est de distinguer les
couples $(\mathcal{C}, F)$ pour lesquels le foncteur (\ref{equation-remember})
est une équivalence. Voici notre définition d'une catégorie galoisienne.

\begin{definition}
\label{definition-galois-category}
\begin{reference}
Cette définition diffère de celle de \cite[Expos\'e V, Définition 5.1]{SGA1}.
Comparer avec \cite[Définition 7.2.1]{BS}.
\end{reference}
Soit $\mathcal{C}$ une catégorie et soit $F : \mathcal{C} \to \textit{Ensembles}$
un foncteur. Le couple $(\mathcal{C}, F)$ est une {\it catégorie galoisienne} si
\begin{enumerate}
\item $\mathcal{C}$ admet les limites finies et les colimites finies,
\item
\label{item-connected-components}
tout objet de $\mathcal{C}$ est un coproduit fini (éventuellement vide)
d'objets connexes,
\item $F(X)$ est fini pour tout $X \in \Ob(\mathcal{C})$, et
\item $F$ est conservatif\footnote{Plus précisément, étant donné un morphisme
$f$ de $\mathcal{C}$, si $F(f)$ est un isomorphisme, alors
$f$ est un isomorphisme.} et est exact\footnote{Cela signifie que
$F$ commute aux limites finies et aux colimites finies ; voir
Catégories, section \ref{categories-section-exact-functor}.}.
\end{enumerate}
On dit ici que $X \in \Ob(\mathcal{C})$ est connexe s'il
n'est pas initial et si, pour tout monomorphisme $Y \to X$,
soit $Y$ est initial, soit $Y \to X$ est un isomorphisme.
\end{definition}

\noindent
{\bf Avertissement :} Cette définition n'est pas la même que celle donnée dans la plupart
des références, même si nous verrons finalement que les deux sont équivalentes.
En effet, dans \cite[Expos\'e V, Définition 5.1]{SGA1}, une catégorie galoisienne est
définie comme une catégorie équivalente à $\textit{Finite-}G\textit{-Ensembles}$
pour un certain groupe profini $G$. Grothendieck caractérise alors les
catégories galoisiennes par une liste d'axiomes (G1) -- (G6) plus faibles
que les nôtres ci-dessus. Notre choix vise à souligner
l'existence des limites finies et des colimites finies, ainsi que l'exactitude du
foncteur $F$. En contrepartie, il nous faudra plus tard travailler
un peu davantage pour appliquer les résultats de cette section.

\begin{lemma}
\label{lemma-epi-mono}
Soit $(\mathcal{C}, F)$ une catégorie galoisienne. Soit
$X \to Y \in \text{Arrows}(\mathcal{C})$. Alors
\begin{enumerate}
\item $F$ est fidèle,
\item $X \to Y$ est un monomorphisme
$\Leftrightarrow F(X) \to F(Y)$ est injective,
\item $X \to Y$ est un épimorphisme
$\Leftrightarrow F(X) \to F(Y)$ est surjective,
\item un objet $A$ de $\mathcal{C}$ est initial si et seulement si
$F(A) = \emptyset$,
\item un objet $Z$ de $\mathcal{C}$ est final si et seulement si
$F(Z)$ est un ensemble à un élément,
\item si $X$ et $Y$ sont connexes, alors $X \to Y$ est un épimorphisme,
\item
\label{item-one-element}
si $X$ est connexe et si $a, b : X \to Y$ sont deux morphismes,
alors $a = b$ dès que $F(a)$ et $F(b)$ coïncident sur un élément de $F(X)$,
\item si $X = \coprod_{i = 1, \ldots, n} X_i$ et
$Y = \coprod_{j = 1, \ldots, m} Y_j$, où les $X_i$, $Y_j$ sont connexes,
alors il existe une application $\alpha : \{1, \ldots, n\} \to \{1, \ldots, m\}$
telle que $X \to Y$ provienne d'une famille de morphismes
$X_i \to Y_{\alpha(i)}$.
\end{enumerate}
\end{lemma}

\begin{proof}
Démonstration de (1). Supposons que $a, b : X \to Y$ vérifient $F(a) = F(b)$.
Soit $E$ l'égalisateur de $a$ et $b$. Alors $F(E) = F(X)$,
et l'on a $E = X$ puisque $F$ reflète les isomorphismes.

\medskip\noindent
Démonstration de (2). Cela vient de ce que $F$ transforme le morphisme $X \to X \times_Y X$
en l'application $F(X) \to F(X) \times_{F(Y)} F(X)$ et que $F$ reflète les isomorphismes.

\medskip\noindent
Démonstration de (3). Cela vient de ce que $F$ transforme le morphisme $Y \amalg_X Y \to Y$
en l'application $F(Y) \amalg_{F(X)} F(Y) \to F(Y)$ et que $F$ reflète les isomorphismes.

\medskip\noindent
Démonstration de (4). Il existe un objet initial $A$ et, bien entendu,
$F(A) = \emptyset$. Réciproquement, si $X$ est un objet tel que
$F(X) = \emptyset$, alors l'unique morphisme $A \to X$ induit une bijection
$F(A) \to F(X)$ ; par conséquent, $A \to X$ est un isomorphisme.

\medskip\noindent
Démonstration de (5). Il existe un objet final $Z$ et, bien entendu,
$F(Z)$ est un ensemble à un élément. Réciproquement, si $X$ est un objet tel que
$F(X)$ soit un ensemble à un élément, alors l'unique morphisme $X \to Z$ induit une bijection
$F(X) \to F(Z)$ ; par conséquent, $X \to Z$ est un isomorphisme.

\medskip\noindent
Démonstration de (6). L'égalisateur $E$ des deux morphismes $Y \to Y \amalg_X Y$ n'est pas
un objet initial de $\mathcal{C}$, car $X \to Y$ se factorise par $E$
et $F(X) \not = \emptyset$. Ainsi $E = Y$, d'où la conclusion.

\medskip\noindent
Démonstration de (\ref{item-one-element}).
L'égalisateur $E$ de $a$ et $b$ est muni d'un monomorphisme
$E \to X$, et $F(E) \subset F(X)$ est l'ensemble des éléments sur lesquels
$F(a)$ et $F(b)$ coïncident. Pour conclure, on utilise le fait que $E$ est initial,
ou que $E = X$.

\medskip\noindent
Démonstration de (8). Pour tous $i, j$, on voit que $E_{ij} = X_i \times_Y Y_j$
est soit initial, soit égal à $X_i$. En choisissant $s \in F(X_i)$,
on voit que $E_{ij} = X_i$ si et seulement si $s$ s'envoie sur un élément
de $F(Y_j) \subset F(Y)$, ce qui arrive pour un unique $j = \alpha(i)$.
\end{proof}

\noindent
Le lemme précédent montre que, pour tout objet connexe $X$ d'une
catégorie galoisienne $(\mathcal{C}, F)$, le groupe d'automorphismes
$\text{Aut}(X)$ est d'ordre au plus $|F(X)|$. En effet, si $s \in F(X)$
et $g \in \text{Aut}(X)$, alors $F(g)(s) = s$ si et seulement
si $g = \text{id}_X$, d'après (\ref{item-one-element}).
On dit que $X$ est {\it galoisien} si l'égalité a lieu.
De manière équivalente, $X$ est galoisien s'il est connexe et si
$\text{Aut}(X)$ agit transitivement sur $F(X)$.

\begin{lemma}
\label{lemma-galois}
Soit $(\mathcal{C}, F)$ une catégorie galoisienne. Pour tout objet connexe $X$
de $\mathcal{C}$, il existe un objet galoisien $Y$ et un morphisme $Y \to X$.
\end{lemma}

\begin{proof}
Nous utiliserons sans autre mention les résultats du lemme \ref{lemma-epi-mono}.
Posons $n = |F(X)|$. Considérons $X^n$ muni de son action naturelle de
$S_n$. Soit
$$
X^n = \coprod\nolimits_{t \in T} Z_t
$$
sa décomposition en objets connexes. Choisissons $t$ de sorte que
$F(Z_t)$ contienne $(s_1, \ldots, s_n)$, où les $s_i$ sont deux à deux distincts.
Si $(s'_1, \ldots, s'_n) \in F(Z_t)$ est un autre élément, nous
affirmons que les $s'_i$ sont eux aussi deux à deux distincts. Dans le cas contraire, disons
$s'_i = s'_j$ ; alors $Z_t$ est l'image d'une composante connexe de
$X^{n - 1}$ par le morphisme diagonal
$$
\Delta_{ij} : X^{n - 1} \longrightarrow X^n
$$
Puisque les morphismes entre objets connexes sont des épimorphismes et induisent
des surjections après application de $F$, on en déduirait que $s_i = s_j$,
ce qui n'est pas le cas.

\medskip\noindent
Soit $G \subset S_n$ le sous-groupe des éléments tels que $g(Z_t) = Z_t$.
En considérant l'action de $S_n$ sur
$$
F(X)^n = F(X^n) = \coprod\nolimits_{t' \in T} F(Z_{t'})
$$
on voit que $G = \{g \in S_n \mid g(s_1, \ldots, s_n) \in F(Z_t)\}$.
Choisissons maintenant un second élément $(s'_1, \ldots, s'_n) \in F(Z_t)$.
Nous avons vu plus haut que les $s'_i$ sont deux à deux distincts. On peut donc
trouver $g \in S_n$ tel que $g(s_1, \ldots, s_n) = (s'_1, \ldots, s'_n)$.
Autrement dit, l'action de $G$ sur $F(Z_t)$ est transitive, ce qui
achève la démonstration.
\end{proof}

\noindent
Voici un lemme essentiel.

\begin{lemma}
\label{lemma-tame}
\begin{reference}
Comparer avec \cite[Définition 7.2.4]{BS}.
\end{reference}
Soit $(\mathcal{C}, F)$ une catégorie galoisienne. Soit $G = \text{Aut}(F)$
comme dans l'exemple \ref{example-from-C-F-to-G-sets}. Pour tout objet connexe
$X$ de $\mathcal{C}$, l'action de $G$ sur $F(X)$ est transitive.
\end{lemma}

\begin{proof}
Nous utiliserons sans autre mention les résultats du lemme \ref{lemma-epi-mono}.
Soit $I$ l'ensemble des classes d'isomorphie d'objets galoisiens de $\mathcal{C}$.
Pour chaque $i \in I$, soit $X_i$ un représentant de la classe d'isomorphie.
Choisissons $\gamma_i \in F(X_i)$ pour tout $i \in I$.
Nous munissons $I$ d'une relation d'ordre en posant $i \geq i'$ si
et seulement s'il existe un morphisme $f_{ii'} : X_i \to X_{i'}$.
Étant donné un tel morphisme, on peut le composer à gauche par un automorphisme
$X_{i'} \to X_{i'}$ de façon que $F(f_{ii'})(\gamma_i) = \gamma_{i'}$.
Avec cette normalisation, le morphisme $f_{ii'}$ est unique.
Remarquons que $I$ est un ensemble ordonné filtrant :
(Catégories, définition \ref{categories-definition-directed-set})
si $i_1, i_2 \in I$, il existe un objet galoisien $Y$ et un morphisme
$Y \to X_{i_1} \times X_{i_2}$ par le lemme \ref{lemma-galois}, appliqué
à une composante connexe de $X_{i_1} \times X_{i_2}$.
Alors $Y \cong X_i$ pour un certain $i \in I$, et $i \geq i_1$, $i \geq i_2$.

\medskip\noindent
Nous affirmons que le foncteur $F$ est isomorphe au foncteur $F'$
qui à $X$ associe
$$
F'(X) = \colim_I \Mor_\mathcal{C}(X_i, X)
$$
au moyen de la transformation de foncteurs $t : F' \to F$ définie comme suit :
pour $f : X_i \to X$, on pose $t_X(f) = F(f)(\gamma_i)$.
En utilisant (\ref{item-one-element}), on constate que $t_X$ est injective.
Pour montrer la surjectivité, soit $\gamma \in F(X)$. On peut aussitôt
se ramener au cas où $X$ est connexe, par la définition d'une
catégorie galoisienne. On peut ensuite supposer $X$ galoisien grâce au
lemme \ref{lemma-galois}. Dans ce cas, $X$ est isomorphe à $X_i$
pour un certain $i$, et l'on peut choisir l'isomorphisme $X_i \to X$ de sorte
que $\gamma_i$ s'envoie sur $\gamma$ (par définition des objets galoisiens).
On conclut que $t$ est un isomorphisme.

\medskip\noindent
Posons $A_i = \text{Aut}(X_i)$.
Nous affirmons que, si $i \geq i'$, il existe une application canonique
$h_{ii'} : A_i \to A_{i'}$ telle que, pour tout $a \in A_i$,
le diagramme
$$
\xymatrix{
X_i \ar[d]_a \ar[r]_{f_{ii'}} & X_{i'} \ar[d]^{h_{ii'}(a)} \\
X_i \ar[r]^{f_{ii'}} & X_{i'}
}
$$
soit commutatif. Il suffit de définir $h_{ii'}(a) = a' : X_{i'} \to X_{i'}$
comme l'unique automorphisme tel que
$F(a')(\gamma_{i'}) = F(f_{ii'} \circ a)(\gamma_i)$.
Comme précédemment, cela rend le diagramme commutatif et, de plus, ce choix
est unique.
Il s'ensuit que
$h_{i'i''} \circ h_{ii'} = h_{ii''}$
si $i \geq i' \geq i''$.
Puisque $F(X_i) \to F(X_{i'})$ est surjective, on voit que
$A_i \to A_{i'}$ est surjective.
En passant à la limite projective, on obtient un groupe
$$
A = \lim_I A_i
$$
C'est un groupe profini, puisque les groupes d'automorphismes sont finis.
L'application $A \to A_i$ est surjective pour tout $i$, d'après
Catégories, lemme \ref{categories-lemma-nonempty-limit}.

\medskip\noindent
Puisque les éléments de $A$ agissent sur le système projectif des $X_i$, on obtient une action de
$A$ (à droite) sur $F'$ par précomposition. Autrement dit, on obtient
un homomorphisme $A^{opp} \to G$. Puisque $A \to A_i$ est surjective, on conclut
que $G$ agit transitivement sur $F(X_i)$ pour tout $i$. Comme tout objet connexe
est dominé par l'un des $X_i$, le lemme en découle.
\end{proof}

\begin{proposition}
\label{proposition-galois}
\begin{reference}
C'est une version faible de \cite[Expos\'e V]{SGA1}.
La démonstration est empruntée à \cite[Théorème 7.2.5]{BS}.
\end{reference}
Soit $(\mathcal{C}, F)$ une catégorie galoisienne. Soit $G = \text{Aut}(F)$
comme dans l'exemple \ref{example-from-C-F-to-G-sets}. Le foncteur
$F : \mathcal{C} \to \textit{Finite-}G\textit{-Ensembles}$
(\ref{equation-remember}) est une équivalence.
\end{proposition}

\begin{proof}
Nous utiliserons sans plus les mentionner les résultats du lemme \ref{lemma-epi-mono}.
En particulier, nous savons que le foncteur est fidèle.
Par le lemme \ref{lemma-tame}, nous savons que pour tout $X$ connexe,
l'action de $G$ sur $F(X)$ est transitive. Ainsi, pour $X$ quelconque, le foncteur
$F$ fait correspondre les composantes connexes de $X$ (dont l'existence résulte
de l'un des axiomes d'une catégorie galoisienne) aux orbites de $G$ sur $F(X)$.
Soient $X$ et $Y$ des objets et soit
$s : F(X) \to F(Y)$ une application équivariante. Alors le graphe
$\Gamma_s \subset F(X) \times F(Y)$ de $s$
est une réunion de composantes connexes. Il existe donc une
réunion de composantes connexes $Z$ de $X \times Y$,
munie d'un monomorphisme $Z \to X \times Y$,
telle que $F(Z) = \Gamma_s$. Puisque $F(Z) \to F(X)$ est bijective,
nous voyons que $Z \to X$ est un isomorphisme et nous concluons
que $s = F(f)$, où $f : X \cong Z \to Y$ est le composé.
Ainsi, $F$ est pleinement fidèle.

\medskip\noindent
Pour achever la démonstration, montrons que $F$ est essentiellement surjectif.
Il suffit de montrer que $G/H$ appartient à l'image essentielle pour
tout sous-groupe ouvert $H \subset G$ d'indice fini.
Par définition de la topologie sur $G$, il existe une famille finie
d'objets $X_i$ telle que
$$
\Ker(G \longrightarrow \prod\nolimits_i \text{Aut}(F(X_i)))
$$
soit contenu dans $H$. Nous pouvons supposer que $X_i$ est connexe
pour tout $i$. Nous pouvons choisir un objet galoisien $Y$ muni d'un morphisme
vers une composante connexe de $\prod X_i$ à l'aide du
lemme \ref{lemma-galois}. Choisissons un isomorphisme $F(Y) = G/U$
dans $G\textit{-sets}$ pour un certain sous-groupe ouvert $U \subset G$.
Comme $Y$ est galoisien, le groupe
$\text{Aut}(Y) = \text{Aut}_{G\textit{-Ensembles}}(G/U)$ opère transitivement
sur $F(Y) = G/U$. Cela implique que $U$ est distingué. Puisque
le morphisme induit de $F(Y)$ vers $F(X_i)$ est surjectif pour tout $i$, nous voyons que
$U \subset H$. Soit $M \subset \text{Aut}(Y)$ le sous-groupe fini
correspondant à
$$
(H/U)^{opp} \subset (G/U)^{opp} = \text{Aut}_{G\textit{-Ensembles}}(G/U)
= \text{Aut}(Y).
$$
Posons $X = Y/M$, c'est-à-dire que $X$ est le coégalisateur
des flèches $m : Y \to Y$, $m \in M$.
Comme $F$ est exact, nous voyons que $F(X) = G/H$, ce qui achève
la démonstration.
\end{proof}

\begin{lemma}
\label{lemma-functoriality-galois}
Soient $(\mathcal{C}, F)$ et $(\mathcal{C}', F')$ des catégories galoisiennes.
Soit $H : \mathcal{C} \to \mathcal{C}'$ un foncteur exact.
Il existe un isomorphisme $t : F' \circ H \to F$.
Le choix de $t$ détermine un homomorphisme continu
$h : G' = \text{Aut}(F') \to \text{Aut}(F) = G$ et
un diagramme $2$-commutatif
$$
\xymatrix{
\mathcal{C} \ar[r]_H \ar[d] & \mathcal{C}' \ar[d] \\
\textit{Finite-}G\textit{-Ensembles} \ar[r]^h &
\textit{Finite-}G'\textit{-Ensembles}
}
$$
L'homomorphisme $h$ ne dépend du choix de $t$ qu'à
un automorphisme intérieur de $G$ près.
Réciproquement, étant donné un homomorphisme continu $h : G' \to G$, il
existe un foncteur exact $H : \mathcal{C} \to \mathcal{C}'$ et un
isomorphisme $t$ permettant de retrouver $h$ comme ci-dessus.
\end{lemma}

\begin{proof}
D'après la proposition \ref{proposition-galois} et
le lemme \ref{lemma-single-out-profinite}, nous pouvons supposer que
$\mathcal{C} = \textit{Finite-}G\textit{-Ensembles}$ et que $F$ est le
foncteur d'oubli, et de même pour $\mathcal{C}'$. L'existence de
$t$ résulte alors du lemme \ref{lemma-second-fundamental-functor}. L'homomorphisme $h$
s'obtient par transport de structure au moyen de $t$. La commutativité du
diagramme est évidente. L'homomorphisme $h$ est unique à automorphisme intérieur
de $G$ près, car le choix de $t$ est
unique à un élément de $G$ près. La dernière assertion est immédiate.
\end{proof}





\section{Foncteurs et homomorphismes}
\label{section-translation}

\noindent
Soient $(\mathcal{C}, F)$, $(\mathcal{C}', F')$, $(\mathcal{C}'', F'')$
des catégories galoisiennes. Posons $G = \text{Aut}(F)$, $G' = \text{Aut}(F')$ et
$G'' = \text{Aut}(F'')$. Soient $H : \mathcal{C} \to \mathcal{C}'$
et $H' : \mathcal{C}' \to \mathcal{C}''$ des foncteurs exacts.
Soient $h : G' \to G$ et $h' : G'' \to G'$ les homomorphismes
continus correspondants comme dans le lemme \ref{lemma-functoriality-galois}.
Dans cette section, nous considérons le diagramme $2$-commutatif correspondant
\begin{equation}
\label{equation-translation}
\vcenter{
\xymatrix{
\mathcal{C} \ar[r]_H \ar[d] &
\mathcal{C}' \ar[r]_{H'} \ar[d] &
\mathcal{C}'' \ar[d] \\
\textit{Finite-}G\textit{-Ensembles} \ar[r]^h &
\textit{Finite-}G'\textit{-Ensembles} \ar[r]^{h'} &
\textit{Finite-}G''\textit{-Ensembles}
}
}
\end{equation}
et nous relions les propriétés d'exactitude de la suite
$1 \to G'' \to G' \to G \to 1$ aux propriétés des foncteurs $H$ et $H'$.

\begin{lemma}
\label{lemma-functoriality-galois-surjective}
Dans le diagramme (\ref{equation-translation}), les assertions suivantes sont équivalentes :
\begin{enumerate}
\item $h : G' \to G$ est surjectif,
\item $H : \mathcal{C} \to \mathcal{C}'$ est pleinement fidèle,
\item si $X \in \Ob(\mathcal{C})$ est connexe, alors $H(X)$ est connexe,
\item si $X \in \Ob(\mathcal{C})$ est connexe et s'il existe
un morphisme $*' \to H(X)$ dans $\mathcal{C}'$, alors
il existe un morphisme $* \to X$, et
\item pour tout objet $X$ de $\mathcal{C}$, l'application
$\Mor_\mathcal{C}(*, X) \to \Mor_{\mathcal{C}'}(*', H(X))$
est bijective.
\end{enumerate}
Ici, $*$ et $*'$ sont des objets finaux de $\mathcal{C}$ et $\mathcal{C}'$.
\end{lemma}

\begin{proof}
Les implications (5) $\Rightarrow$ (4) et (2) $\Rightarrow$ (5) sont évidentes.

\medskip\noindent
Supposons (3). Soit $X$ un objet connexe de $\mathcal{C}$ et soit
$*' \to H(X)$ un morphisme. Comme $H(X)$ est connexe d'après (3),
nous voyons que $*' \to H(X)$ est un isomorphisme. Ainsi, le $G'$-ensemble
correspondant à $H(X)$ possède exactement un élément, ce qui signifie que le
$G$-ensemble correspondant à $X$ possède un seul élément, donc que $X$ est
isomorphe à l'objet final de $\mathcal{C}$ ; en particulier,
il existe un morphisme $* \to X$. Nous avons ainsi montré que (3) $\Rightarrow$ (4).

\medskip\noindent
Si (1) est vraie, alors le foncteur
$\textit{Finite-}G\textit{-Ensembles} \to \textit{Finite-}G'\textit{-Ensembles}$
est pleinement fidèle : en effet, une application entre ensembles munis d'une action de $G$ commute avec
l'action de $G$ si et seulement si elle commute avec l'action de $G'$.
Ainsi, (1) $\Rightarrow$ (2).

\medskip\noindent
Si (1) est vraie, alors, pour un $G$-ensemble $X$, les orbites de $G$ et de $G'$
coïncident. Ainsi, (1) $\Rightarrow$ (3).

\medskip\noindent
Pour achever la démonstration, il suffit de montrer que (4) implique (1).
Si (1) est fausse, c'est-à-dire si $h$ n'est pas surjectif, il existe alors
un sous-groupe ouvert $U \subset G$ contenant $h(G')$ et distinct
de $G$. Le $G$-ensemble fini $M = G/U$ est alors muni d'une action transitive,
mais $G'$ y possède un point fixe. L'objet $X$ de $\mathcal{C}$
correspondant à $M$ contredirait (4). Nous voyons ainsi que
(4) $\Rightarrow$ (1), ce qui achève la démonstration.
\end{proof}

\begin{lemma}
\label{lemma-composition-trivial}
Dans le diagramme (\ref{equation-translation}), les assertions suivantes sont équivalentes :
\begin{enumerate}
\item $h \circ h'$ est trivial, et
\item l'image de $H' \circ H$ est formée d'objets isomorphes à des coproduits
finis d'objets finaux.
\end{enumerate}
\end{lemma}

\begin{proof}
Nous pouvons remplacer $H$ et $H'$ par les foncteurs canoniques
$\textit{Finite-}G\textit{-Ensembles} \to \textit{Finite-}G'\textit{-Ensembles}
\to \textit{Finite-}G''\textit{-Ensembles}$ déterminés par $h$ et $h'$.
Il s'agit alors de dire que l'action de $G''$ sur tout $G$-ensemble est triviale
si et seulement si l'homomorphisme $G'' \to G$ est trivial. Cela est évident.
\end{proof}

\begin{lemma}
\label{lemma-functoriality-galois-ses}
Dans le diagramme (\ref{equation-translation}), les assertions suivantes sont équivalentes :
\begin{enumerate}
\item la suite $G'' \xrightarrow{h'} G' \xrightarrow{h} G \to 1$
est exacte au sens suivant : $h$ est surjectif, $h \circ h'$ est trivial,
et $\Ker(h)$ est le plus petit sous-groupe fermé distingué contenant $\Im(h')$,
\item $H$ est pleinement fidèle et un objet $X'$ de $\mathcal{C}'$ appartient à
l'image essentielle de $H$ si et seulement si $H'(X')$ est isomorphe à un
coproduit fini d'objets finaux, et
\item $H$ est pleinement fidèle, $H' \circ H$ envoie tout objet sur un coproduit
fini d'objets finaux, et, pour tout objet $X'$ de $\mathcal{C}'$
tel que $H'(X')$ soit un coproduit fini d'objets finaux, il existe
un objet $X$ de $\mathcal{C}$ et un épimorphisme $H(X) \to X'$.
\end{enumerate}
\end{lemma}

\begin{proof}
D'après les lemmes \ref{lemma-functoriality-galois-surjective} et
\ref{lemma-composition-trivial}, nous pouvons supposer que
$H$ est pleinement fidèle, que $h$ est surjectif, que $H' \circ H$ envoie les
objets sur des coproduits de copies de l'objet final et que $h \circ h'$
est trivial. Soit $N \subset G'$ le plus petit sous-groupe fermé distingué
contenant l'image de $h'$. Il est clair que
$N \subset \Ker(h)$.
Nous pouvons supposer que les foncteurs $H$ et $H'$ sont les foncteurs canoniques
$\textit{Finite-}G\textit{-Ensembles} \to \textit{Finite-}G'\textit{-Ensembles}
\to \textit{Finite-}G''\textit{-Ensembles}$ déterminés par $h$ et $h'$.

\medskip\noindent
Supposons que (2) soit vraie. Cela signifie que, pour tout $G'$-ensemble fini $X'$
sur lequel $G''$ agit trivialement, l'action de $G'$ se factorise par $G$.
Appliquons ceci à $X' = G'/U'N$, où $U'$ est un sous-groupe ouvert suffisamment petit de $G'$.
Nous voyons alors que $\Ker(h) \subset U'N$ pour tout $U'$. Comme $N$ est fermé,
il s'ensuit que $\Ker(h) \subset N$, c'est-à-dire que (1) est vraie.

\medskip\noindent
Supposons que (1) soit vraie. Cela signifie que $N = \Ker(h)$. Soit $X'$ un
$G'$-ensemble fini sur lequel $G''$ agit trivialement. Cela signifie que
$\Ker(G' \to \text{Aut}(X'))$ est un sous-groupe fermé distingué contenant
$\Im(h')$. Ainsi, $N = \Ker(h)$ y est contenu et l'action de $G'$
sur $X'$ se factorise par $G$, c'est-à-dire que (2) est vraie.

\medskip\noindent
Supposons que (3) soit vraie. Cela signifie que, pour tout $G'$-ensemble fini $X'$
sur lequel $G''$ agit trivialement, il existe une surjection de $G'$-ensembles
$X \to X'$, où $X$ est un $G$-ensemble. Cela signifie manifestement que l'action de
$G'$ sur $X'$ se factorise par $G$, c'est-à-dire que (2) est vraie.

\medskip\noindent
L'implication (2) $\Rightarrow$ (3) est immédiate. Ceci achève la démonstration.
\end{proof}

\begin{lemma}
\label{lemma-functoriality-galois-injective}
Dans le diagramme (\ref{equation-translation}), les assertions suivantes sont équivalentes :
\begin{enumerate}
\item $h'$ est injectif, et
\item pour tout objet connexe $X''$ de $\mathcal{C}''$,
il existe un objet $X'$ de $\mathcal{C}'$ et un diagramme
$$
X'' \leftarrow Y'' \rightarrow H'(X')
$$
dans $\mathcal{C}''$, où $Y'' \to X''$ est un épimorphisme et
$Y'' \to H'(X')$ est un monomorphisme.
\end{enumerate}
\end{lemma}

\begin{proof}
Nous pouvons remplacer $H'$ par le foncteur correspondant entre les catégories
de $G'$-ensembles finis et de $G''$-ensembles finis.

\medskip\noindent
Supposons que $h' : G'' \to G'$ soit injectif. Soit $H'' \subset G''$
un sous-groupe ouvert. Puisque la topologie sur $G''$ est la topologie induite
par celle de $G'$, il existe un sous-groupe ouvert $H' \subset G'$
tel que $(h')^{-1}(H') \subset H''$.
Alors le diagramme voulu est
$$
G''/H'' \leftarrow G''/(h')^{-1}(H') \rightarrow G'/H'
$$
Réciproquement, supposons que (2) soit satisfaite pour le foncteur
$\textit{Finite-}G'\textit{-Ensembles} \to \textit{Finite-}G''\textit{-Ensembles}$.
Soit $g'' \in \Ker(h')$. Choisissons un sous-groupe ouvert quelconque $H'' \subset G''$.
Par hypothèse, il existe un $G'$-ensemble fini $X'$ et un diagramme
$$
G''/H'' \leftarrow Y'' \rightarrow X'
$$
de $G''$-ensembles, dont la flèche de gauche est surjective et celle de droite injective.
Comme $g''$ appartient au noyau de $h'$, nous voyons que $g''$ agit trivialement sur $X'$.
Ainsi, $g''$ agit trivialement sur $Y''$, et donc aussi sur $G''/H''$.
Par conséquent, $g'' \in H''$. Comme ceci vaut pour tout sous-groupe ouvert, nous concluons
que $g''$ est l'élément neutre, comme voulu.
\end{proof}

\begin{lemma}
\label{lemma-functoriality-galois-normal}
Dans le diagramme (\ref{equation-translation}), les assertions suivantes sont équivalentes :
\begin{enumerate}
\item l'image de $h'$ est distinguée, et
\item pour tout objet connexe $X'$ de $\mathcal{C}'$ tel
qu'il existe un morphisme de l'objet final de $\mathcal{C}''$
vers $H'(X')$, l'objet $H'(X')$ est isomorphe à un coproduit fini
d'objets finaux.
\end{enumerate}
\end{lemma}

\begin{proof}
Cela se traduit par l'énoncé suivant pour l'homomorphisme continu
de groupes $h' : G'' \to G'$ : l'image de $h'$ est distinguée
si et seulement si tout sous-groupe ouvert $U' \subset G'$ qui
contient $h'(G'')$ contient aussi chacun des conjugués de $h'(G'')$.
Le résultat s'en déduit aisément ; nous omettons certains détails.
\end{proof}






\section{Morphismes finis \'etales}
\label{section-finite-etale}

\noindent
Dans cette section, nous établissons suffisamment de résultats élémentaires sur les morphismes finis \'etales
pour pouvoir construire le groupe fondamental \'etale.

\medskip\noindent
Soit $X$ un schéma. Nous noterons $\textit{F\'Et}_X$
la catégorie des schémas finis et \'etales sur $X$.
Ainsi,
\begin{enumerate}
\item un objet de $\textit{F\'Et}_X$ est un morphisme fini \'etale
$Y \to X$ ayant pour but $X$, et
\item un morphisme dans $\textit{F\'Et}_X$
de $Y \to X$ vers $Y' \to X$ est un morphisme $Y \to Y'$ qui rend
commutatif le diagramme
$$
\xymatrix{
Y \ar[rr] \ar[rd] & &  Y' \ar[ld] \\
& X
}
$$
ci-dessus.
\end{enumerate}
Nous appellerons souvent un objet de $\textit{F\'Et}_X$ un
{\it revêtement \'etale fini} de $X$ (même si $Y$ est vide).
Il existe en fait un champ $p : \textit{F\'Et} \to \Sch$
sur la catégorie des schémas, dont la fibre au-dessus de $X$ est la catégorie
$\textit{F\'Et}_X$ que nous venons de définir. Voir Exemples de champs, section
\ref{examples-stacks-section-finite-etale}.

\begin{example}
\label{example-finite-etale-geometric-point}
Soient $k$ un corps algébriquement clos et $X = \Spec(k)$. Dans ce cas,
$\textit{F\'Et}_X$ est équivalente à la catégorie des ensembles finis. Ceci vaut
plus généralement lorsque $k$ est séparablement clos. La raison en est que
un schéma \'etale sur $k$ est une réunion disjointe de spectres de
corps qui sont des extensions finies séparables de $k$ ; voir
Morphismes, lemme \ref{morphisms-lemma-etale-over-field}.
\end{example}

\begin{lemma}
\label{lemma-finite-etale-covers-limits-colimits}
Soit $X$ un schéma. La catégorie $\textit{F\'Et}_X$ admet les limites finies et
les colimites finies, et, pour tout morphisme $X' \to X$, le foncteur de changement de base
$\textit{F\'Et}_X \to \textit{F\'Et}_{X'}$ est exact.
\end{lemma}

\begin{proof}
Limites finies et exactitude à gauche. D'après
Catégories, lemme \ref{categories-lemma-finite-limits-exist},
il suffit de montrer que $\textit{F\'Et}_X$ admet un objet final
et des produits fibrés. Cela est clair, car la catégorie de
tous les schémas sur $X$ admet un objet final (à savoir $X$) et des produits fibrés.
De plus, les produits fibrés de schémas finis \'etales sur $X$ sont
finis \'etales sur $X$. Il est en outre clair que le changement
de base commute à ces opérations et que, par conséquent, le changement de base
est exact à gauche (Catégories, lemme
\ref{categories-lemma-characterize-left-exact}).

\medskip\noindent
Colimites finies et exactitude à droite. D'après
Catégories, lemme \ref{categories-lemma-colimits-exist},
il suffit de montrer que $\textit{F\'Et}_X$ admet des sommes
finies et des coégalisateurs. Les sommes finies sont données
par les réunions disjointes (la somme vide est le schéma vide).
Soient $a, b : Z \to Y$ deux morphismes de $\textit{F\'Et}_X$.
Comme $Z \to X$ et $Y \to X$ sont finis \'etales, nous pouvons écrire
$Z = \underline{\Spec}(\mathcal{C})$ et $Y = \underline{\Spec}(\mathcal{B})$
pour certaines $\mathcal{O}_X$-algèbres localement libres de rang fini $\mathcal{C}$
et $\mathcal{B}$. Les morphismes $a, b$ induisent deux homomorphismes
$a^\sharp, b^\sharp : \mathcal{B} \to \mathcal{C}$.
Soit $\mathcal{A} = \text{Eq}(a^\sharp, b^\sharp)$ leur
égalisateur. Si
$$
\underline{\Spec}(\mathcal{A}) \longrightarrow X
$$
est fini \'etale, il est clair qu'il s'agit du coégalisateur
(en effet, tout objet de $\textit{F\'Et}_X$ peut s'écrire
comme le spectre relatif d'un faisceau de $\mathcal{O}_X$-algèbres).
Nous pouvons le vérifier après avoir remplacé $X$ par les éléments d'un recouvrement \'etale
(Descente, lemmes \ref{descent-lemma-descending-property-finite}
et \ref{descent-lemma-descending-property-etale}).
D'après Morphismes \'etales, lemme \ref{etale-lemma-finite-etale-etale-local},
nous pouvons supposer que
$Y = \coprod_{i = 1, \ldots, n} X$ et $Z = \coprod_{j = 1, \ldots, m} X$.
Alors
$$
\mathcal{C} = \prod\nolimits_{1 \leq j \leq m} \mathcal{O}_X
\quad\text{et}\quad
\mathcal{B} = \prod\nolimits_{1 \leq i \leq n} \mathcal{O}_X
$$
Après un nouveau remplacement par les éléments d'un recouvrement ouvert,
nous pouvons supposer que $a, b$ correspondent à des
applications $a_s, b_s : \{1, \ldots, m\} \to \{1, \ldots, n\}$, c'est-à-dire que
la composante $X$ de $Z$ d'indice $j$ est envoyée dans
la composante $X$ de $Y$ d'indice $a_s(j)$, resp.\ $b_s(j)$,
par le morphisme $a$, resp.\ $b$.
Soit $\{1, \ldots, n\} \to T$ le coégalisateur de $a_s, b_s$.
Alors nous voyons que
$$
\mathcal{A} = \prod\nolimits_{t \in T} \mathcal{O}_X
$$
dont le spectre est bien entendu fini \'etale sur $X$. Nous
omettons la vérification de la compatibilité avec le changement de base.
Ainsi, le changement de base est un foncteur exact à droite.
\end{proof}

\begin{remark}
\label{remark-colimits-commute-forgetful}
Soit $X$ un schéma. Considérons les foncteurs naturels
$F_1 : \textit{F\'Et}_X \to \Sch$ et $F_2 : \textit{F\'Et}_X \to \Sch/X$.
Alors
\begin{enumerate}
\item Les foncteurs $F_1$ et $F_2$ commutent aux colimites finies.
\item Le foncteur $F_2$ commute aux limites finies,
\item Le foncteur $F_1$ commute aux limites finies connexes, c'est-à-dire
aux égalisateurs et aux produits fibrés.
\end{enumerate}
Les assertions sur les limites résultent immédiatement de la discussion dans
la démonstration du lemme \ref{lemma-finite-etale-covers-limits-colimits}
et de Catégories, lemme \ref{categories-lemma-connected-limit-over-X}.
Il est clair que $F_1$ et $F_2$ commutent aux sommes finies.
D'après l'énoncé dual de Catégories, lemme
\ref{categories-lemma-characterize-left-exact},
il suffit de montrer que $F_1$ et $F_2$ commutent aux coégalisateurs.
Dans la démonstration du lemme \ref{lemma-finite-etale-covers-limits-colimits},
nous avons vu que les coégalisateurs dans $\textit{F\'Et}_X$ prennent localement pour la topologie \'etale
la forme suivante :
$$
\xymatrix{
\coprod_{j \in J} U \ar@<1ex>[r]^a \ar@<-1ex>[r]_b &
\coprod_{i \in I} U \ar[r] &
\coprod_{t \in \text{Coeq}(a, b)} U
}
$$
ce qui est manifestement un coégalisateur dans la catégorie des schémas.
L'assertion découle donc du fait que la propriété d'être un coégalisateur
est locale pour la topologie fpqc, sous la forme précise donnée dans
Descente, lemme \ref{descent-lemma-coequalizer-fpqc-local}.
\end{remark}

\begin{lemma}
\label{lemma-internal-hom-finite-etale}
Soit $X$ un schéma. Étant donnés $U, V$ finis \'etales sur $X$, il
existe un schéma $W$ fini \'etale sur $X$ tel que
$$
\Mor_X(X, W) = \Mor_X(U, V)
$$
et tel que cette égalité reste vraie après tout changement de base.
\end{lemma}

\begin{proof}
D'après Compléments sur les morphismes, lemme
\ref{more-morphisms-lemma-hom-from-finite-locally-free-separated-lqf},
il existe un schéma $W$ qui représente $\mathit{Mor}_X(U, V)$.
(On utilise le fait qu'un morphisme \'etale est localement quasi-fini d'après
Morphismes, lemme \ref{morphisms-lemma-etale-locally-quasi-finite},
et qu'un morphisme fini est séparé.)
Ce schéma satisfait manifestement la formule après tout changement de base.
Pour achever la démonstration, il reste à montrer que $W \to X$ est fini \'etale.
Nous pouvons le faire après avoir remplacé $X$ par les éléments d'un recouvrement \'etale
(Descente, lemmes \ref{descent-lemma-descending-property-finite}
et \ref{descent-lemma-descending-property-etale}).
D'après Morphismes \'etales, lemme \ref{etale-lemma-finite-etale-etale-local},
nous pouvons supposer que $U = \coprod_{i = 1, \ldots, n} X$
et $V = \coprod_{j = 1, \ldots, m} X$.
Dans ce cas,
$W = \coprod_{\alpha : \{1, \ldots, n\} \to \{1, \ldots, m\}} X$
immédiatement (nous omettons les détails), ce qui achève la démonstration.
\end{proof}

\noindent
Soit $X$ un schéma. Un {\it point géométrique} de $X$ est un morphisme
$\Spec(k) \to X$, où $k$ est algébriquement clos. Un tel point est
généralement noté $\overline{x}$, c'est-à-dire par une lettre minuscule surlignée.
Nous utilisons souvent $\overline{x}$ pour désigner aussi bien le schéma $\Spec(k)$ que
le morphisme, et nous notons $\kappa(\overline{x})$
le corps $k$. Nous disons que $\overline{x}$ {\it est au-dessus de} $x$
pour indiquer que $x \in X$ est l'image de $\overline{x}$.
Nous reviendrons plus en détail sur cette notion dans
Cohomologie \'etale, section \ref{etale-cohomology-section-stalks}.
Étant donnés $\overline{x}$ et un morphisme \'etale $U \to X$, nous pouvons
considérer
$$
|U_{\overline{x}}| : \text{l'ensemble sous-jacent des points du
schéma }U_{\overline{x}} = U \times_X \overline{x}
$$
Puisque le schéma $U_{\overline{x}}$, considéré comme schéma sur $\overline{x}$,
est une réunion disjointe de copies de $\overline{x}$
(Morphismes, lemme \ref{morphisms-lemma-etale-over-field}),
nous pouvons également décrire cet ensemble comme suit :
$$
|U_{\overline{x}}| =
\left\{
\begin{matrix}
\text{diagrammes} \\
\text{commutatifs}
\end{matrix}
\vcenter{
\xymatrix{
\overline{x} \ar[rd]_{\overline{x}} \ar[r]_{\overline{u}} & U \ar[d] \\
& X
}
}
\right\}
$$
L'association $U \mapsto |U_{\overline{x}}|$ est un foncteur
souvent noté $F_{\overline{x}}$.

\begin{lemma}
\label{lemma-finite-etale-connected-galois-category}
Soit $X$ un schéma connexe. Soit $\overline{x}$ un point géométrique.
Le foncteur
$$
F_{\overline{x}} : \textit{F\'Et}_X \longrightarrow \textit{Ensembles},\quad
Y \longmapsto |Y_{\overline{x}}|
$$
définit une catégorie galoisienne (définition \ref{definition-galois-category}).
\end{lemma}

\begin{proof}
Après avoir identifié $\textit{F\'Et}_{\overline{x}}$ à la catégorie des
ensembles finis (exemple \ref{example-finite-etale-geometric-point}),
on voit que notre foncteur $F_{\overline{x}}$ n'est autre que
le foncteur de changement de base pour le morphisme $\overline{x} \to X$.
Ainsi, $\textit{F\'Et}_X$ admet des limites finies et des colimites finies,
et $F_{\overline{x}}$ est exact d'après le
lemme \ref{lemma-finite-etale-covers-limits-colimits}.
Nous utiliserons aussi le fait que les limites finies dans $\textit{F\'Et}_X$
coïncident avec les limites finies correspondantes dans la catégorie
des schémas sur $X$ ; voir la remarque \ref{remark-colimits-commute-forgetful}.

\medskip\noindent
Si $Y' \to Y$ est un monomorphisme dans $\textit{F\'Et}_X$,
alors $Y' \to Y' \times_Y Y'$ est un isomorphisme ; par
conséquent, $Y' \to Y$ est un monomorphisme de schémas. Il s'ensuit que
$Y' \to Y$ est une immersion ouverte
(Morphismes étales, théorème \ref{etale-theorem-etale-radicial-open}). Puisque
$Y'$ est fini sur $X$ et que $Y$ est séparé sur $X$,
le morphisme $Y' \to Y$ est fini
(Morphismes, lemme \ref{morphisms-lemma-finite-permanence}), donc fermé
(Morphismes, lemme \ref{morphisms-lemma-finite-proper}) ;
c'est donc l'inclusion d'un sous-schéma ouvert et fermé de $Y$.
Il s'ensuit que $Y$ est un objet connexe de la catégorie
$\textit{F\'Et}_X$ (au sens de la définition \ref{definition-galois-category})
si et seulement si $Y$ est connexe comme schéma. Le
lemme \ref{topology-lemma-finite-fibre-connected-components} de Topologie
montre alors que $Y$ est le coproduit fini de ses composantes connexes,
tant comme schéma qu'au sens de la
définition \ref{definition-galois-category}.

\medskip\noindent
Soit $Y \to Z$ un morphisme de $\textit{F\'Et}_X$ qui induit une
bijection $F_{\overline{x}}(Y) \to F_{\overline{x}}(Z)$. Il faut
montrer que $Y \to Z$ est un isomorphisme. D'après ce qui précède, on peut
supposer $Z$ connexe. Puisque $Y \to Z$ est fini étale, donc fini
localement libre, il suffit de montrer que $Y \to Z$ est fini localement
libre de degré $1$. Cela est vrai au voisinage de tout point de
$Z$ situé au-dessus de $\overline{x}$ ; comme $Z$ est connexe et que
le degré est localement constant, on conclut.
\end{proof}



\section{Groupes fondamentaux}
\label{section-fundamental-groups}

\noindent
Dans cette section, nous définissons le groupe fondamental algébrique de Grothendieck.
La définition suivante a un sens grâce au
lemme \ref{lemma-finite-etale-connected-galois-category}.

\begin{definition}
\label{definition-fundamental-group}
Soit $X$ un schéma connexe. Soit $\overline{x}$ un point géométrique
de $X$. Le {\it groupe fondamental} de $X$, de
{\it point-base} $\overline{x}$, est le groupe
$$
\pi_1(X, \overline{x}) = \text{Aut}(F_{\overline{x}})
$$
des automorphismes du foncteur fibre
$F_{\overline{x}} : \textit{F\'Et}_X \to \textit{Ensembles}$,
muni de sa topologie profinie canonique définie dans le
lemme \ref{lemma-aut-inverse-limit}.
\end{definition}

\noindent
En combinant ce qui précède avec les résultats de la section \ref{section-galois},
on obtient le théorème suivant.

\begin{theorem}
\label{theorem-fundamental-group}
Soit $X$ un schéma connexe. Soit $\overline{x}$ un point géométrique
de $X$.
\begin{enumerate}
\item Le foncteur fibre $F_{\overline{x}}$ définit une équivalence de
catégories
$$
\textit{F\'Et}_X \longrightarrow
\textit{Finite-}\pi_1(X, \overline{x})\textit{-Ensembles}
$$
\item Étant donné un second point géométrique $\overline{x}'$ de $X$, il
existe un isomorphisme $t : F_{\overline{x}} \to F_{\overline{x}'}$.
Celui-ci donne un isomorphisme $\pi_1(X, \overline{x}) \to \pi_1(X, \overline{x}')$
compatible avec les équivalences de (1). Cet isomorphisme est
indépendant de $t$, à automorphisme intérieur près.
\item Étant donné un morphisme $f : X \to Y$ de schémas connexes, posons
$\overline{y} = f \circ \overline{x}$. Il existe un homomorphisme
continu canonique
$$
f_* : \pi_1(X, \overline{x}) \to \pi_1(Y, \overline{y})
$$
tel que le diagramme
$$
\xymatrix{
\textit{F\'Et}_Y \ar[r]_{\text{changement de base}} \ar[d]_{F_{\overline{y}}} &
\textit{F\'Et}_X \ar[d]^{F_{\overline{x}}} \\
\textit{Finite-}\pi_1(Y, \overline{y})\textit{-Ensembles} \ar[r]^{f_*} &
\textit{Finite-}\pi_1(X, \overline{x})\textit{-Ensembles}
}
$$
soit commutatif.
\end{enumerate}
\end{theorem}

\begin{proof}
L'assertion (1) résulte du lemme \ref{lemma-finite-etale-connected-galois-category}
et de la proposition \ref{proposition-galois}.
L'assertion (2) est un cas particulier du lemme \ref{lemma-functoriality-galois}.
Pour (3), observons que le diagramme
$$
\xymatrix{
\textit{F\'Et}_Y \ar[r] \ar[d]_{F_{\overline{y}}} &
\textit{F\'Et}_X \ar[d]^{F_{\overline{x}}} \\
\textit{Ensembles} \ar@{=}[r] & \textit{Ensembles}
}
$$
est commutatif (au sens strict, et pas seulement $2$-commutatif), car
$\overline{y} = f \circ \overline{x}$. On peut donc
appliquer le lemme \ref{lemma-functoriality-galois}, avec la transformation
de foncteurs qui s'en déduit, pour obtenir (3).
\end{proof}

\begin{lemma}
\label{lemma-fundamental-group-Galois-group}
Soit $K$ un corps et posons $X = \Spec(K)$. Soit $\overline{K}$ une
clôture algébrique et notons $\overline{x} : \Spec(\overline{K}) \to X$
le point géométrique correspondant. Soit $K^{sep} \subset \overline{K}$
la clôture séparable.
\begin{enumerate}
\item Le foncteur du lemme \ref{lemma-sheaves-point} induit une équivalence
$$
\textit{F\'Et}_X \longrightarrow
\textit{Finite-}\text{Gal}(K^{sep}/K)\textit{-Ensembles}.
$$
compatible avec $F_{\overline{x}}$ et avec le foncteur
$\textit{Finite-}\text{Gal}(K^{sep}/K)\textit{-Ensembles} \to \textit{Ensembles}$.
\item Elle induit un isomorphisme canonique
$$
\text{Gal}(K^{sep}/K) \longrightarrow \pi_1(X, \overline{x})
$$
de groupes topologiques profinis.
\end{enumerate}
\end{lemma}

\begin{proof}
Le foncteur du lemme \ref{lemma-sheaves-point} est le même que le foncteur
$F_{\overline{x}}$, car pour tout $Y$ étale sur $X$ on a
$$
\Mor_X(\Spec(\overline{K}), Y) = \Mor_X(\Spec(K^{sep}), Y)
$$
En effet, comme on l'a vu dans la preuve du lemme \ref{lemma-sheaves-point},
$Y = \coprod_{i \in I} \Spec(L_i)$, où $L_i/K$ est une extension séparable finie de $K$.
Par conséquent, tout homomorphisme de $K$-algèbres $L_i \to \overline{K}$ se factorise
par $K^{sep}$. Notons aussi que $F_{\overline{x}}(Y)$ est fini
si et seulement si $I$ est fini, si et seulement si $Y \to X$ est fini étale.
Ceci démontre (1).

\medskip\noindent
L'assertion (2) est une conséquence formelle de (1), du
lemme \ref{lemma-functoriality-galois} et du
lemme \ref{lemma-single-out-profinite}.
(Voir aussi la remarque ci-dessous.)
\end{proof}

\begin{remark}
\label{remark-variance}
Dans la situation du lemme \ref{lemma-fundamental-group-Galois-group},
donnons une construction plus explicite de l'isomorphisme
$\text{Gal}(K^{sep}/K) \to
\pi_1(X, \overline{x}) = \text{Aut}(F_{\overline{x}})$.
Observons que
$\text{Gal}(K^{sep}/K) = \text{Aut}(\overline{K}/K)$,
car $\overline{K}$ est la clôture parfaite de $K^{sep}$.
Puisque $F_{\overline{x}}(Y) = \Mor_X(\Spec(\overline{K}), Y)$,
on peut considérer l'application
$$
\text{Aut}(\overline{K}/K) \times F_{\overline{x}}(Y) \to F_{\overline{x}}(Y),
\quad
(\sigma, \overline{y}) \mapsto
\sigma \cdot \overline{y} = \overline{y} \circ \Spec(\sigma)
$$
Il s'agit d'une action, car
$$
\sigma\tau \cdot \overline{y} =
\overline{y} \circ \Spec(\sigma\tau) =
\overline{y} \circ \Spec(\tau) \circ \Spec(\sigma) =
\sigma \cdot (\tau \cdot \overline{y})
$$
L'action est fonctorielle en $Y \in \textit{F\'Et}_X$, et l'on
obtient ainsi l'application voulue.
\end{remark}





\section{Revêtements galoisiens des schémas connexes}
\label{section-finite-etale-under-galois}

\noindent
Soit $X$ un schéma connexe muni d'un point géométrique $\overline{x}$.
Puisque $F_{\overline{x}} : \textit{F\'Et}_X \to \textit{Ensembles}$ est le foncteur fibre d'une
catégorie galoisienne (lemme \ref{lemma-finite-etale-connected-galois-category}),
les résultats de la section \ref{section-galois} s'appliquent.
Dans cette section, nous transposons explicitement une partie de la terminologie
et des résultats au cadre des schémas et des morphismes finis étales.

\medskip\noindent
Nous dirons qu'un morphisme fini étale $Y \to X$ est un
{\it revêtement galoisien} si $Y$ définit un objet galoisien de
$\textit{F\'Et}_X$.
Pour un morphisme fini étale $Y \to X$, avec $G = \text{Aut}_X(Y)$,
les assertions suivantes sont équivalentes :
\begin{enumerate}
\item $Y$ est un revêtement galoisien de $X$,
\item $Y$ est connexe et $|G|$ est égal au degré de $Y \to X$,
\item $Y$ est connexe et $G$ agit transitivement sur $F_{\overline{x}}(Y)$, et
\item $Y$ est connexe et $G$ agit simplement transitivement sur
$F_{\overline{x}}(Y)$.
\end{enumerate}
Ceci résulte immédiatement de la discussion de la section \ref{section-galois}.

\medskip\noindent
Pour tout morphisme fini étale $f : Y \to X$ avec $Y$ connexe,
il existe un revêtement galoisien fini étale $Y' \to X$ qui domine $Y$
(lemme \ref{lemma-galois}).

\medskip\noindent
Les objets galoisiens de $\textit{F\'Et}_X$ correspondent, par l'équivalence
$$
F_{\overline{x}} : \textit{F\'Et}_X \to
\textit{Finite-}\pi_1(X, \overline{x})\textit{-Ensembles}
$$
du théorème \ref{theorem-fundamental-group},
aux $\pi_1(X, \overline{x})\textit{-Ensembles}$ finis
de la forme $G = \pi_1(X, \overline{x})/H$, où $H$ est un
sous-groupe ouvert distingué. De manière équivalente, si $G$ est un groupe fini
et si $\pi_1(X, \overline{x}) \to G$ est une surjection continue,
alors $G$, considéré comme un $\pi_1(X, \overline{x})$-ensemble, correspond
à un revêtement galoisien.

\medskip\noindent
Si $Y_i \to X$, $i = 1, 2$, sont des revêtements galoisiens finis étales
de groupes de Galois $G_i$, alors il existe un revêtement galoisien
fini étale $Y \to X$ dont le groupe de Galois est un sous-groupe de
$G_1 \times G_2$. En effet, considérons les homomorphismes continus
correspondants $\pi_1(X, \overline{x}) \to G_i$ et
prenons pour $G$ l'image de l'homomorphisme continu induit
$\pi_1(X, \overline{x}) \to G_1 \times G_2$.









\section{Invariance topologique du groupe fondamental}
\label{section-topological-invariance}

\noindent
Le résultat principal de cette section est qu'un homéomorphisme universel
de schémas connexes induit un isomorphisme sur les groupes fondamentaux.
Voir la proposition \ref{proposition-universal-homeomorphism}.

\medskip\noindent
Au lieu de démontrer directement que deux schémas ont le même groupe
fondamental, nous démontrons souvent que leurs catégories de revêtements
finis étales sont les mêmes. Cela implique bien sûr que
leurs groupes fondamentaux sont égaux, pourvu qu'ils soient connexes.

\begin{lemma}
\label{lemma-what-equivalence-gives}
Soit $f : X \to Y$ un morphisme de schémas quasi-compacts et quasi-séparés
tel que le foncteur de changement de base $\textit{F\'Et}_Y \to \textit{F\'Et}_X$
soit une équivalence de catégories. Alors
\begin{enumerate}
\item $f$ induit un homéomorphisme $\pi_0(X) \to \pi_0(Y)$,
\item si $X$, ou de manière équivalente $Y$, est connexe, alors
$\pi_1(X, \overline{x}) = \pi_1(Y, \overline{y})$.
\end{enumerate}
\end{lemma}

\begin{proof}
Soit $Y = Y_0 \amalg Y_1$ une décomposition en sous-schémas ouverts et fermés
non vides. Nous affirmons que $f(X)$ rencontre chacun des $Y_i$. En effet, sinon,
disons que $f(X) \subset Y_1$, nous pouvons considérer le morphisme fini étale
$V = Y_1 \to Y$. Celui-ci n'est pas un isomorphisme,
mais $V \times_Y X \to X$ en est un, ce qui
est une contradiction.

\medskip\noindent
Supposons que $X = X_0 \amalg X_1$ soit une décomposition en sous-schémas
ouverts et fermés. Considérons le morphisme fini étale $U = X_1 \to X$. Alors
$U = X \times_Y V$ pour un certain morphisme fini étale $V \to Y$. Le degré
du morphisme $V \to Y$ est localement constant ; on obtient donc une décomposition
$Y = \coprod_{d \geq 0} Y_d$ en sous-schémas ouverts et fermés
telle que $V \to Y$ soit de degré $d$ au-dessus de $Y_d$. Puisque
$f^{-1}(Y_d) = \emptyset$ pour $d > 1$, on conclut comme ci-dessus que $Y_d = \emptyset$
pour $d > 1$. Et l'on conclut que $f^{-1}(Y_i) = X_i$
pour $i = 0, 1$.

\medskip\noindent
Il s'ensuit que $f^{-1}$ induit une bijection entre l'ensemble des parties
ouvertes et fermées de $Y$ et celui des parties ouvertes et fermées de $X$.
Notons que $X$ et $Y$ sont des espaces spectraux ; voir Propriétés, lemme
\ref{properties-lemma-quasi-compact-quasi-separated-spectral}.
D'après Topologie, lemme \ref{topology-lemma-connected-component-intersection},
le treillis des parties ouvertes et fermées d'un espace spectral
détermine l'ensemble de ses composantes connexes.
Ainsi $\pi_0(X) \to \pi_0(Y)$ est bijective. Puisque $\pi_0(X)$ et
$\pi_0(Y)$ sont des espaces profinis
(Topologie, lemme \ref{topology-lemma-pi0-profinite}),
on conclut que $\pi_0(X) \to \pi_0(Y)$ est un homéomorphisme d'après
Topologie, lemme \ref{topology-lemma-bijective-map}. Ceci démontre (1).
L'assertion (2) est immédiate.
\end{proof}

\noindent
Le lemme suivant nous dit que le groupe fondamental d'un couple hensélien
est le groupe fondamental de la partie fermée.

\begin{lemma}
\label{lemma-gabber}
Soit $(A, I)$ un couple hensélien. Posons $X = \Spec(A)$ et $Z = \Spec(A/I)$.
Le foncteur
$$
\textit{F\'Et}_X \longrightarrow \textit{F\'Et}_Z,\quad
U \longmapsto U \times_X Z
$$
est une équivalence de catégories.
\end{lemma}

\begin{proof}
C'est la traduction géométrique du
lemme \ref{more-algebra-lemma-finite-etale-equivalence} de Compléments d'algèbre.
\end{proof}

\noindent
Le lemme suivant nous dit que le groupe fondamental d'un épaississement
est le même que celui du schéma de départ. Nous l'utiliserons
dans la preuve de la proposition forte sur les homéomorphismes universels
ci-dessous.

\begin{lemma}
\label{lemma-thickening}
Soit $X \subset X'$ un épaississement de schémas. Le foncteur
$$
\textit{F\'Et}_{X'} \longrightarrow \textit{F\'Et}_X,\quad
U' \longmapsto U' \times_{X'} X
$$
est une équivalence de catégories.
\end{lemma}

\begin{proof}
Pour une discussion des épaississements, voir
Compléments sur les morphismes, section \ref{more-morphisms-section-thickenings}.
Soit $U' \to X'$ un morphisme étale tel que $U = U' \times_{X'} X \to X$
soit fini étale. Alors $U' \to X'$ est lui aussi fini étale.
Cela résulte par exemple de Compléments sur les morphismes, lemme
\ref{more-morphisms-lemma-properties-that-extend-over-thickenings}.
Maintenant, si $X \subset X'$ est un épaississement d'ordre fini, cette remarque,
combinée avec Morphismes étales, théorème
\ref{etale-theorem-remarkable-equivalence},
démontre le lemme. Nous allons démontrer ci-dessous le lemme pour les épaississements
généraux, mais nous suggérons au lecteur de sauter la preuve.

\medskip\noindent
Soit $X' = \bigcup X_i'$ un recouvrement ouvert affine. Posons
$X_i = X \times_{X'} X_i'$, $X_{ij}' = X'_i \cap X'_j$,
$X_{ij} = X \times_{X'} X_{ij}'$, $X_{ijk}' = X'_i \cap X'_j \cap X'_k$,
$X_{ijk} = X \times_{X'} X_{ijk}'$.
Supposons que nous puissions démontrer
le résultat pour chacun des épaississements
$X_i \subset X'_i$, $X_{ij} \subset X_{ij}'$ et $X_{ijk} \subset X_{ijk}'$.
Alors le résultat pour $X \subset X'$ découle du recollement relatif des
schémas ; voir
Constructions, section \ref{constructions-section-relative-glueing}.
Observons que les schémas $X_i'$, $X_{ij}'$, $X_{ijk}'$ sont
tous séparés, comme sous-schémas ouverts de schémas affines. En répétant
l'argument une fois de plus, on se ramène au cas où les schémas
$X'_i$, $X_{ij}'$, $X_{ijk}'$ sont affines.

\medskip\noindent
Dans le cas affine, on a $X' = \Spec(A')$ et $X = \Spec(A'/I')$,
où $I'$ est un idéal localement nilpotent. Alors $(A', I')$ est un
couple hensélien (Compléments d'algèbre, lemme
\ref{more-algebra-lemma-locally-nilpotent-henselian}),
et le résultat découle du lemme \ref{lemma-gabber} (qui est
beaucoup plus facile dans ce cas).
\end{proof}

\noindent
La ``bonne'' manière de démontrer la proposition suivante serait de
la déduire de l'invariance du site étale ; voir
Cohomologie étale, théorème
\ref{etale-cohomology-theorem-topological-invariance}.

\begin{proposition}
\label{proposition-universal-homeomorphism}
Soit $f : X \to Y$ un homéomorphisme universel de schémas. Alors
$$
\textit{F\'Et}_Y \longrightarrow \textit{F\'Et}_X,\quad
V \longmapsto V \times_Y X
$$
est une équivalence. Ainsi, si $X$ et $Y$ sont connexes, alors
$f$ induit un isomorphisme $\pi_1(X, \overline{x}) \to \pi_1(Y, \overline{y})$
de groupes fondamentaux.
\end{proposition}

\begin{proof}
Rappelons qu'un homéomorphisme universel est la même chose qu'un morphisme
entier, universellement injectif et surjectif ; voir
Morphismes, lemme \ref{morphisms-lemma-universal-homeomorphism}.
En particulier, la diagonale $\Delta : X \to X \times_Y X$ est un épaississement
d'après Morphismes, lemme \ref{morphisms-lemma-universally-injective}.
Ainsi, d'après le lemme \ref{lemma-thickening},
étant donné un morphisme fini étale $U \to X$,
il existe un unique isomorphisme
$$
\varphi : U \times_Y X \to X \times_Y U
$$
de schémas finis étales sur $X \times_Y X$ dont l'image réciproque par
$\Delta$ est $\text{id} : U \to U$ sur $X$.
Puisque $X \to X \times_Y X \times_Y X$
est également un épaississement (c'est une immersion fermée bijective),
on conclut que $(U, \varphi)$ est une donnée de descente relativement à $X/Y$.
D'après Morphismes étales, proposition \ref{etale-proposition-effective},
on conclut que $U = X \times_Y V$ pour un certain $V \to Y$
quasi-compact, séparé et étale.
Nous omettons la preuve que $V \to Y$ est fini (indications :
le morphisme $U \to V$ est surjectif et $U \to Y$ est entier).
On conclut que $\textit{F\'Et}_Y \to \textit{F\'Et}_X$
est essentiellement surjectif.

\medskip\noindent
En raisonnant de la même manière que ci-dessus, on voit qu'étant donnés
$V_1 \to Y$ et $V_2 \to Y$ dans $\textit{F\'Et}_Y$, tout
morphisme $a : X \times_Y V_1 \to X \times_Y V_2$ sur $X$
est compatible avec les données de descente canoniques. Ainsi $a$
provient d'un morphisme $V_1 \to V_2$ sur $Y$ d'après
Morphismes étales, lemme \ref{etale-lemma-fully-faithful-cases}.
\end{proof}










\section{Revêtements finis étales des schémas propres}
\label{section-finite-etale-over-proper}

\noindent
Dans cette section, nous montrons que le groupe fondamental d'un schéma propre
connexe sur un anneau local hensélien est le même que le groupe fondamental
de sa fibre spéciale. Nous démontrons aussi une variante de ce résultat pour
un couple hensélien.

\medskip\noindent
Nous montrons également que le groupe
fondamental d'un schéma propre connexe sur un corps algébriquement clos $k$
ne change pas si l'on remplace $k$ par une extension algébriquement close.

\medskip\noindent
Au lieu d'énoncer et de démontrer les résultats dans le cas connexe,
nous les démontrons en général et laissons au lecteur le soin d'en déduire
le résultat pour les groupes fondamentaux à l'aide du
lemme \ref{lemma-what-equivalence-gives}.

\begin{lemma}
\label{lemma-finite-etale-on-proper-over-henselian}
Soit $A$ un anneau local hensélien. Soit $X$ un schéma propre sur $A$,
de fibre fermée $X_0$. Alors le foncteur
$$
\textit{F\'Et}_X \to \textit{F\'Et}_{X_0},\quad
U \longmapsto U_0 = U \times_X X_0
$$
est une équivalence de catégories.
\end{lemma}

\begin{proof}
La preuve donnée ici est un exemple d'application de l'algébrisation et de
l'approximation. Nous procédons en plusieurs étapes.

\medskip\noindent
Surjectivité essentielle lorsque $A$ est un anneau local noethérien complet.
Soit $X_n = X \times_{\Spec(A)} \Spec(A/\mathfrak m^{n + 1})$.
D'après Morphismes étales, théorème \ref{etale-theorem-remarkable-equivalence},
les inclusions
$$
X_0 \to X_1 \to X_2 \to \ldots
$$
induisent des équivalences de catégories entre la catégorie
des schémas étales sur $X_0$ et la catégorie des schémas
étales sur $X_n$.
De plus, si $U_n \to X_n$ correspond à un morphisme fini étale
$U_0 \to X_0$, alors $U_n \to X_n$ est également fini, par exemple
d'après Compléments sur les morphismes, lemme
\ref{more-morphisms-lemma-thicken-property-morphisms-cartesian}.
Dans ce cas, le morphisme $U_0 \to \Spec(A/\mathfrak m)$
est propre puisque $X_0$ est propre sur $A/\mathfrak m$. Nous pouvons donc appliquer
le théorème d'algébrisation de Grothendieck
(sous la forme du
lemme de Cohomologie des schémas
\ref{coherent-lemma-algebraize-formal-scheme-finite-over-proper})
pour voir qu'il existe un morphisme fini $U \to X$ dont la restriction
à $X_0$ redonne $U_0$. D'après Compléments sur les morphismes, lemme
\ref{more-morphisms-lemma-check-smoothness-on-infinitesimal-nbhds},
on voit que $U \to X$ est étale en tout point de $U_0$.
Cependant, puisque tout point de $U$ se spécialise en un point de $U_0$
(car $U$ est propre sur $A$), on conclut que $U \to X$ est étale.
On conclut ainsi que le foncteur est essentiellement surjectif.

\medskip\noindent
Pleine fidélité lorsque $A$ est un anneau local noethérien complet.
Soient $U \to X$ et $V \to X$ des morphismes finis étales, et
soit $\varphi_0 : U_0 \to V_0$ un morphisme sur $X_0$. Considérons
le morphisme
$$
\Gamma_{\varphi_0} : U_0 \longrightarrow U_0 \times_{X_0} V_0
$$
Ce morphisme est à la fois fini étale et une immersion fermée.
En appliquant la surjectivité essentielle au schéma $U \times_X V$, on trouve
un morphisme fini étale $W \to U \times_X V$ dont la fibre spéciale
est isomorphe à $\Gamma_{\varphi_0}$. Considérons la projection
$W \to U$. Elle est finie étale et est un isomorphisme au-dessus de $U_0$ par
construction. D'après Morphismes étales, lemme
\ref{etale-lemma-finite-etale-one-point},
$W \to U$ est un isomorphisme dans un voisinage ouvert de $U_0$.
C'est donc un isomorphisme, et la composée $\varphi : U \cong W \to V$
est le relèvement voulu de $\varphi_0$.

\medskip\noindent
Surjectivité essentielle lorsque $A$ est un G-anneau local noethérien hensélien.
Soit $U_0 \to X_0$ un morphisme fini étale.
Soit $A^\wedge$ le complété de $A$ pour la topologie définie par l'idéal maximal.
Soit $X^\wedge$ le changement de base de $X$ à $A^\wedge$.
D'après le résultat précédent, il existe un morphisme fini étale
$V \to X^\wedge$ dont la fibre spéciale est $U_0$.
Écrivons $A^\wedge = \colim A_i$, où $A \to A_i$ est de type fini.
D'après Limites, lemme \ref{limits-lemma-descend-finite-presentation},
il existe un indice $i$ et un morphisme de présentation finie $U_i \to X_{A_i}$
dont le changement de base à $X^\wedge$ est $V$. En augmentant $i$,
on peut supposer que $U_i \to X_{A_i}$ est fini et étale
(Limites, lemmes \ref{limits-lemma-descend-finite-finite-presentation} et
\ref{limits-lemma-descend-etale}). En écrivant
$$
A_i = A[x_1, \ldots, x_n]/(f_1, \ldots, f_m)
$$
l'homomorphisme d'anneaux $A_i \to A^\wedge$ peut se réinterpréter comme une solution
$(a_1, \ldots, a_n)$ dans $A^\wedge$ du système d'équations $f_j = 0$.
D'après Lissage des morphismes d'anneaux, théorème \ref{smoothing-theorem-approximation-property},
on peut approcher cette solution (à l'ordre $11$, par exemple) par une solution
$(b_1, \ldots, b_n)$ dans $A$. En retraduisant, on trouve un homomorphisme de $A$-algèbres
$A_i \to A$ qui donne le même point fermé que l'homomorphisme initial
$A_i \to A^\wedge$ (car $11 > 1$). Le changement de base $U \to X$ de $U_i \to X_{A_i}$
par cet homomorphisme d'anneaux sera donc un morphisme fini étale dont la
fibre spéciale est isomorphe à $U_0$.

\medskip\noindent
Pleine fidélité lorsque $A$ est un G-anneau local noethérien hensélien.
Elle se déduit de la surjectivité essentielle exactement de la même
manière que dans le cas où $A$ est noethérien complet.

\medskip\noindent
Cas général. Soit $(A, \mathfrak m)$ un anneau local hensélien.
Posons $S = \Spec(A)$ et notons $s \in S$ le point fermé. D'après Limites, lemme
\ref{limits-lemma-proper-limit-of-proper-finite-presentation-noetherian},
on peut écrire $X \to \Spec(A)$ comme une limite cofiltrante de
morphismes propres $X_i \to S_i$, où $S_i$ est de type fini sur $\mathbf{Z}$.
Pour chaque $i$, soit $s_i \in S_i$ l'image de $s$.
Puisque $S = \lim S_i$ et $A = \mathcal{O}_{S, s}$, on a
$A = \colim \mathcal{O}_{S_i, s_i}$. L'anneau $A_i = \mathcal{O}_{S_i, s_i}$
est un G-anneau local noethérien (Compléments d'algèbre, proposition
\ref{more-algebra-proposition-ubiquity-G-ring}).
D'après Compléments d'algèbre, lemme \ref{more-algebra-lemma-henselization-colimit},
on voit que $A = \colim A_i^h$. D'après
Compléments d'algèbre, lemme \ref{more-algebra-lemma-henselization-G-ring},
les anneaux $A_i^h$ sont des G-anneaux. Ainsi $A = \colim A_i^h$ et
$$
X = \lim (X_i \times_{S_i} \Spec(A_i^h))
$$
comme schémas. La catégorie des schémas finis étales sur $X$ est la limite
des catégories de schémas finis étales sur
$X_i \times_{S_i} \Spec(A_i^h)$ (d'après
Limites, lemmes
\ref{limits-lemma-descend-finite-presentation},
\ref{limits-lemma-descend-finite-finite-presentation} et
\ref{limits-lemma-descend-etale}).
Il en va de même pour les schémas finis étales sur
$X_0 = \lim (X_i \times_{S_i} s_i)$.
Ainsi, on déduit formellement le résultat pour $X / \Spec(A)$
du résultat pour les $(X_i \times_{S_i} \Spec(A_i^h)) / \Spec(A_i^h)$
traité ci-dessus.
\end{proof}

\begin{lemma}
\label{lemma-finite-etale-on-proper-over-henselian-pair}
Soit $(A, I)$ un couple hensélien. Soit $X$ un schéma propre sur $A$.
Posons $X_0 = X \times_{\Spec(A)} \Spec(A/I)$. Alors le foncteur
$$
\textit{F\'Et}_X \to \textit{F\'Et}_{X_0},\quad
U \longmapsto U_0 = U \times_X X_0
$$
est une équivalence de catégories.
\end{lemma}

\begin{proof}
La preuve de ce lemme est exactement la même que celle du
lemme \ref{lemma-finite-etale-on-proper-over-henselian}.

\medskip\noindent
Surjectivité essentielle lorsque $A$ est noethérien et complet pour la topologie $I$-adique.
Soit $X_n = X \times_{\Spec(A)} \Spec(A/I^{n + 1})$.
D'après Morphismes étales, théorème \ref{etale-theorem-remarkable-equivalence},
les inclusions
$$
X_0 \to X_1 \to X_2 \to \ldots
$$
induisent des équivalences de catégories entre la catégorie
des schémas étales sur $X_0$ et la catégorie des schémas
étales sur $X_n$.
De plus, si $U_n \to X_n$ correspond à un morphisme fini étale
$U_0 \to X_0$, alors $U_n \to X_n$ est également fini, par exemple
d'après Compléments sur les morphismes, lemme
\ref{more-morphisms-lemma-thicken-property-morphisms-cartesian}.
Dans ce cas, le morphisme $U_0 \to \Spec(A/I)$
est propre puisque $X_0$ est propre sur $A/I$. Nous pouvons donc appliquer
le théorème d'algébrisation de Grothendieck
(sous la forme du
lemme de Cohomologie des schémas
\ref{coherent-lemma-algebraize-formal-scheme-finite-over-proper})
pour voir qu'il existe un morphisme fini $U \to X$ dont la restriction
à $X_0$ redonne $U_0$. D'après Compléments sur les morphismes, lemme
\ref{more-morphisms-lemma-check-smoothness-on-infinitesimal-nbhds},
on voit que $U \to X$ est étale en tout point de $U_0$.
Cependant, puisque tout point de $U$ se spécialise en un point de $U_0$
(car $U$ est propre sur $A$), on conclut que $U \to X$ est étale.
On conclut ainsi que le foncteur est essentiellement surjectif.

\medskip\noindent
Pleine fidélité lorsque $A$ est noethérien et complet pour la topologie $I$-adique.
Soient $U \to X$ et $V \to X$ des morphismes finis étales, et
soit $\varphi_0 : U_0 \to V_0$ un morphisme sur $X_0$. Considérons
le morphisme
$$
\Gamma_{\varphi_0} : U_0 \longrightarrow U_0 \times_{X_0} V_0
$$
Ce morphisme est à la fois fini étale et une immersion fermée.
En appliquant la surjectivité essentielle au schéma $U \times_X V$, on trouve
un morphisme fini étale $W \to U \times_X V$ dont la fibre spéciale
est isomorphe à $\Gamma_{\varphi_0}$. Considérons la projection
$W \to U$. Elle est finie étale et est un isomorphisme au-dessus de $U_0$ par
construction. D'après Morphismes étales, lemme
\ref{etale-lemma-finite-etale-one-point},
$W \to U$ est un isomorphisme dans un voisinage ouvert de $U_0$.
C'est donc un isomorphisme, et la composée $\varphi : U \cong W \to V$
est le relèvement voulu de $\varphi_0$.

\medskip\noindent
Surjectivité essentielle lorsque $(A, I)$ est un couple hensélien
et que $A$ est un G-anneau noethérien.
Soit $U_0 \to X_0$ un morphisme fini étale.
Soit $A^\wedge$ le complété de $A$ pour la topologie $I$-adique.
Observons que $A^\wedge$ est un anneau noethérien qui est
complet pour la topologie $IA^\wedge$-adique ; voir Algèbre, lemmes
\ref{algebra-lemma-completion-complete} et
\ref{algebra-lemma-completion-Noetherian-Noetherian}.
Soit $X^\wedge$ le changement de base de $X$ à $A^\wedge$.
D'après le résultat précédent, il existe un morphisme fini étale
$V \to X^\wedge$ dont la fibre spéciale est $U_0$.
Écrivons $A^\wedge = \colim A_i$, où $A \to A_i$ est de type fini.
D'après Limites, lemme \ref{limits-lemma-descend-finite-presentation},
il existe un indice $i$ et un morphisme de présentation finie $U_i \to X_{A_i}$
dont le changement de base à $X^\wedge$ est $V$. En augmentant $i$,
on peut supposer que $U_i \to X_{A_i}$ est fini et étale
(Limites, lemmes \ref{limits-lemma-descend-finite-finite-presentation} et
\ref{limits-lemma-descend-etale}). En écrivant
$$
A_i = A[x_1, \ldots, x_n]/(f_1, \ldots, f_m)
$$
l'homomorphisme d'anneaux $A_i \to A^\wedge$ peut se réinterpréter comme une solution
$(a_1, \ldots, a_n)$ dans $A^\wedge$ du système d'équations $f_j = 0$.
D'après Lissage des morphismes d'anneaux, lemme \ref{smoothing-lemma-henselian-pair},
on peut approcher cette solution (à l'ordre $11$, par exemple) par une solution
$(b_1, \ldots, b_n)$ dans $A$. En retraduisant, on trouve un homomorphisme de $A$-algèbres
$A_i \to A$ qui donne le même point fermé que l'homomorphisme initial
$A_i \to A^\wedge$ (car $11 > 1$). Le changement de base $U \to X$ de $U_i \to X_{A_i}$
par cet homomorphisme d'anneaux sera donc un morphisme fini étale dont la
fibre spéciale est isomorphe à $U_0$.

\medskip\noindent
Pleine fidélité lorsque $(A, I)$ est un couple hensélien et que $A$ est un
G-anneau noethérien. Elle se déduit de la surjectivité essentielle exactement de la même
manière que dans le cas où $A$ est noethérien complet.

\medskip\noindent
Cas général. Soit $(A, I)$ un couple hensélien.
Posons $S = \Spec(A)$ et notons $S_0 = \Spec(A/I)$. D'après Limites, lemme
\ref{limits-lemma-proper-limit-of-proper-finite-presentation-noetherian},
on peut écrire $X \to \Spec(A)$ comme une limite cofiltrante de
morphismes propres $X_i \to S_i$, où $S_i$ est affine et de type fini
sur $\mathbf{Z}$. Écrivons $S_i = \Spec(A_i)$ et notons
$I_i \subset A_i$ l'image réciproque de $I$ par l'application $A_i \to A$.
Posons $S_{i, 0} = \Spec(A_i/I_i)$. Puisque $S = \lim S_i$, on a
$A = \colim A_i$. Ainsi, on a aussi $I = \colim I_i$ et $A/I = \colim A_i/I_i$.
L'anneau $A_i$ est un G-anneau noethérien (Compléments d'algèbre, proposition
\ref{more-algebra-proposition-ubiquity-G-ring}).
Notons $(A_i^h, I_i^h)$ l'hensélisation de la paire $(A_i, I_i)$.
D'après Compléments d'algèbre, lemme \ref{more-algebra-lemma-henselization-colimit},
on voit que $A = \colim A_i^h$. D'après
Compléments d'algèbre, lemme \ref{more-algebra-lemma-henselization-pair-G-ring},
les anneaux $A_i^h$ sont des G-anneaux. Ainsi $A = \colim A_i^h$ et
$$
X = \lim (X_i \times_{S_i} \Spec(A_i^h))
$$
comme schémas. La catégorie des schémas finis étales sur $X$ est la limite
des catégories de schémas finis étales sur
$X_i \times_{S_i} \Spec(A_i^h)$ (d'après
Limites, lemmes
\ref{limits-lemma-descend-finite-presentation},
\ref{limits-lemma-descend-finite-finite-presentation} et
\ref{limits-lemma-descend-etale}).
Il en va de même pour les schémas finis étales sur
$X_0 = \lim (X_i \times_{S_i} S_{i, 0})$.
Ainsi, on déduit formellement le résultat pour $X / \Spec(A)$
du résultat pour les $(X_i \times_{S_i} \Spec(A_i^h)) / \Spec(A_i^h)$
traité ci-dessus.
\end{proof}

\begin{lemma}
\label{lemma-finite-etale-invariant-over-proper}
Soit $k'/k$ une extension de corps algébriquement clos.
Soit $X$ un schéma propre sur $k$. Alors le foncteur
$$
U \longmapsto U_{k'}
$$
est une équivalence de catégories entre les schémas finis étales sur
$X$ et les schémas finis étales sur $X_{k'}$.
\end{lemma}

\begin{proof}
Démontrons que le foncteur est essentiellement surjectif.
Soit $U' \to X_{k'}$ un morphisme fini étale.
Écrivons $k' = \colim A_i$ comme colimite filtrante de $k$-algèbres de type fini.
D'après Limites, lemme \ref{limits-lemma-descend-finite-presentation},
il existe un indice $i$ et un morphisme de présentation finie $U_i \to X_{A_i}$
dont le changement de base à $X_{k'}$ est $U'$. En augmentant $i$,
on peut supposer que $U_i \to X_{A_i}$ est fini et étale
(Limites, lemmes \ref{limits-lemma-descend-finite-finite-presentation} et
\ref{limits-lemma-descend-etale}).
Puisque $k$ est algébriquement clos, on peut trouver un
point à valeurs dans $k$, disons $t$, de $\Spec(A_i)$. Soit $U = (U_i)_t$ la
fibre de $U_i$ au-dessus de $t$. Soit $A_i^h$
l'hensélisé de $(A_i)_{\mathfrak m}$, où $\mathfrak m$ est
l'idéal maximal correspondant au point $t$. D'après
le lemme \ref{lemma-finite-etale-on-proper-over-henselian},
on voit que $(U_i)_{A_i^h} = U \times \Spec(A_i^h)$ comme schémas
sur $X_{A_i^h}$. Maintenant, puisque
$A_i^h$ est algébrique sur $A_i$ (voir par exemple la discussion de
Lissage des morphismes d'anneaux, exemple \ref{smoothing-example-describe-henselian})
et puisque $k'$ est algébriquement clos,
on peut trouver un homomorphisme d'anneaux $A_i^h \to k'$ qui prolonge
l'inclusion donnée $A_i \subset k'$. On conclut donc que $U'$
est isomorphe au changement de base de $U$.
La preuve de la pleine fidélité est exactement la même.
\end{proof}








\section{Connexité locale}
\label{section-unibranch}

\noindent
Dans cette section, nous demandons quand $\pi_1(U) \to \pi_1(X)$ est surjectif
pour $U$ un ouvert dense d'un schéma $X$. Nous verrons que c'est
le cas (en gros) lorsque $U \cap B$ est connexe pour toute petite
``boule'' $B$ autour d'un point $x \in X \setminus U$.

\begin{lemma}
\label{lemma-dense-faithful}
Soit $f : X \to Y$ un morphisme de schémas. Si $f(X)$ est dense dans $Y$,
alors le foncteur de changement de base $\textit{F\'Et}_Y \to \textit{F\'Et}_X$
est fidèle.
\end{lemma}

\begin{proof}
Puisque la catégorie des revêtements finis étales possède un
Hom interne (lemme \ref{lemma-internal-hom-finite-etale}),
il suffit de démontrer ce qui suit : étant donné $W$ fini étale sur $Y$
et un morphisme $s : X \to W$ sur $Y$, il existe au plus une section
$t : Y \to W$ telle que $s = t \circ f$. Considérons deux sections
$t_1, t_2 : Y \to W$ telles que $s = t_1 \circ f = t_2 \circ f$.
Puisque l'égalisateur de $t_1$ et $t_2$ est fermé dans $Y$
(Schémas, lemme \ref{schemes-lemma-where-are-they-equal})
et puisque $f(X)$ est dense dans $Y$, on voit que $t_1$ et $t_2$
coïncident sur $Y_{red}$. Il s'ensuit que $t_1$ et $t_2$ ont
la même image, qui est un sous-schéma ouvert et fermé de $W$ s'envoyant
isomorphiquement sur $Y$
(Morphismes étales, proposition \ref{etale-proposition-properties-sections}) ;
ils sont donc égaux.
\end{proof}

\noindent
La condition du lemme suivant, selon laquelle le spectre épointé
du hensélisé strict est connexe, résulte par exemple de
l'hypothèse que l'anneau local est géométriquement unibranche ; voir
Compléments d'algèbre, lemme \ref{more-algebra-lemma-geometrically-unibranch}.
Il existe une réciproque partielle dans
Propriétés, lemme \ref{properties-lemma-geometrically-unibranch}.

\begin{lemma}
\label{lemma-same-etale-extensions}
Soit $(A, \mathfrak m)$ un anneau local. Posons $X = \Spec(A)$
et soit $U = X \setminus \{\mathfrak m\}$. Si le spectre épointé
du hensélisé strict de $A$ est connexe, alors
$$
\textit{F\'Et}_X \longrightarrow \textit{F\'Et}_U,\quad
Y \longmapsto Y \times_X U
$$
est un foncteur pleinement fidèle.
\end{lemma}

\begin{proof}
Supposons $A$ strictement hensélien. Dans ce cas, tout revêtement fini étale
$Y$ de $X$ est isomorphe à une réunion disjointe finie de
copies de $X$. Il suffit donc de démontrer que tout morphisme
$U \to U \amalg \ldots \amalg U$ sur $U$ se prolonge de manière unique en un morphisme
$X \to X \amalg \ldots \amalg X$ sur $X$.
Si $U$ est connexe (donc en particulier non vide), cela est vrai.

\medskip\noindent
Cas général. Puisque la catégorie des revêtements finis étales possède un
Hom interne (lemme \ref{lemma-internal-hom-finite-etale}),
il suffit de démontrer ce qui suit : étant donné $Y$ fini étale sur $X$,
tout morphisme $s : U \to Y$ sur $X$ se prolonge en un morphisme $t : X \to Y$
sur $X$. Soit $A^{sh}$ l'hensélisé strict de $A$ et notons
$X^{sh} = \Spec(A^{sh})$, $U^{sh} = U \times_X X^{sh}$,
$Y^{sh} = Y \times_X X^{sh}$. D'après le premier paragraphe et notre hypothèse
sur $A$, on peut prolonger le changement de base $s^{sh} : U^{sh} \to Y^{sh}$ de $s$ en
$t^{sh} : X^{sh} \to Y^{sh}$. Posons $A' = A^{sh} \otimes_A A^{sh}$.
Alors les deux images réciproques $t'_1, t'_2$ de $t^{sh}$ sur $X' = \Spec(A')$
sont des prolongements de l'image réciproque $s'$ de $s$ sur $U' = U \times_X X'$.
Comme $A \to A'$ est plat, on voit que $U' \subset X'$ est dense (topologiquement)
par descente des idéaux premiers pour $A \to A'$
(Algèbre, lemme \ref{algebra-lemma-flat-going-down}). Ainsi,
$t'_1 = t'_2$ d'après le lemme \ref{lemma-dense-faithful}.
Par conséquent, $t^{sh}$ provient d'un morphisme $t : X \to Y$,
par exemple d'après
Descente, lemme \ref{descent-lemma-fpqc-universal-effective-epimorphisms}.
\end{proof}

\noindent
Au vu du lemme \ref{lemma-same-etale-extensions},
il est intéressant de savoir quand le
spectre épointé d'un anneau (et celui de son hensélisé strict)
est connexe. Il existe un célèbre lemme dû à Hartshorne
qui donne une condition suffisante ; voir
Cohomologie locale, lemme
\ref{local-cohomology-lemma-depth-2-connected-punctured-spectrum}.

\begin{lemma}
\label{lemma-quasi-compact-dense-open-connected-at-infinity-Noetherian}
Soit $X$ un schéma. Soit $U \subset X$ un ouvert dense. Supposons que
\begin{enumerate}
\item l'espace topologique sous-jacent à $X$ soit noethérien, et
\item pour tout $x \in X \setminus U$, le spectre épointé du
hensélisé strict de $\mathcal{O}_{X, x}$ soit connexe.
\end{enumerate}
Alors $\textit{F\'Et}_X \to \textit{F\'Et}_U$ est pleinement fidèle.
\end{lemma}

\begin{proof}
Soient $Y_1, Y_2$ finis étales sur $X$ et soit
$\varphi : (Y_1)_U \to (Y_2)_U$ un morphisme sur $U$. Il faut montrer que
$\varphi$ se prolonge de manière unique en un morphisme $Y_1 \to Y_2$ sur $X$.
L'unicité découle du lemme \ref{lemma-dense-faithful}.

\medskip\noindent
Soit $x \in X \setminus U$ un point générique d'une composante irréductible
de $X \setminus U$. Posons $V = U \times_X \Spec(\mathcal{O}_{X, x})$.
Par notre choix de $x$, c'est le spectre épointé de
$\Spec(\mathcal{O}_{X, x})$. D'après le
lemme \ref{lemma-same-etale-extensions},
on peut prolonger de manière unique le morphisme $\varphi_V : (Y_1)_V \to (Y_2)_V$
en un morphisme
$(Y_1)_{\Spec(\mathcal{O}_{X, x})} \to (Y_2)_{\Spec(\mathcal{O}_{X, x})}$.
D'après Limites, lemme \ref{limits-lemma-glueing-near-point},
on trouve un ouvert $U \subset U'$ contenant $x$ et un prolongement
$\varphi' : (Y_1)_{U'} \to (Y_2)_{U'}$ de $\varphi$.
Puisque l'espace topologique sous-jacent à $X$ est noethérien,
ceci achève la preuve par récurrence noethérienne sur le complémentaire
de l'ouvert sur lequel $\varphi$ est défini.
\end{proof}

\begin{lemma}
\label{lemma-retrocompact-dense-open-connected-at-infinity-closed}
Soit $X$ un schéma. Soit $U \subset X$ un ouvert dense. Supposons que
\begin{enumerate}
\item $U \to X$ soit quasi-compact,
\item tout point de $X \setminus U$ soit fermé, et
\item pour tout $x \in X \setminus U$, le spectre épointé du
hensélisé strict de $\mathcal{O}_{X, x}$ soit connexe.
\end{enumerate}
Alors $\textit{F\'Et}_X \to \textit{F\'Et}_U$ est pleinement fidèle.
\end{lemma}

\begin{proof}
Soient $Y_1, Y_2$ finis étales sur $X$ et soit
$\varphi : (Y_1)_U \to (Y_2)_U$ un morphisme sur $U$. Il faut montrer que
$\varphi$ se prolonge de manière unique en un morphisme $Y_1 \to Y_2$ sur $X$.
L'unicité découle du lemme \ref{lemma-dense-faithful}.

\medskip\noindent
Soit $x \in X \setminus U$. Posons $V = U \times_X \Spec(\mathcal{O}_{X, x})$.
Puisque tout point de $X \setminus U$ est fermé, $V$ est le spectre épointé
de $\Spec(\mathcal{O}_{X, x})$. D'après le
lemme \ref{lemma-same-etale-extensions},
on peut prolonger de manière unique le morphisme $\varphi_V : (Y_1)_V \to (Y_2)_V$
en un morphisme
$(Y_1)_{\Spec(\mathcal{O}_{X, x})} \to (Y_2)_{\Spec(\mathcal{O}_{X, x})}$.
D'après Limites, lemme \ref{limits-lemma-glueing-near-point},
(ceci utilise que $U$ est rétrocompact dans $X$),
on trouve un ouvert $U \subset U'_x$ contenant $x$ et un prolongement
$\varphi'_x : (Y_1)_{U'_x} \to (Y_2)_{U'_x}$ de $\varphi$.
Notons qu'étant donnés deux points $x, x' \in X \setminus U$, les
morphismes $\varphi'_x$ et $\varphi'_{x'}$ coïncident sur
$U'_x \cap U'_{x'}$, car $U$ est dense dans cet ouvert
(lemme \ref{lemma-dense-faithful}). Ainsi, on peut prolonger $\varphi$
à $\bigcup U'_x = X$, comme souhaité.
\end{proof}

\begin{lemma}
\label{lemma-quasi-compact-dense-open-connected-at-infinity}
Soit $X$ un schéma. Soit $U \subset X$ un ouvert dense. Supposons que
\begin{enumerate}
\item tout ouvert quasi-compact de $X$ possède un nombre fini de
composantes irréductibles,
\item pour tout $x \in X \setminus U$, le spectre épointé du
hensélisé strict de $\mathcal{O}_{X, x}$ soit connexe.
\end{enumerate}
Alors $\textit{F\'Et}_X \to \textit{F\'Et}_U$ est pleinement fidèle.
\end{lemma}

\begin{proof}
Soient $Y_1, Y_2$ finis étales sur $X$ et soit
$\varphi : (Y_1)_U \to (Y_2)_U$ un morphisme sur $U$. Il faut montrer que
$\varphi$ se prolonge de manière unique en un morphisme $Y_1 \to Y_2$ sur $X$.
L'unicité découle du lemme \ref{lemma-dense-faithful}.
Nous allons démontrer l'existence en montrant que nous pouvons agrandir $U$
si $U \not = X$, puis conclure à l'aide du lemme de Zorn.

\medskip\noindent
Soit $x \in X \setminus U$ un point générique d'une composante irréductible
de $X \setminus U$. Posons $V = U \times_X \Spec(\mathcal{O}_{X, x})$.
Par notre choix de $x$, c'est le spectre épointé de
$\Spec(\mathcal{O}_{X, x})$. D'après le
lemme \ref{lemma-same-etale-extensions},
on peut prolonger le morphisme $\varphi_V : (Y_1)_V \to (Y_2)_V$
(de manière unique) en un morphisme
$(Y_1)_{\Spec(\mathcal{O}_{X, x})} \to (Y_2)_{\Spec(\mathcal{O}_{X, x})}$.
Choisissons un voisinage affine $W \subset X$ de $x$.
Puisque $U \cap W$ est dense dans $W$, il contient les points génériques
$\eta_1, \ldots, \eta_n$ de $W$. Choisissons un ouvert affine
$W' \subset W \cap U$ contenant $\eta_1, \ldots, \eta_n$.
Posons $V' = W' \times_X \Spec(\mathcal{O}_{X, x})$.
D'après Limites, lemme \ref{limits-lemma-glueing-near-point},
appliqué à $x \in W \supset W'$,
on trouve un ouvert $W' \subset W'' \subset W$ avec $x \in W''$
et un morphisme $\varphi'' : (Y_1)_{W''} \to (Y_2)_{W''}$
coïncidant avec $\varphi$ sur $W'$. Puisque $W'$ est dense dans
$W'' \cap U$, le lemme \ref{lemma-dense-faithful} montre
que $\varphi$ et $\varphi''$ coïncident sur $U \cap W''$.
Ainsi, $\varphi$ et $\varphi''$ se recollent en un morphisme
$\varphi'$ sur $U' = U \cup W''$ qui coïncide avec $\varphi$ sur $U$.
Observons que $x \in U'$, de sorte que nous avons prolongé $\varphi$
à un ouvert strictement plus grand.

\medskip\noindent
Considérons l'ensemble $\mathcal{S}$ des couples $(U', \varphi')$ où $U \subset U'$
et où $\varphi'$ est un prolongement de $\varphi$. Munissons $\mathcal{S}$
de l'ordre partiel évident. Si les couples $(U'_i, \varphi'_i)$
forment une chaîne, celle-ci admet une borne supérieure $(U', \varphi')$.
Il suffit de prendre $U' = \bigcup U'_i$ et de définir
$\varphi' : (Y_1)_{U'} \to (Y_2)_{U'}$ comme le morphisme
qui coïncide avec $\varphi'_i$ sur $U'_i$. Le lemme de Zorn s'applique donc,
et $\mathcal{S}$ possède un élément maximal. D'après l'argument précédent,
on voit que cet élément maximal est un prolongement de $\varphi$
sur tout $X$.
\end{proof}

\begin{lemma}
\label{lemma-local-exact-sequence}
Soit $(A, \mathfrak m)$ un anneau local. Posons $X = \Spec(A)$ et
$U = X \setminus \{\mathfrak m\}$. Soit $U^{sh}$ le spectre épointé
de l'hensélisé strict $A^{sh}$ de $A$.
Supposons $U$ quasi-compact et $U^{sh}$ connexe. Alors la suite
$$
\pi_1(U^{sh}, \overline{u}) \to \pi_1(U, \overline{u}) \to
\pi_1(X, \overline{u}) \to 1
$$
est exacte au sens du lemme \ref{lemma-functoriality-galois-ses}, partie (1).
\end{lemma}

\begin{proof}
L'application $\pi_1(U) \to \pi_1(X)$ est surjective d'après les
lemmes \ref{lemma-same-etale-extensions} et
\ref{lemma-functoriality-galois-surjective}.

\medskip\noindent
Écrivons $X^{sh} = \Spec(A^{sh})$. Soit $Y \to X$ un morphisme fini étale.
Alors $Y^{sh} = Y \times_X X^{sh} \to X^{sh}$ est un morphisme fini étale.
Puisque $A^{sh}$ est strictement hensélien, on voit que $Y^{sh}$ est isomorphe
à une réunion disjointe de copies de $X^{sh}$. Il en va donc de même pour
$Y \times_X U^{sh}$. Il s'ensuit que la composée
$\pi_1(U^{sh}) \to \pi_1(U) \to \pi_1(X)$ est triviale ; voir
le lemme \ref{lemma-composition-trivial}.

\medskip\noindent
Pour achever la preuve, il suffit, d'après le
lemme \ref{lemma-functoriality-galois-ses},
de montrer ce qui suit : étant donné un morphisme fini étale
$V \to U$ tel que $V \times_U U^{sh}$ soit une réunion
disjointe de copies de $U^{sh}$, on peut trouver un morphisme fini étale
$Y \to X$ avec $V \cong Y \times_X U$ sur $U$.
L'hypothèse implique qu'il existe un morphisme fini étale
$Y^{sh} \to X^{sh}$ et un isomorphisme
$V \times_U U^{sh} \cong Y^{sh} \times_{X^{sh}} U^{sh}$.
Considérons le diagramme suivant
$$
\xymatrix{
U \ar[d] & U^{sh} \ar[d] \ar[l] &
U^{sh} \times_U U^{sh} \ar[d] \ar@<1ex>[l] \ar@<-1ex>[l] &
U^{sh} \times_U U^{sh} \times_U U^{sh}
\ar[d] \ar@<1ex>[l] \ar[l] \ar@<-1ex>[l] \\
X & X^{sh} \ar[l] &
X^{sh} \times_X X^{sh} \ar@<1ex>[l] \ar@<-1ex>[l] &
X^{sh} \times_X X^{sh} \times_X X^{sh} \ar@<1ex>[l] \ar[l] \ar@<-1ex>[l]
}
$$
Puisque l'immersion ouverte $U \subset X$ est quasi-compacte par hypothèse, toutes les
flèches verticales sont des immersions ouvertes quasi-compactes.
Soit $\xi \in X^{sh} \times_X X^{sh}$ un point qui n'appartient pas
à $U^{sh} \times_U U^{sh}$. Alors $\xi$ est au-dessus du point
fermé $x^{sh}$ de $X^{sh}$.
Considérons l'homomorphisme d'anneaux locaux
$$
A^{sh} = \mathcal{O}_{X^{sh}, x^{sh}} \to
\mathcal{O}_{X^{sh} \times_X X^{sh}, \xi}
$$
déterminé par la première projection $X^{sh} \times_X X^{sh}$.
C'est une colimite filtrante d'homomorphismes locaux obtenus par
localisation d'homomorphismes d'anneaux étales.
Puisque $A^{sh}$ est strictement hensélien, on conclut que c'est un
isomorphisme. Comme ceci vaut pour tout $\xi$ du complémentaire,
il n'y a pas de spécialisations entre ces points ;
chacun de ces $\xi$ est donc un point fermé (on peut aussi le démontrer
directement). Puisque l'anneau local en $\xi$ est isomorphe
à $A^{sh}$, il est strictement hensélien et son spectre épointé est connexe.
Il en va de même pour les points $\xi$ de $X^{sh} \times_X X^{sh} \times_X X^{sh}$ qui ne sont pas
dans $U^{sh} \times_U U^{sh} \times_U U^{sh}$. Il résulte du
lemme \ref{lemma-retrocompact-dense-open-connected-at-infinity-closed}
que l'image réciproque par les flèches verticales induit des foncteurs pleinement fidèles sur
les catégories de schémas finis étales. Ainsi, la
donnée de descente canonique sur $V \times_U U^{sh}$ relativement au
recouvrement fpqc $\{U^{sh} \to U\}$ se transporte en une
donnée de descente pour $Y^{sh}$ relativement au recouvrement fpqc $\{X^{sh} \to X\}$.
Puisque $Y^{sh} \to X^{sh}$ est fini, donc affine, cette donnée de descente est
effective (Descente, lemme \ref{descent-lemma-affine}).
On obtient ainsi un morphisme affine $Y \to X$ et un isomorphisme
$Y \times_X X^{sh} \to Y^{sh}$ compatible avec les données de descente.
Par pleine fidélité des données de descente
(comme dans Descente, lemme \ref{descent-lemma-refine-coverings-fully-faithful}),
on obtient un isomorphisme $V \to U \times_X Y$.
Enfin, $Y \to X$ est fini étale puisque $Y^{sh} \to X^{sh}$ l'est ; voir
Descente, lemmes \ref{descent-lemma-descending-property-etale} et
\ref{descent-lemma-descending-property-finite}.
\end{proof}

\noindent
Soit $X$ un schéma irréductible. Soit $\eta \in X$ le point générique.
Le morphisme canonique $\eta \to X$ induit une application canonique
\begin{equation}
\label{equation-inclusion-generic-point}
\text{Gal}(\kappa(\eta)^{sep}/\kappa(\eta)) = \pi_1(\eta, \overline{\eta})
\longrightarrow \pi_1(X, \overline{\eta})
\end{equation}
L'identification du membre de gauche est donnée par le
lemme \ref{lemma-fundamental-group-Galois-group}.

\begin{lemma}
\label{lemma-irreducible-geometrically-unibranch}
Soit $X$ un schéma irréductible et géométriquement unibranche.
Pour tout ouvert non vide $U \subset X$, l'application canonique
$$
\pi_1(U, \overline{u}) \longrightarrow \pi_1(X, \overline{u})
$$
est surjective. L'application (\ref{equation-inclusion-generic-point})
$\pi_1(\eta, \overline{\eta}) \to \pi_1(X, \overline{\eta})$
est également surjective.
\end{lemma}

\begin{proof}
D'après le lemme \ref{lemma-thickening}, on peut remplacer $X$ par son réduit.
On peut donc supposer que $X$ est un schéma intègre. D'après le
lemme \ref{lemma-functoriality-galois-surjective},
l'assertion du lemme revient à affirmer que
les foncteurs $\textit{F\'Et}_X \to \textit{F\'Et}_U$ et
$\textit{F\'Et}_X \to \textit{F\'Et}_\eta$ sont pleinement fidèles.

\medskip\noindent
Le résultat pour $\textit{F\'Et}_X \to \textit{F\'Et}_U$ découle
du lemme \ref{lemma-quasi-compact-dense-open-connected-at-infinity}
et du fait que, pour un anneau local $A$ qui est
géométriquement unibranche, son hensélisé strict possède un
spectre irréductible. Voir
Compléments d'algèbre, lemme \ref{more-algebra-lemma-geometrically-unibranch}.

\medskip\noindent
Observons que le corps résiduel $\kappa(\eta) = \mathcal{O}_{X, \eta}$
est la colimite filtrante des $\mathcal{O}_X(U)$, où $U \subset X$
parcourt les ouverts affines non vides. Ainsi, $\textit{F\'Et}_\eta$ est la colimite des
catégories $\textit{F\'Et}_U$ sur de tels $U$ ; voir
Limites, lemmes \ref{limits-lemma-descend-finite-presentation},
\ref{limits-lemma-descend-finite-finite-presentation} et
\ref{limits-lemma-descend-etale}.
Un argument formel montre alors que la pleine fidélité de
$\textit{F\'Et}_X \to \textit{F\'Et}_\eta$ découle de la
pleine fidélité des foncteurs $\textit{F\'Et}_X \to \textit{F\'Et}_U$.
\end{proof}

\begin{lemma}
\label{lemma-exact-sequence-finite-nr-closed-pts}
Soit $X$ un schéma. Soient $x_1, \ldots, x_n \in X$ des points fermés en nombre fini
tels que
\begin{enumerate}
\item $U = X \setminus \{x_1, \ldots, x_n\}$ soit connexe et soit
un ouvert rétrocompact de $X$, et
\item pour chaque $i$, le spectre épointé $U_i^{sh}$ du
hensélisé strict de $\mathcal{O}_{X, x_i}$ soit connexe.
\end{enumerate}
Alors l'application $\pi_1(U) \to \pi_1(X)$ est surjective et son noyau
est le plus petit sous-groupe distingué fermé de $\pi_1(U)$ qui contienne
l'image de $\pi_1(U_i^{sh}) \to \pi_1(U)$ pour $i = 1, \ldots, n$.
\end{lemma}

\begin{proof}
La surjectivité découle des
lemmes \ref{lemma-retrocompact-dense-open-connected-at-infinity-closed} et
\ref{lemma-functoriality-galois-surjective}.
Nous pouvons considérer la suite d'applications
$$
\pi_1(U)  \to \ldots \to
\pi_1(X \setminus \{x_1, x_2\}) \to \pi_1(X \setminus \{x_1\}) \to \pi_1(X)
$$
Un argument de théorie des groupes montre alors qu'il suffit de démontrer l'assertion sur le
noyau dans le cas $n = 1$ (détails omis). Écrivons
$x = x_1$, $U^{sh} = U_1^{sh}$,
posons $A = \mathcal{O}_{X, x}$ et soit $A^{sh}$ l'hensélisé strict.
Considérons le diagramme
$$
\xymatrix{
U \ar[d] &
\Spec(A) \setminus \{\mathfrak m\} \ar[l] \ar[d] &
U^{sh} \ar[d] \ar[l] \\
X & \Spec(A) \ar[l] & \Spec(A^{sh}) \ar[l]
}
$$
D'après le lemme \ref{lemma-functoriality-galois-ses},
il faut montrer que les morphismes finis étales
$V \to U$ dont les images réciproques sont des revêtements triviaux de $U^{sh}$
se prolongent en des schémas finis étales sur $X$.
D'après le lemme \ref{lemma-local-exact-sequence},
nous connaissons l'énoncé correspondant
pour les schémas finis étales sur le spectre épointé de $A$.
Cependant, d'après Limites, lemme \ref{limits-lemma-glueing-near-closed-point},
les schémas de présentation finie sur $X$ s'identifient aux données formées de
schémas de présentation finie sur $U$ et sur $A$, recollés sur
le spectre épointé de $A$. Ceci achève la preuve.
\end{proof}









\section{Groupes fondamentaux des schémas normaux}
\label{section-normal}

\noindent
Soit $X$ un schéma intègre et géométriquement unibranche. Dans la section précédente,
nous avons vu que le groupe fondamental de $X$ est un quotient du
groupe de Galois du corps des fonctions rationnelles $K$ de $X$. Puisque l'application est continue,
son noyau est un sous-groupe distingué fermé du groupe de Galois. Ce noyau
correspond donc à une extension galoisienne $M/K$ par la théorie de Galois
(Corps, théorème \ref{fields-theorem-inifinite-galois-theory}).
Dans cette section, nous allons déterminer $M$ lorsque $X$ est un schéma intègre normal.

\medskip\noindent
Soit $X$ un schéma intègre normal de corps des fonctions rationnelles $K$.
Soit $L/K$ une extension finie. Considérons la normalisation
$Y \to X$ de $X$ dans le morphisme $\Spec(L) \to X$, telle qu'elle est définie dans
Morphismes, section \ref{morphisms-section-normalization-X-in-Y}.
Nous dirons (dans cette situation) que {\it $X$ est non ramifié dans $L$}
si $Y \to X$ est un morphisme non ramifié de schémas. Dans le
lemme \ref{lemma-unramified}, nous expliciterons cette condition.
Observons que la fibre schématique de $Y \to X$ au-dessus de $\Spec(K)$
est $\Spec(L)$. Ainsi, l'extension de corps $L/K$ est séparable si $X$ est
non ramifié dans $L$ ; voir
Morphismes, lemme \ref{morphisms-lemma-unramified-over-field}.

\begin{lemma}
\label{lemma-unramified-in-L}
Dans la situation ci-dessus, les conditions suivantes sont équivalentes :
\begin{enumerate}
\item $X$ est non ramifié dans $L$,
\item $Y \to X$ est étale, et
\item $Y \to X$ est fini étale.
\end{enumerate}
\end{lemma}

\begin{proof}
Observons que $Y \to X$ est un morphisme entier.
Dans chacun des cas, le morphisme $Y \to X$ est localement de type fini
par définition.
On trouve donc que, dans chacun des cas, $Y \to X$ est fini d'après
Morphismes, lemme \ref{morphisms-lemma-finite-integral}.
En particulier, on voit que (2) est équivalent à (3).
Un morphisme étale est non ramifié ; ainsi (2) implique (1).

\medskip\noindent
Réciproquement, supposons $Y \to X$ non ramifié. Puisqu'un schéma normal est
géométriquement unibranche
(Propriétés, lemme \ref{properties-lemma-normal-geometrically-unibranch}),
on voit que le morphisme $Y \to X$ est étale d'après Compléments sur les morphismes, lemme
\ref{more-morphisms-lemma-global-unramified-dominant-unibranch-is-etale}.
Nous donnons aussi une preuve directe dans le paragraphe suivant.

\medskip\noindent
Soit $x \in X$.
On peut choisir un voisinage étale $(U, u) \to (X, x)$ tel que
$$
Y \times_X U = \coprod V_j \longrightarrow U
$$
est une réunion disjointe d'immersions fermées, voir
Morphismes étales, lemme \ref{etale-lemma-finite-unramified-etale-local}.
En rétrécissant, on peut supposer $U$ quasi-compact.
Alors $U$ possède un nombre fini de composantes irréductibles
(Descente, lemme \ref{descent-lemma-locally-finite-nr-irred-local-fppf}).
Puisque $U$ est normal
(Descente, lemme \ref{descent-lemma-normal-local-smooth}), les
composantes irréductibles de $U$ sont ouvertes et fermées
(Propriétés, lemme \ref{properties-lemma-normal-locally-finite-nr-irreducibles})
et on peut supposer $U$ irréductible. Alors $U$ est un schéma intègre
dont le point générique $\xi$ s'envoie sur le point générique de $X$.
D'autre part, nous savons que $Y \times_X U$
est la normalisation de $U$ dans $\Spec(L) \times_X U$
d'après Compléments sur les morphismes, lemme
\ref{more-morphisms-lemma-normalization-smooth-localization}.
Tout point de $\Spec(L) \times_X U$ s'envoie sur $\xi$.
Ainsi, tout $V_j$ contient un point s'envoyant sur $\xi$ d'après
Morphismes, lemme \ref{morphisms-lemma-normalization-generic}.
Ainsi $V_j \to U$ est un isomorphisme puisque $U = \overline{\{\xi\}}$.
Par conséquent, $Y \times_X U \to U$ est étale. D'après
Descente, lemme \ref{descent-lemma-descending-property-etale},
on conclut que $Y \to X$ est étale au-dessus de
l'image de $U \to X$ (un voisinage ouvert de $x$).
\end{proof}

\begin{lemma}
\label{lemma-finite-etale-covering-normal-unramified}
Soit $X$ un schéma intègre normal de corps des fonctions rationnelles $K$.
Soit $Y \to X$ un morphisme fini étale. Si $Y$ est connexe,
alors $Y$ est un schéma intègre normal et $Y$ est la normalisation
de $X$ dans le corps des fonctions rationnelles de $Y$.
\end{lemma}

\begin{proof}
Le schéma $Y$ est normal d'après
Descente, lemme \ref{descent-lemma-normal-local-smooth}.
Puisque $Y \to X$ est plat, tout point générique de $Y$ s'envoie
sur le point générique de $X$ d'après
Morphismes, lemme \ref{morphisms-lemma-generalizations-lift-flat}.
Puisque $Y \to X$ est fini, on voit que $Y$ possède un nombre fini
de composantes irréductibles. Ainsi $Y$ est la réunion disjointe d'un
nombre fini de schémas intègres normaux d'après
Propriétés, lemme \ref{properties-lemma-normal-locally-finite-nr-irreducibles}.
Ainsi, si $Y$ est connexe, alors $Y$ est un schéma intègre normal.

\medskip\noindent
Soit $L$ le corps des fonctions rationnelles de $Y$ et soit $Y' \to X$ la normalisation
de $X$ dans $L$. D'après
Morphismes, lemme \ref{morphisms-lemma-characterize-normalization},
on obtient une factorisation $Y' \to Y \to X$ et $Y' \to Y$ est
la normalisation de $Y$ dans $L$. Puisque $Y$ est normal, il est clair
que $Y' = Y$ (cela peut aussi se déduire de
Morphismes, lemme \ref{morphisms-lemma-finite-birational-over-normal}).
\end{proof}

\begin{proposition}
\label{proposition-normal}
Soit $X$ un schéma intègre normal de corps des fonctions rationnelles $K$.
Alors l'application canonique (\ref{equation-inclusion-generic-point})
$$
\text{Gal}(K^{sep}/K) = \pi_1(\eta, \overline{\eta})
\longrightarrow \pi_1(X, \overline{\eta})
$$
s'identifie à la projection canonique
$\text{Gal}(K^{sep}/K) \to \text{Gal}(M/K)$ où $M \subset K^{sep}$
est la réunion des sous-extensions finies $L$
telles que $X$ est non ramifié dans $L$.
\end{proposition}

\begin{proof}
Le schéma normal $X$ est géométriquement unibranche
(Propriétés, lemme \ref{properties-lemma-normal-geometrically-unibranch}).
Le lemme \ref{lemma-irreducible-geometrically-unibranch} s'applique donc à $X$.
Ainsi $\pi_1(\eta, \overline{\eta}) \to \pi_1(X, \overline{\eta})$
est surjectif et la flèche horizontale supérieure du diagramme commutatif
$$
\xymatrix{
\textit{F\'Et}_X \ar[r] \ar[d] \ar[rd]_c & \textit{F\'Et}_\eta \ar[d] \\
\textit{Finite-}\pi_1(X, \overline{\eta})\textit{-sets} \ar[r] &
\textit{Finite-}\text{Gal}(K^{sep}/K)\textit{-sets}
}
$$
est pleinement fidèle. La flèche verticale de gauche est l'équivalence du
théorème \ref{theorem-fundamental-group}
et la flèche verticale de droite est l'équivalence du
lemme \ref{lemma-fundamental-group-Galois-group}. La flèche
horizontale inférieure est induite par l'application de la proposition.
D'après les lemmes \ref{lemma-unramified-in-L} et
\ref{lemma-finite-etale-covering-normal-unramified},
on voit que l'image essentielle de $c$
est constituée des $\text{Gal}(K^{sep}/K)\textit{-Ensembles}$ isomorphes
à des ensembles de la forme
$$
S = \Hom_K(\prod\nolimits_{i = 1, \ldots, n} L_i, K^{sep}) =
\coprod\nolimits_{i = 1, \ldots, n} \Hom_K(L_i, K^{sep})
$$
où $L_i/K$ est finie séparable et où $X$ est non ramifié dans $L_i$.
Ainsi, si $M \subset K^{sep}$ est comme dans l'énoncé du lemme,
alors $\text{Gal}(K^{sep}/M)$ est exactement le sous-groupe de
$\text{Gal}(K^{sep}/K)$ qui agit trivialement sur tout objet
de l'image essentielle de $c$. D'autre part, l'image essentielle de $c$
est exactement la catégorie des $S$ tels que l'action de $\text{Gal}(K^{sep}/K)$
se factorise par la surjection
$\text{Gal}(K^{sep}/K) \to \pi_1(X, \overline{\eta})$.
On conclut que $\text{Gal}(K^{sep}/M)$ est le noyau.
Ainsi $\text{Gal}(K^{sep}/M)$ est un sous-groupe distingué, $M/K$ est galoisienne,
et on a une suite exacte courte
$$
1 \to \text{Gal}(K^{sep}/M) \to
\text{Gal}(K^{sep}/K) \to
\text{Gal}(M/K) \to 1
$$
d'après la théorie de Galois (Corps, théorème
\ref{fields-theorem-inifinite-galois-theory} et
lemme \ref{fields-lemma-ses-infinite-galois}). La démonstration est terminée.
\end{proof}

\begin{lemma}
\label{lemma-local-exact-sequence-normal}
Soit $(A, \mathfrak m)$ un anneau local normal.
Posons $X = \Spec(A)$. Soit $A^{sh}$ l'hensélisé strict de $A$.
Soient $K$ et $K^{sh}$ les corps des fractions de $A$ et de $A^{sh}$.
Alors la suite
$$
\pi_1(\Spec(K^{sh})) \to \pi_1(\Spec(K)) \to \pi_1(X) \to 1
$$
est exacte au sens de la partie (1) du lemme \ref{lemma-functoriality-galois-ses}.
\end{lemma}

\begin{proof}
Observons que $A^{sh}$ est un domaine normal, voir
Compléments d'algèbre, lemme \ref{more-algebra-lemma-henselization-normal}.
L'application $\pi_1(\Spec(K)) \to \pi_1(X)$ est surjective d'après la
proposition \ref{proposition-normal}.

\medskip\noindent
Écrivons $X^{sh} = \Spec(A^{sh})$. Soit $Y \to X$ un morphisme fini étale.
Alors $Y^{sh} = Y \times_X X^{sh} \to X^{sh}$ est un morphisme fini étale.
Puisque $A^{sh}$ est strictement hensélien, on voit que $Y^{sh}$ est isomorphe
à une réunion disjointe de copies de $X^{sh}$. Il en va donc de même pour
$Y \times_X \Spec(K^{sh})$. Il s'ensuit que la composée
$\pi_1(\Spec(K^{sh})) \to \pi_1(X)$ est triviale, voir
lemme \ref{lemma-composition-trivial}.

\medskip\noindent
Pour achever la démonstration, il suffit, d'après le
lemme \ref{lemma-functoriality-galois-ses},
de montrer ce qui suit : étant donné un morphisme fini étale
$V \to \Spec(K)$ tel que $V \times_{\Spec(K)} \Spec(K^{sh})$
soit une réunion disjointe de copies de $\Spec(K^{sh})$, on peut trouver un
morphisme fini étale
$Y \to X$ avec $V \cong Y \times_X \Spec(K)$ sur $\Spec(K)$.
Écrivons $V = \Spec(L)$, de sorte que $L$ est un produit fini
d'extensions finies séparables de $K$.
Soit $B \subset L$ la clôture intégrale de $A$ dans $L$.
Si $A \to B$ est étale, alors on peut prendre $Y = \Spec(B)$
et la démonstration est terminée. D'après
Algèbre, lemme \ref{algebra-lemma-integral-closure-commutes-smooth}
(et un passage à la limite que nous omettons),
on voit que $B \otimes_A A^{sh}$ est la clôture intégrale de
$A^{sh}$ dans $L^{sh} = L \otimes_K K^{sh}$.
Notre hypothèse est que $L^{sh}$ est un produit de copies de
$K^{sh}$ et donc $B^{sh}$ est un produit de copies de $A^{sh}$.
Ainsi $A^{sh} \to B^{sh}$ est étale. Comme $A \to A^{sh}$ est
fidèlement plat, il s'ensuit que $A \to B$ est étale
(Descente, lemme \ref{descent-lemma-descending-property-etale}),
comme souhaité.
\end{proof}







\section{Actions de groupes et clôture intégrale}
\label{section-group-actions-integral}

\noindent
Dans cette section, on poursuit la discussion de
Compléments d'algèbre, section \ref{more-algebra-section-group-actions-integral}.
Rappelons qu'un anneau local normal est un domaine par définition.

\begin{lemma}
\label{lemma-get-algebraic-closure}
Soit $A$ un domaine normal dont le corps des fractions $K$ est séparablement
clos. Soit $\mathfrak p \subset A$ un idéal premier non nul.
Alors le corps résiduel $\kappa(\mathfrak p)$ est algébriquement clos.
\end{lemma}

\begin{proof}
Supposons le lemme faux afin d'obtenir une contradiction. Il existe alors un
polynôme unitaire irréductible $P(T) \in \kappa(\mathfrak p)[T]$ de
degré $d > 1$. Après avoir remplacé $P$ par $a^d P(a^{-1}T)$ pour un $a \in A$
convenable (afin d'éliminer les dénominateurs), on peut supposer que $P$ est l'image d'un
polynôme unitaire $Q$ dans $A[T]$. Observons que $Q$ est irréductible dans
$K[T]$. En effet, une factorisation sur $K$ conduit à une factorisation
sur $A$ d'après Algèbre, lemme \ref{algebra-lemma-polynomials-divide},
que l'on pourrait réduire modulo $\mathfrak p$ pour obtenir une factorisation de $P$.
Comme $K$ est séparablement clos, $Q$ n'est pas un polynôme séparable
(Corps, définition \ref{fields-definition-separable}).
La caractéristique de $K$ est alors $p > 0$ et $Q$ a
un terme linéaire nul (Corps, définition \ref{fields-definition-separable}).
Cependant, on peut alors remplacer $Q$ par
$Q + a T$, où $a \in \mathfrak p$ est non nul, pour obtenir une contradiction.
\end{proof}

\begin{lemma}
\label{lemma-normal-local-domain-separablly-closed-fraction-field}
Un anneau local normal dont le corps des fractions est séparablement clos est
strictement hensélien.
\end{lemma}

\begin{proof}
Soit $(A, \mathfrak m, \kappa)$ un anneau local normal dont le corps des fractions
$K$ est séparablement clos. Si $A = K$, alors c'est terminé. Sinon,
le corps résiduel $\kappa$ est algébriquement clos
d'après le lemme \ref{lemma-get-algebraic-closure} et il suffit de
vérifier que $A$ est hensélien.
Soit $f \in A[T]$ unitaire et soit $a_0 \in \kappa$ une racine
de multiplicité $1$ de la réduction $\overline{f} \in \kappa[T]$.
Soit $f = \prod f_i$ la factorisation dans $K[T]$.
D'après Algèbre, lemme \ref{algebra-lemma-polynomials-divide}, on a
$f_i \in A[T]$. Ainsi $a_0$ est une racine de $f_i$ pour un certain $i$.
Après avoir remplacé $f$ par $f_i$, on peut supposer $f$ irréductible.
Alors, puisque la dérivée $f'$ ne peut être nulle dans $A[T]$
car $a_0$ est une racine simple, on conclut que $f$ est linéaire
du fait que $K$ est séparablement clos.
Ainsi $A$ est hensélien, voir
Algèbre, définition \ref{algebra-definition-henselian}.
\end{proof}

\begin{lemma}
\label{lemma-inertia-base-change}
Soit $G$ un groupe fini agissant sur un anneau $R$. Soit $R^G \to A$ un morphisme
d'anneaux. Soit $\mathfrak q' \subset A \otimes_{R^G} R$ un idéal premier au-dessus
de l'idéal premier $\mathfrak q \subset R$. Alors
$$
I_\mathfrak q = \{\sigma \in G \mid
\sigma(\mathfrak q) = \mathfrak q\text{ et }
\sigma \bmod \mathfrak q = \text{id}_{\kappa(\mathfrak q)}\}
$$
est égal à
$$
I_{\mathfrak q'} = \{\sigma \in G \mid
\sigma(\mathfrak q') = \mathfrak q'\text{ et }
\sigma \bmod \mathfrak q' = \text{id}_{\kappa(\mathfrak q')}\}
$$
\end{lemma}

\begin{proof}
Puisque $\mathfrak q$ est l'image réciproque de $\mathfrak q'$
et puisque $\kappa(\mathfrak q) \subset \kappa(\mathfrak q')$,
on obtient $I_{\mathfrak q'} \subset I_\mathfrak q$.
Réciproquement, si $\sigma \in I_\mathfrak q$, alors $\sigma$
agit trivialement sur l'anneau de la fibre $A \otimes_{R^G} \kappa(\mathfrak q)$.
Ainsi $\sigma$ fixe tous les idéaux premiers au-dessus de $\mathfrak q$
et induit l'identité sur leurs corps résiduels.
\end{proof}

\begin{lemma}
\label{lemma-inertia-invariants-etale}
Soit $G$ un groupe fini agissant sur un anneau $R$. Soit $\mathfrak q \subset R$
un idéal premier. Posons
$$
I = \{\sigma \in G \mid \sigma(\mathfrak q) = \mathfrak q
\text{ et } \sigma \bmod \mathfrak q = \text{id}_{\kappa(\mathfrak q)}\}
$$
Alors $R^G \to R^I$ est étale en $R^I \cap \mathfrak q$.
\end{lemma}

\begin{proof}
La stratégie de la démonstration consiste à utiliser une localisation étale pour
se ramener au cas où $R^G \to R^I$ est un isomorphisme local en
$R^I \cap \mathfrak q$.
Soit $R^G \to A$ un morphisme étale d'anneaux. On affirme que si le résultat
vaut pour l'action de $G$ sur $A \otimes_{R^G} R$ et un idéal premier
$\mathfrak q'$ de $A \otimes_{R^G} R$ au-dessus de $\mathfrak q$, alors
le résultat est vrai.

\medskip\noindent
Pour le vérifier, observons que, puisque $R^G \to A$ est plat, on a
$A = (A \otimes_{R^G} R)^G$, voir Compléments d'algèbre,
lemme \ref{more-algebra-lemma-base-change-invariants}.
D'après le lemme \ref{lemma-inertia-base-change}, le groupe $I$ ne change pas.
Une seconde application de Compléments d'algèbre,
lemme \ref{more-algebra-lemma-base-change-invariants},
montre alors que $A \otimes_{R^G} R^I = (A \otimes_{R^G} R)^I$
(parce que $R^I \to A \otimes_{R^G} R^I$ est plat).
Ainsi
$$
\xymatrix{
\Spec((A \otimes_{R^G} R)^I) \ar[d] \ar[r] & \Spec(R^I) \ar[d] \\
\Spec(A) \ar[r] & \Spec(R^G)
}
$$
est cartésien et les flèches horizontales sont étales. Ainsi, si la
flèche verticale de gauche est étale dans un voisinage ouvert $W$ de
$(A \otimes_{R^G} R)^I \cap \mathfrak q'$, alors la flèche verticale de droite
est étale aux points de l'image (ouverte) de $W$ dans
$\Spec(R^I)$, voir
Descente, lemme \ref{descent-lemma-smooth-permanence}. En particulier,
le morphisme $\Spec(R^I) \to \Spec(R^G)$ est étale en $R^I \cap \mathfrak q$.

\medskip\noindent
Soit $\mathfrak p = R^G \cap \mathfrak q$.
D'après Compléments d'algèbre, lemme \ref{more-algebra-lemma-one-orbit},
la fibre de $\Spec(R) \to \Spec(R^G)$ au-dessus de $\mathfrak p$ est
finie. De plus, les extensions de corps résiduels en ces points
sont algébriques, normales, à groupes d'automorphismes finis d'après
Compléments d'algèbre, lemme \ref{more-algebra-lemma-one-orbit-geometric}.
On peut donc appliquer
Compléments sur les morphismes,
lemme \ref{more-morphisms-lemma-etale-makes-integral-split},
au morphisme d'anneaux entier $R^G \to R$ et à l'idéal premier $\mathfrak p$.
Combiné à l'affirmation ci-dessus, cela nous ramène au cas où
$R = A_1 \times \ldots \times A_n$, chaque $A_i$ possédant un unique
idéal premier $\mathfrak q_i$ au-dessus de $\mathfrak p$, tel que les
extensions de corps résiduels $\kappa(\mathfrak q_i)/\kappa(\mathfrak p)$
soient purement inséparables. Bien entendu, $\mathfrak q$ est l'un de
ces idéaux premiers, disons $\mathfrak q = \mathfrak q_1$.

\medskip\noindent
Il se peut que $G$ ne permute pas les facteurs $A_i$
(ce serait vrai si le spectre de $A_i$ était connexe,
par exemple si $R^G$ était local). On peut y remédier comme suit ;
nous invitons le lecteur à le vérifier lui-même, éventuellement
en utilisant des idempotents plutôt que la topologie.
Rappelons que la décomposition en produit donne une décomposition correspondante
de $\Spec(R)$ en réunion disjointe de sous-ensembles ouverts et fermés
$U_i$. Puisque $G$ est fini, on peut raffiner ce recouvrement
en une décomposition en réunion disjointe finie
$\Spec(R) = \coprod_{j \in J} W_j$ par des sous-ensembles ouverts
et fermés $W_j$, telle que, pour tout $j \in J$, il existe
un $j' \in J$ avec $\sigma(W_j) = W_{j'}$. La réunion des
$W_j$ qui ne rencontrent pas $\{\mathfrak q_1, \ldots, \mathfrak q_n\}$
est un sous-ensemble fermé qui ne rencontre pas la fibre au-dessus de $\mathfrak p$,
donc s'envoie sur un sous-ensemble fermé de $\Spec(R^G)$ qui ne contient pas
$\mathfrak p$, car $\Spec(R) \to \Spec(R^G)$ est fermé.
Ainsi, après avoir remplacé $R^G$ par une localisation principale
(ce qui est permis par l'affirmation), on peut supposer que chaque $W_j$ rencontre
l'un des points $\mathfrak q_i$. On pose alors $U_i = W_j$
si $\mathfrak q_i \in W_j$. La décomposition en produit correspondante
$R = A_1 \times \ldots \times A_n$ est telle
que $G$ permute les facteurs $A_i$.

\medskip\noindent
On peut donc supposer que l'on dispose d'une décomposition en produit
$R = A_1 \times \ldots \times A_n$ compatible avec l'action de $G$,
où chaque $A_i$ possède un unique idéal premier $\mathfrak q_i$ au-dessus
de $\mathfrak p$ et les extensions de corps
$\kappa(\mathfrak q_i)/\kappa(\mathfrak p)$ sont purement inséparables.
Écrivons $A' = A_2 \times \ldots \times A_n$, de sorte que
$$
R = A_1 \times A'
$$
Puisque $\mathfrak q = \mathfrak q_1$, on trouve que tout
$\sigma \in I$ préserve la décomposition en produit ci-dessus.
Par conséquent,
$$
R^I = (A_1)^I \times (A')^I
$$
Observons que $I = D = \{\sigma \in G \mid \sigma(\mathfrak q) = \mathfrak q\}$
parce que $\kappa(\mathfrak q)/\kappa(\mathfrak p)$ est purement inséparable.
Puisque l'action de $G$ sur les idéaux premiers au-dessus de $\mathfrak p$ est transitive
(Compléments d'algèbre, lemme \ref{more-algebra-lemma-one-orbit}),
on conclut que l'indice de $I$ dans $G$ est $n$ et que l'on peut écrire
$G = eI \amalg \sigma_2I \amalg \ldots \amalg \sigma_nI$, de sorte que
$A_i = \sigma_i(A_1)$ pour $i = 2, \ldots, n$. Il s'ensuit que
$$
R^G = (A_1)^I.
$$
Ainsi, l'application $R^G \to R^I$ est étale en $R^I \cap \mathfrak q$
et la démonstration est terminée.
\end{proof}

\noindent
Le lemme suivant généralise
Compléments d'algèbre, lemme \ref{more-algebra-lemma-inertial-invariants-unramified}.

\begin{lemma}
\label{lemma-inertial-invariants-unramified}
Soit $A$ un domaine normal de corps des fractions $K$.
Soit $L/K$ une extension galoisienne (éventuellement infinie).
Soit $G = \text{Gal}(L/K)$ et soit
$B$ la clôture intégrale de $A$ dans $L$.
Soit $\mathfrak q \subset B$. Posons
$$
I = \{\sigma \in G \mid
\sigma(\mathfrak q) = \mathfrak q \text{ et }
\sigma \bmod \mathfrak q = \text{id}_{\kappa(\mathfrak q)}\}
$$
Alors $(B^I)_{B^I \cap \mathfrak q}$ est une colimite filtrante
d'$A$-algèbres étales.
\end{lemma}

\begin{proof}
On peut écrire $L$ comme colimite filtrante d'extensions galoisiennes finies
de $K$. Il suffit donc de démontrer ce lemme lorsque $L/K$ est
une extension galoisienne finie, voir
Algèbre, lemme \ref{algebra-lemma-colimit-colimit-etale}.
Puisque $A = B^G$, car $A$ est intégralement
clos dans $K = L^G$, le résultat découle du
lemme \ref{lemma-inertia-invariants-etale}.
\end{proof}







\section{Théorie de la ramification}
\label{section-ramification}

\noindent
Dans cette section, on poursuit la discussion de
Compléments d'algèbre, section \ref{more-algebra-section-ramification},
et on la relie à notre discussion des groupes fondamentaux des schémas.

\medskip\noindent
Soit $(A, \mathfrak m, \kappa)$ un anneau local normal de
corps des fractions $K$. Choisissons une clôture séparable $K^{sep}$. Soit
$A^{sep}$ la clôture intégrale de $A$ dans $K^{sep}$.
Choisissons un idéal maximal $\mathfrak m^{sep} \subset A^{sep}$.
Soient $A \subset A^h \subset A^{sh}$ l'hensélisé et
l'hensélisé strict. Observons que $A^h$ et $A^{sh}$ sont eux aussi des anneaux normaux
(Compléments d'algèbre, lemme \ref{more-algebra-lemma-henselization-normal}).
Notons $K^h$ et $K^{sh}$ leurs corps des fractions.
Puisque $(A^{sep})_{\mathfrak m^{sep}}$ est strictement hensélien d'après le
lemme \ref{lemma-normal-local-domain-separablly-closed-fraction-field},
on peut choisir un morphisme de $A$-algèbres $A^{sh} \to (A^{sep})_{\mathfrak m^{sep}}$.
En effet, choisissons d'abord un $\kappa$-plongement\footnote{Cela est possible
parce que $\kappa(\mathfrak m^{sh})$ est une clôture séparable
de $\kappa$ et que $\kappa(\mathfrak m^{sep})$ est une clôture algébrique
de $\kappa$ d'après le lemme \ref{lemma-get-algebraic-closure}.}
$\kappa(\mathfrak m^{sh}) \to \kappa(\mathfrak m^{sep})$ et
prolongeons-le ensuite (de manière unique) en un homomorphisme de $A$-algèbres d'après
Algèbre, lemme \ref{algebra-lemma-strictly-henselian-functorial}.
On obtient le diagramme suivant
$$
\xymatrix{
K^{sep} & K^{sh} \ar[l] & K^h \ar[l] & K \ar[l] \\
(A^{sep})_{\mathfrak m^{sep}} \ar[u] &
A^{sh} \ar[u] \ar[l] &
A^h \ar[u] \ar[l] &
A \ar[u] \ar[l]
}
$$
On peut prendre les groupes fondamentaux des spectres de ces anneaux.
Bien entendu, puisque $K^{sep}$, $(A^{sep})_{\mathfrak m^{sep}}$ et
$A^{sh}$ sont strictement henséliens, on obtient pour eux des groupes triviaux.
La partie intéressante est donc la suivante
\begin{equation}
\label{equation-inertia-diagram-pione}
\vcenter{
\xymatrix{
\pi_1(U^{sh}) \ar[r] \ar[rd]_1 & \pi_1(U^h) \ar[d] \ar[r] & \pi_1(U) \ar[d] \\
& \pi_1(X^h) \ar[r] & \pi_1(X)
}
}
\end{equation}
Ici, $X^h$ et $X$ sont les spectres de $A^h$ et $A$, et
$U^{sh}$, $U^h$, $U$ sont les spectres de $K^{sh}$, $K^h$ et $K$.
L'indice $1$ signifie que l'application est triviale ; cela résulte
de ce qu'elle se factorise par le groupe trivial $\pi_1(X^{sh})$.
D'autre part, le groupe profini $G = \text{Gal}(K^{sep}/K)$
agit sur $A^{sep}$ et on peut poser les définitions suivantes
$$
D = \{\sigma \in G \mid \sigma(\mathfrak m^{sep}) = \mathfrak m^{sep}\}
\supset
I = \{\sigma \in D \mid \sigma \bmod \mathfrak m^{sep} =
\text{id}_{\kappa(\mathfrak m^{sep})}\}
$$
Ces groupes sont parfois appelés le
{\it groupe de décomposition} et le {\it groupe d'inertie},
notamment lorsque $A$ est un anneau de valuation discrète.

\begin{lemma}
\label{lemma-identify-inertia}
Dans la situation décrite ci-dessus, via l'isomorphisme
$\pi_1(U) = \text{Gal}(K^{sep}/K)$, le diagramme
(\ref{equation-inertia-diagram-pione})
se traduit par le diagramme
$$
\xymatrix{
I \ar[r] \ar[rd]_1 & D \ar[d] \ar[r] & \text{Gal}(K^{sep}/K) \ar[d] \\
& \text{Gal}(\kappa(\mathfrak m^{sh})/\kappa) \ar[r] & \text{Gal}(M/K)
}
$$
où $K^{sep}/M/K$ est la sous-extension maximale non ramifiée
relativement à $A$. De plus, les flèches verticales sont surjectives,
le noyau de la flèche verticale de gauche est $I$ et le noyau de la
flèche verticale de droite est
le plus petit sous-groupe distingué fermé de $\text{Gal}(K^{sep}/K)$
qui contient $I$.
\end{lemma}

\begin{proof}
Par construction, le groupe $D$ agit sur $(A^{sep})_{\mathfrak m^{sep}}$
au-dessus de $A$. Par l'unicité de $A^{sh} \to (A^{sep})_{\mathfrak m^{sep}}$
une fois donnée l'application sur les corps résiduels
(Algèbre, lemme \ref{algebra-lemma-strictly-henselian-functorial}),
on voit que l'image de $A^{sh} \to (A^{sep})_{\mathfrak m^{sep}}$
est contenue dans $((A^{sep})_{\mathfrak m^{sep}})^I$.
D'autre part, le
lemme \ref{lemma-inertial-invariants-unramified}
montre que $((A^{sep})_{\mathfrak m^{sep}})^I$
est une colimite filtrante d'extensions étales de $A$.
Puisque $A^{sh}$ est l'extension maximale de ce type, on conclut
que $A^{sh} = ((A^{sep})_{\mathfrak m^{sep}})^I$.
Ainsi $K^{sh} = (K^{sep})^I$.

\medskip\noindent
Rappelons que $I$ est le noyau d'une application surjective
$D \to \text{Aut}(\kappa(\mathfrak m^{sep})/\kappa)$, voir
Compléments d'algèbre, lemme \ref{more-algebra-lemma-one-orbit-geometric-galois}.
On a $\text{Aut}(\kappa(\mathfrak m^{sep})/\kappa) =
\text{Gal}(\kappa(\mathfrak m^{sh})/\kappa)$,
car on a vu ci-dessus que ces corps sont respectivement la clôture algébrique
et la clôture séparable de $\kappa$.
D'autre part, tout automorphisme de $A^{sh}$ au-dessus de $A$
est un automorphisme de $A^{sh}$ au-dessus de $A^h$ par l'unicité
dans Algèbre, lemme \ref{algebra-lemma-henselian-functorial}.
En outre, $A^{sh}$ est la colimite des extensions finies étales
$A^h \subset A'$ qui correspondent $1$-à-$1$
aux extensions finies séparables $\kappa'/\kappa$, voir
Algèbre, remarque \ref{algebra-remark-construct-sh-from-h}.
Ainsi
$$
\text{Aut}(A^{sh}/A) = \text{Aut}(A^{sh}/A^h) =
\text{Gal}(\kappa(\mathfrak m^{sh})/\kappa)
$$
Soit $\kappa'/\kappa$ une extension galoisienne finie de
groupe de Galois $G$. Soit $A^h \subset A'$ l'extension finie étale
qui correspond à $\kappa \subset \kappa'$ d'après
Algèbre, lemme \ref{algebra-lemma-henselian-cat-finite-etale}.
Il s'ensuit alors que
$(A')^G = A^h$ en considérant les corps des fractions et les degrés
(on omet ce petit détail). En prenant la colimite, on conclut que
$(A^{sh})^{\text{Gal}(\kappa(\mathfrak m^{sh})/\kappa)} = A^h$.
En combinant tout ce qui précède, on trouve $A^h = ((A^{sep})_{\mathfrak m^{sep}})^D$.
Ainsi $K^h = (K^{sep})^D$.

\medskip\noindent
Puisque $U$, $U^h$, $U^{sh}$ sont les spectres des corps
$K$, $K^h$, $K^{sh}$, on voit que les lignes supérieures des diagrammes
se correspondent via le
lemme \ref{lemma-fundamental-group-Galois-group}.
D'après le lemme \ref{lemma-gabber}, on a
$\pi_1(X^h) = \text{Gal}(\kappa(\mathfrak m^{sh})/\kappa)$.
L'exactitude de la suite
$1 \to I \to D \to \text{Gal}(\kappa(\mathfrak m^{sh})/\kappa) \to 1$
a été signalée ci-dessus.
D'après la proposition \ref{proposition-normal},
on voit que $\pi_1(X) = \text{Gal}(M/K)$.
Enfin, l'assertion sur le noyau de
$\text{Gal}(K^{sep}/K) \to \text{Gal}(M/K) = \pi_1(X)$
découle du lemme \ref{lemma-local-exact-sequence-normal}.
Cela achève la démonstration.
\end{proof}

\noindent
Soit $X$ un schéma intègre normal de corps des fonctions rationnelles $K$.
Soit $K^{sep}$ une clôture séparable de $K$.
Soit $X^{sep} \to X$ la normalisation de $X$ dans $K^{sep}$.
Puisque $G = \text{Gal}(K^{sep}/K)$ agit sur $K^{sep}$,
on obtient une action à droite de $G$ sur $X^{sep}$.
Pour $y \in X^{sep}$, posons
$$
D_y = \{\sigma \in G \mid \sigma(y) = y\} \supset
I_y = \{\sigma \in D_y \mid \sigma \bmod \mathfrak m_y =
\text{id}_{\kappa(y)} \}
$$
comme ci-dessus. D'autre part, pour $x \in X$,
soit $\mathcal{O}_{X, x}^{sh}$ un hensélisé strict,
soit $K_x^{sh}$ le corps des fractions de $\mathcal{O}_{X, x}^{sh}$
et choisissons un $K$-plongement $K_x^{sh} \to K^{sep}$.

\begin{lemma}
\label{lemma-normal-pione-quotient-inertia}
Soit $X$ un schéma intègre normal de corps des fonctions rationnelles $K$.
Avec les notations ci-dessus, les trois sous-groupes suivants de
$\text{Gal}(K^{sep}/K) = \pi_1(\Spec(K))$
sont égaux :
\begin{enumerate}
\item le noyau de la surjection
$\text{Gal}(K^{sep}/K) \longrightarrow \pi_1(X)$,
\item le plus petit sous-groupe distingué fermé qui contient $I_y$
pour tout $y \in X^{sep}$, et
\item le plus petit sous-groupe distingué fermé qui contient
$\text{Gal}(K^{sep}/K_x^{sh})$ pour tout $x \in  X$.
\end{enumerate}
\end{lemma}

\begin{proof}
L'équivalence de (2) et (3) découle du
lemme \ref{lemma-identify-inertia},
qui nous dit que $I_y$ est conjugué à $\text{Gal}(K^{sep}/K_x^{sh})$
si $y$ est au-dessus de $x$. D'après le lemme \ref{lemma-local-exact-sequence-normal},
on voit que $\text{Gal}(K^{sep}/K_x^{sh})$ s'envoie trivialement dans
$\pi_1(\Spec(\mathcal{O}_{X, x}))$ et que, par conséquent, le sous-groupe
$N \subset G = \text{Gal}(K^{sep}/K)$
de (2) et (3) est contenu dans le noyau de
$G \longrightarrow \pi_1(X)$.

\medskip\noindent
Pour démontrer l'autre inclusion, puisque $N$ est distingué, il suffit de montrer ceci :
étant donné $N \subset U \subset G$ avec $U$ ouvert normal,
la projection canonique $G \to G/U$ se factorise par $\pi_1(X)$.
Autrement dit, si $L/K$ est l'extension galoisienne correspondant
à $U$, alors il faut montrer que $X$ est non ramifié dans $L$
(section \ref{section-normal}, en particulier
proposition \ref{proposition-normal}).
Il suffit de le faire lorsque $X$ est affine (on procède ainsi
afin de pouvoir invoquer des résultats algébriques dans la suite de la démonstration).
Soit $Y \to X$ la normalisation de $X$ dans $L$.
L'inclusion $L \subset K^{sep}$ induit un morphisme
$\pi : X^{sep} \to Y$. Pour $y \in X^{sep}$,
le groupe d'inertie de $\pi(y)$ dans $\text{Gal}(L/K)$
est l'image de $I_y$ dans $\text{Gal}(L/K)$ ; cela découle
de Compléments d'algèbre, lemme
\ref{more-algebra-lemma-one-orbit-geometric-galois-compare}.
Puisque $N \subset U$, tous ces groupes d'inertie sont triviaux.
On conclut que $Y \to X$ est étale en appliquant le
lemme \ref{lemma-inertia-invariants-etale}.
(Autre possibilité : on peut utiliser le lemme \ref{lemma-local-exact-sequence-normal}
pour voir que le changement de base de $Y$ à $\Spec(\mathcal{O}_{X, x})$ est
étale pour tout $x \in X$, puis conclure à partir de là
avec un peu plus de travail.)
\end{proof}

\begin{example}
\label{example-bigger-codim}
Soit $X$ un schéma noethérien intègre normal de corps des fonctions rationnelles $K$.
La pureté du lieu de branchement (voir plus bas) nous dit que si $X$ est régulier, alors
il suffit, dans le lemme \ref{lemma-normal-pione-quotient-inertia},
de considérer les groupes d'inertie $I = \pi_1(\Spec(K_x^{sh}))$
pour les points $x$ de codimension $1$ dans $X$.
En général, cela ne suffit cependant pas. En effet, soit
$Y = \mathbf{A}_k^n = \Spec(k[t_1, \ldots, t_n])$,
où $k$ est un corps de caractéristique différente de $2$.
Soit $G = \{\pm 1\}$ le groupe d'ordre $2$ agissant sur $Y$
par multiplication sur les coordonnées. Posons
$$
X = \Spec(k[t_it_j, i, j \in \{1, \ldots, n\}])
$$
Le plongement $k[t_it_j] \subset k[t_1, \ldots, t_n]$
définit un morphisme de degré $2$, à savoir $Y \to X$, qui est non ramifié partout
sauf au-dessus de l'idéal maximal $\mathfrak m = (t_it_j)$,
qui est un point de codimension $n$ dans $X$.
\end{example}

\begin{lemma}
\label{lemma-unramified}
Soit $X$ un schéma intègre normal de corps des fonctions rationnelles $K$.
Soit $L/K$ une extension finie. Soit $Y \to X$ la normalisation
de $X$ dans $L$. Les conditions suivantes sont équivalentes :
\begin{enumerate}
\item $X$ est non ramifié dans $L$ au sens de la section \ref{section-normal},
\item $Y \to X$ est un morphisme non ramifié de schémas,
\item $Y \to X$ est un morphisme étale de schémas,
\item $Y \to X$ est un morphisme fini étale de schémas,
\item pour $x \in X$, la projection
$Y \times_X \Spec(\mathcal{O}_{X, x}) \to \Spec(\mathcal{O}_{X, x})$
est non ramifiée,
\item la même condition que (5), mais avec $\mathcal{O}_{X, x}^h$,
\item la même condition que (5), mais avec $\mathcal{O}_{X, x}^{sh}$,
\item pour $x \in X$, la fibre schématique $Y_x$
est étale sur $x$ de degré $\geq [L : K]$.
\end{enumerate}
Si $L/K$ est galoisienne de groupe de Galois $G$, ces conditions sont aussi
équivalentes à
\begin{enumerate}
\item[(9)] pour $y \in Y$, le groupe
$I_y = \{g \in G \mid g(y) = y\text{ et }
g \bmod \mathfrak m_y = \text{id}_{\kappa(y)}\}$ est trivial.
\end{enumerate}
\end{lemma}

\begin{proof}
L'équivalence de (1) et (2) est la définition de (1).
L'équivalence de (2), (3) et (4) est le lemme \ref{lemma-unramified-in-L}.
Il est immédiat de démontrer que (4) $\Rightarrow$ (5),
(5) $\Rightarrow$ (6), (6) $\Rightarrow$ (7).

\medskip\noindent
Supposons (7). Observons que $\mathcal{O}_{X, x}^{sh}$ est un domaine local normal
(Compléments d'algèbre, lemme \ref{more-algebra-lemma-henselization-normal}).
Soit $L^{sh} = L \otimes_K K_x^{sh}$, où $K_x^{sh}$ est le corps des fractions
de $\mathcal{O}_{X, x}^{sh}$. Alors $L^{sh} = \prod_{i = 1, \ldots, n} L_i$,
avec $L_i/K_x^{sh}$ finie séparable. D'après
Algèbre, lemme \ref{algebra-lemma-integral-closure-commutes-smooth}
(et un passage à la limite que nous omettons),
on voit que $Y \times_X \Spec(\mathcal{O}_{X, x}^{sh})$
est la clôture intégrale de $\Spec(\mathcal{O}_{X, x}^{sh})$ dans $L^{sh}$.
Par conséquent, d'après le lemme \ref{lemma-unramified-in-L} (appliqué aux facteurs
$L_i$ de $L^{sh}$), on voit que
$Y \times_X \Spec(\mathcal{O}_{X, x}^{sh}) \to \Spec(\mathcal{O}_{X, x}^{sh})$
est fini étale. En considérant le point générique, on voit que
le degré est égal à $[L : K]$ et on voit donc que (8) est vraie.

\medskip\noindent
Supposons (8). Supposons que $x \in X$ et que la fibre schématique $Y_x$
est étale sur $x$ de degré $\geq [L : K]$. Observons que cela signifie
que $Y$ possède $\geq [L : K]$ points géométriques au-dessus de $x$.
Nous allons montrer que $Y \to X$ est fini étale au-dessus d'un voisinage de $x$.
Cela démontrera que (1) est vraie.
Pour le démontrer, on peut supposer $X = \Spec(R)$, que le point $x$ correspond à
l'idéal premier $\mathfrak p \subset R$ et que $Y = \Spec(S)$. On applique
Compléments sur les morphismes,
lemme \ref{more-morphisms-lemma-etale-makes-integral-split}, et on trouve un
voisinage étale $(U, u) \to (X, x)$ tel que
$Y \times_X U = V_1 \amalg \ldots \amalg V_m$, où $V_i$
possède un unique point $v_i$ au-dessus de $u$ avec $\kappa(v_i)/\kappa(u)$
purement inséparable. En rétrécissant $U$ si nécessaire, on peut supposer que $U$ est
un schéma intègre normal de point générique $\xi$ (utiliser
Descente, lemmes \ref{descent-lemma-locally-finite-nr-irred-local-fppf} et
\ref{descent-lemma-normal-local-smooth}, et
Propriétés, lemme \ref{properties-lemma-normal-locally-finite-nr-irreducibles}).
D'après notre remarque sur les points géométriques, on voit que $m \geq [L : K]$.
D'autre part, d'après Compléments sur les morphismes, lemme
\ref{more-morphisms-lemma-normalization-smooth-localization},
on voit que $\coprod V_i \to U$ est la normalisation de $U$ dans
$\Spec(L) \times_X U$. Comme $K \subset \kappa(\xi)$ est finie séparable,
on peut écrire $\Spec(L) \times_X U = \Spec(\prod_{i = 1, \ldots, n} L_i)$,
avec $L_i/\kappa(\xi)$ finie et $[L : K] = \sum [L_i : \kappa(\xi)]$.
Puisque $V_j$ est non vide pour tout $j$ et que $m \geq [L : K]$,
on conclut que $m = n$ et que $[L_i : \kappa(\xi)] = 1$
pour tout $i$. Alors $V_j \to U$ est en particulier un isomorphisme
étale, et donc $Y \times_X U \to U$ est étale. D'après
Descente, lemme \ref{descent-lemma-descending-property-etale},
on conclut que $Y \to X$ est étale au-dessus de
l'image de $U \to X$ (un voisinage ouvert de $x$).

\medskip\noindent
Supposons $L/K$ galoisienne et (9) vérifiée. Alors $Y \to X$ est étale
d'après le lemme \ref{lemma-inertial-invariants-unramified}.
On omet la démonstration de ce que (1) implique (9).
\end{proof}

\noindent
Dans le cas d'extensions galoisiennes infinies d'anneaux de valuation discrète,
on peut en dire un tout petit peu plus. Pour cela, introduisons la notation suivante.
Un sous-ensemble $S \subset \mathbf{N}$ d'entiers est {\it filtrant pour la divisibilité}
si $1 \in S$ et si, pour $n, m \in S$, il existe $k \in S$ tel que
$n | k$ et $m | k$. Définissons sur $S$ un ordre partiel par la règle
$n \geq_S m$ si et seulement si $m | n$. Étant donné un corps $\kappa$, on obtient
un système projectif de groupes finis $\{\mu_n(\kappa)\}_{n \in S}$
dont les applications de transition sont
$$
\mu_n(\kappa) \longrightarrow \mu_m(\kappa),\quad
\zeta \longmapsto \zeta^{n/m}
$$
pour $n \geq_S m$. On peut alors former le groupe profini
$$
\lim_{n \in S} \mu_n(\kappa)
$$
Observons que la catégorie d'indices de cette limite est cofiltrante (car $S$ est filtrant).
La construction est fonctorielle en $\kappa$. En particulier,
$\text{Aut}(\kappa)$ agit sur ce groupe profini.
Par exemple, si $S = \{1, n\}$, on obtient alors $\mu_n(\kappa)$.
Si $S = \{1, \ell, \ell^2, \ell^3, \ldots\}$ pour un nombre premier
$\ell$ différent de la caractéristique de $\kappa$, cela produit
$\lim_n \mu_{\ell^n}(\kappa)$,
que l'on appelle parfois le module de Tate $\ell$-adique du groupe
multiplicatif de $\kappa$ (comparer avec
Compléments d'algèbre, exemple
\ref{more-algebra-example-spectral-sequence-principal}).

\begin{lemma}
\label{lemma-structure-decomposition}
Soit $A$ un anneau de valuation discrète de corps des fractions $K$.
Soit $L/K$ une extension galoisienne (éventuellement infinie).
Soit $B$ la clôture intégrale de $A$ dans $L$.
Soit $\mathfrak m$ un idéal maximal de $B$ et posons $\kappa_A = A/(\mathfrak m \cap A)$.
Soient $G = \text{Gal}(L/K)$,
$D = \{\sigma \in G \mid \sigma(\mathfrak m) = \mathfrak m\}$, et
$I = \{\sigma \in D \mid \sigma \bmod \mathfrak m =
\text{id}_{\kappa(\mathfrak m)}\}$.
Le groupe de décomposition $D$ s'insère dans une suite exacte canonique
$$
1 \to I \to D \to \text{Aut}(\kappa(\mathfrak m)/\kappa_A) \to 1
$$
Le groupe d'inertie $I$ s'insère dans une suite exacte canonique
$$
1 \to P \to I \to I_t \to 1
$$
telle que
\begin{enumerate}
\item $P$ est un sous-groupe distingué de $D$,
\item $P$ est un pro-$p$-groupe si la caractéristique de
$\kappa_A$ est $p > 1$, et $P = \{1\}$ si la caractéristique de $\kappa_A$
est nulle (avec la convention $p = 1$ dans ce dernier cas),
\item il existe un $S \subset \mathbf{N}$ filtrant pour la divisibilité
tel que $\kappa(\mathfrak m)$ contient une racine $n$-ième primitive de l'unité
pour tout $n \in S$ (les éléments de $S$ sont premiers à $p$),
\item il existe une application surjective canonique
$$
\theta_{can} : I \to \lim_{n \in S} \mu_n(\kappa(\mathfrak m))
$$
dont le noyau est $P$, qui vérifie
$\theta_{can}(\tau \sigma \tau^{-1}) = \tau(\theta_{can}(\sigma))$
pour $\tau \in D$, $\sigma \in I$, et qui induit un isomorphisme
$I_t \to \lim_{n \in S} \mu_n(\kappa(\mathfrak m))$.
\end{enumerate}
\end{lemma}

\begin{proof}
Il s'agit essentiellement d'une reformulation des résultats sur les extensions galoisiennes finies
démontrés dans Compléments d'algèbre, section \ref{more-algebra-section-ramification}.
La surjectivité de l'application $D \to \text{Aut}(\kappa(\mathfrak m)/\kappa_A)$ est le
lemme \ref{more-algebra-lemma-one-orbit-geometric-galois} de Compléments d'algèbre.
Cela donne la première suite exacte.

\medskip\noindent
Pour construire la seconde suite exacte courte, soit $\Lambda$ l'ensemble
des sous-extensions galoisiennes finies, c'est-à-dire que $\lambda \in \Lambda$ correspond
à $L/L_\lambda/K$. Posons $G_\lambda = \text{Gal}(L_\lambda/K)$.
Rappelons que les $G_\lambda$ forment un système projectif de groupes finis à applications
de transition surjectives et que $G = \lim_{\lambda \in \Lambda} G_\lambda$, voir
Corps, lemme \ref{fields-lemma-infinite-galois-limit}.
On note $B_\lambda$ la clôture intégrale de $A$ dans $L_\lambda$.
On pose ensuite $\mathfrak m_\lambda = \mathfrak m \cap B_\lambda$
et on note $P_\lambda, I_\lambda, D_\lambda$ les groupes
d'inertie sauvage, d'inertie et de décomposition de
$\mathfrak m_\lambda$, voir Compléments d'algèbre, lemme
\ref{more-algebra-lemma-galois-inertia}.
Pour $\lambda \geq \lambda'$, la restriction définit
un diagramme commutatif
$$
\xymatrix{
P_\lambda \ar[d] \ar[r] &
I_\lambda \ar[d] \ar[r] &
D_\lambda \ar[d] \ar[r] &
G_\lambda \ar[d] \\
P_{\lambda'} \ar[r] &
I_{\lambda'} \ar[r] &
D_{\lambda'} \ar[r] &
G_{\lambda'}
}
$$
dont les flèches verticales sont surjectives, voir
Compléments d'algèbre, lemme \ref{more-algebra-lemma-compare-inertia}.

\medskip\noindent
Il résulte immédiatement des définitions
que $I = \lim I_\lambda$ et $D = \lim D_\lambda$
sous l'isomorphisme $G = \lim G_\lambda$ ci-dessus.
Puisque $L = \colim L_\lambda$, on a $B = \colim B_\lambda$
et $\kappa(\mathfrak m) = \colim \kappa(\mathfrak m_\lambda)$.
Puisque les applications de transition du système $D_\lambda$
sont compatibles avec les applications
$D_\lambda \to \text{Aut}(\kappa(\mathfrak m_\lambda)/\kappa)$
(voir Compléments d'algèbre, lemme \ref{more-algebra-lemma-compare-inertia}),
on voit que l'application $D \to \text{Aut}(\kappa(\mathfrak m)/\kappa)$
est la limite des applications
$D_\lambda \to \text{Aut}(\kappa(\mathfrak m_\lambda)/\kappa)$.

\medskip\noindent
Il existe des applications canoniques
$$
\theta_{\lambda, can} :
I_\lambda
\longrightarrow
\mu_{n_\lambda}(\kappa(\mathfrak m_\lambda))
$$
où $n_\lambda = |I_\lambda|/|P_\lambda|$, où
$\mu_{n_\lambda}(\kappa(\mathfrak m_\lambda))$ est
d'ordre $n_\lambda$, telles que
$\theta_{\lambda, can}(\tau \sigma \tau^{-1}) =
\tau(\theta_{\lambda, can}(\sigma))$ pour
$\tau \in D_\lambda$ et $\sigma \in I_\lambda$, et telles que
l'on obtient des diagrammes commutatifs
$$
\xymatrix{
I_\lambda \ar[r]_-{\theta_{\lambda, can}} \ar[d] &
\mu_{n_\lambda}(\kappa(\mathfrak m_\lambda))
\ar[d]^{(-)^{n_\lambda/n_{\lambda'}}} \\
I_{\lambda'} \ar[r]^-{\theta_{\lambda', can}} &
\mu_{n_{\lambda'}}(\kappa(\mathfrak m_{\lambda'}))
}
$$
voir
Compléments d'algèbre, remarque \ref{more-algebra-remark-canonical-inertia-character}.

\medskip\noindent
Soit $S \subset \mathbf{N}$ l'ensemble des entiers $n_\lambda$.
Puisque $\Lambda$ est filtrant, on voit que $S$ est filtrant pour la divisibilité.
À l'aide des diagrammes commutatifs affichés ci-dessus, on peut prendre les limites des
applications $\theta_{\lambda, can}$ afin d'obtenir
$$
\theta_{can} : I \to \lim_{n \in S} \mu_n(\kappa(\mathfrak m)).
$$
Cette application est continue (on omet ce petit détail). Puisque les applications de transition
du système des $I_\lambda$ sont surjectives
et que $\Lambda$ est filtrant, les projections $I \to I_\lambda$
sont surjectives. Pour tout $\lambda$, le diagramme
$$
\xymatrix{
I \ar[d] \ar[r]_-{\theta_{can}} &
\lim_{n \in S} \mu_n(\kappa(\mathfrak m)) \ar[d] \\
I_{\lambda} \ar[r]^-{\theta_{\lambda, can}} &
\mu_{n_\lambda}(\kappa(\mathfrak m_\lambda))
}
$$
est commutatif. L'image de $\theta_{can}$ se projette donc sur le groupe fini
$\mu_{n_\lambda}(\kappa(\mathfrak m)) =
\mu_{n_\lambda}(\kappa(\mathfrak m_\lambda))$ d'ordre $n_\lambda$
(voir ci-dessus). Il s'ensuit que l'image de $\theta_{can}$ est dense.
D'autre part, $\theta_{can}$ est continue et la
source est un groupe profini. Ainsi $\theta_{can}$ est surjective
d'après un argument topologique.

\medskip\noindent
La propriété $\theta_{can}(\tau \sigma \tau^{-1}) = \tau(\theta_{can}(\sigma))$
pour $\tau \in D$, $\sigma \in I$ découle des propriétés correspondantes
des applications $\theta_{\lambda, can}$ et de la compatibilité de l'application
$D \to \text{Aut}(\kappa(\mathfrak m))$ avec les applications
$D_\lambda \to \text{Aut}(\kappa(\mathfrak m_\lambda))$.
En posant $P = \Ker(\theta_{can})$, cela implique
que $P$ est un sous-groupe distingué de $D$. En posant $I_t = I/P$,
on obtient l'isomorphisme $I_t \to \lim_{n \in S} \mu_n(\kappa(\mathfrak m))$
par la surjectivité de $\theta_{can}$.

\medskip\noindent
Pour achever la démonstration, montrons que $P = \lim P_\lambda$, ce qui prouve
que $P$ est un pro-$p$-groupe. Rappelons que le groupe d'inertie modérée
$I_{\lambda, t} = I_\lambda/P_\lambda$ est d'ordre $n_\lambda$.
Puisque les applications de transition $P_\lambda \to P_{\lambda'}$ sont surjectives
et que $\Lambda$ est filtrant, on obtient une suite exacte courte
$$
1 \to \lim P_\lambda \to I \to \lim I_{\lambda, t} \to 1
$$
(on omet les détails). Puisque, pour tout $\lambda$, l'application $\theta_{\lambda, can}$
induit un isomorphisme
$I_{\lambda, t} \cong \mu_{n_\lambda}(\kappa(\mathfrak m))$,
le résultat souhaité en découle.
\end{proof}

\begin{lemma}
\label{lemma-structure-decomposition-separable-closure}
Soit $A$ un anneau de valuation discrète de corps des fractions $K$.
Soit $K^{sep}$ une clôture séparable de $K$.
Soit $A^{sep}$ la clôture intégrale de $A$ dans $K^{sep}$.
Soit $\mathfrak m^{sep}$ un idéal maximal de $A^{sep}$.
Soit $\mathfrak m = \mathfrak m^{sep} \cap A$, soient
$\kappa = A/\mathfrak m$, et
$\overline{\kappa} = A^{sep}/\mathfrak m^{sep}$.
Alors $\overline{\kappa}$ est une clôture algébrique de $\kappa$.
Soient $G = \text{Gal}(K^{sep}/K)$,
$D = \{\sigma \in G \mid \sigma(\mathfrak m^{sep}) = \mathfrak m^{sep}\}$, et
$I = \{\sigma \in D \mid \sigma \bmod \mathfrak m^{sep} =
\text{id}_{\kappa(\mathfrak m^{sep})}\}$.
Le groupe de décomposition $D$ s'insère dans une suite exacte canonique
$$
1 \to I \to D \to \text{Gal}(\kappa^{sep}/\kappa) \to 1
$$
où $\kappa^{sep} \subset \overline{\kappa}$ est la clôture
séparable de $\kappa$.
Le groupe d'inertie $I$ s'insère dans une suite exacte canonique
$$
1 \to P \to I \to I_t \to 1
$$
telle que
\begin{enumerate}
\item $P$ est un sous-groupe distingué de $D$,
\item $P$ est un pro-$p$-groupe si la caractéristique de
$\kappa$ est $p > 1$, et $P = \{1\}$ si la caractéristique de $\kappa$
est nulle (avec la convention $p = 1$ dans ce dernier cas),
\item il existe une application surjective canonique
$$
\theta_{can} : I \to \lim_{n\text{ premier à }p} \mu_n(\kappa^{sep})
$$
dont le noyau est $P$, qui vérifie
$\theta_{can}(\tau \sigma \tau^{-1}) = \tau(\theta_{can}(\sigma))$
pour $\tau \in D$, $\sigma \in I$, et qui induit un isomorphisme
$I_t \to \lim_{n\text{ premier à }p} \mu_n(\kappa^{sep})$.
\end{enumerate}
\end{lemma}

\begin{proof}
Le corps $\overline{\kappa}$ est la clôture algébrique de $\kappa$ d'après le
lemme \ref{lemma-get-algebraic-closure}.
La plupart des assertions découlent immédiatement des parties correspondantes
du lemme \ref{lemma-structure-decomposition}. Par exemple, puisque
$\text{Aut}(\overline{\kappa}/\kappa) = \text{Gal}(\kappa^{sep}/\kappa)$,
on obtient la première suite.
Alors la seule autre assertion qui requiert une démonstration est le fait que,
avec $S$ comme dans le lemme \ref{lemma-structure-decomposition}, la
limite $\lim_{n \in S} \mu_n(\overline{\kappa})$ est égale à
$\lim_{n\text{ premier à }p} \mu_n(\kappa^{sep})$. Pour le voir,
il suffit de montrer que tout entier $n$ premier à $p$
divise un élément de $S$.
Soit $\pi \in A$ une uniformisante et considérons le corps de décomposition
$L$ du polynôme $X^n - \pi$. Puisque ce polynôme
est séparable, on voit que $L$ est une extension galoisienne finie de $K$.
Choisissons un plongement $L \to K^{sep}$.
Observons que si $B$ est la clôture intégrale de $A$ dans $L$,
alors l'indice de ramification de $A \to B_{\mathfrak m^{sep} \cap B}$
est divisible par $n$ (car $\pi$ admet une racine $n$-ième dans $B$ ; en fait,
l'indice de ramification est égal à $n$, mais nous n'en avons pas besoin).
Il résulte alors de la construction de l'ensemble $S$ dans la démonstration du
lemme \ref{lemma-structure-decomposition}
que $n$ divise un élément de $S$.
\end{proof}








\section{Groupes fondamentaux géométriques et arithmétiques}
\label{section-galois-action}

\noindent
Dans cette section, nous étudions ce qui se passe lorsque l'on compare le
groupe fondamental d'un schéma $X$ sur un corps $k$ au
groupe fondamental de $X_{\overline{k}}$, où $\overline{k}$
est la clôture algébrique de $k$.

\begin{lemma}
\label{lemma-limit}
Soit $I$ un ensemble ordonné filtrant. Soit $X_i$ un
système projectif de schémas quasi-compacts et quasi-séparés
indexé par $I$, dont les morphismes de transition sont affines.
Soit $X = \lim X_i$ comme dans Limites, section \ref{limits-section-limits}.
Il existe alors une équivalence de catégories
$$
\colim \textit{F\'Et}_{X_i} = \textit{F\'Et}_X
$$
Si $X_i$ est connexe pour tout $i$ suffisamment grand et si $\overline{x}$
est un point géométrique de $X$, alors
$$
\pi_1(X, \overline{x}) = \lim \pi_1(X_i, \overline{x})
$$
\end{lemma}

\begin{proof}
L'équivalence de catégories résulte de Limites, lemmes
\ref{limits-lemma-descend-finite-presentation},
\ref{limits-lemma-descend-finite-finite-presentation}, et
\ref{limits-lemma-descend-etale}.
La seconde assertion découle formellement de l'assertion sur les
catégories.
\end{proof}

\begin{lemma}
\label{lemma-perfection}
Soit $k$ un corps et notons $k^{perf}$ sa clôture parfaite. Soit $X$ un schéma connexe
sur $k$. Alors $X_{k^{perf}}$ est connexe et
$\pi_1(X_{k^{perf}}) \to \pi_1(X)$ est un isomorphisme.
\end{lemma}

\begin{proof}
C'est un cas particulier de l'invariance topologique du groupe fondamental.
Voir la proposition \ref{proposition-universal-homeomorphism}.
Pour voir que $\Spec(k^{perf}) \to \Spec(k)$ est un
homéomorphisme universel, on peut utiliser
Algèbre, lemme \ref{algebra-lemma-radicial-integral-bijective}.
\end{proof}

\begin{lemma}
\label{lemma-ses-field}
Soit $k$ un corps dont $\overline{k}$ est une clôture algébrique.
Soit $X$ un schéma quasi-compact et quasi-séparé sur $k$.
Si le changement de base $X_{\overline{k}}$ est connexe, alors
il existe une suite exacte courte
$$
1 \to \pi_1(X_{\overline{k}}) \to \pi_1(X) \to \pi_1(\Spec(k)) \to 1
$$
de groupes topologiques profinis.
\end{lemma}

\begin{proof}
Les objets connexes de $\textit{F\'Et}_{\Spec(k)}$ sont de la forme
$\Spec(k') \to \Spec(k)$, où $k'/k$ est une extension finie séparable.
Alors $X_{\Spec(k')}$ est connexe, puisque le morphisme
$X_{\overline{k}} \to X_{\Spec(k')}$ est surjectif et
$X_{\overline{k}}$ est connexe par hypothèse. Ainsi
$\pi_1(X) \to \pi_1(\Spec(k))$ est surjectif d'après le
lemme \ref{lemma-functoriality-galois-surjective}.

\medskip\noindent
Avant de poursuivre, remarquons que l'on peut supposer que $k$ est parfait.
En effet, on a $\pi_1(X_{k^{perf}}) = \pi_1(X)$ et
$\pi_1(\Spec(k^{perf})) = \pi_1(\Spec(k))$ d'après le lemme \ref{lemma-perfection}.

\medskip\noindent
Il est clair que le composé des foncteurs
$\textit{F\'Et}_{\Spec(k)} \to \textit{F\'Et}_X \to
\textit{F\'Et}_{X_{\overline{k}}}$ envoie les objets sur des réunions disjointes
de copies de $X_{\Spec(\overline{k})}$. Par conséquent, le composé
$\pi_1(X_{\overline{k}}) \to \pi_1(X) \to \pi_1(\Spec(k))$
est l'homomorphisme trivial d'après le lemme \ref{lemma-composition-trivial}.

\medskip\noindent
Soit $U \to X$ un morphisme fini \'etale, avec $U$ connexe.
Observons que $U \times_X X_{\overline{k}} = U_{\overline{k}}$.
Supposons que $U_{\overline{k}} \to X_{\overline{k}}$
admette une section $s : X_{\overline{k}} \to U_{\overline{k}}$.
Alors $s(X_{\overline{k}})$ est une composante connexe ouverte de
$U_{\overline{k}}$. Pour $\sigma \in \text{Gal}(\overline{k}/k)$,
notons $s^\sigma$ le changement de base de $s$ par $\Spec(\sigma)$.
Puisque $U_{\overline{k}} \to X_{\overline{k}}$ est fini \'etale,
il n'a qu'un nombre fini de sections. Ainsi
$$
\overline{T} = \bigcup_{\sigma \in \text{Gal}(\overline{k}/k)} s^\sigma(X_{\overline{k}})
$$
est une réunion finie et l'on voit que $\overline{T}$ est un
sous-ensemble ouvert et fermé stable par $\text{Gal}(\overline{k}/k)$.
D'après Variétés, lemme \ref{varieties-lemma-closed-fixed-by-Galois},
on voit que $\overline{T}$ est l'image réciproque d'un sous-ensemble fermé
$T \subset U$. Puisque $U_{\overline{k}} \to U$ est ouvert
(Morphismes, lemme \ref{morphisms-lemma-scheme-over-field-universally-open}),
on conclut que $T$ est également ouvert. Comme $U$ est connexe, on
voit que $T = U$. Par conséquent, $U_{\overline{k}}$ est une réunion disjointe (finie)
de copies de $X_{\overline{k}}$. D'après le
lemme \ref{lemma-functoriality-galois-normal}, on conclut que l'image de
$\pi_1(X_{\overline{k}}) \to \pi_1(X)$ est un sous-groupe distingué.

\medskip\noindent
Soit $V \to X_{\overline{k}}$ un revêtement fini \'etale. Rappelons que
$\overline{k}$ est la réunion des extensions finies séparables de $k$.
D'après le lemme \ref{lemma-limit}, il existe une extension finie séparable $k'/k$
et un morphisme fini \'etale $U \to X_{k'}$ tels que
$V = X_{\overline{k}} \times_{X_{k'}} U =
U \times_{\Spec(k')} \Spec(\overline{k})$.
Alors le composé $U \to X_{k'} \to X$ est fini \'etale
et $U \times_{\Spec(k)} \Spec(\overline{k})$
contient $V = U \times_{\Spec(k')} \Spec(\overline{k})$
comme sous-schéma ouvert et fermé. (En effet, $\Spec(\overline{k})$
est un sous-schéma ouvert et fermé de
$\Spec(k') \times_{\Spec(k)} \Spec(\overline{k})$ via
l'application de multiplication $k' \otimes_k \overline{k} \to \overline{k}$.) D'après le
lemme \ref{lemma-functoriality-galois-injective},
on conclut que $\pi_1(X_{\overline{k}}) \to \pi_1(X)$ est injectif.

\medskip\noindent
Enfin, nous devons montrer que, pour tout morphisme fini \'etale
$U \to X$ tel que $U_{\overline{k}}$ soit une réunion disjointe
de copies de $X_{\overline{k}}$, il existe un morphisme fini \'etale
$V \to \Spec(k)$ et une surjection $V \times_{\Spec(k)} X \to U$.
Voir le lemme \ref{lemma-functoriality-galois-ses}.
En raisonnant comme ci-dessus à l'aide du lemme \ref{lemma-limit},
il existe une extension finie séparable $k'/k$
pour laquelle il existe un isomorphisme
$U_{k'} \cong \coprod_{i = 1, \ldots, n} X_{k'}$.
Ainsi, en posant $V = \coprod_{i = 1, \ldots, n} \Spec(k')$,
on conclut.
\end{proof}






\section{Suite exacte d'homotopie}
\label{section-homotopy-exact-sequence}

\noindent
Dans cette section, nous examinons le résultat suivant.
Soit $f : X \to S$ un morphisme propre et plat de
présentation finie dont les
fibres géométriques sont connexes et réduites.
Supposons $S$ connexe et soit $\overline{s}$
un point géométrique de $S$. Il existe alors une suite
exacte
$$
\pi_1(X_{\overline{s}}) \to \pi_1(X) \to \pi_1(S) \to 1
$$
de groupes fondamentaux. Voir la
proposition \ref{proposition-first-homotopy-sequence}.

\begin{lemma}
\label{lemma-stein-factorization-etale}
\begin{reference}
\cite[Exposé X, proposition 1.2, p. 262]{SGA1}.
\end{reference}
Soit $f : X \to S$ un morphisme propre de schémas.
Soit $X \to S' \to S$ la factorisation de Stein de $f$ ; voir
Compléments sur les morphismes, théorème
\ref{more-morphisms-theorem-stein-factorization-general}.
Si $f$ est de présentation finie et plat, à fibres géométriquement
réduites, alors $S' \to S$ est fini \'etale.
\end{lemma}

\begin{proof}
Cela résulte de Catégories dérivées des schémas,
lemme \ref{perfect-lemma-proper-flat-geom-red}
et des informations contenues dans
Compléments sur les morphismes, théorème
\ref{more-morphisms-theorem-stein-factorization-general}.
\end{proof}

\begin{proposition}
\label{proposition-first-homotopy-sequence}
Soit $f : X \to S$ un morphisme propre et plat de présentation finie dont les
fibres géométriques sont connexes et réduites. Supposons $S$ connexe et
soit $\overline{s}$ un point géométrique de $S$. Alors il existe une suite
exacte
$$
\pi_1(X_{\overline{s}}) \to \pi_1(X) \to \pi_1(S) \to 1
$$
de groupes fondamentaux.
\end{proposition}

\begin{proof}
Soit $Y \to X$ un morphisme fini \'etale. Considérons la factorisation de Stein
$$
\xymatrix{
Y \ar[d] \ar[r] & X \ar[d] \\
T \ar[r] & S
}
$$
de $Y \to S$. D'après le lemme \ref{lemma-stein-factorization-etale},
le morphisme $T \to S$ est fini \'etale. On obtient ainsi
un foncteur $\textit{F\'Et}_X \to \textit{F\'Et}_S$.
Pour tout morphisme fini \'etale $U \to S$, se donner un morphisme
$Y \to U \times_S X$ au-dessus de $X$ revient à se donner un morphisme
$Y \to U$ au-dessus de $S$ ; un tel morphisme se factorise de manière unique par
la factorisation de Stein, c'est-à-dire correspond à un unique
morphisme $T \to U$
(par la construction de la factorisation de Stein comme
normalisation relative dans Compléments sur les morphismes, lemme
\ref{more-morphisms-lemma-stein-universally-closed}
et par la factorisation donnée par
Morphismes, lemme \ref{morphisms-lemma-characterize-normalization}).
Ainsi, on voit que les foncteurs
$\textit{F\'Et}_X \to \textit{F\'Et}_S$ et
$\textit{F\'Et}_S \to \textit{F\'Et}_X$ sont adjoints l'un de l'autre.
Remarquons que la factorisation de Stein de $U \times_S X \to S$ est
$U$, car les fibres de $U \times_S X \to U$ sont géométriquement connexes.

\medskip\noindent
D'après la discussion ci-dessus et
Catégories, lemme \ref{categories-lemma-adjoint-fully-faithful},
on conclut que
$\textit{F\'Et}_S \to \textit{F\'Et}_X$
est pleinement fidèle, c'est-à-dire que $\pi_1(X) \to \pi_1(S)$ est surjectif
(lemme \ref{lemma-functoriality-galois-surjective}).

\medskip\noindent
Il est immédiat que le composé
$\textit{F\'Et}_S \to \textit{F\'Et}_X \to \textit{F\'Et}_{X_{\overline{s}}}$
envoie tout $U$ sur une réunion disjointe de copies de $X_{\overline{s}}$.
Ainsi la composée $\pi_1(X_{\overline{s}}) \to \pi_1(X) \to \pi_1(S)$ est triviale
d'après le lemme \ref{lemma-composition-trivial}.

\medskip\noindent
Soit $Y \to X$ un morphisme fini \'etale, où $Y$ est connexe, tel que
$Y \times_X X_{\overline{s}}$ contienne une composante connexe $Z$
isomorphe à $X_{\overline{s}}$. Considérons la factorisation de Stein $T$
ci-dessus. Soit $\overline{t} \in T_{\overline{s}}$ le point correspondant
à la fibre $Z$. Observons que $T$ est connexe (en tant qu'image d'un schéma
connexe) et que, par la surjectivité ci-dessus, $T \times_S X$ est connexe.
Considérons maintenant la factorisation
$$
\pi : Y \longrightarrow T \times_S X
$$
Soit $\overline{x} \in X_{\overline{s}}$ un point fermé quelconque. Notons que
$\kappa(\overline{t}) = \kappa(\overline{s}) = \kappa(\overline{x})$
est un corps algébriquement clos.
Alors la fibre de $\pi$ au-dessus de $(\overline{t}, \overline{x})$ est constituée
d'un unique point, à savoir l'unique point $\overline{z} \in Z$
correspondant à $\overline{x} \in X_{\overline{s}}$ par
l'isomorphisme $Z \to X_{\overline{s}}$. On conclut que le morphisme fini
\'etale $\pi$ est de degré $1$ dans un voisinage de
$(\overline{t}, \overline{x})$. Puisque $T \times_S X$ est connexe,
ce morphisme est partout de degré $1$, et l'on trouve que $Y \cong T \times_S X$.
Ainsi $Y \times_X X_{\overline{s}}$ se scinde complètement.
En combinant tout ce qui précède, on voit que les
lemmes \ref{lemma-functoriality-galois-ses} et
\ref{lemma-functoriality-galois-normal}
s'appliquent tous deux, ce qui achève la démonstration.
\end{proof}




\section{Homomorphismes de spécialisation}
\label{section-specialization-map}

\noindent
Dans cette section, nous construisons des homomorphismes de spécialisation.
Soit $f : X \to S$ un morphisme propre de schémas
à fibres géométriquement connexes.
Soit $s' \leadsto s$ une spécialisation de points de $S$.
Soient $\overline{s}$ et $\overline{s}'$ des points géométriques
au-dessus de $s$ et $s'$. Il existe alors un homomorphisme de spécialisation
$$
sp : \pi_1(X_{\overline{s}'}) \longrightarrow \pi_1(X_{\overline{s}})
$$
La construction de cet homomorphisme est la suivante. Soit $A$ le
hensélisé strict de $\mathcal{O}_{S, s}$ relativement à
$\kappa(s) \subset \kappa(s)^{sep} \subset \kappa(\overline{s})$, voir
Algèbre, définition \ref{algebra-definition-henselization}.
Puisque $s' \leadsto s$, le point $s'$ correspond à un point de
$\Spec(\mathcal{O}_{S, s})$ et il existe donc au moins un point
(et éventuellement plusieurs points)
de $\Spec(A)$ au-dessus de $s'$ dont le corps résiduel est une extension
algébrique séparable de $\kappa(s')$.
Puisque $\kappa(\overline{s}')$ est algébriquement clos, on peut choisir
un morphisme $\varphi : \overline{s}' \to \Spec(A)$ donnant lieu à un diagramme
commutatif
$$
\xymatrix{
\overline{s}' \ar[r]_-\varphi \ar[rd] &
\Spec(A) \ar[d] &
\overline{s} \ar[l] \ar[ld] \\
& S
}
$$
L'homomorphisme de spécialisation est la composée
$$
\pi_1(X_{\overline{s}'}) \longrightarrow
\pi_1(X_A) =
\pi_1(X_{\kappa(s)^{sep}}) =
\pi_1(X_{\overline{s}})
$$
où la première égalité est donnée par le
lemme \ref{lemma-finite-etale-on-proper-over-henselian}
et la seconde découle des
lemmes \ref{lemma-perfection} et
\ref{lemma-finite-etale-invariant-over-proper}.
Par construction, l'homomorphisme de spécialisation s'insère dans un diagramme
commutatif
$$
\xymatrix{
\pi_1(X_{\overline{s}'}) \ar[rr]_{sp} \ar[rd] & &
\pi_1(X_{\overline{s}}) \ar[ld] \\
& \pi_1(X)
}
$$
pourvu que $X$ soit connexe. L'homomorphisme de spécialisation dépend du
choix de $\varphi : \overline{s}' \to \Spec(A)$ ci-dessus et nous
écrirons $sp_\varphi$ si nous voulons préciser ce choix.

\begin{lemma}
\label{lemma-specialization-map-base-change}
Considérons un diagramme commutatif
$$
\xymatrix{
Y \ar[d]_g \ar[r] & X \ar[d]^f \\
T \ar[r] & S
}
$$
de schémas où $f$ et $g$ sont des morphismes propres à fibres géométriquement connexes.
Soit $t' \leadsto t$ une spécialisation de points de $T$
et considérons un homomorphisme de spécialisation
$sp : \pi_1(Y_{\overline{t}'}) \to \pi_1(Y_{\overline{t}})$ comme ci-dessus.
Il existe alors un diagramme commutatif
$$
\xymatrix{
\pi_1(Y_{\overline{t}'}) \ar[r]_{sp} \ar[d] & \pi_1(Y_{\overline{t}}) \ar[d] \\
\pi_1(X_{\overline{s}'}) \ar[r]^{sp} & \pi_1(X_{\overline{s}})
}
$$
d'homomorphismes de spécialisation, où $\overline{s}$ et $\overline{s}'$
sont les images de $\overline{t}$ et $\overline{t}'$.
\end{lemma}

\begin{proof}
Soit $B$ le hensélisé strict de $\mathcal{O}_{T, t}$ relativement à
$\kappa(t) \subset \kappa(t)^{sep} \subset \kappa(\overline{t})$.
Choisissons $\psi : \overline{t}' \to \Spec(B)$ relevant $\overline{t}' \to T$
comme dans la construction de l'homomorphisme de spécialisation.
Notons $s$ et $s'$ les images de $t$ et $t'$ dans $S$.
Soit $A$ le hensélisé strict de $\mathcal{O}_{S, s}$
relativement à
$\kappa(s) \subset \kappa(s)^{sep} \subset \kappa(\overline{s})$.
Puisque $\kappa(\overline{s}) = \kappa(\overline{t})$,
la fonctorialité de l'hensélisation stricte
(Algèbre, lemme \ref{algebra-lemma-strictly-henselian-functorial})
nous donne un morphisme d'anneaux $A \to B$ s'insérant dans le diagramme commutatif
$$
\xymatrix{
\overline{t}' \ar[r]_-\psi \ar[d] & \Spec(B) \ar[d] \ar[r] & T \ar[d] \\
\overline{s}' \ar[r]^-\varphi & \Spec(A) \ar[r] & S
}
$$
Ici, le morphisme $\varphi : \overline{s}' \to \Spec(A)$ est simplement pris
égal à la composée $\overline{t}' \to \Spec(B) \to \Spec(A)$.
En effectuant le changement de base, on obtient un diagramme commutatif
$$
\xymatrix{
Y_{\overline{t}'} \ar[r] \ar[d] & Y_B \ar[d] \\
X_{\overline{s}'} \ar[r] & X_A
}
$$
et, d'après la construction de l'homomorphisme de spécialisation, la commutativité
de ce diagramme implique celle du diagramme du lemme.
\end{proof}

\begin{lemma}
\label{lemma-specialization-map-composition}
Soit $f : X \to S$ un morphisme propre à fibres géométriquement connexes.
Soient $s'' \leadsto s' \leadsto s$ des spécialisations de points de $S$.
Une composée d'homomorphismes de spécialisation
$\pi_1(X_{\overline{s}''}) \to \pi_1(X_{\overline{s}'}) \to
\pi_1(X_{\overline{s}})$ est un homomorphisme de spécialisation
$\pi_1(X_{\overline{s}''}) \to \pi_1(X_{\overline{s}})$.
\end{lemma}

\begin{proof}
Soit $\mathcal{O}_{S, s} \to A$ l'hensélisé strict
construit à l'aide de $\kappa(s) \to \kappa(\overline{s})$.
Soit $A \to \kappa(\overline{s}')$ le morphisme utilisé pour construire
le premier homomorphisme de spécialisation. Soit $\mathcal{O}_{S, s'} \to A'$
l'hensélisé strict construit à l'aide de
$\kappa(s') \subset \kappa(\overline{s}')$.
Par fonctorialité de l'hensélisation stricte, il existe un morphisme
$A \to A'$ tel que sa composée avec $A' \to \kappa(\overline{s}')$
soit le morphisme donné
(Algèbre, lemme \ref{algebra-lemma-map-into-henselian-colimit}).
Ensuite, soit $A' \to \kappa(\overline{s}'')$ le morphisme utilisé pour
construire le second homomorphisme de spécialisation. Il est alors clair que
la composée du premier et du second homomorphisme de spécialisation
est l'homomorphisme de spécialisation
$\pi_1(X_{\overline{s}''}) \to \pi_1(X_{\overline{s}})$
construit à l'aide de $A \to A' \to \kappa(\overline{s}'')$.
\end{proof}

\noindent
Soit $X \to S$ un morphisme propre à fibres géométriquement connexes.
Soit $R$ un anneau de valuation strictement hensélien dont le corps des fractions
est algébriquement clos et soit $\Spec(R) \to S$
un morphisme. Soient $\eta, s \in \Spec(R)$ les points générique et fermé.
On peut alors considérer l'homomorphisme de spécialisation
$$
sp_R : \pi_1(X_\eta) \to \pi_1(X_s)
$$
pour le changement de base $X_R/\Spec(R)$. Notons que cela a un sens puisque
$\eta$ et $s$ ont tous deux un corps résiduel algébriquement clos.

\begin{lemma}
\label{lemma-specialization-map-valuation-ring}
Soit $f : X \to S$ un morphisme propre à fibres géométriquement connexes.
Soit $s' \leadsto s$ une spécialisation de points de $S$ et soit
$sp : \pi_1(X_{\overline{s}'}) \to \pi_1(X_{\overline{s}})$
un homomorphisme de spécialisation. Il existe alors un anneau de valuation strictement hensélien
$R$ au-dessus de $S$ dont le corps des fractions est algébriquement clos,
tel que $sp$ soit isomorphe à $sp_R$ défini ci-dessus.
\end{lemma}

\begin{proof}
Soit $\mathcal{O}_{S, s} \to A$ l'hensélisé strict
construit à l'aide de $\kappa(s) \to \kappa(\overline{s})$.
Soit $A \to \kappa(\overline{s}')$ le morphisme utilisé pour construire $sp$.
Soit $R \subset \kappa(\overline{s}')$ un anneau de valuation de
corps des fractions $\kappa(\overline{s}')$ qui domine l'image de $A$.
Voir Algèbre, lemme \ref{algebra-lemma-dominate}.
Observons que $R$ est strictement hensélien, par exemple d'après le
lemme \ref{lemma-normal-local-domain-separablly-closed-fraction-field}
et Algèbre, lemme \ref{algebra-lemma-valuation-ring-normal}.
Le lemme est alors clair.
\end{proof}

\noindent
Soit $X \to S$ un morphisme propre à fibres géométriquement connexes.
Soit $R$ un anneau de valuation discrète strictement hensélien et soit
$\Spec(R) \to S$ un morphisme. Soient $\eta, s \in \Spec(R)$ les points
générique et fermé. On peut alors considérer l'homomorphisme de spécialisation
$$
sp_R : \pi_1(X_{\overline{\eta}}) \to \pi_1(X_s)
$$
pour le changement de base $X_R/\Spec(R)$. Notons que cela a un sens puisque $s$
a un corps résiduel algébriquement clos.

\begin{lemma}
\label{lemma-specialization-map-discrete-valuation-ring}
Soit $f : X \to S$ un morphisme propre à fibres géométriquement connexes.
Soit $s' \leadsto s$ une spécialisation de points de $S$ et soit
$sp : \pi_1(X_{\overline{s}'}) \to \pi_1(X_{\overline{s}})$
un homomorphisme de spécialisation. Si $S$ est noethérien, alors
il existe un anneau de valuation discrète strictement hensélien
$R$ au-dessus de $S$ tel que $sp$ soit isomorphe à $sp_R$
défini ci-dessus.
\end{lemma}

\begin{proof}
Soit $\mathcal{O}_{S, s} \to A$ l'hensélisé strict
construit à l'aide de $\kappa(s) \to \kappa(\overline{s})$.
Soit $A \to \kappa(\overline{s}')$ le morphisme utilisé pour construire $sp$.
Soit $R \subset \kappa(\overline{s}')$ un anneau de valuation discrète
qui domine l'image de $A$, voir Algèbre, lemme \ref{algebra-lemma-exists-dvr}.
Choisissons un diagramme de corps
$$
\xymatrix{
\kappa(\overline{s}) \ar[r] & k \\
A/\mathfrak m_A \ar[r] \ar[u] & R/\mathfrak m_R \ar[u]
}
$$
où $k$ est algébriquement clos. Soit $R^{sh}$ l'hensélisé
strict de $R$ construit à l'aide de $R \to k$. On obtient
un morphisme $A \to R^{sh}$ d'après Algèbre, lemme
\ref{algebra-lemma-strictly-henselian-functorial}.
L'anneau $R^{sh}$ est un anneau de valuation discrète d'après
Compléments d'algèbre, lemme \ref{more-algebra-lemma-henselization-dvr}.
Notons $\eta, o$ les points générique et fermé de $\Spec(R^{sh})$.
Comme le diagramme de schémas
$$
\xymatrix{
\overline{\eta} \ar[d] \ar[r] & \Spec(R^{sh}) \ar[d] &
\Spec(k) \ar[d] \ar[l] \\
\overline{s}' \ar[r] & \Spec(A) & \overline{s} \ar[l]
}
$$
est commutatif, on obtient un diagramme commutatif
$$
\xymatrix{
\pi_1(X_{\overline{\eta}}) \ar[d] \ar[r]_{sp_{R^{sh}}} & \pi_1(X_o) \ar[d] \\
\pi_1(X_{\overline{s}'}) \ar[r]^{sp} & \pi_1(X_{\overline{s}})
}
$$
d'homomorphismes de spécialisation, par la construction de ces homomorphismes.
Comme les flèches verticales sont des isomorphismes
(lemme \ref{lemma-finite-etale-invariant-over-proper}), cela démontre le lemme.
\end{proof}








\section{Restriction à un sous-schéma fermé}
\label{section-lefschetz}

\noindent
Dans cette section, nous démontrons quelques résultats sur le foncteur de restriction
$$
\textit{F\'Et}_X \longrightarrow \textit{F\'Et}_Y,\quad
U \longmapsto V = U \times_X Y
$$
où $X$ est un schéma et $Y$ un sous-schéma fermé. Grâce à
l'invariance topologique du groupe fondamental, on peut ramener
l'étude de ce foncteur à celle du foncteur de complétion sur les
modules finis localement libres.

\medskip\noindent
Dans les lemmes suivants, nous utilisons la notion de modules formels cohérents
définie dans
Cohomologie des schémas, section \ref{coherent-section-coherent-formal}.
Étant donnés un schéma noethérien et un faisceau quasi-cohérent d'idéaux
$\mathcal{I} \subset \mathcal{O}_X$, nous dirons
qu'un objet $(\mathcal{F}_n)$ de $\textit{Coh}(X, \mathcal{I})$
est {\it fini localement libre} si chaque $\mathcal{F}_n$ est un
$\mathcal{O}_X/\mathcal{I}^n$-module fini localement libre.

\begin{lemma}
\label{lemma-restriction-fully-faithful}
Soit $X$ un schéma noethérien et soit $Y \subset X$ un sous-schéma fermé
de faisceau d'idéaux $\mathcal{I} \subset \mathcal{O}_X$.
Supposons que le foncteur de complétion
$$
\textit{Coh}(\mathcal{O}_X)
\longrightarrow 
\textit{Coh}(X, \mathcal{I}),\quad
\mathcal{F} \longmapsto \mathcal{F}^\wedge
$$
soit pleinement fidèle sur la sous-catégorie pleine des objets finis
localement libres (voir ci-dessus).
Alors le foncteur de restriction $\textit{F\'Et}_X \to \textit{F\'Et}_Y$
est pleinement fidèle.
\end{lemma}

\begin{proof}
Comme la catégorie des revêtements finis étales possède un
Hom interne (lemme \ref{lemma-internal-hom-finite-etale}),
il suffit de démontrer ceci : étant donnés $U$ fini étale sur $X$
et un morphisme $t : Y \to U$ sur $X$, il existe une unique section
$s : X \to U$ telle que $t = s|_Y$. Voici le diagramme :
$$
\xymatrix{
& U \ar[d]^f \\
Y \ar[r] \ar[ru] & X \ar@{..>}@/^1em/[u]
}
$$
Trouver la flèche pointillée $s$ revient à trouver un morphisme de
$\mathcal{O}_X$-algèbres
$$
s^\sharp : f_*\mathcal{O}_U \longrightarrow \mathcal{O}_X
$$
qui, modulo le faisceau d'idéaux de $Y$, se réduise au morphisme d'algèbres donné
$t^\sharp : f_*\mathcal{O}_U \to \mathcal{O}_Y$.
D'après le lemme \ref{lemma-thickening}, on peut relever $t$ de façon unique en un système
compatible de morphismes $t_n : Y_n \to U$, et donc en un morphisme
$$
\lim t_n^\sharp : f_*\mathcal{O}_U \longrightarrow \lim \mathcal{O}_{Y_n}
$$
de faisceaux d'algèbres sur $X$.
Comme $f_*\mathcal{O}_U$ est un $\mathcal{O}_X$-module fini localement libre,
on en déduit un unique morphisme de $\mathcal{O}_X$-modules
$\sigma : f_*\mathcal{O}_U \to \mathcal{O}_X$ dont la complétion
est $\lim t_n^\sharp$. Pour voir que $\sigma$ est un homomorphisme d'algèbres,
il faut vérifier que le diagramme
$$
\xymatrix{
f_*\mathcal{O}_U \otimes_{\mathcal{O}_X} f_*\mathcal{O}_U
\ar[r] \ar[d]_{\sigma \otimes \sigma} &
f_*\mathcal{O}_U \ar[d]^\sigma \\
\mathcal{O}_X \otimes_{\mathcal{O}_X} \mathcal{O}_X \ar[r] &
\mathcal{O}_X
}
$$
est commutatif. Pour tout $n$, on sait que ce diagramme est commutatif après restriction
à $Y_n$, c'est-à-dire après application du foncteur de complétion.
La fidélité du foncteur de complétion permet donc de conclure.
\end{proof}

\begin{lemma}
\label{lemma-restriction-equivalence}
Soit $X$ un schéma noethérien et soit $Y \subset X$ un sous-schéma fermé
de faisceau d'idéaux $\mathcal{I} \subset \mathcal{O}_X$.
Supposons que le foncteur de complétion
$$
\textit{Coh}(\mathcal{O}_X)
\longrightarrow 
\textit{Coh}(X, \mathcal{I}),\quad
\mathcal{F} \longmapsto \mathcal{F}^\wedge
$$
soit une équivalence sur les sous-catégories pleines d'objets finis
localement libres (voir ci-dessus).
Alors le foncteur de restriction $\textit{F\'Et}_X \to \textit{F\'Et}_Y$
est une équivalence.
\end{lemma}

\begin{proof}
Le foncteur de restriction est pleinement fidèle d'après le
lemme \ref{lemma-restriction-fully-faithful}.

\medskip\noindent
Soit $U_1 \to Y$ un morphisme fini étale. Pour achever la démonstration,
nous allons montrer que $U_1$ appartient à l'image essentielle du
foncteur de restriction.

\medskip\noindent
Pour $n \geq 1$, soit $Y_n$ le $n$-ième voisinage infinitésimal de $Y$.
D'après le lemme \ref{lemma-thickening},
il existe un unique morphisme fini étale
$\pi_n : U_n \to Y_n$ dont le changement de base à $Y = Y_1$
redonne $U_1 \to Y_1$.
Considérons les faisceaux $\mathcal{F}_n = \pi_{n, *}\mathcal{O}_{U_n}$.
Nous pouvons et allons considérer $\mathcal{F}_n$ comme un $\mathcal{O}_X$-module sur $X$
qui est localement isomorphe à
$(\mathcal{O}_X/\mathcal{I}^n)^{\oplus r}$.
Ce $(\mathcal{F}_n)$ est un objet fini localement libre de
$\textit{Coh}(X, \mathcal{I})$.
Par hypothèse, il existe un $\mathcal{O}_X$-module fini localement libre
$\mathcal{F}$ et un système compatible d'isomorphismes
$$
\mathcal{F}/\mathcal{I}^n\mathcal{F} \to \mathcal{F}_n
$$
de $\mathcal{O}_X$-modules.

\medskip\noindent
Pour munir $\mathcal{F}$ d'une structure d'algèbre, considérons les
morphismes de multiplication
$\mathcal{F}_n \otimes_{\mathcal{O}_X} \mathcal{F}_n \to \mathcal{F}_n$
provenant du fait que les $\mathcal{F}_n = \pi_{n, *}\mathcal{O}_{U_n}$
sont des faisceaux d'algèbres. Ceux-ci définissent un morphisme
$$
(\mathcal{F}\otimes_{\mathcal{O}_X} \mathcal{F})^\wedge
\longrightarrow
\mathcal{F}^\wedge
$$
dans la catégorie $\textit{Coh}(X, \mathcal{I})$. Par hypothèse,
on peut donc supposer qu'il existe un morphisme
$\mu : \mathcal{F}\otimes_{\mathcal{O}_X} \mathcal{F} \to \mathcal{F}$
dont la restriction à $Y_n$ donne les morphismes de multiplication ci-dessus.
La fidélité du foncteur de l'énoncé du lemme montre alors
que $\mu$ définit une structure de $\mathcal{O}_X$-algèbre commutative
sur $\mathcal{F}$, compatible avec les structures d'algèbres données
sur les $\mathcal{F}_n$.
En posant
$$
U = \underline{\Spec}_X((\mathcal{F}, \mu))
$$
on obtient un schéma fini localement libre $\pi : U \to X$ dont la restriction
à $Y$ est isomorphe à $U_1$. Le discriminant de $\pi$ est le lieu des zéros
de la section
$$
\det(Q_\pi) :
\mathcal{O}_X
\longrightarrow
\wedge^{top}(\pi_*\mathcal{O}_U)^{\otimes -2}
$$
construite dans
Discriminants, section \ref{discriminant-section-discriminant}.
Comme sa restriction à $Y_n$ est un isomorphisme pour tout $n$
d'après Discriminants, lemme \ref{discriminant-lemma-discriminant},
on en conclut que c'est un isomorphisme. Ainsi $\pi$ est étale d'après
Discriminants, lemme \ref{discriminant-lemma-discriminant}.
\end{proof}

\begin{lemma}
\label{lemma-restriction-fully-faithful-general}
Soit $X$ un schéma noethérien et soit $Y \subset X$ un sous-schéma fermé
de faisceau d'idéaux $\mathcal{I} \subset \mathcal{O}_X$.
Soit $\mathcal{V}$ l'ensemble des sous-schémas ouverts $V \subset X$ contenant $Y$,
ordonné par inclusion inverse. Supposons que le foncteur de complétion
$$
\colim_\mathcal{V} \textit{Coh}(\mathcal{O}_V)
\longrightarrow
\textit{Coh}(X, \mathcal{I}),
\quad
\mathcal{F} \longmapsto \mathcal{F}^\wedge
$$
soit pleinement fidèle sur la sous-catégorie pleine des
objets finis localement libres (voir ci-dessus).
Alors le foncteur de restriction
$\colim_\mathcal{V} \textit{F\'Et}_V \to \textit{F\'Et}_Y$
est pleinement fidèle.
\end{lemma}

\begin{proof}
Observons que $\mathcal{V}$ est un ensemble filtrant, de sorte que les colimites sont
celles de Catégories, section \ref{categories-section-directed-colimits}.
La suite de l'argument est presque exactement la même que dans la démonstration
du lemme \ref{lemma-restriction-fully-faithful}; nous invitons le lecteur
à l'omettre.

\medskip\noindent
Comme la catégorie des revêtements finis étales possède un
Hom interne (lemme \ref{lemma-internal-hom-finite-etale}),
il suffit de démontrer ceci : étant donnés $U$ fini étale sur
$V \in \mathcal{V}$
et un morphisme $t : Y \to U$ sur $V$, il existe $V' \geq V$
et un morphisme $s : V' \to U$ sur $V$ tel que $t = s|_Y$. Voici le diagramme :
$$
\xymatrix{
& & U \ar[d]^f \\
Y \ar[r] \ar[rru] & V' \ar@{..>}[ru] \ar[r] & V
}
$$
Trouver la flèche pointillée $s$ revient à trouver un morphisme de
$\mathcal{O}_{V'}$-algèbres
$$
s^\sharp : f_*\mathcal{O}_U|_{V'} \longrightarrow \mathcal{O}_{V'}
$$
qui, modulo le faisceau d'idéaux de $Y$, se réduise au morphisme d'algèbres donné
$t^\sharp : f_*\mathcal{O}_U \to \mathcal{O}_Y$.
D'après le lemme \ref{lemma-thickening}, on peut relever $t$ de façon unique en un système
compatible de morphismes $t_n : Y_n \to U$, et donc un morphisme
$$
\lim t_n^\sharp : f_*\mathcal{O}_U \longrightarrow \lim \mathcal{O}_{Y_n}
$$
de faisceaux d'algèbres sur $V$.
Observons que $f_*\mathcal{O}_U$ est un
$\mathcal{O}_V$-module fini localement libre. Il existe donc $V' \geq V$ et un morphisme
$\sigma : f_*\mathcal{O}_U|_{V'} \to \mathcal{O}_{V'}$
dont la complétion est $\lim t_n^\sharp$.
Pour voir que $\sigma$ est un homomorphisme d'algèbres, il faut vérifier
que le diagramme
$$
\xymatrix{
(f_*\mathcal{O}_U \otimes_{\mathcal{O}_V} f_*\mathcal{O}_U)|_{V'}
\ar[r] \ar[d]_{\sigma \otimes \sigma} &
f_*\mathcal{O}_U|_{V'} \ar[d]^\sigma \\
\mathcal{O}_{V'} \otimes_{\mathcal{O}_{V'}} \mathcal{O}_{V'} \ar[r] &
\mathcal{O}_{V'}
}
$$
est commutatif. Pour tout $n$, on sait que ce diagramme est commutatif après restriction
à $Y_n$, c'est-à-dire après application du foncteur de complétion.
La fidélité du foncteur de complétion permet donc
d'en déduire qu'il existe $V'' \geq V'$ tel que
$\sigma|_{V''}$ soit un homomorphisme d'algèbres, comme souhaité.
\end{proof}

\begin{lemma}
\label{lemma-restriction-equivalence-general}
Soit $X$ un schéma noethérien et soit $Y \subset X$ un sous-schéma fermé
de faisceau d'idéaux $\mathcal{I} \subset \mathcal{O}_X$.
Soit $\mathcal{V}$ l'ensemble des sous-schémas ouverts $V \subset X$ contenant $Y$,
ordonné par inclusion inverse. Supposons que le foncteur de complétion
$$
\colim_\mathcal{V} \textit{Coh}(\mathcal{O}_V)
\longrightarrow
\textit{Coh}(X, \mathcal{I}),
\quad
\mathcal{F} \longmapsto \mathcal{F}^\wedge
$$
induise une équivalence des sous-catégories pleines d'objets
finis localement libres (voir l'explication ci-dessus).
Alors le foncteur de restriction
$$
\colim_\mathcal{V} \textit{F\'Et}_V \to \textit{F\'Et}_Y
$$
est une équivalence.
\end{lemma}

\begin{proof}
Observons que $\mathcal{V}$ est un ensemble filtrant, de sorte que les colimites sont
celles de Catégories, section \ref{categories-section-directed-colimits}.
La suite de l'argument est presque exactement la même que dans la démonstration
du lemme \ref{lemma-restriction-equivalence}; nous invitons le lecteur
à l'omettre.

\medskip\noindent
Le foncteur de restriction est pleinement fidèle d'après le
lemme \ref{lemma-restriction-fully-faithful-general}.

\medskip\noindent
Soit $U_1 \to Y$ un morphisme fini étale. Pour achever la démonstration,
nous allons montrer que $U_1$ appartient à l'image essentielle du
foncteur de restriction.

\medskip\noindent
Pour $n \geq 1$, soit $Y_n$ le $n$-ième voisinage infinitésimal de $Y$.
D'après le lemme \ref{lemma-thickening},
il existe un unique morphisme fini étale
$\pi_n : U_n \to Y_n$ dont le changement de base à $Y = Y_1$
redonne $U_1 \to Y_1$.
Considérons les faisceaux $\mathcal{F}_n = \pi_{n, *}\mathcal{O}_{U_n}$.
Nous pouvons et allons considérer $\mathcal{F}_n$ comme un $\mathcal{O}_X$-module sur $X$
qui est localement isomorphe à
$(\mathcal{O}_X/\mathcal{I}^n)^{\oplus r}$.
Ce $(\mathcal{F}_n)$ est un objet fini localement libre de
$\textit{Coh}(X, \mathcal{I})$.
Par hypothèse, il existe $V \in \mathcal{V}$,
un $\mathcal{O}_V$-module fini localement libre $\mathcal{F}$
et un système compatible d'isomorphismes
$$
\mathcal{F}/\mathcal{I}^n\mathcal{F} \to \mathcal{F}_n
$$
de $\mathcal{O}_V$-modules.

\medskip\noindent
Pour munir $\mathcal{F}$ d'une structure d'algèbre, considérons les
morphismes de multiplication
$\mathcal{F}_n \otimes_{\mathcal{O}_V} \mathcal{F}_n \to \mathcal{F}_n$
provenant du fait que les $\mathcal{F}_n = \pi_{n, *}\mathcal{O}_{U_n}$
sont des faisceaux d'algèbres. Ceux-ci définissent un morphisme
$$
(\mathcal{F}\otimes_{\mathcal{O}_V} \mathcal{F})^\wedge
\longrightarrow
\mathcal{F}^\wedge
$$
dans la catégorie $\textit{Coh}(X, \mathcal{I})$. Par hypothèse,
après avoir rétréci $V$, on peut supposer qu'il existe un morphisme
$\mu : \mathcal{F}\otimes_{\mathcal{O}_V} \mathcal{F} \to \mathcal{F}$
dont la restriction à $Y_n$ donne les morphismes de multiplication ci-dessus.
Après l'avoir éventuellement rétréci davantage, on peut supposer que $\mu$
définit une structure de $\mathcal{O}_V$-algèbre commutative
sur $\mathcal{F}$, compatible avec les structures d'algèbres données
sur les $\mathcal{F}_n$.
En posant
$$
U = \underline{\Spec}_V((\mathcal{F}, \mu))
$$
on obtient un schéma fini localement libre sur $V$ dont la restriction
à $Y$ est isomorphe à $U_1$. Il s'ensuit que $U \to V$
est étale en tous les points situés au-dessus de $Y$; voir
Compléments sur les morphismes, lemme
\ref{more-morphisms-lemma-check-smoothness-on-infinitesimal-nbhds}.
Ainsi, après avoir rétréci $V$ une fois encore, on peut supposer que $U \to V$ est
fini étale. Cela achève la démonstration.
\end{proof}

\begin{lemma}
\label{lemma-restriction-faithful}
Soit $X$ un schéma et soit $Y \subset X$ un sous-schéma fermé.
Si toute composante connexe de $X$ rencontre $Y$, alors
le foncteur de restriction $\textit{F\'Et}_X \to \textit{F\'Et}_Y$
est fidèle.
\end{lemma}

\begin{proof}
Soient $a, b : U \to U'$ deux morphismes de schémas finis étales sur $X$
dont les restrictions à $Y$ sont égales. L'image d'une composante connexe
de $U$ est une composante connexe de $X$; cela résulte de
Topologie, lemme \ref{topology-lemma-finite-fibre-connected-components},
appliqué à la restriction de $U \to X$ à une composante connexe de $U$.
Par conséquent, l'image de toute composante connexe de $U$ rencontre $Y$
par hypothèse. On en conclut que $a = b$ après restriction à chaque
composante connexe de $U$, d'après Morphismes étales, proposition
\ref{etale-proposition-equality}. Comme l'égalisateur de $a$ et $b$
est un sous-schéma ouvert de $U$ (puisque la diagonale de $U'$ sur $X$ est ouverte),
on conclut.
\end{proof}

\begin{lemma}
\label{lemma-restriction-fully-faithful-special}
Soit $X$ un schéma noethérien et soit $Y \subset X$ un sous-schéma fermé.
Soit $Y_n \subset X$ le $n$-ième voisinage infinitésimal de $Y$ dans $X$.
Supposons que l'une des conditions suivantes soit satisfaite :
\begin{enumerate}
\item $X$ est quasi-affine et
$\Gamma(X, \mathcal{O}_X) \to \lim \Gamma(Y_n, \mathcal{O}_{Y_n})$
est un isomorphisme, ou bien
\item $X$ possède un module inversible ample $\mathcal{L}$ et
$\Gamma(X, \mathcal{L}^{\otimes m}) \to
\lim \Gamma(Y_n, \mathcal{L}^{\otimes m}|_{Y_n})$
est un isomorphisme pour tout $m \gg 0$, ou bien
\item pour tout $\mathcal{O}_X$-module fini localement libre
$\mathcal{E}$, le morphisme
$\Gamma(X, \mathcal{E}) \to \lim \Gamma(Y_n, \mathcal{E}|_{Y_n})$
est un isomorphisme.
\end{enumerate}
Alors le foncteur de restriction $\textit{F\'Et}_X \to \textit{F\'Et}_Y$
est pleinement fidèle.
\end{lemma}

\begin{proof}
Ce lemme résulte formellement du
lemme \ref{lemma-restriction-fully-faithful} et de
Géométrie algébrique et formelle, lemme
\ref{algebraization-lemma-completion-fully-faithful}.
\end{proof}

\begin{lemma}
\label{lemma-restriction-fully-faithful-general-special}
Soit $X$ un schéma noethérien et soit $Y \subset X$ un sous-schéma fermé.
Soit $Y_n \subset X$ le $n$-ième voisinage infinitésimal de $Y$ dans $X$.
Soit $\mathcal{V}$ l'ensemble des sous-schémas ouverts $V \subset X$ contenant $Y$,
ordonné par inclusion inverse. Supposons que l'une des conditions suivantes soit satisfaite :
\begin{enumerate}
\item $X$ est quasi-affine et
$$
\colim_\mathcal{V} \Gamma(V, \mathcal{O}_V)
\longrightarrow
\lim \Gamma(Y_n, \mathcal{O}_{Y_n})
$$
est un isomorphisme, ou bien
\item $X$ possède un module inversible ample $\mathcal{L}$ et
$$
\colim_\mathcal{V} \Gamma(V, \mathcal{L}^{\otimes m})
\longrightarrow
\lim \Gamma(Y_n, \mathcal{L}^{\otimes m}|_{Y_n})
$$
est un isomorphisme pour tout $m \gg 0$, ou bien
\item pour tout $V \in \mathcal{V}$ et tout
$\mathcal{O}_V$-module fini localement libre $\mathcal{E}$, le morphisme
$$
\colim_{V' \geq V} \Gamma(V', \mathcal{E}|_{V'})
\longrightarrow
\lim \Gamma(Y_n, \mathcal{E}|_{Y_n})
$$
est un isomorphisme.
\end{enumerate}
Alors le foncteur
$$
\colim_\mathcal{V} \textit{F\'Et}_V \to \textit{F\'Et}_Y
$$
est pleinement fidèle.
\end{lemma}

\begin{proof}
Ce lemme résulte formellement du
lemme \ref{lemma-restriction-fully-faithful-general} et de
Géométrie algébrique et formelle, lemme
\ref{algebraization-lemma-completion-fully-faithful-general}.
\end{proof}






\section{Sommes amalgamées et groupes fondamentaux}
\label{section-pushouts}

\noindent
Voici le résultat principal.

\begin{lemma}
\label{lemma-pushout-along-closed-immersion-and-integral}
Dans Compléments sur les morphismes, situation
\ref{more-morphisms-situation-pushout-along-closed-immersion-and-integral},
par exemple si $Z \to Y$ et $Z \to X$ sont des immersions fermées de schémas,
il existe une équivalence de catégories
$$
\textit{F\'Et}_{Y \amalg_Z X}
\longrightarrow
\textit{F\'Et}_Y
\times_{\textit{F\'Et}_Z}
\textit{F\'Et}_X
$$
\end{lemma}

\begin{proof}
La somme amalgamée existe d'après
Compléments sur les morphismes, proposition
\ref{more-morphisms-proposition-pushout-along-closed-immersion-and-integral}.
Le foncteur associe à un schéma $U$ fini étale sur la
somme amalgamée les changements de base $Y' = U \times_{Y \amalg_Z X} Y$
et $X' = U \times_{Y \amalg_Z X} X$, munis de l'isomorphisme naturel
$Y' \times_Y Z \to X' \times_X Z$ sur $Z$. Pour montrer que ce foncteur
est une équivalence, nous utilisons
Compléments sur les morphismes, lemme
\ref{more-morphisms-lemma-pushout-functor-equivalence-flat}
pour construire un foncteur quasi-inverse.
Il reste seulement à montrer que, si l'on se donne un morphisme
$U \to Y \amalg_Z X$ séparé, quasi-fini et étale
tel que $X' \to X$ et $Y' \to Y$ soient finis,
alors $U \to Y \amalg_Z X$ est fini.
Cela se déduit soit du fait algébrique correspondant
(Compléments d'algèbre, lemme
\ref{more-algebra-lemma-finite-module-over-fibre-product}),
soit du fait que
$$
X' \amalg Y' \to U
$$
est surjectif et que $X'$ et $Y'$ sont propres sur $Y \amalg_Z X$
(on utilise ici la description de la somme amalgamée dans Compléments sur les morphismes, proposition
\ref{more-morphisms-proposition-pushout-along-closed-immersion-and-integral}),
après quoi on peut appliquer
Morphismes, lemme \ref{morphisms-lemma-scheme-theoretic-image-is-proper},
pour conclure que $U$ est propre sur $Y \amalg_Z X$.
Comme un morphisme quasi-fini et propre est fini
(Compléments sur les morphismes, lemme \ref{more-morphisms-lemma-characterize-finite}),
la démonstration est achevée.
\end{proof}





\section{Revêtements finis étales des spectres épointés, I}
\label{section-pi1-punctured-spec}

\noindent
Nous démontrons d'abord quelques résultats à la Lefschetz.

\begin{situation}
\label{situation-local-lefschetz}
Soient $(A, \mathfrak m)$ un anneau local noethérien et $f \in \mathfrak m$.
Posons $X = \Spec(A)$ et $X_0 = \Spec(A/fA)$, et
notons $U = X \setminus \{\mathfrak m\}$ et
$U_0 = X_0 \setminus \{\mathfrak m\}$ les spectres épointés de
$A$ et $A/fA$.
\end{situation}

\noindent
Rappelons que, pour un schéma $X$, la catégorie des schémas finis
étales sur $X$ est notée $\textit{F\'Et}_X$ ; voir la
section \ref{section-finite-etale}.
Dans la situation \ref{situation-local-lefschetz},
nous étudierons les foncteurs de changement de base
$$
\xymatrix{
\textit{F\'Et}_X \ar[d] \ar[r] & \textit{F\'Et}_U \ar[d] \\
\textit{F\'Et}_{X_0} \ar[r] & \textit{F\'Et}_{U_0}
}
$$
Dans de nombreux cas, la flèche verticale de droite est fidèle.

\begin{lemma}
\label{lemma-faithful}
Dans la situation \ref{situation-local-lefschetz},
supposons que l'une des conditions suivantes soit satisfaite :
\begin{enumerate}
\item $\dim(A/\mathfrak p) \geq 2$ pour tout idéal premier minimal
$\mathfrak p \subset A$ tel que $f \not \in \mathfrak p$, ou bien
\item toute composante connexe de $U$ rencontre $U_0$.
\end{enumerate}
Alors
$$
\textit{F\'Et}_U \longrightarrow \textit{F\'Et}_{U_0},\quad
V \longmapsto V_0 = V \times_U U_0
$$
est un foncteur fidèle.
\end{lemma}

\begin{proof}
Le cas (2) résulte immédiatement du lemme \ref{lemma-restriction-faithful}.
L'hypothèse (1) implique que toute composante irréductible de $U$ rencontre $U_0$ ; voir
Algèbre, lemme \ref{algebra-lemma-one-equation}.
Ainsi (1) résulte de (2).
\end{proof}

\noindent
Avant de démontrer un résultat plus intéressant, il nous faut quelques lemmes.

\begin{lemma}
\label{lemma-fill-in-missing}
Dans la situation \ref{situation-local-lefschetz}, soit $V \to U$ un morphisme
fini. Soit $A^\wedge$ le complété $\mathfrak m$-adique de $A$,
soit $X' = \Spec(A^\wedge)$, et soient $U'$ et $V'$ les changements de base de
$U$ et $V$ à $X'$. S'il existe un morphisme fini $Y' \to X'$ tel que
$V' = Y' \times_{X'} U'$, alors il existe un morphisme fini $Y \to X$
tel que $V = Y \times_X U$ et $Y' = Y \times_X X'$.
\end{lemma}

\begin{proof}
C'est une application directe de
Compléments d'algèbre, proposition \ref{more-algebra-proposition-equivalence}.
Choisissons en effet des générateurs $f_1, \ldots, f_t$ de $\mathfrak m$.
Pour tout $i$, écrivons $V \times_U D(f_i) = \Spec(B_i)$.
Pour $1 \leq i, j \leq t$, on obtient un isomorphisme
$\alpha_{ij} : (B_i)_{f_j} \to (B_j)_{f_i}$ de $A_{f_if_j}$-algèbres,
car les spectres de ces deux algèbres représentent $V \times_U D(f_if_j)$.
Écrivons $Y' = \Spec(B')$. Comme $V \times_U U' = Y' \times_{X'} U'$
on obtient des isomorphismes $\alpha_i : B'_{f_i} \to B_i \otimes_A A^\wedge$.
Un argument direct montre que $(B', B_i, \alpha_i, \alpha_{ij})$
est un objet de $\text{Glue}(A \to A^\wedge, f_1, \ldots, f_t)$ ; voir
Compléments d'algèbre, remarque \ref{more-algebra-remark-glueing-data}.
En appliquant la proposition citée ci-dessus (et en utilisant
Compléments d'algèbre, remarque \ref{more-algebra-remark-formal-glueing-algebras},
pour obtenir la structure d'algèbre), on trouve une $A$-algèbre $B$ telle que
$\text{Can}(B)$ soit isomorphe à $(B', B_i, \alpha_i, \alpha_{ij})$.
En posant $Y = \Spec(B)$, on voit que $Y \to X$ est un morphisme
muni d'isomorphismes compatibles
$V \cong Y \times_X U$ et $Y' = Y \times_X X'$, comme souhaité.
\end{proof}

\begin{lemma}
\label{lemma-fully-faithful-henselian-completion}
Dans la situation \ref{situation-local-lefschetz}, supposons $A$ hensélien
ou, plus généralement, que $(A, (f))$ soit un couple hensélien.
Soit $A^\wedge$ le complété $\mathfrak m$-adique de $A$,
soit $X' = \Spec(A^\wedge)$, et soient $U'$ et $U'_0$ les changements de base de
$U$ et $U_0$ à $X'$. Si $\textit{F\'Et}_{U'} \to \textit{F\'Et}_{U'_0}$
est pleinement fidèle, alors $\textit{F\'Et}_U \to \textit{F\'Et}_{U_0}$
est pleinement fidèle.
\end{lemma}

\begin{proof}
Supposons $\textit{F\'Et}_{U'} \longrightarrow \textit{F\'Et}_{U'_0}$
pleinement fidèle. Comme $X' \to X$ est fidèlement plat, il est
immédiat que le foncteur $V \to V_0 = V \times_U U_0$ est fidèle.
Comme la catégorie des revêtements finis étales possède un Hom interne
(lemme \ref{lemma-internal-hom-finite-etale}),
il suffit de démontrer ceci : pour $V$ fini étale sur $U$,
on a
$$
\Mor_U(U, V) = \Mor_{U_0}(U_0, V_0)
$$
Supposons donc donné un morphisme $s_0 : U_0 \to V_0$ sur $U_0$ ; nous allons
construire un morphisme $s : U \to V$ sur $U$.

\medskip\noindent
Par hypothèse, il existe bien un morphisme $s' : U' \to V'$
dont la restriction à $U'_0$ est le changement de base $s'_0$ de $s_0$.
Comme $V' \to U'$ est fini étale, cela signifie que $V' = s'(U') \amalg W'$
pour un certain morphisme $W' \to U'$ fini et étale.
Choisissons un morphisme fini $Z' \to X'$ tel que $W' = Z' \times_{X'} U'$.
C'est possible grâce au théorème principal de Zariski, sous la forme énoncée dans
Compléments sur les morphismes, lemme
\ref{more-morphisms-lemma-quasi-finite-separated-pass-through-finite}
(nous omettons un petit détail).
Alors
$$
V' = s'(U') \amalg W' \longrightarrow X' \amalg Z' = Y'
$$
est une immersion ouverte telle que $V' = Y' \times_{X'} U'$.
D'après le lemme \ref{lemma-fill-in-missing}, il existe $Y \to X$ fini
tel que $V = Y \times_X U$ et $Y' = Y \times_X X'$.
Écrivons $Y = \Spec(B)$, de sorte que $Y' = \Spec(B \otimes_A A^\wedge)$.
Alors $B \otimes_A A^\wedge$ possède un idempotent $e'$
correspondant au sous-schéma ouvert et fermé $X'$ de $Y' = X' \amalg Z'$.

\medskip\noindent
Supposons d'abord $A$ hensélien (ce cas est légèrement plus simple). L'image $\overline{e}$
de $e'$ dans $B \otimes_A \kappa(\mathfrak m) = B/\mathfrak mB$ se relève en un
idempotent $e$ de $B$, puisque $A$ est hensélien (car $B$ est un produit d'anneaux
locaux d'après Algèbre, lemme \ref{algebra-lemma-characterize-henselian}).
On voit alors que $e$ s'envoie sur $e'$ par unicité du relèvement des idempotents
(en utilisant que $B \otimes_A A^\wedge$ est un produit d'anneaux locaux).
Soit $Y_1 \subset Y$ le sous-schéma ouvert et fermé correspondant à $e$.
Alors $Y_1 \times_X X' = s'(X')$, ce qui implique que $Y_1 \to X$ est
un isomorphisme (par descente fidèlement plate) et fournit la section voulue.

\medskip\noindent
Supposons maintenant que $(A, (f))$ soit un couple hensélien. Nous utilisons ici que $s'$ est
un relèvement de $s'_0$. Soit en effet $Y_{0, 1} \subset Y_0 = Y \times_X X_0$
l'adhérence de $s_0(U_0) \subset V_0 = Y_0 \times_{X_0} U_0$.
Comme $X' \to X$ est plat, le changement de base $Y'_{0, 1} \subset Y'_0$
est l'adhérence de $s'_0(U'_0)$, laquelle est égale à $X'_0 \subset Y'_0$
(voir Morphismes, lemme
\ref{morphisms-lemma-flat-base-change-scheme-theoretic-image}).
Comme $Y'_0 \to Y_0$ est submersif
(Morphismes, lemme \ref{morphisms-lemma-fpqc-quotient-topology}),
on conclut que $Y_{0, 1}$ est ouvert et fermé dans $Y_0$.
Soit $e_0 \in B/fB$ l'idempotent correspondant.
D'après Compléments d'algèbre, lemme
\ref{more-algebra-lemma-characterize-henselian-pair},
on peut relever $e_0$ en un idempotent $e \in B$.
On conclut alors comme précédemment.
\end{proof}

\noindent
Dans la situation \ref{situation-local-lefschetz},
la pleine fidélité du foncteur de restriction
$\textit{F\'Et}_U \longrightarrow \textit{F\'Et}_{U_0}$
vaut sous des hypothèses assez faibles.
En particulier, ces hypothèses n'impliquent souvent pas que
$U$ soit un schéma connexe, mais la conclusion garantit
que $U$ et $U_0$ ont le même nombre de composantes connexes.

\begin{lemma}
\label{lemma-fully-faithful-simple}
Dans la situation \ref{situation-local-lefschetz}, supposons que
\begin{enumerate}
\item[(a)] $A$ possède un complexe dualisant,
\item[(b)] le couple $(A, (f))$ est hensélien,
\item[(c)] l'une des conditions suivantes est satisfaite :
\begin{enumerate}
\item[(i)] $A_f$ vérifie $(S_2)$ et toute composante irréductible de $X$
non contenue dans $X_0$ est de dimension $\geq 3$, ou bien
\item[(ii)] pour tout idéal premier
$\mathfrak p \subset A$ tel que $f \not \in \mathfrak p$, on ait
$\text{depth}(A_\mathfrak p) + \dim(A/\mathfrak p) > 2$.
\end{enumerate}
\end{enumerate}
Alors le foncteur de restriction
$\textit{F\'Et}_U \longrightarrow \textit{F\'Et}_{U_0}$
est pleinement fidèle.
\end{lemma}

\begin{proof}
Soit $A'$ le complété $\mathfrak m$-adique de $A$. Nous allons montrer que
les hypothèses restent vraies pour $A'$. C'est clair pour les conditions
(a) et (b). La condition (c)(ii) est préservée par
Cohomologie locale, lemme \ref{local-cohomology-lemma-change-completion}.
Supposons ensuite que (c)(i) soit satisfaite. Comme $A$ est universellement caténaire
(Complexes dualisants, lemme \ref{dualizing-lemma-universally-catenary}),
on voit que toute composante irréductible de $\Spec(A')$ non contenue dans $V(f)$
est de dimension $\geq 3$ ; voir
Compléments d'algèbre, proposition \ref{more-algebra-proposition-ratliff}.
Comme $A \to A'$ est plat à fibres de Gorenstein,
le fait que $A_f$ vérifie $(S_2)$ implique que $A'_f$ vérifie $(S_2)$.
Références utilisées :
Complexes dualisants, section \ref{dualizing-section-formal-fibres},
Compléments d'algèbre, section \ref{more-algebra-section-properties-formal-fibres},
et Algèbre, lemme \ref{algebra-lemma-Sk-goes-up}.
Ainsi, d'après le lemme \ref{lemma-fully-faithful-henselian-completion},
on peut supposer que $A$ est un anneau local noethérien complet.

\medskip\noindent
Supposons, en plus des autres hypothèses, que $A$ soit un anneau local complet.
D'après le lemme \ref{lemma-restriction-fully-faithful}, le résultat découle de
Géométrie algébrique et formelle, lemme
\ref{algebraization-lemma-fully-faithful-simple-two}.
\end{proof}

\begin{lemma}
\label{lemma-fully-faithful-minimal}
\begin{reference}
\cite[corollaire 1.11]{Bhatt-local}
\end{reference}
Dans la situation \ref{situation-local-lefschetz}, supposons que
\begin{enumerate}
\item $H^1_\mathfrak m(A)$ et $H^2_\mathfrak m(A)$ sont
annulés par une puissance de $f$, et
\item $A$ est hensélien ou, plus généralement, $(A, (f))$ est un couple hensélien.
\end{enumerate}
Alors le foncteur de restriction
$\textit{F\'Et}_U \longrightarrow \textit{F\'Et}_{U_0}$
est pleinement fidèle.
\end{lemma}

\begin{proof}
D'après le lemme \ref{lemma-fully-faithful-henselian-completion},
on peut supposer que $A$ est un anneau local noethérien complet.
(Les hypothèses se transportent ; utiliser
Complexes dualisants, lemme \ref{dualizing-lemma-torsion-change-rings}.)
D'après le lemme \ref{lemma-restriction-fully-faithful},
le résultat découle du
lemme de Géométrie algébrique et formelle
\ref{algebraization-lemma-fully-faithful-alternative}.
\end{proof}

\begin{lemma}
\label{lemma-fully-faithful}
Dans la situation \ref{situation-local-lefschetz}, supposons que $A$ soit de profondeur $\geq 3$
et que $A$ soit hensélien ou, plus généralement, que $(A, (f))$ soit un couple hensélien. Alors
le foncteur de restriction
$\textit{F\'Et}_U \to \textit{F\'Et}_{U_0}$
est pleinement fidèle.
\end{lemma}

\begin{proof}
L'hypothèse sur la profondeur impose
$H^1_\mathfrak m(A) = H^2_\mathfrak m(A) = 0$ ; voir
Complexes dualisants, lemme \ref{dualizing-lemma-depth}.
Le lemme \ref{lemma-fully-faithful-minimal} s'applique donc.
\end{proof}



























\section{Pureté dans le cas local, I}
\label{section-local-purity}

\noindent
Soit $(A, \mathfrak m)$ un anneau local noethérien. Posons $X = \Spec(A)$
et soit $U = X \setminus \{\mathfrak m\}$ le spectre épointé.
Nous disons que {\it la pureté vaut pour $(A, \mathfrak m)$} si le foncteur de restriction
$$
\textit{F\'Et}_X \longrightarrow \textit{F\'Et}_U
$$
est essentiellement surjectif. Dans cette section, nous cherchons à comprendre comment la
question se transforme lorsque l'on passe de $X$ à une hypersurface $X_0$ de $X$ ;
autrement dit, nous étudions une sorte de propriété de Lefschetz locale pour les
groupes fondamentaux des spectres épointés.
Ces résultats seront utiles pour procéder par récurrence sur la dimension
dans les démonstrations de nos principaux résultats sur la pureté locale, à savoir le
lemme \ref{lemma-local-purity},
la proposition \ref{proposition-purity-complete-intersection} et la
proposition \ref{proposition-purity-smooth-over-depth2}.

\begin{lemma}
\label{lemma-sections-over-punctured-spec}
Soit $(A, \mathfrak m)$ un anneau local noethérien. Posons $X = \Spec(A)$
et soit $U = X \setminus \{\mathfrak m\}$.
Soit $\pi : Y \to X$ un morphisme fini tel que
$\text{depth}(\mathcal{O}_{Y, y}) \geq 2$ pour tout point fermé
$y \in Y$.
Alors $Y$ est le spectre de $B = \mathcal{O}_Y(\pi^{-1}(U))$.
\end{lemma}

\begin{proof}
Posons $V = \pi^{-1}(U)$ et notons $\pi' : V \to U$ la restriction de $\pi$.
Considérons le morphisme de $\mathcal{O}_X$-modules
$$
\pi_*\mathcal{O}_Y \longrightarrow j_*\pi'_*\mathcal{O}_V
$$
où $j : U \to X$ est le morphisme d'inclusion. Nous affirmons que le
lemme de Diviseurs \ref{divisors-lemma-depth-2-hartog}
s'applique à ce morphisme. Si tel est le cas, alors $B = \Gamma(Y, \mathcal{O}_Y)$
et le lemme en résulte. Soit $x \in X$ le point fermé.
Il suffit de montrer que
$\text{depth}((\pi_*\mathcal{O}_Y)_x) \geq 2$.
Soient $y_1, \ldots, y_n \in Y$ les points au-dessus de $x$.
D'après Algèbre, lemme \ref{algebra-lemma-depth-goes-down-finite},
il suffit de montrer que
$\text{depth}(\mathcal{O}_{Y, y_i}) \geq 2$ pour $i = 1, \ldots, n$.
Comme c'est précisément l'hypothèse du lemme, la démonstration est achevée.
\end{proof}

\begin{lemma}
\label{lemma-reformulate-purity}
Soit $(A, \mathfrak m)$ un anneau local noethérien. Posons $X = \Spec(A)$
et soit $U = X \setminus \{\mathfrak m\}$.
Soit $V$ fini étale
sur $U$. Supposons $A$ de profondeur $\geq 2$. Les assertions suivantes sont équivalentes :
\begin{enumerate}
\item $V = Y \times_X U$ pour un certain morphisme fini étale $Y \to X$,
\item $B = \Gamma(V, \mathcal{O}_V)$ est fini étale sur $A$.
\end{enumerate}
\end{lemma}

\begin{proof}
Notons $\pi : V \to U$ le morphisme fini étale donné.
Supposons qu'il existe $Y$ comme dans (1). Soit $x \in X$ le point
correspondant à $\mathfrak m$.
Soit $y \in Y$ un point au-dessus de $x$. Nous affirmons que
$\text{depth}(\mathcal{O}_{Y, y}) \geq 2$.
Cela est vrai parce que $Y \to X$ est étale et que
$A = \mathcal{O}_{X, x}$ et $\mathcal{O}_{Y, y}$ ont donc
la même profondeur (Algèbre, lemme \ref{algebra-lemma-apply-grothendieck}).
Le lemme \ref{lemma-sections-over-punctured-spec}
s'applique donc et $Y = \Spec(B)$.

\medskip\noindent
L'implication (2) $\Rightarrow$ (1) est plus facile et nous en
omettons les détails.
\end{proof}

\begin{lemma}
\label{lemma-reformulate-purity-normal}
Soit $(A, \mathfrak m)$ un anneau local noethérien. Posons $X = \Spec(A)$
et soit $U = X \setminus \{\mathfrak m\}$. Supposons $A$ normal
de dimension $\geq 2$. Le foncteur
$$
\textit{F\'Et}_U \longrightarrow
\left\{
\begin{matrix}
A\text{-algèbres finies normales }B\text{ telles que} \\
\Spec(B) \to X\text{ soit étale au-dessus de }U
\end{matrix}
\right\},
\quad
V \longmapsto \Gamma(V, \mathcal{O}_V)
$$
est une équivalence. De plus, $V = Y \times_X U$ pour un certain morphisme $Y \to X$
fini étale si et seulement si $B = \Gamma(V, \mathcal{O}_V)$
est fini étale sur $A$.
\end{lemma}

\begin{proof}
Remarquons que $\text{depth}(A) \geq 2$ puisque $A$ est normal
(critère de normalité de Serre, Algèbre, lemme
\ref{algebra-lemma-criterion-normal}).
La dernière assertion découle donc du lemme \ref{lemma-reformulate-purity}.
Étant donné $\pi : V \to U$ fini étale, posons $B = \Gamma(V, \mathcal{O}_V)$.
Si nous montrons que $B$ est normal et fini sur $A$, nous
obtenons le foncteur affiché. Comme il existe un foncteur quasi-inverse
évident, c'est aussi tout ce qu'il reste à montrer.

\medskip\noindent
Comme $A$ est normal, le schéma $V$ est normal
(Descente, lemme \ref{descent-lemma-normal-local-smooth}).
Par conséquent, $V$ est une réunion disjointe finie de schémas intègres
(Propriétés, lemme \ref{properties-lemma-normal-Noetherian}).
On peut donc supposer $V$ intègre.
Dans ce cas, le corps des fonctions $L$ de $V$
(Morphismes, section \ref{morphisms-section-rational-maps})
est une extension finie séparable du corps des fractions de $A$
(car on l'obtient en considérant la fibre générique
de $V \to U$ et en utilisant Morphismes, lemme
\ref{morphisms-lemma-etale-over-field}).
D'après Algèbre, lemme
\ref{algebra-lemma-Noetherian-normal-domain-finite-separable-extension},
la clôture intégrale $B' \subset L$ de $A$ dans $L$ est finie sur $A$.
D'après Compléments d'algèbre, lemme \ref{more-algebra-lemma-integral-closure-reflexive},
$B'$ est un $A$-module réflexif, ce qui implique à son tour
que $\text{depth}_A(B') \geq 2$ d'après
Compléments d'algèbre, lemme \ref{more-algebra-lemma-reflexive-over-normal}.

\medskip\noindent
Soit $f \in \mathfrak m$. Alors $B_f = \Gamma(V \times_U D(f), \mathcal{O}_V)$
(Propriétés, lemme \ref{properties-lemma-invert-f-sections}).
On a donc $B'_f = B_f$, puisque $B_f$ est normal (voir ci-dessus),
fini sur $A_f$ et de corps des fractions $L$.
Il s'ensuit que $V = \Spec(B') \times_X U$.
Nous concluons alors que $B = B'$ grâce au
lemme \ref{lemma-sections-over-punctured-spec}
appliqué à $\Spec(B') \to X$.
Ce lemme s'applique parce que les localisés $B'_{\mathfrak m'}$
de $B'$ aux idéaux maximaux $\mathfrak m' \subset B'$ situés au-dessus de
$\mathfrak m$ sont de profondeur $\geq 2$ d'après
Algèbre, lemme \ref{algebra-lemma-depth-goes-down-finite},
et la remarque sur la profondeur du paragraphe précédent.
\end{proof}

\begin{lemma}
\label{lemma-purity-and-completion}
Soit $(A, \mathfrak m)$ un anneau local noethérien. Posons $X = \Spec(A)$
et soit $U = X \setminus \{\mathfrak m\}$.
Soit $V$ fini étale sur $U$.
Soit $A^\wedge$ le complété $\mathfrak m$-adique de $A$,
soit $X' = \Spec(A^\wedge)$, et soient $U'$ et $V'$ les changements de base de
$U$ et $V$ à $X'$. Les assertions suivantes sont équivalentes :
\begin{enumerate}
\item $V = Y \times_X U$ pour un certain morphisme fini étale $Y \to X$, et
\item $V' = Y' \times_{X'} U'$ pour un certain morphisme fini étale $Y' \to X'$.
\end{enumerate}
\end{lemma}

\begin{proof}
L'implication (1) $\Rightarrow$ (2) s'obtient en effectuant le changement de base
d'une solution $Y \to X$. Soit $Y' \to X'$ comme dans (2).
D'après le lemme \ref{lemma-fill-in-missing}, il existe un morphisme fini $Y \to X$
tel que $V = Y \times_X U$ et $Y' = Y \times_X X'$.
Par descente, le morphisme $Y \to X$ est fini étale
(Algèbre, lemmes \ref{algebra-lemma-descend-properties-modules} et
\ref{algebra-lemma-etale}). Cela achève la démonstration.
\end{proof}

\noindent
L'intérêt des deux lemmes suivants est que leurs hypothèses n'imposent pas à
$A$ d'être de profondeur $\geq 3$. Par exemple, si $A$ est un anneau local complet
normal et intègre de dimension $\geq 3$ et si $f \in \mathfrak m$ est non nul,
alors les hypothèses sont satisfaites.

\begin{lemma}
\label{lemma-lift-simple}
Dans la situation \ref{situation-local-lefschetz}, soit $V$ fini
étale sur $U$. Supposons que
\begin{enumerate}
\item[(a)] $A$ possède un complexe dualisant,
\item[(b)] le couple $(A, (f))$ soit hensélien,
\item[(c)] l'une des conditions suivantes soit satisfaite :
\begin{enumerate}
\item[(i)] $A_f$ vérifie $(S_2)$ et toute composante irréductible de $X$
non contenue dans $X_0$ est de dimension $\geq 3$, ou bien
\item[(ii)] pour tout idéal premier $\mathfrak p \subset A$, $f \not \in \mathfrak p$
on ait $\text{depth}(A_\mathfrak p) + \dim(A/\mathfrak p) > 2$.
\end{enumerate}
\item[(d)] $V_0 = V \times_U U_0$ soit égal à $Y_0 \times_{X_0} U_0$
pour un certain morphisme fini étale $Y_0 \to X_0$.
\end{enumerate}
Alors $V = Y \times_X U$ pour un certain morphisme fini étale $Y \to X$.
\end{lemma}

\begin{proof}
On se ramène au cas complet à l'aide du lemme \ref{lemma-purity-and-completion}.
(Les hypothèses se transportent ; voir la démonstration du
lemme \ref{lemma-fully-faithful-simple}.)

\medskip\noindent
Dans le cas complet, on peut relever $Y_0 \to X_0$ en un morphisme
fini étale $Y \to X$ d'après le
lemme de Compléments d'algèbre \ref{more-algebra-lemma-finite-etale-equivalence} ;
remarquons que $(A, fA)$ est un couple hensélien d'après le
lemme de Compléments d'algèbre \ref{more-algebra-lemma-complete-henselian}.
On peut alors utiliser le lemme \ref{lemma-fully-faithful-simple}
pour voir que $V$ est isomorphe à $Y \times_X U$ et
la démonstration est achevée.
\end{proof}

\begin{lemma}
\label{lemma-lift-purity-general}
Dans la situation \ref{situation-local-lefschetz},
soit $V$ fini étale sur $U$. Supposons que
\begin{enumerate}
\item $H^1_\mathfrak m(A)$ et $H^2_\mathfrak m(A)$
soient annulés par une puissance de $f$,
\item $V_0 = V \times_U U_0$ soit égal à $Y_0 \times_{X_0} U_0$
pour un certain morphisme fini étale $Y_0 \to X_0$.
\end{enumerate}
Alors $V = Y \times_X U$ pour un certain morphisme fini étale $Y \to X$.
\end{lemma}

\begin{proof}
On se ramène au cas complet à l'aide du lemme \ref{lemma-purity-and-completion}.
(Les hypothèses se transportent ; utiliser Complexes dualisants, lemme
\ref{dualizing-lemma-torsion-change-rings}.)

\medskip\noindent
Dans le cas complet, on peut relever $Y_0 \to X_0$ en un morphisme
fini étale $Y \to X$ d'après le
lemme de Compléments d'algèbre \ref{more-algebra-lemma-finite-etale-equivalence} ;
remarquons que $(A, fA)$ est un couple hensélien d'après le
lemme de Compléments d'algèbre \ref{more-algebra-lemma-complete-henselian}.
On peut alors utiliser le lemme \ref{lemma-fully-faithful-minimal}
pour voir que $V$ est isomorphe à $Y \times_X U$ et
la démonstration est achevée.
\end{proof}

\begin{lemma}
\label{lemma-lift-purity}
Dans la situation \ref{situation-local-lefschetz},
soit $V$ fini étale sur $U$. Supposons que
\begin{enumerate}
\item $A$ soit de profondeur $\geq 3$,
\item $V_0 = V \times_U U_0$ soit égal à $Y_0 \times_{X_0} U_0$
pour un certain morphisme fini étale $Y_0 \to X_0$.
\end{enumerate}
Alors $V = Y \times_X U$ pour un certain morphisme fini étale $Y \to X$.
\end{lemma}

\begin{proof}
L'hypothèse sur la profondeur impose
$H^1_\mathfrak m(A) = H^2_\mathfrak m(A) = 0$ ; voir
Complexes dualisants, lemme \ref{dualizing-lemma-depth}.
Le lemme \ref{lemma-lift-purity-general} s'applique donc.
\end{proof}










\section{Pureté du lieu de branchement}
\label{section-purity}

\noindent
Nous utiliserons le discriminant d'un morphisme fini localement libre. Voir
Discriminants, section \ref{discriminant-section-discriminant}.

\begin{lemma}
\label{lemma-find-point-codim-1}
Soit $(A, \mathfrak m)$ un anneau local noethérien tel que $\dim(A) \geq 1$.
Soit $f \in \mathfrak m$. Alors il existe $\mathfrak p \in V(f)$ tel que
$\dim(A_\mathfrak p) = 1$.
\end{lemma}

\begin{proof}
Raisonnons par récurrence sur $\dim(A)$. Si $\dim(A) = 1$, alors $\mathfrak p = \mathfrak m$
convient. Si $\dim(A) > 1$, soit $Z \subset \Spec(A)$ une composante irréductible
de dimension $> 1$. Alors $V(f) \cap Z$ est de dimension $> 0$
(Algèbre, lemme \ref{algebra-lemma-one-equation}). Choisissons un idéal premier
$\mathfrak q \in V(f) \cap Z$, $\mathfrak q \not = \mathfrak m$,
correspondant à un point fermé du spectre épointé de $A$ ;
c'est possible d'après
Propriétés, lemme \ref{properties-lemma-complement-closed-point-Jacobson}.
Alors $\mathfrak q$ n'est pas le point générique de $Z$. Par conséquent,
$0 < \dim(A_\mathfrak q) < \dim(A)$ et $f \in \mathfrak q A_\mathfrak q$.
Par récurrence sur la dimension, on peut trouver
$f \in \mathfrak p \subset A_\mathfrak q$ tel que
$\dim((A_\mathfrak q)_\mathfrak p) = 1$.
Alors $\mathfrak p \cap A$ convient.
\end{proof}

\begin{lemma}
\label{lemma-ramification-quasi-finite-flat}
Soit $f : X \to Y$ un morphisme de schémas localement noethériens.
Soit $x \in X$. Supposons que
\begin{enumerate}
\item $f$ soit plat,
\item $f$ soit quasi-fini en $x$,
\item $x$ ne soit pas le point générique d'une composante irréductible de $X$,
\item pour les spécialisations $x' \leadsto x$ telles que
$\dim(\mathcal{O}_{X, x'}) = 1$, le morphisme $f$ soit non ramifié en $x'$.
\end{enumerate}
Alors $f$ est étale en $x$.
\end{lemma}

\begin{proof}
Remarquons que l'ensemble des points où $f$ est non ramifié coïncide avec
l'ensemble des points où $f$ est étale, et que cet ensemble est ouvert.
Voir Morphismes, définitions \ref{morphisms-definition-unramified}
et \ref{morphisms-definition-etale}, ainsi que le
lemme \ref{morphisms-lemma-flat-unramified-etale}.
Pour vérifier que $f$ est étale en $x$, on peut travailler localement pour la topologie étale
sur la base et sur la
source (Descente, lemmes \ref{descent-lemma-descending-property-etale} et
\ref{descent-lemma-etale-etale-local-source}).
On peut donc appliquer Compléments sur les morphismes, lemme
\ref{more-morphisms-lemma-etale-makes-quasi-finite-finite-at-point},
et supposer que $f : X \to Y$ est fini et que $x$ est l'unique
point de $X$ situé au-dessus de $y = f(x)$.
Il s'ensuit que $f$ est fini localement libre
(Morphismes, lemme \ref{morphisms-lemma-finite-flat}).

\medskip\noindent
Supposons $f$ fini localement libre et $x$ l'unique point de
$X$ situé au-dessus de $y = f(x)$. D'après
Discriminants, lemme \ref{discriminant-lemma-discriminant},
il existe un sous-schéma fermé localement principal $D_f \subset Y$
tel que $y' \in D_f$ si et seulement s'il existe $x' \in X$
avec $f(x') = y'$ et $f$ ramifié en $x'$. Il faut donc montrer
que $y \not \in D_f$. Supposons $y \in D_f$ et cherchons une contradiction.

\medskip\noindent
D'après la condition (3), on a $\dim(\mathcal{O}_{X, x}) \geq 1$.
On a $\dim(\mathcal{O}_{X, x}) = \dim(\mathcal{O}_{Y, y})$ d'après
Algèbre, lemme \ref{algebra-lemma-dimension-base-fibre-equals-total}.
D'après le lemme \ref{lemma-find-point-codim-1},
on peut trouver $y' \in D_f$ se spécialisant en $y$
tel que $\dim(\mathcal{O}_{Y, y'}) = 1$.
Choisissons $x' \in X$ tel que $f(x') = y'$ et où $f$ est ramifié. Comme $f$
est fini, il est fermé, et par conséquent $x' \leadsto x$.
On a $\dim(\mathcal{O}_{X, x'}) = \dim(\mathcal{O}_{Y, y'}) = 1$
comme précédemment. Cela contredit la propriété (4).
\end{proof}

\begin{lemma}
\label{lemma-local-purity}
Soit $(A, \mathfrak m)$ un anneau local régulier de dimension $d \geq 2$.
Posons $X = \Spec(A)$ et $U = X \setminus \{\mathfrak m\}$. Alors
\begin{enumerate}
\item le foncteur $\textit{F\'Et}_X \to \textit{F\'Et}_U$
est essentiellement surjectif, c'est-à-dire que la pureté vaut pour $A$,
\item tout morphisme fini $A \to B$ avec $B$ normal qui
induit un morphisme fini étale sur les spectres épointés est étale.
\end{enumerate}
\end{lemma}

\begin{proof}
Rappelons qu'un anneau local régulier est normal d'après
Algèbre, lemme \ref{algebra-lemma-regular-normal}.
Par conséquent, (1) et (2) sont équivalents d'après le
lemme \ref{lemma-reformulate-purity-normal}.
Démontrons le lemme par récurrence sur $d$.

\medskip\noindent
Le cas $d = 2$. Dans ce cas, $A \to B$ est plat.
En effet, la descente des idéaux premiers vaut pour $A \to B$ d'après
Algèbre, proposition \ref{algebra-proposition-going-down-normal-integral}.
Alors $\dim(B_{\mathfrak m'}) = 2$ pour tout idéal maximal
$\mathfrak m' \subset B$ d'après
Algèbre, lemme \ref{algebra-lemma-dimension-base-fibre-equals-total}.
Puis $B_{\mathfrak m'}$ est un anneau de Cohen--Macaulay d'après
Algèbre, lemme \ref{algebra-lemma-criterion-normal}.
Par conséquent, et c'est l'étape importante, le
lemme d'Algèbre \ref{algebra-lemma-CM-over-regular-flat}
s'applique et montre que $A \to B_{\mathfrak m'}$ est plat.
Le lemme d'Algèbre \ref{algebra-lemma-flat-localization}
montre alors que $A \to B$ est plat. On peut donc appliquer le
lemme \ref{lemma-ramification-quasi-finite-flat}
(ou raisonner directement à l'aide du plus simple
lemme de Discriminants \ref{discriminant-lemma-discriminant})
pour voir que $A \to B$ est étale.

\medskip\noindent
Le cas $d \geq 3$. Soit $V \to U$ fini étale.
Soit $f \in \mathfrak m_A$, $f \not \in \mathfrak m_A^2$.
Alors $A/fA$ est un anneau local régulier de dimension $d - 1 \geq 2$ ; voir
Algèbre, lemme \ref{algebra-lemma-regular-ring-CM}.
Soit $U_0$ le spectre épointé de $A/fA$ et posons
$V_0 = V \times_U U_0$.
D'après le lemme \ref{lemma-lift-purity},
il suffit de montrer que $V_0$ appartient à l'image essentielle
de $\textit{F\'Et}_{\Spec(A/fA)} \to \textit{F\'Et}_{U_0}$.
Cela résulte de l'hypothèse de récurrence.
\end{proof}

\begin{lemma}[Pureté du lieu de branchement]
\label{lemma-purity}
\begin{reference}
\cite{Nagata-Purity} et \cite[exposé X, théorème 3.1]{SGA1}
\end{reference}
\begin{history}
Ce résultat fut d'abord énoncé et démontré par Zariski sous
forme géométrique dans \cite{Zariski-Purity}.
La généralisation aux corps de base non parfaits, due à Nagata,
fut publiée comme article suivant dans le même volume des
Comptes rendus de l'Académie nationale des sciences des États-Unis d'Amérique,
dans \cite{Nagata-Remarks-Purity}. L'année suivante, Nagata
démontra le résultat pour les anneaux locaux noethériens dans \cite{Nagata-Purity}.
Sa démonstration utilise un résultat de Chow qui est un théorème de Bertini pour les
anneaux locaux complets et intègres ; voir \cite{Chow-Bertini} ;
l'histoire des théorèmes de Bertini est étudiée dans
l'article historique de Kleiman \cite{Kleiman-Bertini}.
Quelques années plus tard, Auslander trouva une démonstration entièrement différente ;
voir \cite{Auslander-Purity}.
\end{history}
Soit $f : X \to Y$ un morphisme de schémas localement noethériens.
Soit $x \in X$ et posons $y = f(x)$. Supposons que
\begin{enumerate}
\item $\mathcal{O}_{X, x}$ soit normal,
\item $\mathcal{O}_{Y, y}$ soit régulier,
\item $f$ soit quasi-fini en $x$,
\item $\dim(\mathcal{O}_{X, x}) = \dim(\mathcal{O}_{Y, y}) \geq 1$
\item pour les spécialisations $x' \leadsto x$ telles que
$\dim(\mathcal{O}_{X, x'}) = 1$, le morphisme $f$ soit non ramifié en $x'$.
\end{enumerate}
Alors $f$ est étale en $x$.
\end{lemma}

\begin{proof}
Nous allons démontrer le lemme par récurrence sur
$d = \dim(\mathcal{O}_{X, x}) = \dim(\mathcal{O}_{Y, y})$.

\medskip\noindent
Un cas sans intérêt est celui où $d = 1$.
Dans ce cas, nous supposons que $f$ est non ramifié en $x$
et que $\mathcal{O}_{Y, y}$ est un anneau de valuation discrète
(Algèbre, lemme \ref{algebra-lemma-characterize-dvr}).
Alors $\mathcal{O}_{X, x}$ est plat sur $\mathcal{O}_{Y, y}$
(sinon le morphisme ne serait pas quasi-fini en $x$),
et l'on voit que $f$ est plat en $x$. Comme plat $+$
non ramifié implique étale, on conclut (certains détails sont omis).

\medskip\noindent
Le cas $d \geq 2$. On va utiliser la récurrence sur $d$ pour se ramener
au cas étudié dans le lemme \ref{lemma-local-purity}.
Pour vérifier que $f$ est étale en $x$, on peut travailler localement pour la topologie étale
sur la base et sur la source
(Descente, lemmes \ref{descent-lemma-descending-property-etale} et
\ref{descent-lemma-etale-etale-local-source}).
On peut donc appliquer Compléments sur les morphismes, lemme
\ref{more-morphisms-lemma-etale-makes-quasi-finite-finite-at-point},
et supposer que $f : X \to Y$ est fini et que $x$ est l'unique
point de $X$ situé au-dessus de $y$. On utilise ici le fait que les extensions étales entre
anneaux locaux ne changent ni la dimension, ni la normalité, ni la régularité ; voir
Compléments d'algèbre, section \ref{more-algebra-section-permanence-etale},
et
Morphismes étales, section \ref{etale-section-properties-permanence}.

\medskip\noindent
Ensuite, on peut effectuer le changement de base par $\Spec(\mathcal{O}_{Y, y})$
et supposer que $Y$ est le spectre d'un anneau local régulier.
Il s'ensuit que $X = \Spec(\mathcal{O}_{X, x})$, puisque
tout point de $X$ se spécialise nécessairement en $x$.

\medskip\noindent
Le morphisme d'anneaux $\mathcal{O}_{Y, y} \to \mathcal{O}_{X, x}$ est
fini et nécessairement injectif (par égalité des dimensions).
On conclut que la descente des idéaux premiers vaut pour
$\mathcal{O}_{Y, y} \to \mathcal{O}_{X, x}$ d'après
Algèbre, proposition \ref{algebra-proposition-going-down-normal-integral}
(et le fait qu'un anneau régulier est normal d'après
Algèbre, lemme \ref{algebra-lemma-regular-normal}).
Choisissons $x' \in X$, $x' \not = x$, d'image $y' = f(x')$.
Alors $\mathcal{O}_{X, x'}$ est normal comme localisé
d'un anneau normal intègre. De même, $\mathcal{O}_{Y, y'}$ est
régulier (voir Algèbre, lemme
\ref{algebra-lemma-localization-of-regular-local-is-regular}).
On a $\dim(\mathcal{O}_{X, x'}) = \dim(\mathcal{O}_{Y, y'})$ d'après
Algèbre, lemme \ref{algebra-lemma-dimension-base-fibre-equals-total}
(on a vérifié ci-dessus la descente des idéaux premiers).
Bien entendu, ces dimensions sont strictement inférieures à $d$, puisque $x' \not = x$,
et l'hypothèse de récurrence sur $d$ montre que $f$ est étale en $x'$.

\medskip\noindent
On arrive donc à la situation suivante : on a un morphisme
local fini $A \to B$ d'anneaux locaux noethériens
de dimension $d \geq 2$, avec $A$ régulier et $B$ normal, qui
induit un morphisme fini étale $V \to U$ sur les spectres épointés.
Notre but est de montrer que $A \to B$ est étale.
Cela résulte du lemme \ref{lemma-local-purity},
et la démonstration est achevée.
\end{proof}

\noindent
Le lemme suivant est parfois utile pour déterminer le plus grand
ouvert sur lequel un morphisme fini étale se prolonge.

\begin{lemma}
\label{lemma-extend-S2}
Soit $j : U \to X$ une immersion ouverte de schémas localement noethériens
telle que $\text{depth}(\mathcal{O}_{X, x}) \geq 2$ pour $x \not \in U$.
Soit $\pi : V \to U$ fini étale. Alors
\begin{enumerate}
\item $\mathcal{B} = j_*\pi_*\mathcal{O}_V$ est une algèbre cohérente
réflexive sur $\mathcal{O}_X$ ; posons $Y = \underline{\Spec}_X(\mathcal{B})$,
\item $Y \to X$ est l'unique morphisme fini tel que
$V = Y \times_X U$ et $\text{depth}(\mathcal{O}_{Y, y}) \geq 2$
pour $y \in Y \setminus V$,
\item $Y \to X$ est étale en $y$ si et seulement si $Y \to X$ est plat en $y$, et
\item $Y \to X$ est étale si et seulement si $\mathcal{B}$
est fini localement libre comme $\mathcal{O}_X$-module.
\end{enumerate}
De plus, (a) la construction de $\mathcal{B}$ et de $Y \to X$ commute
au changement de base par les morphismes plats $X' \to X$ de schémas
localement noethériens, et (b), si $V' \to U'$ est un morphisme fini étale avec
$U \subset U' \subset X$ ouvert qui se restreint à $V \to U$ au-dessus de $U$,
alors il existe un unique isomorphisme $Y \times_X U' = V'$ au-dessus de $U'$.
\end{lemma}

\begin{proof}
Remarquons que $\pi_*\mathcal{O}_V$ est un module fini localement libre
sur $\mathcal{O}_U$, donc en particulier réflexif.
D'après Diviseurs, lemme \ref{divisors-lemma-reflexive-S2-extend},
le module $j_*\pi_*\mathcal{O}_V$ est l'unique
module cohérent réflexif sur $X$ dont la restriction à $U$
est $\pi_*\mathcal{O}_V$. Cela démontre (1).

\medskip\noindent
Par construction, $Y \times_X U = V$.
Comme $\mathcal{B}$ est cohérente, on voit que $Y \to X$ est fini.
On a $\text{depth}(\mathcal{B}_x) \geq 2$ pour $x \in X \setminus U$
d'après Diviseurs, lemme \ref{divisors-lemma-reflexive-S2}.
Par conséquent, $\text{depth}(\mathcal{O}_{Y, y}) \geq 2$ pour $y \in Y \setminus V$
d'après Algèbre, lemme \ref{algebra-lemma-depth-goes-down-finite}.
Réciproquement, supposons que $\pi' : Y' \to X$ soit un morphisme fini tel que
$V = Y' \times_X U$ et $\text{depth}(\mathcal{O}_{Y', y'}) \geq 2$
pour $y' \in Y' \setminus V$. Alors $\pi'_*\mathcal{O}_{Y'}$
a pour restriction $\pi_*\mathcal{O}_V$ sur $U$ et vérifie
$\text{depth}((\pi'_*\mathcal{O}_{Y'})_x) \geq 2$ pour
$x \in X \setminus U$ d'après le
lemme d'Algèbre \ref{algebra-lemma-depth-goes-down-finite}.
Alors $\pi'_*\mathcal{O}_{Y'}$ est canoniquement isomorphe
à $j_*\pi_*\mathcal{O}_V$, par exemple d'après
Diviseurs, lemme \ref{divisors-lemma-depth-2-hartog}.
Cela démontre (2).

\medskip\noindent
Si $Y \to X$ est étale en $y$, alors $Y \to X$ est plat en $y$.
Réciproquement, supposons que $Y \to X$ soit plat en $y$.
Si $y \in V$, alors $Y \to X$ est étale en $y$.
Si $y \not \in V$, vérifions que (1), (2), (3) et (4) du
lemme \ref{lemma-ramification-quasi-finite-flat} sont satisfaits
pour voir que $Y \to X$ est étale en $y$. Les parties (1) et (2)
sont claires, de même que (3), puisque $\text{depth}(\mathcal{O}_{Y, y}) \geq 2$.
Si $y' \leadsto y$ est une spécialisation et $\dim(\mathcal{O}_{Y, y'}) = 1$,
alors $y' \in V$, car sinon la profondeur de cet anneau local
serait au moins égale à $2$, ce qui contredirait le
lemme d'Algèbre \ref{algebra-lemma-bound-depth}.
Ainsi, $Y \to X$ est étale en $y'$, et l'on conclut que (4) du
lemme \ref{lemma-ramification-quasi-finite-flat} est également satisfait.
Cela achève la démonstration de (3).

\medskip\noindent
La partie (4) résulte de (3) et du fait que $((Y \to X)_*\mathcal{O}_Y)_x$
est un $\mathcal{O}_{X, x}$-module plat si et seulement si $\mathcal{O}_{Y, y}$
est un $\mathcal{O}_{X, x}$-module plat pour tout $y \in Y$ au-dessus de $x$ ; voir
Algèbre, lemme \ref{algebra-lemma-flat-localization}. On utilise aussi ici
le fait qu'un module fini plat sur un anneau noethérien est fini localement
libre ; voir Algèbre, lemme \ref{algebra-lemma-finite-projective}
(et
Algèbre, lemme
\ref{algebra-lemma-Noetherian-finite-type-is-finite-presentation}).

\medskip\noindent
Quant aux dernières assertions du lemme, la partie (a) résulte du
changement de base plat ; voir
Cohomologie des schémas, lemme \ref{coherent-lemma-flat-base-change-cohomology},
et la partie (b) découle de l'unicité dans (2), appliquée à la restriction
$Y \times_X U'$.
\end{proof}

\begin{lemma}
\label{lemma-extend-pure}
Soit $j : U \to X$ une immersion ouverte de schémas noethériens
telle que la pureté vaille pour $\mathcal{O}_{X, x}$ pour tout $x \not \in U$.
Alors
$$
\textit{F\'Et}_X \longrightarrow \textit{F\'Et}_U
$$
est essentiellement surjectif.
\end{lemma}

\begin{proof}
Soit $V \to U$ un morphisme fini étale. Par récurrence noethérienne,
il suffit de prolonger $V \to U$ en un morphisme fini étale
au-dessus d'un ouvert strictement plus grand de $X$.
Soit $x \in X \setminus U$ le point générique d'une
composante irréductible de $X \setminus U$.
Alors l'image inverse $U_x$ de $U$ dans $\Spec(\mathcal{O}_{X, x})$
est le spectre épointé de $\mathcal{O}_{X, x}$.
Par hypothèse, $V_x = V \times_U U_x$ est la restriction
d'un morphisme fini étale $Y_x \to \Spec(\mathcal{O}_{X, x})$
à $U_x$.
D'après Limites, lemme \ref{limits-lemma-glueing-near-point},
il existe un sous-schéma ouvert $U \subset U' \subset X$
contenant $x$ et un morphisme de présentation finie $V' \to U'$
dont la restriction à $U$ redonne $V \to U$ et
dont la restriction à $\Spec(\mathcal{O}_{X, x})$ redonne $Y_x$.
Enfin, le morphisme $V' \to U'$ est fini étale
après avoir éventuellement remplacé $U'$ par un ouvert plus petit, d'après
Limites, lemme \ref{limits-lemma-glueing-near-point-properties}.
\end{proof}
















\section{Revêtements finis étales des spectres épointés, II}
\label{section-pi1-punctured-spec-II}

\noindent
Dans cette section, nous démontrons quelques variantes des résultats étudiés
dans la section \ref{section-pi1-punctured-spec}. Supposons
donné un anneau local noethérien $(A, \mathfrak m)$ et un élément $f \in \mathfrak m$.
Posons $X = \Spec(A)$ et $X_0 = \Spec(A/fA)$, et
posons $U = X \setminus \{\mathfrak m\}$ et
$U_0 = X_0 \setminus \{\mathfrak m\}$, les spectres épointés de
$A$ et $A/fA$. Tout ceci est exactement comme dans la
situation \ref{situation-local-lefschetz}.
La différence est que nous considérerons le foncteur de restriction
$$
\colim_{U_0 \subset U' \subset U\text{ ouvert}} \textit{F\'Et}_{U'}
\longrightarrow
\textit{F\'Et}_{U_0}
$$
Autrement dit, nous ne chercherons pas à relever les revêtements finis étales
de $U_0$ à tout $U$, mais seulement à un voisinage ouvert
$U'$ de $U_0$ dans $U$.

\begin{lemma}
\label{lemma-faithful-general}
Dans la situation \ref{situation-local-lefschetz}, soit $U' \subset U$
un ouvert contenant $U_0$. Supposons que, pour tout idéal premier minimal $\mathfrak p \subset A$
tel que $\mathfrak p \in U'$ et $\mathfrak p \not \in U_0$, on ait
$\dim(A/\mathfrak p) \geq 2$. Alors
$$
\textit{F\'Et}_{U'} \longrightarrow \textit{F\'Et}_{U_0},\quad
V' \longmapsto V_0 = V' \times_{U'} U_0
$$
est un foncteur fidèle. De plus, il existe un $U'$ satisfaisant
l'hypothèse, et tout ouvert plus petit $U'' \subset U'$ contenant
$U_0$ satisfait également cette hypothèse. En particulier, le
foncteur de restriction
$$
\colim_{U_0 \subset U' \subset U\text{ ouvert}} \textit{F\'Et}_{U'}
\longrightarrow
\textit{F\'Et}_{U_0}
$$
est fidèle.
\end{lemma}

\begin{proof}
D'après Algèbre, lemme \ref{algebra-lemma-one-equation},
on voit que $V(\mathfrak p)$ rencontre $U_0$ pour
tout idéal premier $\mathfrak p$ de $A$ tel que $\dim(A/\mathfrak p) \geq 2$.
Ainsi, le foncteur affiché est fidèle pour un $U'$ comme dans l'énoncé,
d'après le lemme \ref{lemma-restriction-faithful}.
Pour voir qu'un tel $U'$ existe, remarquons que, pour
$\mathfrak p \subset A$ tel que $\mathfrak p \in U$,
$\mathfrak p \not \in U_0$ et $\dim(A/\mathfrak p) = 1$,
alors $\mathfrak p$ correspond à un point fermé de $U$
et donc $V(\mathfrak p) \cap U_0 = \emptyset$.
On peut ainsi prendre pour $U'$ le complémentaire des composantes
irréductibles de $X$ qui ne rencontrent pas $U_0$ et sont de dimension $1$.
\end{proof}

\begin{lemma}
\label{lemma-fully-faithful-general-better}
Dans la situation \ref{situation-local-lefschetz}, supposons que
\begin{enumerate}
\item $A$ possède un complexe dualisant et soit complet pour la topologie $f$-adique,
\item toute composante irréductible de $X$ non contenue dans $X_0$
soit de dimension $\geq 3$.
\end{enumerate}
Alors le foncteur de restriction
$$
\colim_{U_0 \subset U' \subset U\text{ ouvert}} \textit{F\'Et}_{U'}
\longrightarrow
\textit{F\'Et}_{U_0}
$$
est pleinement fidèle.
\end{lemma}

\begin{proof}
Pour le démontrer, on peut remplacer $A$ par son réduit, grâce à l'invariance topologique
du groupe fondamental ; voir le lemme \ref{lemma-thickening}.
Le résultat découle alors du
lemme \ref{lemma-restriction-fully-faithful-general}
et du lemme de Géométrie algébrique et formelle
\ref{algebraization-lemma-fully-faithful-general}.
\end{proof}

\begin{lemma}
\label{lemma-fully-faithful-general}
Dans la situation \ref{situation-local-lefschetz}, supposons que
\begin{enumerate}
\item $A$ soit complet pour la topologie $f$-adique,
\item $f$ soit non diviseur de zéro,
\item $H^1_\mathfrak m(A/fA)$ soit un $A$-module fini.
\end{enumerate}
Alors le foncteur de restriction
$$
\colim_{U_0 \subset U' \subset U\text{ ouvert}} \textit{F\'Et}_{U'}
\longrightarrow
\textit{F\'Et}_{U_0}
$$
est pleinement fidèle.
\end{lemma}

\begin{proof}
Cela résulte du
lemme \ref{lemma-restriction-fully-faithful-general} et du
lemme de Géométrie algébrique et formelle
\ref{algebraization-lemma-fully-faithful-general-alternative}.
\end{proof}








\section{Revêtements finis étales des spectres épointés, III}
\label{section-pi1-punctured-spec-III}

\noindent
Dans cette section, nous étudions, dans la situation \ref{situation-local-lefschetz}, quand
le foncteur de restriction
$$
\colim_{U_0 \subset U' \subset U\text{ ouvert}} \textit{F\'Et}_{U'}
\longrightarrow
\textit{F\'Et}_{U_0}
$$
est une équivalence de catégories.

\begin{lemma}
\label{lemma-essentially-surjective-general-better}
Dans la situation \ref{situation-local-lefschetz}, supposons que
\begin{enumerate}
\item $A$ possède un complexe dualisant et soit complet pour la topologie $f$-adique,
\item l'une des conditions suivantes soit satisfaite :
\begin{enumerate}
\item $A_f$ vérifie $(S_2)$ et toute composante irréductible de $X$
non contenue dans $X_0$ soit de dimension $\geq 4$, ou bien
\item si $\mathfrak p \not \in V(f)$ et
$V(\mathfrak p) \cap V(f) \not = \{\mathfrak m\}$, alors
$\text{depth}(A_\mathfrak p) + \dim(A/\mathfrak p) > 3$.
\end{enumerate}
\end{enumerate}
Alors le foncteur de restriction
$$
\colim_{U_0 \subset U' \subset U\text{ ouvert}} \textit{F\'Et}_{U'}
\longrightarrow
\textit{F\'Et}_{U_0}
$$
est une équivalence.
\end{lemma}

\begin{proof}
Cela résulte du lemme \ref{lemma-restriction-equivalence-general}
et du
lemme de Géométrie algébrique et formelle
\ref{algebraization-lemma-equivalence-better}.
\end{proof}

\begin{lemma}
\label{lemma-essentially-surjective-general}
Dans la situation \ref{situation-local-lefschetz}, supposons que
\begin{enumerate}
\item $A$ soit complet pour la topologie $f$-adique,
\item $f$ soit non diviseur de zéro,
\item $H^1_\mathfrak m(A/fA)$ et $H^2_\mathfrak m(A/fA)$
soient des $A$-modules finis.
\end{enumerate}
Alors le foncteur de restriction
$$
\colim_{U_0 \subset U' \subset U\text{ ouvert}} \textit{F\'Et}_{U'}
\longrightarrow
\textit{F\'Et}_{U_0}
$$
est une équivalence.
\end{lemma}

\begin{proof}
Cela résulte du lemme \ref{lemma-restriction-equivalence-general}
et du
lemme de Géométrie algébrique et formelle
\ref{algebraization-lemma-equivalence}.
\end{proof}

\begin{remark}
\label{remark-combine}
Soient $(A, \mathfrak m)$ un anneau local noethérien complet et
$f \in \mathfrak m$ non nul. Supposons que $A_f$ vérifie $(S_2)$ et que
toute composante irréductible de $\Spec(A)$ soit de dimension $\geq 4$.
Alors le lemme \ref{lemma-essentially-surjective-general-better}
nous dit que la catégorie
$$
\colim\nolimits_{U' \subset U\text{ ouvert, }U_0 \subset U'}
\text{ catégorie des schémas finis étales sur }U'
$$
est équivalente à la catégorie des schémas finis étales sur $U_0$.
Cela vaut par exemple si $A$ est un anneau normal intègre de dimension $\geq 4$ !
\end{remark}




\section{Revêtements finis étales des spectres épointés, IV}
\label{section-pi1-punctured-spec-IV}

\noindent
Soient $X, X_0, U, U_0$ comme dans la situation \ref{situation-local-lefschetz}.
Dans cette section, nous demandons quand le foncteur de restriction
$$
\textit{F\'Et}_U
\longrightarrow
\textit{F\'Et}_{U_0}
$$
est essentiellement surjectif.
Nous procéderons en reprenant les résultats de la
section \ref{section-pi1-punctured-spec-III},
puis en comblant les lacunes à l'aide de la pureté. Rappelons que
nous disons que {\it la pureté vaut} pour un anneau local noethérien
$(A, \mathfrak m)$ si le foncteur de restriction
$\textit{F\'Et}_X \to \textit{F\'Et}_U$ est essentiellement
surjectif, où $X = \Spec(A)$ et $U = X \setminus \{\mathfrak m\}$.

\begin{lemma}
\label{lemma-equivalence-better}
Dans la situation \ref{situation-local-lefschetz}, supposons que
\begin{enumerate}
\item $A$ possède un complexe dualisant et soit complet pour la topologie $f$-adique,
\item l'une des conditions suivantes soit satisfaite :
\begin{enumerate}
\item $A_f$ vérifie $(S_2)$ et toute composante irréductible de $X$
non contenue dans $X_0$ soit de dimension $\geq 4$, ou bien
\item si $\mathfrak p \not \in V(f)$ et
$V(\mathfrak p) \cap V(f) \not = \{\mathfrak m\}$, alors
$\text{depth}(A_\mathfrak p) + \dim(A/\mathfrak p) > 3$.
\end{enumerate}
\item pour tout idéal maximal $\mathfrak p \subset A_f$,
la pureté vaille pour $(A_f)_\mathfrak p$.
\end{enumerate}
Alors le foncteur de restriction $\textit{F\'Et}_U \to \textit{F\'Et}_{U_0}$
est essentiellement surjectif.
\end{lemma}

\begin{proof}
Soit $V_0 \to U_0$ un morphisme fini étale. D'après le
lemme \ref{lemma-essentially-surjective-general-better},
il existe un ouvert $U' \subset U$ contenant $U_0$ et
un morphisme fini étale $V' \to U'$ dont le changement de base à $U_0$
est isomorphe à $V_0 \to U_0$. Comme $U' \supset U_0$,
on voit que $U \setminus U'$ est formé de points correspondant
aux idéaux premiers $\mathfrak p_1, \ldots, \mathfrak p_n$ comme dans (3).
Par hypothèse, on peut trouver des morphismes finis étales
$V'_i \to \Spec(A_{\mathfrak p_i})$ qui coïncident avec
$V' \to U'$ sur $U' \times_U \Spec(A_{\mathfrak p_i})$.
D'après Limites, lemme \ref{limits-lemma-glueing-near-closed-point},
appliqué $n$ fois, on voit que $V' \to U'$ se prolonge en un morphisme fini étale
$V \to U$.
\end{proof}

\begin{lemma}
\label{lemma-equivalence}
Soit $(A, \mathfrak m)$ un anneau local noethérien.
Soit $f \in \mathfrak m$. Supposons que
\begin{enumerate}
\item $A$ soit complet pour la topologie $f$-adique,
\item $f$ soit non diviseur de zéro,
\item $H^1_\mathfrak m(A/fA)$ et $H^2_\mathfrak m(A/fA)$ soient des
$A$-modules finis,
\item pour tout idéal maximal $\mathfrak p \subset A_f$,
la pureté vaille pour $(A_f)_\mathfrak p$.
\end{enumerate}
Alors le foncteur de restriction $\textit{F\'Et}_U \to \textit{F\'Et}_{U_0}$
est essentiellement surjectif.
\end{lemma}

\begin{proof}
La démonstration est identique à celle du
lemme \ref{lemma-equivalence-better},
en utilisant le
lemme \ref{lemma-essentially-surjective-general}
au lieu du
lemme \ref{lemma-essentially-surjective-general-better}.
\end{proof}



\section{Pureté dans le cas local, II}
\label{section-local-purity-II}

\noindent
Cette section prolonge la section \ref{section-local-purity}.
Rappelons que nous disons que {\it la pureté vaut} pour un anneau local noethérien
$(A, \mathfrak m)$ si le foncteur de restriction
$\textit{F\'Et}_X \to \textit{F\'Et}_U$ est essentiellement
surjectif, où $X = \Spec(A)$ et $U = X \setminus \{\mathfrak m\}$.

\begin{lemma}
\label{lemma-purity-inherited-by-hypersurface-better}
Soit $(A, \mathfrak m)$ un anneau local noethérien.
Soit $f \in \mathfrak m$. Supposons que
\begin{enumerate}
\item $A$ possède un complexe dualisant et soit complet pour la topologie $f$-adique,
\item l'une des conditions suivantes soit satisfaite :
\begin{enumerate}
\item $A_f$ vérifie $(S_2)$ et toute composante irréductible de $X$
non contenue dans $X_0$ soit de dimension $\geq 4$, ou bien
\item si $\mathfrak p \not \in V(f)$ et
$V(\mathfrak p) \cap V(f) \not = \{\mathfrak m\}$, alors
$\text{depth}(A_\mathfrak p) + \dim(A/\mathfrak p) > 3$.
\end{enumerate}
\item pour tout idéal maximal $\mathfrak p \subset A_f$,
la pureté vaille pour $(A_f)_\mathfrak p$, et
\item la pureté vaille pour $A$.
\end{enumerate}
Alors la pureté vaut pour $A/fA$.
\end{lemma}

\begin{proof}
Notons $X = \Spec(A)$ et $U = X \setminus \{\mathfrak m\}$
le spectre épointé. De même, on a $X_0 = \Spec(A/fA)$
et $U_0 = X_0 \setminus \{\mathfrak m\}$.
Soit $V_0 \to U_0$ un morphisme fini étale. D'après le
lemme \ref{lemma-equivalence-better},
il existe un morphisme fini étale $V \to U$
dont le changement de base à $U_0$
est isomorphe à $V_0 \to U_0$.
D'après l'hypothèse (4), le morphisme $V \to U$ se prolonge
en un morphisme fini étale $Y \to X$. La restriction de
$Y$ à $X_0$ est alors le prolongement cherché de $V_0 \to U_0$.
\end{proof}

\begin{lemma}
\label{lemma-purity-inherited-by-hypersurface}
Soit $(A, \mathfrak m)$ un anneau local noethérien.
Soit $f \in \mathfrak m$. Supposons que
\begin{enumerate}
\item $A$ soit complet pour la topologie $f$-adique,
\item $f$ soit non diviseur de zéro,
\item $H^1_\mathfrak m(A/fA)$ et $H^2_\mathfrak m(A/fA)$ soient des
$A$-modules finis,
\item pour tout idéal maximal $\mathfrak p \subset A_f$,
la pureté vaille pour $(A_f)_\mathfrak p$,
\item la pureté vaille pour $A$.
\end{enumerate}
Alors la pureté vaut pour $A/fA$.
\end{lemma}

\begin{proof}
La démonstration est identique à celle du
lemme \ref{lemma-purity-inherited-by-hypersurface-better},
en utilisant le
lemme \ref{lemma-equivalence}
au lieu du
lemme \ref{lemma-equivalence-better}.
\end{proof}

\noindent
On peut maintenant amplifier les résultats précédents pour démontrer que
la pureté vaut pour les intersections complètes de dimension $\geq 3$.
Rappelons qu'un anneau local noethérien est appelé une intersection
complète si son complété est le quotient d'un
anneau local régulier par l'idéal engendré par une suite régulière.
Voir la discussion dans Algèbre à puissances divisées, section \ref{dpa-section-lci}.

\begin{proposition}
\label{proposition-purity-complete-intersection}
Soit $(A, \mathfrak m)$ un anneau local noethérien. Si $A$ est une
intersection complète de dimension $\geq 3$, alors la pureté
vaut pour $A$, au sens où tout revêtement fini étale du
spectre épointé se prolonge.
\end{proposition}

\begin{proof}
D'après le lemme \ref{lemma-purity-and-completion}, on peut supposer $A$
local complet. Par hypothèse, on peut écrire
$A = B/(f_1, \ldots, f_r)$, où $B$ est un anneau local régulier complet
et $f_1, \ldots, f_r$ une suite régulière.
On achève la démonstration par récurrence sur $r$.
Le cas initial est $r = 0$ et résulte du
lemme \ref{lemma-local-purity}, qui s'applique aux
anneaux réguliers de dimension $\geq 2$.

\medskip\noindent
Supposons que $A = B/(f_1, \ldots, f_r)$ et que la proposition
vaille pour $r - 1$. Posons $A' = B/(f_1, \ldots, f_{r - 1})$ et appliquons le
lemme \ref{lemma-purity-inherited-by-hypersurface} à $f_r \in A'$.
C'est permis :
la condition (2) est satisfaite puisque $f_1, \ldots, f_r$ est une suite régulière,
la condition (1) est satisfaite puisque $B$, et donc $A'$, est complet,
la condition (3) est satisfaite puisque $A = A'/f_r A'$ est de Cohen--Macaulay de dimension
$\dim(A) \geq 3$ ; voir Complexes dualisants, lemme \ref{dualizing-lemma-depth},
la condition (4) est satisfaite par l'hypothèse de récurrence, car
$\dim((A'_{f_r})_\mathfrak p) \geq 3$ pour un idéal premier maximal
$\mathfrak p$ de $A'_{f_r}$, et car
$(A'_{f_r})_\mathfrak p = B_\mathfrak q/(f_1, \ldots, f_{r - 1})$
pour un certain $\mathfrak q \subset B$,
la condition (5) est satisfaite par l'hypothèse de récurrence.
\end{proof}




\section{Pureté dans le cas local, III}
\label{section-local-purity-III}

\noindent
Cette section prolonge la discussion des
sections \ref{section-local-purity} et \ref{section-local-purity-II}.

\begin{lemma}
\label{lemma-fully-faithful-power-series-over-depth2}
Soit $(A, \mathfrak m)$ un anneau local noethérien de profondeur $\geq 2$.
Soit $B = A[[x_1, \ldots, x_d]]$ avec $d \geq 1$.
Posons $Y = \Spec(B)$ et $Y_0 = V(x_1, \ldots, x_d)$.
Pour tout sous-schéma ouvert $V \subset Y$ tel que
$V_0 = V \cap Y_0$ soit égal à $Y_0 \setminus \{\mathfrak m_B\}$,
le foncteur de restriction
$$
\textit{F\'Et}_V \longrightarrow \textit{F\'Et}_{V_0}
$$
est pleinement fidèle.
\end{lemma}

\begin{proof}
Posons $I = (x_1, \ldots, x_d)$. Posons $X = \Spec(A)$.
Si l'on utilise le morphisme $Y \to X$ pour identifier $Y_0$ à $X$,
alors $V_0$ s'identifie au spectre épointé $U$ de $A$.
En prenant les images directes des modules par ce morphisme affine, on obtient
\begin{align*}
\lim_n \Gamma(V_0, \mathcal{O}_V/I^n\mathcal{O}_V)
& =
\lim_n \Gamma(V_0, \mathcal{O}_Y/I^n\mathcal{O}_Y) \\
& =
\lim_n \Gamma(U, \mathcal{O}_U[x_1, \ldots, x_d]/(x_1, \ldots, x_d)^n) \\
& =
\lim_n A[x_1, \ldots, x_d]/(x_1, \ldots, x_d)^n \\
& =
B
\end{align*}
En effet, puisque la profondeur de $A$ est $\geq 2$, on a $\Gamma(U, \mathcal{O}_U) = A$ ;
voir Cohomologie locale, lemme
\ref{local-cohomology-lemma-finiteness-pushforwards-and-H1-local}.
Ainsi, pour tout ouvert $V \subset Y$ comme dans le lemme, on obtient
$$
B = \Gamma(Y, \mathcal{O}_Y) \to \Gamma(V, \mathcal{O}_V) \to
\lim_n \Gamma(V_0, \mathcal{O}_Y/I^n\mathcal{O}_Y) = B
$$
ce qui implique que les deux flèches sont des isomorphismes (on omet un petit détail).
D'après Géométrie algébrique et formelle, lemme
\ref{algebraization-lemma-completion-fully-faithful},
on conclut que
$\textit{Coh}(\mathcal{O}_V) \to \textit{Coh}(V, I\mathcal{O}_V)$
est pleinement fidèle sur la sous-catégorie pleine des objets finis localement libres.
On conclut donc par le
lemme \ref{lemma-restriction-fully-faithful}.
\end{proof}

\begin{lemma}
\label{lemma-purity-power-series-over-depth2}
\begin{slogan}
Ramanujam--Samuel pour les revêtements finis étales
\end{slogan}
Soit $(A, \mathfrak m)$ un anneau local noethérien de profondeur $\geq 2$. Soit
$B = A[[x_1, \ldots, x_d]]$ avec $d \geq 1$. Pour tout ouvert
$V \subset Y = \Spec(B)$ qui contient
\begin{enumerate}
\item tout idéal premier $\mathfrak q \subset B$ tel que
$\mathfrak q \cap A \not = \mathfrak m$,
\item l'idéal premier $\mathfrak m B$,
\end{enumerate}
le foncteur
$
\textit{F\'Et}_Y
\to
\textit{F\'Et}_V
$
est une équivalence. En particulier, la pureté vaut pour $B$.
\end{lemma}

\begin{proof}
Un idéal premier $\mathfrak q \subset B$ qui n'appartient pas à $V$
est au-dessus de $\mathfrak m$. Dans ce cas, $A \to B_\mathfrak q$
est un homomorphisme local plat, et donc $\text{depth}(B_\mathfrak q) \geq 2$
(Algèbre, lemme \ref{algebra-lemma-apply-grothendieck}).
Le foncteur est donc pleinement fidèle d'après le
lemme \ref{lemma-quasi-compact-dense-open-connected-at-infinity-Noetherian},
combiné au lemme de Cohomologie locale
\ref{local-cohomology-lemma-depth-2-connected-punctured-spectrum}.

\medskip\noindent
Soit $W \to V$ un morphisme fini étale. Soit $B \to C$ l'unique morphisme fini
d'anneaux tel que $\Spec(C) \to Y$ soit le morphisme fini prolongeant
$W \to V$ construit dans le lemme \ref{lemma-extend-S2}.
Remarquons que $C = \Gamma(W, \mathcal{O}_W)$.

\medskip\noindent
Posons $Y_0 = V(x_1, \ldots, x_d)$ et $V_0 = V \cap Y_0$. Posons $X = \Spec(A)$.
Si l'on utilise le morphisme $Y \to X$ pour identifier $Y_0$ à $X$,
alors $V_0$ s'identifie au spectre épointé $U$ de $A$.
On peut ainsi considérer $W_0 = W \times_Y Y_0$ comme un schéma fini étale
sur $U$. Alors
$$
W_0 \times_U (U \times_X Y)
\quad\text{et}\quad
W \times_V (U \times_X Y)
$$
sont des schémas finis étales sur $U \times_X Y$ dont les restrictions sont
des schémas finis étales isomorphes sur $V_0$. D'après le
lemme \ref{lemma-fully-faithful-power-series-over-depth2},
appliqué à l'ouvert $U \times_X Y$, on obtient un isomorphisme
$$
W_0 \times_U (U \times_X Y) \longrightarrow W \times_V (U \times_X Y)
$$
sur $U \times_X Y$.

\medskip\noindent
Remarquons que $C_0 = \Gamma(W_0, \mathcal{O}_{W_0})$ est une $A$-algèbre finie
d'après le lemme \ref{lemma-extend-S2} appliqué à $W_0 \to U \subset X$ (exactement
comme ci-dessus pour $B \to C$). Comme la construction du
lemme \ref{lemma-extend-S2} est compatible au changement de base plat
et au changement d'ouverts, l'isomorphisme ci-dessus induit
un isomorphisme
$$
\Psi : C \longrightarrow C_0 \otimes_A B
$$
de $B$-algèbres finies. Cependant, on sait que $\Spec(C) \to Y$
est étale en tout point situé au-dessus d'au moins un point de $Y$ qui est
au-dessus de $\mathfrak m \in X$. Comme $\Psi$ est un isomorphisme, on
conclut que $\Spec(C_0) \to X$
est étale au-dessus de $\mathfrak m$ (on omet un petit détail).
Cela signifie bien sûr que $A \to C_0$ est fini
étale, et donc que $B \to C$ est fini étale.
\end{proof}

\begin{lemma}
\label{lemma-purity-smooth-over-depth2}
Soit $f : X \to S$ un morphisme de schémas. Soit $U \subset X$
un sous-schéma ouvert. Supposons que
\begin{enumerate}
\item $f$ soit lisse,
\item $S$ soit noethérien,
\item pour $s \in S$ tel que $\text{depth}(\mathcal{O}_{S, s}) \leq 1$,
on ait $X_s = U_s$,
\item $U_s \subset X_s$ soit dense pour tout $s \in S$.
\end{enumerate}
Alors $\textit{F\'Et}_X \to \textit{F\'Et}_U$ est une équivalence.
\end{lemma}

\begin{proof}
Le foncteur est pleinement fidèle d'après le
lemme \ref{lemma-quasi-compact-dense-open-connected-at-infinity-Noetherian},
combiné au lemme de Cohomologie locale
\ref{local-cohomology-lemma-depth-2-connected-punctured-spectrum}
(ainsi qu'à une application du
lemme d'Algèbre \ref{algebra-lemma-apply-grothendieck}
pour vérifier la condition de profondeur).

\medskip\noindent
Soit $\pi : V \to U$ un morphisme fini étale. Soit $Y \to X$
le morphisme fini construit dans le lemme \ref{lemma-extend-S2}.
Il faut montrer que $Y \to X$ est fini étale.
Pour le montrer en tout point $x \in X$ s'envoyant sur un
point donné $s \in S$, on peut effectuer le changement de base par un morphisme plat
$S' \to S$ de schémas noethériens tel que $s$ appartienne
à l'image. Cela résulte de la compatibilité de la
construction du lemme \ref{lemma-extend-S2} au changement de base plat.

\medskip\noindent
Après avoir agrandi $U$, on peut supposer que $U \subset X$ est
le plus grand ouvert au-dessus duquel $Y \to X$ est fini étale.
Soit $Z \subset X$ le complémentaire de $U$.
Pour obtenir une contradiction, supposons $Z \not = \emptyset$.
Soit $s \in S$ un point de l'image de $Z \to S$
tel qu'aucune générisation stricte de $s$ n'appartienne à cette image.
Après changement de base à $\Spec(\mathcal{O}_{S, s})$,
on voit alors que $S = \Spec(A)$, où $(A, \mathfrak m, \kappa)$ est
un anneau local noethérien de profondeur $\geq 2$, et que $Z$ est
contenu dans la fibre fermée $X_s$
et n'est dense nulle part dans $X_s$. Choisissons un point fermé $z \in Z$.
Alors $\kappa(z)/\kappa$ est fini (par le théorème des zéros de Hilbert ; voir
Algèbre, théorème \ref{algebra-theorem-nullstellensatz}).
Choisissons un morphisme fini plat $(S', s') \to (S, s)$ de schémas locaux
réalisant l'extension résiduelle $\kappa(z)/\kappa$ ; voir
Algèbre, lemme \ref{algebra-lemma-finite-free-given-residue-field-extension}.
Après avoir effectué le changement de base par $S' \to S$, on se ramène au cas
où $\kappa(z) = \kappa$.

\medskip\noindent
D'après Compléments sur les morphismes, lemme \ref{more-morphisms-lemma-slice-smooth},
il existe un sous-schéma localement fermé $S' \subset X$ passant par $z$
tel que $S' \to S$ soit étale en $z$. Après changement de base
par $S' \to S$, on peut supposer qu'il existe une section $\sigma : S \to X$
telle que $\sigma(s) = z$. Choisissons un voisinage affine
$\Spec(B) \subset X$ de $z$. Alors $A \to B$ est un morphisme lisse
d'anneaux qui possède une section $\sigma : B \to A$. Notons $I = \Ker(\sigma)$
et notons $B^\wedge$ le complété $I$-adique de $B$.
Alors $B^\wedge \cong A[[x_1, \ldots, x_d]]$ pour un certain $d \geq 0$ ; voir
Algèbre, lemme \ref{algebra-lemma-section-smooth}.
Remarquons que $d > 0$, car sinon on verrait que $X \to S$
est étale en $z$, ce qui impliquerait que $z$ est un point générique de
$X_s$, et donc que $z \in U$ d'après l'hypothèse (4).
De même, si $d > 0$, alors $\mathfrak m B^\wedge$ s'envoie dans
$U$ par le morphisme $\Spec(B^\wedge) \to X$.
Il suffit de démontrer que $Y \to X$ est fini étale après changement de base
à $\Spec(B^\wedge)$. Comme $B \to B^\wedge$ est plat
(Algèbre, lemme \ref{algebra-lemma-completion-flat}),
cela résulte du lemme \ref{lemma-purity-power-series-over-depth2}
et de l'unicité dans la construction de $Y \to X$.
\end{proof}

\begin{proposition}
\label{proposition-purity-smooth-over-depth2}
Soit $A \to B$ un homomorphisme local d'anneaux locaux noethériens.
Supposons $A$ de profondeur $\geq 2$, le morphisme $A \to B$ formellement lisse pour la
topologie $\mathfrak m_B$-adique, et $\dim(B) > \dim(A)$. Pour tout ouvert
$V \subset Y = \Spec(B)$ qui contient
\begin{enumerate}
\item tout idéal premier $\mathfrak q \subset B$ tel que
$\mathfrak q \cap A \not = \mathfrak m_A$,
\item l'idéal premier $\mathfrak m_A B$,
\end{enumerate}
le foncteur $\textit{F\'Et}_Y \to \textit{F\'Et}_V$
est une équivalence. En particulier, la pureté vaut pour $B$.
\end{proposition}

\begin{proof}
Un idéal premier $\mathfrak q \subset B$ qui n'appartient pas à $V$
est au-dessus de $\mathfrak m_A$. Dans ce cas, $A \to B_\mathfrak q$
est un homomorphisme local plat, et donc $\text{depth}(B_\mathfrak q) \geq 2$
(Algèbre, lemme \ref{algebra-lemma-apply-grothendieck}).
Le foncteur est donc pleinement fidèle d'après le
lemme \ref{lemma-quasi-compact-dense-open-connected-at-infinity-Noetherian},
combiné au lemme de Cohomologie locale
\ref{local-cohomology-lemma-depth-2-connected-punctured-spectrum}.

\medskip\noindent
Notons $A^\wedge$ et $B^\wedge$ les complétés de $A$ et $B$
pour leurs idéaux maximaux. Remarquons que les hypothèses
de la proposition sont satisfaites par $A^\wedge \to B^\wedge$ ; voir
Compléments d'algèbre, lemmes
\ref{more-algebra-lemma-completion-dimension},
\ref{more-algebra-lemma-completion-depth} et
\ref{more-algebra-lemma-formally-smooth-completion}.
Par l'unicité et la compatibilité au changement de base plat
de la construction du lemme \ref{lemma-extend-S2},
il suffit de démontrer la surjectivité essentielle pour
$A^\wedge \to B^\wedge$ et l'image inverse de $V$
(détails omis ; comparer au lemme \ref{lemma-purity-and-completion}
dans le cas où $V$ est le spectre épointé).
D'après Compléments d'algèbre, proposition \ref{more-algebra-proposition-fs-regular},
cela signifie que l'on peut supposer $A \to B$ régulier.

\medskip\noindent
Soit $W \to V$ un morphisme fini étale.
D'après le théorème de Popescu
(Lissage des morphismes d'anneaux, théorème \ref{smoothing-theorem-popescu}),
on peut écrire $B = \colim B_i$ comme colimite filtrante
de $A$-algèbres lisses. On peut choisir un $i$ et un
ouvert $V_i \subset \Spec(B_i)$ dont l'image inverse est $V$
(Limites, lemme \ref{limits-lemma-descend-opens}).
Après avoir augmenté $i$, on peut supposer qu'il existe un morphisme
fini étale $W_i \to V_i$ dont le changement de base à $V$
est $W \to V$ ; voir
Limites, lemmes \ref{limits-lemma-descend-finite-presentation},
\ref{limits-lemma-descend-finite-finite-presentation} et
\ref{limits-lemma-descend-etale}.
On peut supposer que le complémentaire de $V_i$ est contenu
dans la fibre fermée de $\Spec(B_i) \to \Spec(A)$, puisque cela
est vrai pour $V$ (soit on choisit ainsi $V_i$, soit on utilise
le lemme précédent pour montrer que cela vaut pour $i$ assez grand).
Soit $\eta$ le point générique de la fibre fermée
de $\Spec(B) \to \Spec(A)$. Comme $\eta \in V$, l'image de
$\eta$ appartient à $V_i$. Ainsi, après avoir remplacé $V_i$ par un
voisinage ouvert affine de l'image du point fermé
de $\Spec(B)$, on peut supposer que la fibre fermée
de $\Spec(B_i) \to \Spec(A)$ est irréductible et que
son point générique appartient à $V_i$ (détails omis ; en effet,
un schéma lisse sur un corps est réunion disjointe de schémas irréductibles).
On peut alors appliquer le lemme \ref{lemma-purity-smooth-over-depth2}
pour voir que $W_i \to V_i$ se prolonge en un morphisme fini étale
$\Spec(C_i) \to \Spec(B_i)$ et, en le ramenant
à $\Spec(B)$, on conclut que $W$ appartient à
l'image essentielle du foncteur
$\textit{F\'Et}_Y \to \textit{F\'Et}_V$,
comme souhaité.
\end{proof}





\section{Lefschetz pour le groupe fondamental}
\label{section-global-lefschetz}

\noindent
Nous avons bien sûr déjà démontré plusieurs résultats de ce type
dans le cas local. Dans cette section, nous étudions le cas projectif.

\begin{proposition}
\label{proposition-lefschetz-fully-faithful}
Soit $k$ un corps. Soit $X$ un schéma propre sur $k$.
Soit $\mathcal{L}$ un $\mathcal{O}_X$-module inversible ample.
Soit $s \in \Gamma(X, \mathcal{L})$. Soit $Y = Z(s)$ le
schéma des zéros de $s$. Supposons que, pour tout $x \in X \setminus Y$,
on ait
$$
\text{depth}(\mathcal{O}_{X, x}) + \dim(\overline{\{x\}}) > 1
$$
Alors le foncteur de restriction $\textit{F\'Et}_X \to \textit{F\'Et}_Y$
est pleinement fidèle. En fait, pour tout sous-schéma ouvert $V \subset X$
contenant $Y$, le foncteur de restriction
$\textit{F\'Et}_V \to \textit{F\'Et}_Y$
est pleinement fidèle.
\end{proposition}

\begin{proof}
La première assertion est une conséquence formelle du
lemme \ref{lemma-restriction-fully-faithful-special}
et de la
proposition de Géométrie algébrique et formelle
\ref{algebraization-proposition-lefschetz}.
La seconde assertion découle du
lemme \ref{lemma-restriction-fully-faithful-special}
et du
lemme de Géométrie algébrique et formelle
\ref{algebraization-lemma-lefschetz-addendum}.
\end{proof}

\begin{proposition}
\label{proposition-lefschetz-equivalence-general}
Soit $k$ un corps. Soit $X$ un schéma propre sur $k$.
Soit $\mathcal{L}$ un $\mathcal{O}_X$-module inversible ample.
Soit $s \in \Gamma(X, \mathcal{L})$. Soit $Y = Z(s)$ le
schéma des zéros de $s$. Soit $\mathcal{V}$ l'ensemble des
sous-schémas ouverts de $X$ contenant $Y$, ordonné par inclusion inverse.
Supposons que, pour tout $x \in X \setminus Y$, on ait
$$
\text{depth}(\mathcal{O}_{X, x}) + \dim(\overline{\{x\}}) > 2
$$
Alors le foncteur de restriction
$$
\colim_\mathcal{V} \textit{F\'Et}_V \to \textit{F\'Et}_Y
$$
est une équivalence.
\end{proposition}

\begin{proof}
C'est une conséquence formelle du
lemme \ref{lemma-restriction-equivalence-general} et de la
proposition de Géométrie algébrique et formelle
\ref{algebraization-proposition-lefschetz-equivalence}.
\end{proof}

\begin{proposition}
\label{proposition-lefschetz-equivalence}
Soit $k$ un corps. Soit $X$ un schéma propre sur $k$.
Soit $\mathcal{L}$ un $\mathcal{O}_X$-module inversible ample.
Soit $s \in \Gamma(X, \mathcal{L})$. Soit $Y = Z(s)$ le
schéma des zéros de $s$.
Supposons que, pour tout $x \in X \setminus Y$, on ait
$$
\text{depth}(\mathcal{O}_{X, x}) + \dim(\overline{\{x\}}) > 2
$$
et que, lorsque $x \in X \setminus Y$ est fermé, la pureté vaille pour
$\mathcal{O}_{X, x}$. Alors le foncteur de restriction
$\textit{F\'Et}_X \to \textit{F\'Et}_Y$
est une équivalence. Si $X$, ou de manière équivalente $Y$, est connexe, alors
$$
\pi_1(Y, \overline{y}) \to \pi_1(X, \overline{y})
$$
est un isomorphisme pour tout point géométrique $\overline{y}$ de $Y$.
\end{proposition}

\begin{proof}
La pleine fidélité résulte de la
proposition \ref{proposition-lefschetz-fully-faithful}.
D'après la proposition \ref{proposition-lefschetz-equivalence-general},
tout objet de $\textit{F\'Et}_Y$
est isomorphe au produit fibré $U \times_V Y$ pour un certain
morphisme fini étale $U \to V$, où $V \subset X$
est un sous-schéma ouvert contenant $Y$.
Le complémentaire $T = X \setminus V$
est\footnote{En effet, $T$ est propre sur $k$ (car fermé dans $X$)
et affine (car fermé dans le schéma affine $X \setminus Y$ ; voir
Morphismes, lemme \ref{morphisms-lemma-proper-ample-delete-affine}),
et donc fini sur $k$
(Morphismes, lemme \ref{morphisms-lemma-finite-proper}).
Ainsi, $T$ est un ensemble fini de points fermés.}
un ensemble fini de points fermés de $X \setminus Y$.
Écrivons $T = \{x_1, \ldots, x_n\}$.
Par hypothèse, on peut trouver des morphismes finis étales
$V'_i \to \Spec(\mathcal{O}_{X, x_i})$ qui coïncident avec
$U \to V$ sur $V \times_X \Spec(\mathcal{O}_{X, x_i})$.
D'après Limites, lemme \ref{limits-lemma-glueing-near-closed-point},
appliqué $n$ fois, on voit que $U \to V$ se prolonge en un morphisme fini étale
$U' \to X$, comme souhaité.
Voir le lemme \ref{lemma-what-equivalence-gives} pour la dernière assertion.
\end{proof}










\section{Pureté du lieu de ramification}
\label{section-purity-ramification}

\noindent
Dans cette section, nous examinons l'analogue de la pureté du lieu de branchement pour
les morphismes génériquement finis. Ce résultat est apparemment dû à Gabber.
Un cas particulier est le théorème de pureté de van der Waerden pour le lieu où
un morphisme birationnel d'une variété normale vers une variété lisse n'est pas
un isomorphisme.

\begin{lemma}
\label{lemma-characterize-rational-singularity}
Soit $A$ un domaine local noethérien normal de dimension $2$.
Supposons que $A$ soit un anneau de Nagata, possède un module dualisant $\omega_A$ et admette une
résolution des singularités $f : X \to \Spec(A)$.
Soit $\omega_X$ comme dans Résolution des surfaces,
remarque \ref{resolve-remark-dualizing-setup}.
Si $\omega_X \cong \mathcal{O}_X(E)$ pour un certain diviseur de
Cartier effectif $E \subset X$ porté par la fibre exceptionnelle,
alors $A$ définit une singularité rationnelle.
Si $f$ est une résolution minimale, alors $E = 0$.
\end{lemma}

\begin{proof}
Il existe un morphisme trace $Rf_*\omega_X \to \omega_A$, voir
Dualité pour les schémas, section \ref{duality-section-trace}.
D'après Grauert--Riemenschneider
(Résolution des surfaces,
proposition \ref{resolve-proposition-Grauert-Riemenschneider}),
on a $R^1f_*\omega_X = 0$.
Ainsi, le morphisme trace est un morphisme $f_*\omega_X \to \omega_A$.
On peut alors considérer
$$
\mathcal{O}_{\Spec(A)} = f_*\mathcal{O}_X \to f_*\omega_X \to \omega_A
$$
où le premier morphisme provient du morphisme
$\mathcal{O}_X \to \mathcal{O}_X(E) = \omega_X$ dont l'existence
est supposée dans l'énoncé du lemme.
La composée est un isomorphisme d'après Diviseurs, lemme
\ref{divisors-lemma-check-isomorphism-via-depth-and-ass},
car c'est un isomorphisme sur le spectre épointé de $A$
(par l'hypothèse du lemme et le fait que $f$ est un isomorphisme
sur le spectre épointé) et que $A$ et $\omega_A$
sont des $A$-modules de profondeur $2$ (d'après
Algèbre, lemme \ref{algebra-lemma-criterion-normal} et
Complexes dualisants, lemme \ref{dualizing-lemma-depth-dualizing-module}).
Par conséquent, $f_*\omega_X \to \omega_A$ est surjectif, donc est un isomorphisme.
Ainsi, $Rf_*\omega_X = \omega_A$, ce qui, par dualité, implique
$Rf_*\mathcal{O}_X = \mathcal{O}_{\Spec(A)}$.
Il s'ensuit que $H^1(X, \mathcal{O}_X) = 0$, ce qui implique que $A$
définit une singularité rationnelle (voir la discussion dans
Résolution des surfaces, section
\ref{resolve-section-bounded}, en particulier les
lemmes \ref{resolve-lemma-regular-rational} et
\ref{resolve-lemma-exact-sequence}).
Si $f$ est minimal, alors $E = 0$, car le morphisme
$f^*\omega_A \to \omega_X$ est surjectif d'après
une application répétée de Résolution des surfaces, lemme
\ref{resolve-lemma-dualizing-blow-up-rational},
et $\omega_A \cong A$, comme on l'a vu ci-dessus.
\end{proof}

\begin{lemma}
\label{lemma-key-purity-ramification}
Posons $Y = \Spec(A)$ et soit $f : X \to Y$ un morphisme de type fini.
Soit $x \in X$ un point. Supposons que
\begin{enumerate}
\item $A$ soit un anneau local régulier excellent,
\item $\mathcal{O}_{X, x}$ soit normal de dimension $2$,
\item $f$ soit \'etale hors de $\overline{\{x\}}$.
\end{enumerate}
Alors $f$ est \'etale en $x$.
\end{lemma}

\begin{proof}
On remplace d'abord $X$ par un voisinage ouvert affine de $x$.
Notons que $\mathcal{O}_{X, x}$ est un anneau local excellent
(Compléments d'algèbre, lemme \ref{more-algebra-lemma-finite-type-over-excellent}).
On peut donc choisir une résolution minimale des singularités
$W \to \Spec(\mathcal{O}_{X, x})$, voir
Résolution des surfaces, théorème \ref{resolve-theorem-resolve}.
Après avoir éventuellement remplacé $X$ par un voisinage ouvert affine de $x$,
on peut trouver un morphisme propre $b : X' \to X$ tel que
$X' \times_X \Spec(\mathcal{O}_{X, x}) = W$, voir
Limites, lemme \ref{limits-lemma-glueing-near-closed-point}.
Après avoir encore rétréci $X$, on peut supposer que $X'$ est régulier.
En effet, on sait que $W$ est régulier et que $X'$ est excellent,
et le lieu régulier du spectre d'un anneau excellent est ouvert.
Puisque $W \to \Spec(\mathcal{O}_{X, x})$ est projectif
(comme composé d'une suite d'éclatements normalisés), on peut supposer,
après avoir rétréci $X$, que $b$ est projectif (détails omis).
Soit $U = X \setminus \overline{\{x\}}$.
Puisque $W \to \Spec(\mathcal{O}_{X, x})$ est un isomorphisme sur le
spectre épointé, on peut supposer que $b : X' \to X$ est un isomorphisme sur $U$.
Ainsi, on peut et on va aussi considérer $U$ comme un sous-schéma ouvert de $X'$.
Posons $f' = f \circ b : X' \to \Spec(A)$.

\medskip\noindent
Puisque $A$ est régulier, on voit que $\mathcal{O}_Y$ est un complexe dualisant pour $Y$.
Par conséquent, $f^!\mathcal{O}_Y$ est un complexe dualisant sur $X$
(Dualité pour les schémas, lemme \ref{duality-lemma-shriek-dualizing}).
Le lieu de Cohen--Macaulay de $X$ est ouvert d'après
Dualité pour les schémas, lemme \ref{duality-lemma-dualizing-module-CM-scheme}
(cela peut aussi se démontrer à l'aide de l'excellence).
Puisque $\mathcal{O}_{X, x}$ est de Cohen--Macaulay, après avoir rétréci
$X$, on peut supposer que $X$ est de Cohen--Macaulay.
Notons qu'un morphisme \'etale est une intersection complète locale.
Ainsi,
Dualité pour les schémas, lemme \ref{duality-lemma-fundamental-class-almost-lci},
s'applique avec $r = 0$ et donne un morphisme
$$
\mathcal{O}_X \longrightarrow \omega_{X/Y} = H^0(f^!\mathcal{O}_Y)
$$
qui est un isomorphisme sur $X \setminus \overline{\{x\}}$.
Puisque $\omega_{X/Y}$ est $(S_2)$ d'après
Dualité pour les schémas, lemme \ref{duality-lemma-shriek-over-CM},
on constate que ce morphisme est un isomorphisme d'après
Diviseurs, lemme \ref{divisors-lemma-check-isomorphism-via-depth-and-ass}.
Cela montre déjà que $X$, et en particulier $\mathcal{O}_{X, x}$, est
de Gorenstein.

\medskip\noindent
Posons $\omega_{X'/Y} = H^0((f')^!\mathcal{O}_Y)$. En raisonnant exactement de la
même manière que ci-dessus, on voit que $(f')^!\mathcal{O}_Y = \omega_{X'/Y}[0]$
est un complexe dualisant pour $X'$. Puisque $X'$ est régulier,
le morphisme $X' \to Y$ est un morphisme d'intersection complète locale, voir
Compléments sur les morphismes, lemme
\ref{more-morphisms-lemma-morphism-regular-schemes-is-lci}.
D'après Dualité pour les schémas, lemme \ref{duality-lemma-fundamental-class-lci},
il existe un morphisme
$$
\mathcal{O}_{X'} \longrightarrow \omega_{X'/Y}
$$
qui est un isomorphisme sur $U$. On en déduit que
$\omega_{X'/Y} = \mathcal{O}_{X'}(E)$ pour un certain diviseur de
Cartier effectif $E \subset X'$ disjoint de $U$.

\medskip\noindent
Puisque $\omega_{X/Y} = \mathcal{O}_X$, on voit que
$\omega_{X'/Y} = b^! f^!\mathcal{O}_Y = b^!\mathcal{O}_X$.
En revenant à $W \to \Spec(\mathcal{O}_{X, x})$,
on voit que $\omega_W = \mathcal{O}_W(E|_W)$.
D'après le lemme \ref{lemma-characterize-rational-singularity},
on obtient $E|_W = 0$.
Cela signifie que $f' : X' \to Y$ est \'etale d'après le lemme (déjà utilisé)
de Dualité pour les schémas, lemme \ref{duality-lemma-fundamental-class-lci}.
Cela achève immédiatement la démonstration, car le caractère \'etale
de $f'$ force $b$ à être un isomorphisme.
\end{proof}

\begin{lemma}[Pureté du lieu de ramification]
\label{lemma-purity-ramification}
\begin{reference}
Ce résultat pour les espaces complexes se trouve à la page 170 de \cite{Fischer}.
Dans le cas général, il s'agit de \cite[théorème 2.4]{Zong}, attribué à Gabber.
\end{reference}
Soit $f : X \to Y$ un morphisme de schémas localement noethériens.
Soit $x \in X$ et posons $y = f(x)$. Supposons que
\begin{enumerate}
\item $\mathcal{O}_{X, x}$ soit normal de dimension $\geq 1$,
\item $\mathcal{O}_{Y, y}$ soit régulier,
\item $f$ soit localement de type fini, et
\item pour les spécialisations $x' \leadsto x$ telles que
$\dim(\mathcal{O}_{X, x'}) = 1$, le morphisme $f$ soit \'etale en $x'$.
\end{enumerate}
Alors $f$ est \'etale en $x$.
\end{lemma}

\begin{proof}
On va démontrer le lemme par récurrence sur $d = \dim(\mathcal{O}_{X, x})$.

\medskip\noindent
Le cas $d = 1$ est immédiat, puisque dans ce cas le morphisme
$f$ est \'etale en $x$ par hypothèse. Supposons $d \geq 2$.

\medskip\noindent
On peut effectuer le changement de base par $\Spec(\mathcal{O}_{Y, y}) \to Y$
sans modifier la conclusion du lemme, voir
Morphismes, lemme \ref{morphisms-lemma-set-points-where-fibres-etale}.
Ainsi, on peut supposer $Y = \Spec(A)$, où $A$ est un anneau local régulier,
et que $y$ correspond à l'idéal maximal $\mathfrak m$ de $A$.

\medskip\noindent
Soit $x' \leadsto x$ une spécialisation avec $x' \not = x$.
Alors $\mathcal{O}_{X, x'}$ est normal, comme localisation de
$\mathcal{O}_{X, x}$. Si $x'$ n'est pas un point générique de $X$,
alors $1 \leq \dim(\mathcal{O}_{X, x'}) < d$ et on conclut que
$f$ est \'etale en $x'$ par l'hypothèse de récurrence.
On peut donc supposer que $f$ est \'etale en tout point se spécialisant en
$x$. Puisque l'ensemble des points où $f$ est \'etale est ouvert dans $X$
(par définition), on peut, après avoir remplacé $X$ par un voisinage ouvert de $x$,
supposer que $f$ est \'etale hors de $\overline{\{x\}}$.
En particulier, on voit que $f$ est \'etale sauf en des points
situés au-dessus du point fermé $y \in Y = \Spec(A)$.

\medskip\noindent
Soit $X' = X \times_{\Spec(A)} \Spec(A^\wedge)$. Soit $x' \in X'$
l'unique point au-dessus de $x$. D'après ce qui précède, on voit que
$X'$ est \'etale sur $\Spec(A^\wedge)$ hors de la fibre fermée et, par conséquent,
$X'$ est normal hors de la fibre fermée. Puisque $X$ est normal,
on en déduit que $X'$ est normal d'après
Résolution des surfaces, lemme \ref{resolve-lemma-normalization-completion}.
Alors, si l'on peut montrer que $X' \to \Spec(A^\wedge)$ est \'etale en $x'$,
le morphisme $f$ est \'etale en $x$ (d'après le lemme déjà cité de
Morphismes, lemme \ref{morphisms-lemma-set-points-where-fibres-etale}).
Ainsi, on peut et on va supposer que $A$ est un anneau local régulier complet.

\medskip\noindent
Le cas $d = 2$ résulte alors du lemme \ref{lemma-key-purity-ramification}.

\medskip\noindent
Supposons $d > 2$. Soit $t \in \mathfrak m$, $t \not \in \mathfrak m^2$.
Posons $Y_0 = \Spec(A/tA)$ et $X_0 = X \times_Y Y_0$.
Alors $X_0 \to Y_0$ est \'etale hors de la fibre au-dessus du point fermé.
Puisque $d > 2$, on a $\dim(\mathcal{O}_{X_0, x}) = d - 1$, qui est $\geq 2$.
La normalisation $X_0' \to X_0$ est surjective et finie
(car on travaille sur un anneau local complet et que de tels anneaux sont de Nagata).
Soit $x' \in X_0'$ un point dont l'image est $x$. Par hypothèse de récurrence, le
morphisme $X'_0 \to Y_0$ est \'etale en $x'$. Des inclusions
$\kappa(y) \subset \kappa(x) \subset \kappa(x')$,
on déduit que $\kappa(x)$ est fini sur $\kappa(y)$.
Par conséquent, $x$ est un point fermé de la fibre de $X \to Y$
au-dessus de $y$. Mais comme $x$ est également un point générique de
cette fibre, on conclut que $f$ est quasi-fini en $x$,
et on se ramène au cas de la pureté du lieu de branchement, voir le
lemme \ref{lemma-purity}.
\end{proof}




\section{Affinité du complémentaire du lieu de ramification}
\label{section-stronger-purity}

\noindent
Soit $f : X \to Y$ un morphisme de type fini de schémas noethériens,
avec $X$ normal et $Y$ régulier. Soit $V \subset X$ le plus grand
sous-schéma ouvert sur lequel $f$ est \'etale.
La discussion dans \cite[chapitre IV, section 21.12]{EGA}
laisse penser que $V \to X$ pourrait être un morphisme affine.
Notons que, si $V \to X$ est affine, on en déduit la pureté du
lieu de ramification (lemme \ref{lemma-purity-ramification}) en
utilisant Diviseurs, lemme \ref{divisors-lemma-complement-affine-open-immersion}.
Ainsi, l'affinité de $V \to X$ est une forme ``forte'' de la pureté
du lieu de ramification.
Dans cette section, nous démontrons que $V \to X$ est affine lorsque
$X$ et $Y$ sont équicaractéristiques et excellents, voir le
théorème \ref{theorem-global}. Il paraît raisonnable
de conjecturer que le résultat reste vrai lorsque $X$ et $Y$ sont
de caractéristique mixte (mais toujours excellents).

\begin{lemma}
\label{lemma-structure-cohomology}
Soit $(A, \mathfrak m)$ un anneau local régulier contenant un corps.
Soit $f : V \to \Spec(A)$ un morphisme \'etale et quasi-compact.
Supposons que $\mathfrak m \not \in f(V)$ et que
$g : V \to \Spec(A) \setminus \{\mathfrak m\}$ soit affine.
Alors $H^i(V, \mathcal{O}_V)$, $i > 0$, est isomorphe à une somme
directe de copies de l'enveloppe injective du corps résiduel de $A$.
\end{lemma}

\begin{proof}
Notons $U = \Spec(A) \setminus \{\mathfrak m\}$ le spectre épointé.
Ainsi, $g : V \to U$ est affine.
On a $H^i(V, \mathcal{O}_V) = H^i(U, g_*\mathcal{O}_V)$ d'après
Cohomologie des schémas, lemme \ref{coherent-lemma-relative-affine-cohomology}.
Le $\mathcal{O}_U$-module $g_*\mathcal{O}_V$ est quasi-cohérent d'après
Schémas, lemme \ref{schemes-lemma-push-forward-quasi-coherent}.
Pour tout $\mathcal{O}_U$-module quasi-cohérent $\mathcal{F}$,
la cohomologie $H^i(U, \mathcal{F})$, $i > 0$,
est de $\mathfrak m$-torsion, voir par exemple
Cohomologie locale, lemme \ref{local-cohomology-lemma-local-cohomology}.
En particulier, les $A$-modules $H^i(V, \mathcal{O}_V)$, $i > 0$,
sont de $\mathfrak m$-torsion.
Pour tout homomorphisme plat d'anneaux $A \to A'$, on a
$H^i(V, \mathcal{O}_V) \otimes_A A' = H^i(V', \mathcal{O}_{V'})$
où $V' = V \times_{\Spec(A)} \Spec(A')$, par changement de base plat, d'après
Cohomologie des schémas, lemme \ref{coherent-lemma-flat-base-change-cohomology}.
Si l'on prend pour $A'$ le complété de $A$ (qui est plat d'après
Compléments d'algèbre, section \ref{more-algebra-section-permanence-completion}),
alors on voit que
$$
H^i(V, \mathcal{O}_V) = H^i(V, \mathcal{O}_V) \otimes_A A' =
H^i(V', \mathcal{O}_{V'}),\quad\text{pour } i > 0
$$
La première égalité résulte de la propriété de torsion que nous venons d'établir et du
lemme de Compléments d'algèbre \ref{more-algebra-lemma-neighbourhood-equivalence}.
De plus, l'enveloppe injective du corps résiduel $k$
est la même pour $A$ et $A'$, voir
Complexes dualisants, lemme \ref{dualizing-lemma-compare}.
On se ramène ainsi au cas $A = k[[x_1, \ldots, x_d]]$, voir
Algèbre, section \ref{algebra-section-cohen-structure-theorem}.

\medskip\noindent
Supposons que la caractéristique de $k$ soit $p > 0$. Puisque $F : A \to A$,
$a \mapsto a^p$, est plat (Cohomologie locale, lemme
\ref{local-cohomology-lemma-frobenius-flat-regular})
et puisque $V \times_{\Spec(A), \Spec(F)} \Spec(A) \cong V$ comme schémas
sur $\Spec(A)$ d'après
Morphismes \'etales, lemme \ref{etale-lemma-relative-frobenius-etale},
ce qui précède donne
$H^i(V, \mathcal{O}_V) \otimes_{A, F} A \cong H^i(V, \mathcal{O}_V)$.
Le résultat découle donc du
lemme de Cohomologie locale
\ref{local-cohomology-lemma-structure-torsion-Frobenius-regular}.

\medskip\noindent
Supposons que la caractéristique de $k$ soit $0$. D'après
Cohomologie locale, lemme \ref{local-cohomology-lemma-etale-derivation},
il existe des opérateurs additifs $D_j$, $j = 1, \ldots, d$, sur
$H^i(V, \mathcal{O}_V)$ qui satisfont la règle de Leibniz par
rapport à $\partial_j = \partial/\partial x_j$.
Le résultat découle donc du lemme de Cohomologie locale
\ref{local-cohomology-lemma-structure-torsion-D-module-regular}.
\end{proof}

\begin{lemma}
\label{lemma-conclude}
Dans la situation du lemme \ref{lemma-structure-cohomology},
supposons que $H^i(V, \mathcal{O}_V) = 0$ pour $i \geq \dim(A) - 1$.
Alors $V$ est affine.
\end{lemma}

\begin{proof}
Soit $k = A/\mathfrak m$. Puisque $V \times_{\Spec(A)} \Spec(k) = \emptyset$,
la cohomologie et le changement de base donnent
$$
R\Gamma(V, \mathcal{O}_V) \otimes_A^\mathbf{L} k = 0
$$
Voir
Catégories dérivées de schémas, lemme \ref{perfect-lemma-compare-base-change}.
Il existe donc une suite spectrale
(Compléments d'algèbre, exemple \ref{more-algebra-example-tor})
$$
E_2^{p, q} = \text{Tor}_{-p}(k, H^q(V, \mathcal{O}_V)),\quad
d_2^{p, q} : E_2^{p, q} \to E_2^{p + 2, q - 1}
$$
et $d_r^{p, q} : E_r^{p, q} \to E_r^{p + r, q - r + 1}$,
qui converge vers zéro. D'après le lemme \ref{lemma-structure-cohomology},
le lemme de Complexes dualisants \ref{dualizing-lemma-tor-injective-hull},
et notre hypothèse $H^i(V, \mathcal{O}_V) = 0$ pour $i \geq \dim(A) - 1$,
on conclut qu'aucune différentielle non nulle
n'aboutit à la case $(p, q) = (0, 0)$ ni n'en part. Ainsi,
$H^0(V, \mathcal{O}_V) \otimes_A k = 0$. Cela signifie que,
si $\mathfrak m = (x_1, \ldots, x_d)$, on dispose
d'un recouvrement ouvert $V = \bigcup V \times_{\Spec(A)} \Spec(A_{x_i})$
par les sous-schémas ouverts affines $V \times_{\Spec(A)} \Spec(A_{x_i})$
(car $V$ est affine sur le spectre épointé de $A$)
tel que $x_1, \ldots, x_d$ engendrent
l'idéal unité de $\Gamma(V, \mathcal{O}_V)$.
Cela implique que $V$ est affine d'après
Propriétés, lemme \ref{properties-lemma-characterize-affine}.
\end{proof}

\begin{theorem}
\label{theorem-global}
Soit $Y$ un schéma régulier excellent sur un corps. Soit $f : X \to Y$
un morphisme de schémas de type fini, avec $X$ normal. Soit $V \subset X$
le plus grand sous-schéma ouvert où $f$ est \'etale. Alors le morphisme
d'inclusion $V \to X$ est affine.
\end{theorem}

\begin{proof}
Soit $x \in X$ d'image $y \in Y$. Il suffit de démontrer que
$V \cap W$ est affine pour un voisinage ouvert affine $W$ de $x$.
Puisque $\Spec(\mathcal{O}_{X, x})$ est la limite des schémas $W$,
cela équivaut à dire que
$$
V_x = V \times_X \Spec(\mathcal{O}_{X, x})
$$
est affine (Limites, lemme \ref{limits-lemma-limit-affine}).
Ainsi, si le théorème vaut pour le morphisme
$X \times_Y \Spec(\mathcal{O}_{Y, y}) \to \Spec(\mathcal{O}_{Y, y})$,
alors il vaut dans la situation initiale. En particulier, on peut supposer $Y$
régulier de dimension finie, ce qui permet de raisonner par récurrence
sur la dimension $d = \dim(Y)$. En combinant ceci avec le même argument,
on peut supposer que $Y$ est local, de point fermé $y$, et que
$V \cap (X \setminus f^{-1}(\{y\})) \to X \setminus f^{-1}(\{y\})$
est affine.

\medskip\noindent
Soit $x \in X$ un point au-dessus de $y$. Si $x \in V$, il n'y a
rien à démontrer. Observons que $f^{-1}(\{y\}) \cap V$
est un ensemble fini de points fermés (les fibres d'un morphisme \'etale
sont discrètes). Ainsi, quitte à remplacer $X$
par un voisinage ouvert affine de $x$, on peut supposer que
$y \not \in f(V)$. Il faut démontrer que $V$ est affine.

\medskip\noindent
Soit $e(V)$ le plus grand $i$ tel que $H^i(V, \mathcal{O}_V) \not = 0$.
Comme $X$ est affine, l'entier $e(V)$ est le maximum des nombres $e(V_x)$
pour $x \in X \setminus V$, voir
Cohomologie locale, lemme \ref{local-cohomology-lemma-cd-local},
et la caractérisation de la dimension cohomologique dans
Cohomologie locale, lemme \ref{local-cohomology-lemma-cd}.
On a $e(V_x) \leq \dim(\mathcal{O}_{X, x}) - 1$ d'après
Cohomologie locale, lemme \ref{local-cohomology-lemma-cd-dimension}.
Si $\dim(\mathcal{O}_{X, x}) \geq 2$, la pureté du
lieu de ramification (lemme \ref{lemma-purity-ramification})
montre que $V_x$ est strictement contenu dans le spectre épointé de
$\mathcal{O}_{X, x}$. Puisque $\mathcal{O}_{X, x}$ est
normal et excellent, cela implique
$e(V_x) \leq \dim(\mathcal{O}_{X, x}) - 2$ d'après le théorème
d'annulation de Hartshorne--Lichtenbaum
(Cohomologie locale, lemme \ref{local-cohomology-lemma-affine-complement}).
D'autre part, puisque $X \to Y$ est de type fini
et que $V \subset X$ est dense (quitte à remplacer $X$
par l'adhérence de $V$), on voit que $\dim(\mathcal{O}_{X, x}) \leq d$
d'après la formule des dimensions
(Morphismes, lemme \ref{morphisms-lemma-dimension-formula}).
Par conséquent, $e(V) \leq \max(0, d -  2)$.
Ainsi, $V$ est affine d'après le lemme \ref{lemma-conclude}
si $d \geq 2$. Si $d = 1$ ou $d = 0$, le spectre épointé de
$\mathcal{O}_{Y, y}$ est affine, donc $V$ est affine.
\end{proof}




\section{Homomorphismes de spécialisation dans le cas propre et lisse}
\sectionmark{Spécialisation dans le cas propre et lisse}
\label{section-specialization-smooth-proper}

\noindent
Dans cette section, nous examinons le résultat suivant.
Soit $f : X \to S$ un morphisme propre et lisse de schémas.
Soit $s \leadsto s'$ une spécialisation de points de $S$.
Alors l'homomorphisme de spécialisation
$$
sp : \pi_1(X_{\overline{s}}) \longrightarrow \pi_1(X_{\overline{s}'})
$$
de la section \ref{section-specialization-map}
est surjectif et
\begin{enumerate}
\item si la caractéristique de $\kappa(s')$ est nulle, c'est
un isomorphisme, ou bien
\item si la caractéristique de $\kappa(s')$ est $p > 0$, il
induit un isomorphisme sur les quotients maximaux premiers à $p$.
\end{enumerate}

\begin{lemma}
\label{lemma-specialization-map-surjective}
Soit $f : X \to S$ un morphisme plat et propre à fibres géométriquement
connexes. Soit $s' \leadsto s$ une spécialisation.
Si $X_s$ est géométriquement réduit, alors l'homomorphisme de
spécialisation $sp : \pi_1(X_{\overline{s}'}) \to \pi_1(X_{\overline{s}})$
est surjectif.
\end{lemma}

\begin{proof}
Puisque $X_s$ est géométriquement réduit, on peut supposer que toutes les
fibres sont géométriquement réduites, quitte à rétrécir $S$, voir
Compléments sur les morphismes, lemme \ref{more-morphisms-lemma-geometrically-reduced-open}.
Soit $\mathcal{O}_{S, s} \to A \to \kappa(\overline{s}')$ comme
dans la construction de l'homomorphisme de spécialisation, voir la
section \ref{section-specialization-map}.
Il suffit donc de montrer que
$$
\pi_1(X_{\overline{s}'}) \to \pi_1(X_A)
$$
est surjective. Cela résulte de la
proposition \ref{proposition-first-homotopy-sequence}
et de $\pi_1(\Spec(A)) = \{1\}$.
\end{proof}

\begin{proposition}
\label{proposition-specialization-map-isomorphism}
Soit $f : X \to S$ un morphisme propre et lisse à fibres géométriquement
connexes. Soit $s' \leadsto s$ une spécialisation.
Si la caractéristique de $\kappa(s)$ est nulle, alors l'homomorphisme de
spécialisation
$$
sp : \pi_1(X_{\overline{s}'}) \to \pi_1(X_{\overline{s}})
$$
est un isomorphisme.
\end{proposition}

\begin{proof}
L'homomorphisme est surjectif d'après le
lemme \ref{lemma-specialization-map-surjective}.
Il reste donc à montrer qu'il est injectif.

\medskip\noindent
On peut supposer que $S$ est affine. Alors $S$ est une limite cofiltrante de schémas
affines de type fini sur $\mathbf{Z}$.
On peut donc supposer que $X \to S$ s'obtient par
changement de base à partir de $X_0 \to S_0$, où $S_0$ est le spectre d'une
$\mathbf{Z}$-algèbre de type fini et $X_0 \to S_0$ est propre et lisse.
Voir Limites, lemme \ref{limits-lemma-descend-finite-presentation},
\ref{limits-lemma-descend-smooth}, et
\ref{limits-lemma-eventually-proper}. D'après le
lemme \ref{lemma-specialization-map-base-change},
on se ramène au cas où la base est noethérienne.

\medskip\noindent
En appliquant le lemme \ref{lemma-specialization-map-discrete-valuation-ring},
on se ramène au cas où la base $S$ est le spectre d'un
anneau de valuation discrète strictement hensélien $A$ et où l'on
considère l'homomorphisme de spécialisation sur $A$.
Soit $K$ le corps des fractions de $A$.
Choisissons une clôture algébrique $\overline{K}$ qui
correspond à un point géométrique générique $\overline{\eta}$ de $\Spec(A)$.
Pour toute sous-extension finie séparable $L/K$ de $\overline{K}/K$, soit $B \subset L$ la
clôture intégrale de $A$ dans $L$. C'est un anneau de
valuation discrète d'après
Compléments d'algèbre, remarque \ref{more-algebra-remark-finite-separable-extension}.

\medskip\noindent
Soit $X \to \Spec(A)$ comme dans le paragraphe précédent.
Pour montrer l'injectivité de l'homomorphisme de spécialisation,
il suffit de prouver que tout revêtement fini
\'etale $V$ de $X_{\overline{\eta}}$ provient par changement
de base d'un revêtement fini \'etale $Y \to X$.
En effet, alors $\pi_1(X_{\overline{\eta}}) \to \pi_1(X) = \pi_1(X_s)$
est injective d'après le lemme \ref{lemma-functoriality-galois-injective}.

\medskip\noindent
Étant donné $V$, on peut d'abord descendre $V$ en $V' \to X_{K^{sep}}$ d'après le
lemme \ref{lemma-perfection}, puis en
$V'' \to X_L$ d'après le lemme \ref{lemma-limit}.
Soit $Z \to X_B$ la normalisation de $X_B$ dans $V''$.
Observons que $Z$ est normal et que $Z_L = V''$ comme schémas
sur $X_L$. Ainsi, $Z \to X_B$ est fini \'etale sur
la fibre générique. Le problème est que nous ne savons pas si
$Z \to X_B$ est partout \'etale. Puisque $X \to \Spec(A)$
a des fibres lisses géométriquement connexes, on voit que
la fibre spéciale $X_s$ est géométriquement irréductible.
Ainsi, la fibre spéciale de $X_B \to \Spec(B)$ est irréductible ;
soit $\xi_B$ son point générique. Soient
$\xi_1, \ldots, \xi_r$ les points de $Z$ qui s'envoient sur
$\xi_B$. Notre premier problème (et, comme on le verra, le seul)
est alors que les extensions
$$
\mathcal{O}_{X_B, \xi_B} \subset \mathcal{O}_{Z, \xi_i}
$$
d'anneaux de valuation discrète peuvent être ramifiées. Soit $e_i$
l'indice de ramification de cette extension. Notons que, puisque la
caractéristique de $\kappa(s)$ est nulle, la ramification est modérée !

\medskip\noindent
Pour éliminer la ramification, nous allons choisir une nouvelle extension finie
séparable $K^{sep}/L'/L/K$ telle que l'indice de ramification
$e$ des extensions induites $B'/B$ soit divisible par $e_i$.
Considérons le changement de base normalisé $Z'$ de $Z$ relativement à
$\Spec(B') \to \Spec(B)$, voir la discussion dans
Compléments sur les morphismes, section \ref{more-morphisms-section-reduced-fibre-theorem}.
Soient $\xi_{i, j}$ les points de $Z'$ qui s'envoient sur $\xi_{B'}$
et sur $\xi_i$ dans $Z$. Alors les anneaux locaux
$$
\mathcal{O}_{Z', \xi_{i, j}}
$$
sont des localisés de la clôture intégrale de $\mathcal{O}_{Z, \xi_i}$
dans $L' \otimes_L F_i$, où $F_i$ est le corps des fractions de
$\mathcal{O}_{Z, \xi_i}$ ; les détails sont omis. Le lemme d'Abhyankar
(Compléments d'algèbre, lemme \ref{more-algebra-lemma-abhyankar})
nous dit donc que
$$
\mathcal{O}_{X_{B'}, \xi_{B'}} \subset \mathcal{O}_{Z', \xi_{i, j}}
$$
est non ramifiée. On conclut que le morphisme $Z' \to X_{B'}$
est \'etale en codimension $1$. La pureté du
lieu de branchement (lemme \ref{lemma-purity})
montre donc que $Z' \to X_{B'}$ est fini \'etale !

\medskip\noindent
Cependant, puisque l'extension de corps résiduels induite par $A \to B'$
est triviale (le corps résiduel de $A$ étant algébriquement clos,
car séparablement clos de caractéristique nulle),
on conclut que $Z'$ provient par changement de base d'un revêtement fini \'etale
$Y \to X$ en appliquant deux fois le
lemme \ref{lemma-finite-etale-on-proper-over-henselian}
(d'abord pour obtenir $Y$ sur $A$, puis pour démontrer que
le changement de base à $B'$ est isomorphe à $Z'$).
Ceci achève la démonstration.
\end{proof}

\noindent
Soit $G$ un groupe profini. Soit $p$ un nombre premier.
Le {\it quotient maximal premier à $p$} est par définition
$$
G' = \lim_{U \subset G\text{ ouvert, distingué, d'indice premier à }p} G/U
$$
Si $X$ est un schéma connexe et si $p$ est fixé, le quotient maximal
premier à $p$ de $\pi_1(X)$ est noté $\pi'_1(X)$.

\begin{theorem}
\label{theorem-specialization-map-isomorphism-prime-to-p}
Soit $f : X \to S$ un morphisme propre et lisse à fibres géométriquement
connexes. Soit $s' \leadsto s$ une spécialisation.
Si la caractéristique de $\kappa(s)$ est $p$, alors l'application de
spécialisation
$$
sp : \pi_1(X_{\overline{s}'}) \to \pi_1(X_{\overline{s}})
$$
est surjective et induit un isomorphisme
$$
\pi'_1(X_{\overline{s}'}) \cong \pi'_1(X_{\overline{s}})
$$
des quotients maximaux premiers à $p$.
\end{theorem}

\begin{proof}
Cela se démontre exactement de la même manière que la
proposition \ref{proposition-specialization-map-isomorphism},
avec les différences suivantes :
\begin{enumerate}
\item Étant donné $X/A$, on ne montre plus que le foncteur
$\textit{F\'Et}_X \to \textit{F\'Et}_{X_{\overline{\eta}}}$
est essentiellement surjectif. On montre seulement que les objets galoisiens
dont le groupe de Galois est d'ordre premier à $p$ appartiennent à l'image
essentielle. Cela suffit pour conclure à l'injectivité de
$\pi'_1(X_{\overline{s}'}) \to \pi'_1(X_{\overline{s}})$ par
exactement le même argument.
\item Les extensions
$\mathcal{O}_{X_B, \xi_B} \subset \mathcal{O}_{Z, \xi_i}$
sont modérément ramifiées, car l'extension associée des corps de
fractions est galoisienne, de groupe d'ordre premier à $p$. Voir
Compléments d'algèbre, lemme \ref{more-algebra-lemma-galois-conclusion}.
\item L'extension $\kappa_B/\kappa_A$ n'est plus
nécessairement triviale, mais elle est purement inséparable.
Ainsi, le morphisme $X_{\kappa_B} \to X_{\kappa_A}$
est un homéomorphisme universel et induit un isomorphisme
de groupes fondamentaux d'après la proposition \ref{proposition-universal-homeomorphism}.
\end{enumerate}
\end{proof}













\section{Ramification modérée}
\label{section-tame}

\noindent
Soit $f : X \to Y$ un morphisme fini \'etale de schémas de type fini
sur $\mathbf{Z}$. Il existe de nombreuses façons de définir ce que signifie pour $f$
d'être modérément ramifié à $\infty$. L'article \cite{Kerz-Schmidt}
examine dans quelle mesure ces notions coïncident.

\medskip\noindent
Dans cette section, nous examinons une question différente et plus élémentaire qui
précède la notion de ramification modérée à l'infini.
On comparera avec la discussion (légèrement différente)
de \cite{Grothendieck-Murre}.
Supposons donnés
\begin{enumerate}
\item un schéma localement noethérien $X$,
\item un ouvert dense $U \subset X$,
\item un morphisme fini \'etale $f : Y \to U$
\end{enumerate}
tels que, pour tout diviseur premier $Z \subset X$
tel que $Z \cap U = \emptyset$, l'anneau local $\mathcal{O}_{X, \xi}$
de $X$ au point générique $\xi$ de $Z$ soit un anneau de valuation discrète.
En prenant pour $K_\xi$ le corps des fractions de $\mathcal{O}_{X, \xi}$,
on obtient un carré cartésien
$$
\xymatrix{
\Spec(K_\xi) \ar[r] \ar[d] & U \ar[d] \\
\Spec(\mathcal{O}_{X, \xi}) \ar[r] & X
}
$$
de schémas. En particulier, on voit que $Y \times_U \Spec(K_\xi)$
est le spectre d'une algèbre finie séparable $L_\xi/K_\xi$.
On dit alors que
{\it $Y$ est non ramifié sur $X$ en codimension $1$},
resp.\ {\it $Y$ est modérément ramifié sur $X$ en codimension $1$},
si $L_\xi/K_\xi$ est non ramifiée, resp.\ modérément ramifiée,
relativement à $\mathcal{O}_{X, \xi}$ pour tout $(Z, \xi)$
comme ci-dessus ; voir Compléments d'algèbre, définition
\ref{more-algebra-definition-types-of-extensions}.
Plus précisément, on décompose $L_\xi$ en un produit d'extensions
de corps finies séparables de $K_\xi$ et l'on exige que chacune soit
non ramifiée, resp.\ modérément ramifiée, relativement à
$\mathcal{O}_{X, \xi}$.

\begin{lemma}
\label{lemma-pullback-tame-codim1}
Soit $g : X' \to X$ un morphisme de schémas localement noethériens.
Soit $U \subset X$ un ouvert dense. Supposons que
\begin{enumerate}
\item $U' = g^{-1}(U)$ soit un ouvert dense de $X'$,
\item pour tout diviseur premier $Z \subset X$ tel que $Z \cap U = \emptyset$,
l'anneau local $\mathcal{O}_{X, \xi}$ de $X$ au point générique $\xi$
de $Z$ soit un anneau de valuation discrète,
\item pour tout diviseur premier $Z' \subset X'$
tel que $Z' \cap U' = \emptyset$, l'anneau local $\mathcal{O}_{X', \xi'}$
de $X'$ au point générique $\xi'$ de $Z'$ soit un anneau de valuation discrète,
\item si $\xi' \in X'$ est comme en (3), alors $\xi = g(\xi')$ soit comme en (2).
\end{enumerate}
Alors, si $f : Y \to U$ est fini \'etale et si
$Y$ est non ramifié, resp.\ modérément ramifié, sur $X$
en codimension $1$, alors $Y' = Y \times_X X' \to U'$ est fini \'etale
et $Y'$ est non ramifié, resp.\ modérément ramifié, sur $X'$ en codimension $1$.
\end{lemma}

\begin{proof}
Le seul fait intéressant de ce lemme est le résultat d'algèbre commutative
donné dans Compléments d'algèbre, lemme \ref{more-algebra-lemma-tame-goes-up}.
\end{proof}

\noindent
Avec la terminologie introduite ci-dessus, on peut reformuler les
résultats de pureté obtenus plus haut sous la forme agréable suivante.

\begin{lemma}
\label{lemma-purity-one-divisor-general}
Soit $X$ un schéma localement noethérien. Soit $U \subset X$ ouvert et dense.
Soit $Y \to U$ un morphisme fini \'etale. Supposons que
\begin{enumerate}
\item $Y$ soit non ramifié sur $X$ en codimension $1$, et
\item $\mathcal{O}_{X, x}$ soit régulier pour tout $x \in X \setminus U$.
\end{enumerate}
Alors il existe un morphisme fini \'etale $Y' \to X$
dont la restriction à $U$ est $Y$.
\end{lemma}

\begin{proof}
Soit $\xi \in X \setminus U$ le point générique d'une composante irréductible
de $X \setminus U$ de codimension $1$. Alors $\mathcal{O}_{X, \xi}$ est un
anneau de valuation discrète. Comme dans la discussion ci-dessus, écrivons
$Y \times_U \Spec(K_\xi) = \Spec(L_\xi)$.
Notons $B_\xi$ la clôture intégrale de $\mathcal{O}_{X, \xi}$ dans
$L_\xi$. Notre hypothèse que $Y$ est non ramifié sur $X$ en codimension $1$
signifie que $\mathcal{O}_{X, \xi} \to B_\xi$ est fini \'etale.
On obtient ainsi $Y_\xi \to \Spec(\mathcal{O}_{X, \xi})$ fini
\'etale et un isomorphisme
$$
Y \times_U \Spec(K_\xi) \cong
Y_\xi \times_{\Spec(\mathcal{O}_{X, \xi})} \Spec(K_\xi)
$$
sur $\Spec(K_\xi)$.
D'après Limites, lemme \ref{limits-lemma-glueing-near-point},
on trouve un sous-schéma ouvert $U \subset U' \subset X$
contenant $\xi$ et un morphisme $Y' \to U'$ de présentation finie
dont la restriction à $U$ redonne $Y$ et
dont la restriction à $\Spec(\mathcal{O}_{X, \xi})$ redonne $Y_\xi$.
Enfin, le morphisme $Y' \to U'$ est fini \'etale
après avoir éventuellement remplacé $U'$ par un ouvert plus petit, d'après
Limites, lemme \ref{limits-lemma-glueing-near-point-properties}.
En répétant l'argument avec les autres points génériques de
$X \setminus U$ de codimension $1$, on peut supposer que l'on dispose
d'un morphisme fini \'etale $Y' \to U'$ prolongeant $Y \to U$
à un sous-schéma ouvert $U' \subset X$ contenant $U$
et tous les points de codimension $1$ de $X \setminus U$.
On termine en appliquant le lemme \ref{lemma-extend-pure}
à $Y' \to U'$. En effet, tous les anneaux locaux $\mathcal{O}_{X, x}$
pour $x \in X \setminus U'$ sont réguliers et
vérifient $\dim(\mathcal{O}_{X, x}) \geq 2$. La pureté vaut donc
pour $\mathcal{O}_{X, x}$ d'après le lemme \ref{lemma-local-purity}.
\end{proof}

\begin{lemma}
\label{lemma-purity-one-divisor}
Soit $X$ un schéma localement noethérien. Soit $D \subset X$
un diviseur de Cartier effectif tel que $D$ soit un schéma régulier.
Soit $Y \to X \setminus D$ un morphisme fini \'etale.
Si $Y$ est non ramifié sur $X$ en codimension $1$, alors
il existe un morphisme fini \'etale $Y' \to X$
dont la restriction à $X \setminus D$ est $Y$.
\end{lemma}

\begin{proof}
C'est un cas particulier du lemme \ref{lemma-purity-one-divisor-general}.
Premièrement, $D$ n'est dense nulle part dans $X$ (voir la discussion dans
Diviseurs, section \ref{divisors-section-effective-Cartier-divisors}),
donc $X \setminus D$ est dense dans $X$. Deuxièmement, l'anneau
$\mathcal{O}_{X, x}$ est un anneau local régulier pour tout $x \in D$ d'après
Algèbre, lemme \ref{algebra-lemma-regular-mod-x},
et notre hypothèse selon laquelle $\mathcal{O}_{D, x}$ est régulier.
\end{proof}

\begin{example}[Morphisme standard modérément ramifié]
\label{example-tamely-ramified}
Soit $A$ un anneau noethérien. Soit $f \in A$ un non diviseur de zéro
tel que $A/fA$ soit réduit. Cela implique que $A_\mathfrak p$
est un anneau de valuation discrète d'uniformisante $f$ pour tout
idéal premier minimal $\mathfrak p$ contenant $f$. Soit $e \geq 1$ un entier
inversible dans $A$. Posons
$$
C = A[x]/(x^e - f)
$$
Alors $\Spec(C) \to \Spec(A)$ est un morphisme fini localement libre
qui est \'etale sur le spectre de $A_f$. Le morphisme fini \'etale
$$
\Spec(C_f) \longrightarrow \Spec(A_f)
$$
est modérément ramifié sur $\Spec(A)$ en codimension $1$. La
modération résulte immédiatement de la caractérisation des extensions
modérément ramifiées dans
Compléments d'algèbre, lemme \ref{more-algebra-lemma-characterize-tame}.
\end{example}

\noindent
Voici une version du lemme d'Abhyankar pour les diviseurs réguliers.

\begin{lemma}[Lemme d'Abhyankar pour un diviseur régulier]
\label{lemma-abhyankar-one-divisor}
Soit $X$ un schéma localement noethérien. Soit $D \subset X$
un diviseur de Cartier effectif tel que $D$ soit un schéma régulier.
Soit $Y \to X \setminus D$ un morphisme fini \'etale.
Si $Y$ est modérément ramifié sur $X$ en codimension $1$, alors,
localement pour la topologie \'etale sur $X$, le morphisme $Y \to X \setminus D$ est la restriction
à $X \setminus D$ d'une réunion disjointe finie de morphismes standard modérément ramifiés
tels que ceux décrits dans l'exemple \ref{example-tamely-ramified}.
\end{lemma}

\begin{proof}
Pour tout $x \in X$, nous allons trouver un voisinage \'etale
$(U, u) \to (X, x)$, avec $U = \Spec(A)$, tel que le changement de base
$Y \times_{X \setminus D} (U \setminus D) \to U \setminus D$ soit une réunion disjointe finie de restrictions de morphismes standard
modérément ramifiés comme dans l'exemple \ref{example-tamely-ramified}.
On supposera $x \in D$ ; le cas $x \not \in D$ résulte de
Morphismes \'etales, lemme \ref{etale-lemma-finite-etale-etale-local},
en prenant $e = 1$ et $f = 1$ dans l'exemple \ref{example-tamely-ramified}.

\medskip\noindent
Dans ce paragraphe, on se ramène au cas où $X$ est le spectre d'un
anneau local strictement hensélien et où $x$ est le point fermé. En effet, quitte à rétrécir
$X$, on peut supposer que $X = \Spec(A)$ et que $D \subset X$ est défini par un
non diviseur de zéro $f \in A$. Alors $Y$ est affine, car c'est un revêtement fini
\'etale de $\Spec(A_f)$. Écrivons $Y = \Spec(B)$.
Soit $A^{sh}$ l'hensélisé strict de $\mathcal{O}_{X, x}$.
Puisque $A \to A^{sh}$ est plat et que $x \in D$, on voit que $f$ s'envoie
sur un non diviseur de zéro appartenant à l'idéal maximal de $A^{sh}$.
Observons que $A^{sh}/fA^{sh} = (A/fA)^{sh}$
(Algèbre, lemme \ref{algebra-lemma-quotient-strict-henselization})
est un anneau local régulier, car c'est l'hensélisé strict de $\mathcal{O}_{D, x}$ ;
voir Compléments d'algèbre, lemme \ref{more-algebra-lemma-henselization-regular}.
D'après le lemme \ref{lemma-pullback-tame-codim1},
le changement de base de $Y$ à $\Spec(A^{sh})$ est modérément ramifié en
codimension $1$. Supposons l'assertion démontrée pour le changement de base
de $Y$ à $\Spec(A^{sh})$. Puisque $A^{sh}$ est strictement hensélien, tout
voisinage \'etale possède une section, et l'on conclut que l'on a un isomorphisme
$$
A^{sh} \otimes_A B \cong
\prod\nolimits_{i \in I} A^{sh}_f[x]/(x^{e_i} - f)
$$
où $I$ est un ensemble fini et chaque $e_i$ un entier inversible dans $A^{sh}$.
Écrivons $A^{sh} = \colim A_\lambda$ comme une limite inductive filtrante d'algèbres \'etales
sur $A$, notées $A_\lambda$. L'isomorphisme affiché
descend en un isomorphisme
$$
A_\lambda \otimes_A B \cong \prod\nolimits_{i \in I}
(A_\lambda)_f[x]/(x^{e_i} - f)
$$
pour $\lambda$ assez grand, voir par exemple
Algèbre, lemme \ref{algebra-lemma-colimit-category-fp-algebras}.
Après avoir encore un peu augmenté $\lambda$,
on peut supposer que $e_i$ est inversible dans $A_\lambda$. Alors
$\Spec(A_\lambda) \to \Spec(A)$ est le voisinage \'etale recherché
de $x$.

\medskip\noindent
Supposons que $X = \Spec(A)$, où $A$ est un anneau local strictement hensélien,
et que $x \in X$ corresponde à l'idéal maximal $\mathfrak m$ de $A$.
Soit $f \in \mathfrak m$ le non diviseur de zéro qui définit $D$.
Alors $A/fA$ est régulier et, d'après Algèbre, lemme \ref{algebra-lemma-regular-mod-x},
on voit que $A$ est lui aussi régulier. Nous utiliserons certaines propriétés des anneaux
locaux réguliers, par exemple le fait qu'ils sont des anneaux intègres normaux ; voir
Algèbre, lemmes \ref{algebra-lemma-regular-domain} et
\ref{algebra-lemma-regular-normal}. En particulier, $A$ est intègre.
Comme ci-dessus, on voit que $Y = \Spec(B)$ est affine et que $A_f \to B$
est fini \'etale. Ainsi, $B$ est lui aussi normal
(Algèbre, lemme \ref{algebra-lemma-normal-goes-up}),
et l'on conclut que $Y$ est une réunion disjointe finie de spectres d'anneaux intègres normaux d'après
Algèbre, lemme \ref{algebra-lemma-characterize-reduced-ring-normal}.
On peut donc supposer que $B$ est lui aussi intègre.
Puisque $A$ et $A/fA$ sont intègres, on conclut que $f$ est un élément
premier de $A$ qui engendre un idéal premier de hauteur $1$, $\mathfrak p = fA$.
Observons que $A_\mathfrak p$ est un anneau de valuation discrète d'uniformisante $f$.

\medskip\noindent
Soit $K$ le corps des fractions de $A$ et soit $L$ le corps des
fractions de $B$. L'hypothèse de ramification modérée signifie que
$L$ est modérément ramifié relativement à $A_\mathfrak p$.
Soit $e \geq 1$ le produit des indices de ramification de $L$
sur $A_\mathfrak p$, comme dans
Compléments d'algèbre, remarque \ref{more-algebra-remark-finite-separable-extension}.
Alors $e$ est inversible dans $\kappa(\mathfrak p)$, mais
à ce stade, nous ne savons pas si $e$ est inversible dans $A/fA$
ou dans $A$.

\medskip\noindent
Considérons la $A$-algèbre finie et libre
$$
A' = A[x]/(x^e - f)
$$
Il s'agit d'une extension locale finie d'un anneau local
strictement hensélien ; elle est donc strictement hensélienne.
Observons que $f' = x$ est un non diviseur de zéro dans $A'$ et que
$A'/f'A' \cong A/fA$ est un anneau régulier. Comme précédemment, $A'$
est donc régulier et, a fortiori, intègre. Posons
$B' = B \otimes_A A' = B \otimes_{A_f} A'_{f'}$.
D'après le lemme d'Abhyankar
(Compléments d'algèbre, lemme \ref{more-algebra-lemma-abhyankar}),
on voit que $\Spec(B')$ est non ramifié sur $\Spec(A')$
en codimension $1$. En effet, d'après le lemme \ref{lemma-pullback-tame-codim1},
on voit que $\Spec(B')$ est encore au moins modérément ramifié
sur $\Spec(A')$ en codimension $1$. Mais le lemme d'Abhyankar
nous dit que les indices de ramification sont tous devenus égaux à $1$.
D'après le lemme \ref{lemma-purity-one-divisor}, on conclut que
$\Spec(B') \to \Spec(A'_{f'})$ se prolonge en un morphisme fini \'etale
$\Spec(C) \to \Spec(A')$. Cependant, puisque $A'$ est strictement hensélien,
on conclut que $\Spec(C)$ est une réunion disjointe finie de copies
de $\Spec(A')$. Conclusion : il existe au moins un morphisme
$\Spec(A'_{f'}) \to \Spec(B')$. On conclut qu'il existe
une inclusion $B \subset A'_{f'}$ de $A_f$-algèbres.

\medskip\noindent
Il s'ensuit que l'on a des inclusions $K \subset L \subset K[f^{1/e}]$.
D'après Compléments d'algèbre, lemme \ref{more-algebra-lemma-pull-root-uniformizer},
on conclut que $L$ est égal à $K[f^{1/n}]$ pour un diviseur $n$ de $e$.
La normalité de $B$ implique alors que $B \cong A_f[y]/(y^n - f)$ ; les détails
sont omis.
Cependant, le lieu de ramification de l'homomorphisme $A_f \to A_f[y]/(y^n - f)$
est défini par $ny^{n - 1}$, ce qui implique nécessairement que $n$ est une unité de $A_f$.
Puisque $f$ est un élément premier, cela signifie que $n = u f^s$ pour
une unité $u$ de $A$ et $s \in \mathbf{Z}$. Puisque $n$ s'envoie sur une unité de
$\kappa(\mathfrak p)$, on obtient $s = 0$, ce qui achève la démonstration.
\end{proof}

\begin{lemma}
\label{lemma-extend-tame-covering-normal}
Dans la situation du lemme \ref{lemma-abhyankar-one-divisor},
la normalisation de $X$ dans $Y$ est un morphisme fini localement libre
$\pi : Y' \to X$ tel que
\begin{enumerate}
\item la restriction de $Y'$ à $X \setminus D$ soit isomorphe à $Y$,
\item $D' = \pi^{-1}(D)_{red}$ soit un diviseur de Cartier effectif sur $Y'$, et
\item $D'$ soit un schéma régulier.
\end{enumerate}
De plus, localement pour la topologie \'etale sur $X$, le morphisme $Y' \to X$ est une réunion
disjointe finie de morphismes
$$
\Spec(A[x]/(x^e - f)) \to \Spec(A)
$$
où $A$ est un anneau noethérien, $f \in A$ un non diviseur de zéro tel que
$A/fA$ soit régulier, et $e \geq 1$ est inversible dans $A$.
\end{lemma}

\begin{proof}
Ce n'est qu'un complément au lemme \ref{lemma-abhyankar-one-divisor},
et, en fait, l'énoncé du présent lemme découle presque immédiatement de
la démonstration de ce lemme. Mais on peut également déduire
le présent lemme du résultat du lemme \ref{lemma-abhyankar-one-divisor}.
En effet, la normalisation de $X$ dans $Y$ commute au
changement de base \'etale, voir Compléments sur les morphismes, lemme
\ref{more-morphisms-lemma-normalization-smooth-localization}.
On voit donc que l'on peut démontrer les assertions sur la structure locale
de $Y' \to X$ localement pour la topologie \'etale sur $X$. Ainsi, d'après le
lemme \ref{lemma-abhyankar-one-divisor}, on peut supposer que
$X = \Spec(A)$, où $A$ est un anneau noethérien, que l'on a un
non diviseur de zéro $f\in A$ tel que $A/fA$ soit régulier, et que $Y$
est une réunion disjointe finie de spectres d'anneaux $A_f[x]/(x^e - f)$
où $e$ est inversible dans $A$. Nous omettons la vérification que
la clôture intégrale de $A$ dans $A_f[x]/(x^e - f)$ est
égale à $A' = A[x]/(x^e - f)$. (Pour le voir, montrer que
les localisés de $A'$ aux idéaux premiers au-dessus de $(f)$ sont réguliers.)
Nous omettons les détails.
\end{proof}

\begin{lemma}
\label{lemma-tame-covering-split}
Dans la situation du lemme \ref{lemma-abhyankar-one-divisor},
soit $Y' \to X$ comme dans le lemme \ref{lemma-extend-tame-covering-normal}.
Soit $R$ un anneau de valuation discrète de corps des fractions $K$.
Soit
$$
t : \Spec(R) \to X
$$
un morphisme tel que l'image réciproque schématique
$t^{-1}D$ soit le point fermé réduit de $\Spec(R)$.
\begin{enumerate}
\item Si $t|_{\Spec(K)}$ se relève en un point de $Y$, alors
on obtient un relèvement $t' : \Spec(R) \to Y'$ tel que $Y' \to X$
soit \'etale le long de $t'(\Spec(R))$.
\item Si $\Spec(K) \times_X Y$ est isomorphe à une réunion disjointe
de copies de $\Spec(K)$, alors $Y' \to X$ est fini \'etale
sur un voisinage ouvert de $t(\Spec(R))$.
\end{enumerate}
\end{lemma}

\begin{proof}
En appliquant le critère valuatif de propreté au
morphisme fini $Y' \to X$, on voit que les points à valeurs dans $\Spec(K)$
de $Y$ qui coïncident avec $t|_{\Spec(K)}$ comme morphismes vers $X$
se relèvent de manière unique en morphismes $t' : \Spec(R) \to Y'$.
Ainsi, l'énoncé (1) a un sens.

\medskip\noindent
Choisissons un voisinage \'etale $(U, u) \to (X, t(\mathfrak m_R))$
tel que $U = \Spec(A)$ et que $Y' \times_X U \to U$
admette une description comme dans le lemme \ref{lemma-extend-tame-covering-normal}
pour un certain $f \in A$. Alors $\Spec(R) \times_X U \to \Spec(R)$ est \'etale
et surjectif. Si $R'$ désigne l'anneau local de
$\Spec(R) \times_X U$ situé au-dessus du point fermé de $\Spec(R)$,
alors $R'$ est un anneau de valuation discrète et $R \subset R'$
est une extension non ramifiée d'anneaux de valuation discrète
(Compléments d'algèbre, lemme \ref{more-algebra-lemma-Dedekind-etale-extension}).
L'hypothèse sur $t$ signifie que l'homomorphisme $A \to R'$
correspondant à
$$
\Spec(R') \to \Spec(R) \times_X U \to U
$$
envoie $f$ sur une uniformisante $\pi \in R'$. Supposons maintenant que
$$
Y' \times_X U =
\coprod\nolimits_{i \in I} \Spec(A[x]/(x^{e_i} - f))
$$
pour certains $e_i \geq 1$. Alors on voit que
$$
\Spec(R') \times_U (Y' \times_X U) =
\coprod\nolimits_{i \in I} \Spec(R'[x]/(x^{e_i} - \pi))
$$
Les anneaux $R'[x]/(x^{e_i} - \pi)$ sont des anneaux de valuation discrète
(Compléments d'algèbre, lemme \ref{more-algebra-lemma-pull-root-uniformizer})
et n'admettent donc aucun morphisme vers le corps des fractions de $R'$, sauf si $e_i = 1$.

\medskip\noindent
Démonstration de (1). Dans ce cas, le changement de base de l'application $t' : \Spec(R) \to Y'$
détermine un morphisme correspondant $t'' : \Spec(R') \to Y' \times_X U$
dont l'image doit être contenue dans une composante correspondant à $i \in I$ avec $e_i = 1$
d'après la discussion ci-dessus. On voit donc clairement que $Y' \times_X U \to U$
est \'etale le long de l'image de $t''$. Comme la propriété d'être \'etale
se vérifie après un changement de base \'etale, ceci démontre (1).

\medskip\noindent
Démonstration de (2). Dans ce cas, l'hypothèse implique que $e_i = 1$
pour tout $i \in I$. Ainsi, $Y' \times_X U \to U$ est fini \'etale
et l'on conclut comme ci-dessus.
\end{proof}

\begin{lemma}
\label{lemma-extend-covering}
Soit $S$ un schéma noethérien intègre normal de point générique $\eta$.
Soit $f : X \to S$ un morphisme lisse à fibres géométriquement connexes.
Soit $\sigma : S \to X$ une section de $f$. Soit $Z \to X_\eta$ un
revêtement galoisien fini \'etale (section \ref{section-finite-etale-under-galois})
de groupe $G$, dont l'ordre est inversible sur $S$, tel que
$Z$ possède un point $\kappa(\eta)$-rationnel qui s'envoie sur $\sigma(\eta)$.
Alors il existe un revêtement galoisien fini \'etale $Y \to X$ de groupe $G$
dont la restriction à $X_\eta$ est $Z$.
\end{lemma}

\begin{proof}
Supposons d'abord que $S = \Spec(R)$ soit le spectre d'un anneau de valuation discrète $R$
de point fermé $s \in S$. Alors $X_s$ est un diviseur de Cartier effectif
dans $X$, et $X_s$ est régulier, car c'est un schéma lisse sur un corps.
De plus, la fibre générique $X_\eta$ est le sous-schéma ouvert $X \setminus X_s$.
Il résulte de
Compléments d'algèbre, lemme \ref{more-algebra-lemma-galois-conclusion},
et de l'hypothèse sur $G$ que $Z$ est modérément ramifié
sur $X$ en codimension $1$. Soit $Z' \to X$ comme dans le
lemme \ref{lemma-extend-tame-covering-normal}. Observons que
l'action de $G$ sur $Z$ se prolonge en une action de $G$ sur $Z'$.
D'après le lemme \ref{lemma-tame-covering-split},
on voit que $Z' \to X$ est fini \'etale sur un voisinage
ouvert de $\sigma(s)$.
Puisque $X_s$ est irréductible, cela implique que $Z \to X_\eta$
est non ramifié sur $X$ en codimension $1$.
On obtient alors un morphisme fini \'etale $Y \to X$
dont la restriction à $X_\eta$ est $Z$, d'après le
lemme \ref{lemma-purity-one-divisor}.
Bien sûr, $Y \cong Z'$ (détails omis ; indication : calculer localement pour la topologie \'etale),
et $Y$ est donc un revêtement galoisien de groupe $G$.

\medskip\noindent
Cas général. Soit $U \subset S$ un sous-schéma ouvert maximal
tel qu'il existe un revêtement galoisien fini \'etale
$Y \to X \times_S U$ de groupe $G$
dont la restriction à $X_\eta$ soit isomorphe à $Z$.
Supposons que $U \not = S$ afin d'obtenir une contradiction.
Soit $s \in S \setminus U$ un point générique d'une composante irréductible
de $S \setminus U$. Alors l'image inverse $U_s$
de $U$ dans $\Spec(\mathcal{O}_{S, s})$
est le spectre épointé de $\mathcal{O}_{S, s}$.
Nous affirmons que $Y \times_S U_s \to X \times_S U_s$
est la restriction d'un revêtement galoisien fini \'etale
$Y'_s \to X \times_S \Spec(\mathcal{O}_{S, s})$
de groupe $G$.

\medskip\noindent
Montrons d'abord que l'affirmation donne la contradiction recherchée.
D'après Limites, lemme \ref{limits-lemma-glueing-near-point},
on trouve un sous-schéma ouvert $U \subset U' \subset S$
contenant $s$ et un morphisme $Y'' \to U'$ de présentation finie
dont la restriction à $U$ redonne $Y \to U$ et
dont la restriction à $\Spec(\mathcal{O}_{S, s})$ redonne $Y'_s$.
De plus, par l'équivalence de catégories donnée dans le
lemme, on peut supposer, après avoir rétréci $U'$,
qu'il existe un morphisme $Y'' \to U' \times_S X$
et une action de $G$ sur $Y''$ au-dessus de $U' \times_S X$,
compatibles avec les morphismes et actions donnés après
changement de base à $U$ et à $\Spec(\mathcal{O}_{S, s})$.
Après avoir encore rétréci $U'$ si nécessaire, on peut
supposer que $Y'' \to U' \times_S X$ est fini \'etale, voir
Limites, lemme \ref{limits-lemma-glueing-near-point-properties}.
Cela signifie que nous avons trouvé un ouvert strictement plus grand de $S$ sur lequel
$Y$ se prolonge en un revêtement galoisien fini \'etale de groupe $G$,
ce qui donne la contradiction recherchée.

\medskip\noindent
Démonstration de l'affirmation. On peut remplacer, et l'on remplace, $S$ par $\Spec(\mathcal{O}_{S, s})$.
Alors $S = \Spec(A)$, où $(A, \mathfrak m)$ est un anneau intègre local normal.
De plus, $U \subset S$ est le spectre épointé et l'on dispose d'un revêtement
galoisien fini \'etale $Y \to X \times_S U$ de groupe $G$.
Si $\dim(A) = 1$, on peut construire le prolongement de $Y$ en
un revêtement galoisien de $X$ par le premier paragraphe de la démonstration.
On peut donc supposer $\dim(A) \geq 2$, et donc $\text{depth}(A) \geq 2$,
puisque $S$ est normal, voir Algèbre, lemme \ref{algebra-lemma-criterion-normal}.
Puisque $X \to S$ est plat, on conclut que
$\text{depth}(\mathcal{O}_{X, x}) \geq 2$ pour tout point $x \in X$
qui s'envoie sur $s$, voir
Algèbre, lemme \ref{algebra-lemma-apply-grothendieck}.
Soit
$$
Y' \longrightarrow X
$$
le morphisme fini construit dans le lemme \ref{lemma-extend-S2}
à partir de $Y \to X \times_S U$. Observons que l'on obtient une action canonique
de $G$ sur $Y'$. Il reste donc seulement à montrer que $Y'$
est \'etale sur $X$. En fait, d'après le
lemme \ref{lemma-purity-smooth-over-depth2} (par exemple),
il suffit même de montrer que $Y' \to X$ est \'etale au-dessus du
point générique (unique) de la fibre $X_s$. Nous le faisons par
un calcul local dans un voisinage (formel) de $\sigma(s)$.

\medskip\noindent
Choisissons un ouvert affine $\Spec(B) \subset X$ contenant $\sigma(s)$.
Alors $A \to B$ est un homomorphisme d'anneaux lisse qui admet une section $\sigma : B \to A$.
Notons $I = \Ker(\sigma)$ et notons $B^\wedge$ le
complété $I$-adique de $B$. Alors $B^\wedge \cong A[[x_1, \ldots, x_d]]$
pour un certain $d \geq 0$, voir
Algèbre, lemme \ref{algebra-lemma-section-smooth}.
Bien sûr, $B \to B^\wedge$ est plat
(Algèbre, lemme \ref{algebra-lemma-completion-flat}),
et l'image de $\Spec(B^\wedge) \to X$ contient le point générique de $X_s$.
Soit $V \subset \Spec(B^\wedge)$ l'image inverse de $U$.
Considérons le morphisme fini \'etale
$$
W = Y \times_{(X \times_S U)} V \longrightarrow V
$$
Par la compatibilité de la construction de $Y'$
au changement de base plat dans le lemme \ref{lemma-extend-S2},
on voit que le changement de base
$Y' \times_X \Spec(B^\wedge) \to \Spec(B^\wedge)$
est construit à partir de $W \to V$ sur $\Spec(B^\wedge)$
par le procédé du lemme \ref{lemma-extend-S2}.
Posons $V_0 = V \cap V(x_1, \ldots, x_d) \subset V$ et $W_0 = W \times_V V_0$.
Le schéma $V_0$ est intègre normal, s'envoie dans $\sigma(S)$
par le morphisme $\Spec(B^\wedge) \to X$ et s'identifie en fait
à $\sigma(U)$. On sait donc que $W_0 \to V_0 = U$ se scinde complètement,
car c'est le cas de sa fibre générique d'après notre hypothèse
selon laquelle $Z \to X_\eta$ possède un point $\kappa(\eta)$-rationnel au-dessus de
$\sigma(\eta)$ (et, bien sûr, l'action de $G$ implique alors que toute la
fibre $Z_{\sigma(\eta)}$ est une réunion disjointe de copies du
schéma $\eta = \Spec(\kappa(\eta))$).
Enfin, d'après le
lemme \ref{lemma-fully-faithful-power-series-over-depth2},
on a
$$
W_0 \times_U V \cong W
$$
Cela montre que $W$ est une réunion disjointe de copies de $V$
et donc que $Y' \times_X \Spec(B^\wedge)$ est une réunion disjointe
de copies de $\Spec(B^\wedge)$, ce qui achève la démonstration.
\end{proof}

\begin{lemma}
\label{lemma-extend-covering-general}
Soit $S$ un schéma intègre normal quasi-compact et quasi-séparé
de point générique $\eta$. Soit $f : X \to S$ un morphisme lisse quasi-compact et
quasi-séparé à fibres géométriquement connexes.
Soit $\sigma : S \to X$ une section de $f$. Soit $Z \to X_\eta$ un
revêtement galoisien fini \'etale (section \ref{section-finite-etale-under-galois})
de groupe $G$, dont l'ordre est inversible sur $S$, tel que
$Z$ possède un point $\kappa(\eta)$-rationnel qui s'envoie sur $\sigma(\eta)$.
Alors il existe un revêtement galoisien fini \'etale $Y \to X$ de groupe $G$
dont la restriction à $X_\eta$ est $Z$.
\end{lemma}

\begin{proof}
Si $S$ est noethérien, c'est le résultat du
lemme \ref{lemma-extend-covering}. Le cas général
en découle par un argument standard de passage à la limite.
Nous recommandons vivement au lecteur de sauter la démonstration.

\medskip\noindent
On peut écrire $S = \lim S_i$ comme limite projective d'un système filtrant
de schémas dont les morphismes de transition sont affines et dont les $S_i$ sont
de type fini sur $\mathbf{Z}$, voir
Limites, proposition \ref{limits-proposition-approximate}.
Pour chaque $i$, soit $S \to S'_i \to S_i$ la normalisation
de $S_i$ dans $S$, voir
Morphismes, section \ref{morphisms-section-normalization-X-in-Y}.
En combinant Algèbre, proposition \ref{algebra-proposition-ubiquity-nagata},
avec Morphismes, lemmes \ref{morphisms-lemma-nagata-normalization-finite} et
\ref{morphisms-lemma-normal-normalization},
on conclut que $S'_i$ est de type fini sur $\mathbf{Z}$,
fini sur $S_i$, et que $S'_i$ est un schéma intègre normal
tel que $S \to S'_i$ soit dominant.
D'après Morphismes, lemme \ref{morphisms-lemma-functoriality-normalization},
on obtient des morphismes de transition $S'_{i'} \to S'_i$ compatibles
avec les morphismes de transition $S_{i'} \to S_i$ et avec
les morphismes de source $S$.
Nous affirmons que $S = \lim S'_i$. Démonstration de l'affirmation omise (indication :
considérer les ouverts affines situés au-dessus d'un ouvert affine choisi de $S_i$
pour un certain $i$, afin de traduire ceci en un simple problème
d'algèbre). On conclut que l'on peut écrire $S = \lim S_i$
comme limite projective d'un système filtrant de schémas intègres normaux $S_i$
dont les morphismes de transition sont affines et dont les $S_i$ sont
de type fini sur $\mathbf{Z}$.

\medskip\noindent
Pour un certain $i$, on peut trouver un morphisme lisse $X_i \to S_i$
de présentation finie dont le changement de base à $S$ est $X \to S$.
Voir Limites, lemmes
\ref{limits-lemma-descend-finite-presentation} et
\ref{limits-lemma-descend-smooth}.
En augmentant $i$, on peut supposer que la section
$\sigma$ se relève en une section $\sigma_i : S_i \to X_i$
(par l'équivalence de catégories de
Limites, lemme \ref{limits-lemma-descend-finite-presentation}).
On peut remplacer $X_i$ par son sous-schéma ouvert $X_i^0$
étudié dans Compléments sur les morphismes, section
\ref{more-morphisms-section-connected-components}
car l'image de $X \to X_i$ s'y envoie manifestement
(l'ouverture résulte de Compléments sur les morphismes, lemme
\ref{more-morphisms-lemma-connected-along-section-open}).
On peut donc supposer que les fibres de $X_i \to S_i$ sont
géométriquement connexes.
En augmentant $i$, on peut supposer que $|G|$ est inversible
sur $S_i$.
Soit $\eta_i \in S_i$ le point générique.
Puisque $X_\eta$ est la limite des schémas
$X_{i, \eta_i}$, on peut employer exactement les mêmes arguments
pour faire descendre $Z \to X_\eta$ en un revêtement galoisien fini \'etale
$Z_i \to X_{i, \eta_i}$, quitte à augmenter $i$.
Voir le lemme \ref{lemma-limit}.
Quitte à augmenter encore une fois $i$, on peut
supposer que $Z_i$ possède un point $\kappa(\eta_i)$-rationnel
s'envoyant sur $\sigma_i(\eta_i)$.
On applique alors le lemme dans le cas noethérien
et on effectue le changement de base à $X$ pour conclure.
\end{proof}











\section{Astuces en caractéristique positive}
\label{section-achinger}

\noindent
Dans l'article de Piotr Achinger \cite{Achinger}, il est montré qu'un
schéma affine en caractéristique positive est toujours un $K(\pi, 1)$.
Dans cette section, nous expliquons les parties les plus élémentaires de \cite{Achinger}.
Plus précisément, nous montrons que, pour un corps $k$ de caractéristique positive,
un schéma affine \'etale sur $\mathbf{A}^n_k$ est en fait fini \'etale
sur $\mathbf{A}^n_k$ (par un autre morphisme). Nous montrons aussi
qu'une immersion fermée entre schémas affines connexes
en caractéristique positive induit un homomorphisme injectif entre les
groupes fondamentaux \'etales.

\medskip\noindent
Soit $k$ un corps de caractéristique $p > 0$.
Considérons
$$
k[x_1, \ldots, x_n] \longrightarrow A
$$
une surjection de $k$-algèbres de type fini dont la source est
l'algèbre de polynômes en $x_1, \ldots, x_n$. Notons
$I \subset k[x_1, \ldots, x_n]$ le noyau, de sorte que l'on ait
$A = k[x_1, \ldots, x_n]/I$. Nous ne supposons pas que $A$ soit non nul
(autrement dit, nous autorisons le cas où $A$ est l'anneau nul
et $I = k[x_1, \ldots, x_n]$). Enfin, nous supposons donné un
homomorphisme d'anneaux fini \'etale $\pi : A \to B$.

\medskip\noindent
Supposons donnés $k, n, k[x_1, \ldots, x_n] \to A, I, \pi : A \to B$.
Soit $C$ une $k$-algèbre. Considérons les
diagrammes commutatifs
$$
\xymatrix{
& B \\
C \ar[r] & C/\varphi(I)C \ar[u]^\tau \\
k[x_1, \ldots, x_n] \ar[u]^\varphi \ar[r] &
A \ar[u] \ar@/_3em/[uu]_\pi
}
$$
où $\varphi$ est un homomorphisme \'etale de $k$-algèbres et $\tau$ un
homomorphisme surjectif de $k$-algèbres.
Fixons $C, \varphi, \tau$. Pour tout $r \geq 0$ et tous
$y_1, \ldots, y_r \in C$ qui engendrent $C$ comme algèbre sur
$\Im(\varphi)$, soit $s = s(r, y_1, \ldots, y_r) \in \{0, \ldots, r\}$
le plus grand entier tel que $y_i$ soit entier sur $\Im(\varphi)$
pour $1 \leq i \leq s$. On définit $NF(C, \varphi, \tau)$
comme la valeur minimale de
$r - s = r - s(r, y_1, \ldots, y_r)$ pour tous les choix de
$r$ et $y_1, \ldots, y_r$ comme ci-dessus.
Observons que $NF(C, \varphi, \tau)$ vaut $0$ si et seulement si $\varphi$ est fini.

\begin{lemma}
\label{lemma-lower-invariant}
Dans la situation ci-dessus, si $NF(C, \varphi, \tau) > 0$, alors il existe
un homomorphisme \'etale de $k$-algèbres $\varphi'$ et un homomorphisme surjectif de $k$-algèbres
$\tau'$ s'insérant dans le diagramme commutatif
$$
\xymatrix{
& B \\
C \ar[r] & C/\varphi'(I)C \ar[u]_{\tau'} \\
k[x_1, \ldots, x_n] \ar[u]^{\varphi'} \ar[r] &
A \ar[u] \ar@/_3em/[uu]_\pi
}
$$
tel que $NF(C, \varphi', \tau') < NF(C, \varphi, \tau)$.
\end{lemma}

\begin{proof}
Choisissons $r \geq 0$ et $y_1, \ldots, y_r \in C$ qui engendrent
$C$ sur $\Im(\varphi)$, et soit $0 \leq s \leq r$ tel que
$y_1, \ldots, y_s$ soient entiers sur $\Im(\varphi)$
et que $r - s = NF(C, \varphi, \tau) > 0$.
Comme $B$ est fini sur $A$, l'image de $y_{s + 1}$ dans $B$
annule un polynôme unitaire à coefficients dans $A$. On peut donc trouver
un entier $d \geq 1$ et des éléments $f_1, \ldots, f_d \in k[x_1, \ldots, x_n]$
tels que
$$
z = y_{s + 1}^d + \varphi(f_1) y_{s + 1}^{d - 1} + \ldots + \varphi(f_d) \in
J = \Ker(C \to C/\varphi(I)C \xrightarrow{\tau} B)
$$
Comme $\varphi : k[x_1, \ldots, x_n] \to C$ est \'etale, il existe un
polynôme non nul et non constant $g \in k[T_1, \ldots, T_{n + 1}]$
tel que
$$
g(\varphi(x_1), \ldots, \varphi(x_n), z) = 0
\quad\text{dans}\quad
C
$$
Pour le voir, on peut utiliser par exemple le fait que
$C \otimes_{\varphi, k[x_1, \ldots, x_n]} k(x_1, \ldots, x_n)$
est un produit fini d'extensions de corps finies séparables
de $k(x_1, \ldots, x_n)$
(voir Algèbre, lemme \ref{algebra-lemma-etale-over-field})
et que, par conséquent, $z$ annule un polynôme unitaire
à coefficients dans $k(x_1, \ldots, x_n)$. En chassant les dénominateurs,
on obtient $g$.

\medskip\noindent
L'existence de $g$ et le
lemme \ref{algebra-lemma-helper-polynomial} d'Algèbre
fournissent des entiers $e_1, e_2, \ldots, e_n \geq 1$ tels que
$z$ soit entier sur le sous-anneau $C'$ de $C$ engendré par
$t_1 = \varphi(x_1) + z^{pe_1}, \ldots, t_n = \varphi(x_n) + z^{pe_n}$.
Bien entendu, les éléments $\varphi(x_1), \ldots, \varphi(x_n)$
sont eux aussi entiers sur $C'$, de même que les éléments
$y_1, \ldots, y_s$. Enfin, par notre choix de $z$, l'élément
$y_{s + 1}$ est lui aussi entier sur $C'$.

\medskip\noindent
Considérons l'homomorphisme d'anneaux
$$
\varphi' : k[x_1, \ldots, x_n] \longrightarrow C, \quad
x_i \longmapsto t_i
$$
dont l'image est $C'$. Comme
$\text{d}(\varphi(x_i)) = \text{d}(t_i) = \text{d}(\varphi'(x_i))$
dans $\Omega_{C/k}$ (et c'est ici que nous utilisons l'hypothèse que la caractéristique de $k$
est $p > 0$), on conclut que $\varphi'$ est \'etale puisque $\varphi$ est
\'etale, voir
Algèbre, lemme \ref{algebra-lemma-characterize-etale-over-polynomial-ring}.
Observons que $\varphi'(x_i) - \varphi(x_i) = t_i - \varphi(x_i) = z^{pe_i}$
appartient au noyau $J$ de l'application $C \to C/\varphi(I)C \to B$ par notre
choix de $z$ comme élément de $J$.
Par conséquent, pour $f \in I$, l'élément
$$
\varphi'(f) =
f(t_1, \ldots, t_n) =
f(\varphi(x_1) + z^{pe_1}, \ldots, \varphi(x_n) + z^{pe_n}) =
\varphi(f) + \text{élément de }(z)
$$
appartient également à $J$. Autrement dit, $\varphi'(I)C \subset J$ et
nous obtenons une surjection
$$
\tau' : C/\varphi'(I)C \longrightarrow C/J \cong B
$$
d'algèbres \'etales sur $A$. Enfin, l'algèbre
$C$ est engendrée par les éléments
$\varphi(x_1), \ldots, \varphi(x_n), y_1, \ldots, y_r$
sur $C' = \Im(\varphi')$, avec
$\varphi(x_1), \ldots, \varphi(x_n), y_1, \ldots, y_{s + 1}$
entiers sur $C' = \Im(\varphi')$. Ainsi
$NF(C, \varphi', \tau') < r - s = NF(C, \varphi, \tau)$.
Ceci achève la démonstration.
\end{proof}

\begin{lemma}
\label{lemma-affine-etale-over-affine-space}
Soient $k$ un corps de caractéristique $p > 0$ et $X \to \mathbf{A}^n_k$ un
morphisme \'etale, où $X$ est affine. Alors il existe un morphisme fini \'etale
$X \to \mathbf{A}^n_k$.
\end{lemma}

\begin{proof}
Écrivons $X = \Spec(C)$. Posons $A = 0$ et notons $I = k[x_1, \ldots, x_n]$.
Par hypothèse, il existe un homomorphisme de $k$-algèbres \'etale
$\varphi : k[x_1, \ldots, x_n] \to C$. Notons $\tau : C/\varphi(I)C \to 0$
l'unique surjection. Nous pouvons choisir $\varphi$ et $\tau$ de telle sorte
que $NF(C, \varphi, \tau)$ soit minimal. Par le lemme \ref{lemma-lower-invariant},
nous obtenons $NF(C, \varphi, \tau) = 0$. Ainsi $\varphi$ est fini \'etale.
\end{proof}

\begin{lemma}
\label{lemma-dominate-affine-space}
Soient $k$ un corps de caractéristique $p > 0$ et $Z \subset \mathbf{A}^n_k$
un sous-schéma fermé. Soit $Y \to Z$ un morphisme fini \'etale. Il existe un
morphisme fini \'etale $f : U \to \mathbf{A}^n_k$ tel qu'
il existe une immersion ouverte et fermée $Y \to f^{-1}(Z)$ au-dessus de $Z$.
\end{lemma}

\begin{proof}
Traduisons le problème en algèbre. Écrivons
$\mathbf{A}^n_k = \Spec(k[x_1, \ldots, x_n])$.
Alors $Z = \Spec(A)$, où $A = k[x_1, \ldots, x_n]/I$
pour un certain idéal $I \subset k[x_1, \ldots, x_n]$.
Écrivons $Y = \Spec(B)$, de sorte que $Y \to Z$ corresponde à
l'homomorphisme fini \'etale de $k$-algèbres $A \to B$.

\medskip\noindent
Par Algèbre, lemme \ref{algebra-lemma-lift-etale},
il existe un homomorphisme d'anneaux \'etale
$$
\varphi : k[x_1, \ldots, x_n] \to C
$$
et un homomorphisme surjectif de $A$-algèbres $\tau : C/\varphi(I)C \to B$.
(Nous pouvons même choisir $C, \varphi, \tau$ de sorte que $\tau$ soit un isomorphisme,
mais nous ne nous en servirons pas.) Nous pouvons choisir $\varphi$ et $\tau$ de telle sorte
que $NF(C, \varphi, \tau)$ soit minimal. Par le lemme \ref{lemma-lower-invariant},
nous obtenons $NF(C, \varphi, \tau) = 0$. Ainsi $\varphi$ est fini \'etale.

\medskip\noindent
Soit $f : U = \Spec(C) \to \mathbf{A}^n_k$ le morphisme fini \'etale
correspondant à $\varphi$. Le morphisme
$Y \to f^{-1}(Z) = \Spec(C/\varphi(I)C)$ induit par $\tau$
est une immersion fermée, car $\tau$ est surjective, et une immersion ouverte, car
il s'agit d'un morphisme \'etale d'après Morphismes, lemme
\ref{morphisms-lemma-etale-permanence}. Ceci achève la démonstration.
\end{proof}

\noindent
Voici le résultat principal.

\begin{proposition}
\label{proposition-injective}
Soient $p$ un nombre premier et $i : Z \to X$ une immersion fermée
de schémas affines connexes sur $\mathbf{F}_p$. Pour tout point géométrique
$\overline{z}$ de $Z$, l'application
$$
\pi_1(Z, \overline{z}) \to \pi_1(X, \overline{z})
$$
est injective.
\end{proposition}

\begin{proof}
Soit $Y \to Z$ un morphisme fini \'etale. Il suffit de construire
un morphisme fini \'etale $f : U \to X$ tel que $Y$ soit isomorphe à
un sous-schéma ouvert et fermé de $f^{-1}(Z)$, voir
le lemme \ref{lemma-functoriality-galois-injective}.
Écrivons $Z = \Spec(A)$ et $X = \Spec(R)$, de sorte que l'immersion fermée
$Z \to X$ soit donnée par une surjection $R \to A$. Nous pouvons écrire
$A = \colim A_i$ comme la limite inductive filtrante de ses sous-$\mathbf{F}_p$-algèbres
de type fini. Par le lemme \ref{lemma-limit},
nous pouvons trouver un $i$ et un morphisme fini \'etale $Y_i \to Z_i = \Spec(A_i)$
tel que $Y = Z \times_{Z_i} Y_i$.

\medskip\noindent
Choisissons une surjection $\mathbf{F}_p[x_1, \ldots, x_n] \to A_i$.
Elle détermine une immersion fermée
$$
Z_i = \Spec(A_i)
\longrightarrow
X_i = \mathbf{A}^n_{\mathbf{F}_p} = \Spec(\mathbf{F}_p[x_1, \ldots, x_n])
$$
Par la propriété universelle des algèbres de polynômes et puisque
$R \to A$ est surjective, nous pouvons trouver un diagramme commutatif
$$
\xymatrix{
\mathbf{F}_p[x_1, \ldots, x_n] \ar[r] \ar[d] &
A_i \ar[d] \\
R \ar[r] &
A
}
$$
de $\mathbf{F}_p$-algèbres. Nous avons donc un diagramme commutatif
$$
\xymatrix{
Y_i \ar[r] &
Z_i \ar[r] &
X_i \\
Y \ar[u] \ar[r] &
Z \ar[u] \ar[r] &
X \ar[u]
}
$$
dont le carré de gauche est cartésien. Il est clair que si nous pouvons
trouver $f_i : U_i \to X_i$ fini \'etale tel que $Y_i$ soit
isomorphe à un sous-schéma ouvert et fermé de $f_i^{-1}(Z_i)$,
alors le changement de base $f : U \to X$ de $f_i$ par $X \to X_i$
fournit une solution à notre problème. Nous concluons donc en appliquant
le lemme \ref{lemma-dominate-affine-space} à
$Y_i \to Z_i \to X_i = \mathbf{A}^n_{\mathbf{F}_p}$.
\end{proof}











\begin{multicols}{2}[\section{Autres chapitres}]
\noindent
Préliminaires
\begin{enumerate}
\item \hyperref[introduction-section-phantom]{Introduction}
\item \hyperref[conventions-section-phantom]{Conventions}
\item \hyperref[sets-section-phantom]{Théorie des ensembles}
\item \hyperref[categories-section-phantom]{Catégories}
\item \hyperref[topology-section-phantom]{Topologie}
\item \hyperref[sheaves-section-phantom]{Faisceaux sur les espaces}
\item \hyperref[sites-section-phantom]{Sites et faisceaux}
\item \hyperref[stacks-section-phantom]{Champs}
\item \hyperref[fields-section-phantom]{Corps}
\item \hyperref[algebra-section-phantom]{Algèbre commutative}
\item \hyperref[brauer-section-phantom]{Groupes de Brauer}
\item \hyperref[homology-section-phantom]{Algèbre homologique}
\item \hyperref[derived-section-phantom]{Catégories dérivées}
\item \hyperref[simplicial-section-phantom]{Méthodes simpliciales}
\item \hyperref[more-algebra-section-phantom]{Compléments d'algèbre}
\item \hyperref[smoothing-section-phantom]{Lissage des morphismes d'anneaux}
\item \hyperref[modules-section-phantom]{Faisceaux de modules}
\item \hyperref[sites-modules-section-phantom]{Modules sur les sites}
\item \hyperref[injectives-section-phantom]{Objets injectifs}
\item \hyperref[cohomology-section-phantom]{Cohomologie des faisceaux}
\item \hyperref[sites-cohomology-section-phantom]{Cohomologie sur les sites}
\item \hyperref[dga-section-phantom]{Algèbre différentielle graduée}
\item \hyperref[dpa-section-phantom]{Algèbre à puissances divisées}
\item \hyperref[sdga-section-phantom]{Faisceaux différentiels gradués}
\item \hyperref[hypercovering-section-phantom]{Hyperrecouvrements}
\end{enumerate}
Schémas
\begin{enumerate}
\setcounter{enumi}{25}
\item \hyperref[schemes-section-phantom]{Schémas}
\item \hyperref[constructions-section-phantom]{Constructions de schémas}
\item \hyperref[properties-section-phantom]{Propriétés des schémas}
\item \hyperref[morphisms-section-phantom]{Morphismes de schémas}
\item \hyperref[coherent-section-phantom]{Cohomologie des schémas}
\item \hyperref[divisors-section-phantom]{Diviseurs}
\item \hyperref[limits-section-phantom]{Limites de schémas}
\item \hyperref[varieties-section-phantom]{Variétés}
\item \hyperref[topologies-section-phantom]{Topologies sur les schémas}
\item \hyperref[descent-section-phantom]{Descente}
\item \hyperref[perfect-section-phantom]{Catégories dérivées de schémas}
\item \hyperref[more-morphisms-section-phantom]{Compléments sur les morphismes}
\item \hyperref[flat-section-phantom]{Compléments sur la platitude}
\item \hyperref[groupoids-section-phantom]{Schémas en groupoïdes}
\item \hyperref[more-groupoids-section-phantom]{Compléments sur les schémas en groupoïdes}
\item \hyperref[etale-section-phantom]{Morphismes étales de schémas}
\end{enumerate}
Thèmes de théorie des schémas
\begin{enumerate}
\setcounter{enumi}{41}
\item \hyperref[chow-section-phantom]{Homologie de Chow}
\item \hyperref[intersection-section-phantom]{Théorie de l'intersection}
\item \hyperref[pic-section-phantom]{Schémas de Picard des courbes}
\item \hyperref[weil-section-phantom]{Théories cohomologiques de Weil}
\item \hyperref[adequate-section-phantom]{Modules adéquats}
\item \hyperref[dualizing-section-phantom]{Complexes dualisants}
\item \hyperref[duality-section-phantom]{Dualité pour les schémas}
\item \hyperref[discriminant-section-phantom]{Discriminants et différentes}
\item \hyperref[derham-section-phantom]{Cohomologie de de Rham}
\item \hyperref[local-cohomology-section-phantom]{Cohomologie locale}
\item \hyperref[algebraization-section-phantom]{Géométrie algébrique et formelle}
\item \hyperref[curves-section-phantom]{Courbes algébriques}
\item \hyperref[resolve-section-phantom]{Résolution des surfaces}
\item \hyperref[models-section-phantom]{Réduction semi-stable}
\item \hyperref[functors-section-phantom]{Foncteurs et morphismes}
\item \hyperref[equiv-section-phantom]{Catégories dérivées de variétés}
\item \hyperref[pione-section-phantom]{Groupes fondamentaux de schémas}
\item \hyperref[etale-cohomology-section-phantom]{Cohomologie étale}
\item \hyperref[crystalline-section-phantom]{Cohomologie cristalline}
\item \hyperref[proetale-section-phantom]{Cohomologie pro-étale}
\item \hyperref[relative-cycles-section-phantom]{Cycles relatifs}
\item \hyperref[more-etale-section-phantom]{Compléments de cohomologie étale}
\item \hyperref[trace-section-phantom]{La formule des traces}
\end{enumerate}
Espaces algébriques
\begin{enumerate}
\setcounter{enumi}{64}
\item \hyperref[spaces-section-phantom]{Espaces algébriques}
\item \hyperref[spaces-properties-section-phantom]{Propriétés des espaces algébriques}
\item \hyperref[spaces-morphisms-section-phantom]{Morphismes d'espaces algébriques}
\item \hyperref[decent-spaces-section-phantom]{Espaces algébriques décents}
\item \hyperref[spaces-cohomology-section-phantom]{Cohomologie des espaces algébriques}
\item \hyperref[spaces-limits-section-phantom]{Limites d'espaces algébriques}
\item \hyperref[spaces-divisors-section-phantom]{Diviseurs sur les espaces algébriques}
\item \hyperref[spaces-over-fields-section-phantom]{Espaces algébriques sur des corps}
\item \hyperref[spaces-topologies-section-phantom]{Topologies sur les espaces algébriques}
\item \hyperref[spaces-descent-section-phantom]{Descente et espaces algébriques}
\item \hyperref[spaces-perfect-section-phantom]{Catégories dérivées d'espaces}
\item \hyperref[spaces-more-morphisms-section-phantom]{Compléments sur les morphismes d'espaces}
\item \hyperref[spaces-flat-section-phantom]{Platitude sur les espaces algébriques}
\item \hyperref[spaces-groupoids-section-phantom]{Groupoïdes dans les espaces algébriques}
\item \hyperref[spaces-more-groupoids-section-phantom]{Compléments sur les groupoïdes dans les espaces}
\item \hyperref[bootstrap-section-phantom]{Amorçage}
\item \hyperref[spaces-pushouts-section-phantom]{Sommes amalgamées d'espaces algébriques}
\end{enumerate}
Thèmes de géométrie
\begin{enumerate}
\setcounter{enumi}{81}
\item \hyperref[spaces-chow-section-phantom]{Groupes de Chow des espaces}
\item \hyperref[groupoids-quotients-section-phantom]{Quotients de groupoïdes}
\item \hyperref[spaces-more-cohomology-section-phantom]{Compléments sur la cohomologie des espaces}
\item \hyperref[spaces-simplicial-section-phantom]{Espaces simpliciaux}
\item \hyperref[spaces-duality-section-phantom]{Dualité pour les espaces}
\item \hyperref[formal-spaces-section-phantom]{Espaces algébriques formels}
\item \hyperref[restricted-section-phantom]{Algébrisation des espaces formels}
\item \hyperref[spaces-resolve-section-phantom]{Résolution des surfaces revisitée}
\end{enumerate}
Théorie des déformations
\begin{enumerate}
\setcounter{enumi}{89}
\item \hyperref[formal-defos-section-phantom]{Théorie formelle des déformations}
\item \hyperref[defos-section-phantom]{Théorie des déformations}
\item \hyperref[cotangent-section-phantom]{Le complexe cotangent}
\item \hyperref[examples-defos-section-phantom]{Problèmes de déformation}
\end{enumerate}
Champs algébriques
\begin{enumerate}
\setcounter{enumi}{93}
\item \hyperref[algebraic-section-phantom]{Champs algébriques}
\item \hyperref[examples-stacks-section-phantom]{Exemples de champs}
\item \hyperref[stacks-sheaves-section-phantom]{Faisceaux sur les champs algébriques}
\item \hyperref[criteria-section-phantom]{Critères de représentabilité}
\item \hyperref[artin-section-phantom]{Axiomes d'Artin}
\item \hyperref[quot-section-phantom]{Espaces de Quot et de Hilbert}
\item \hyperref[stacks-properties-section-phantom]{Propriétés des champs algébriques}
\item \hyperref[stacks-morphisms-section-phantom]{Morphismes de champs algébriques}
\item \hyperref[stacks-limits-section-phantom]{Limites de champs algébriques}
\item \hyperref[stacks-cohomology-section-phantom]{Cohomologie des champs algébriques}
\item \hyperref[stacks-perfect-section-phantom]{Catégories dérivées de champs}
\item \hyperref[stacks-introduction-section-phantom]{Introduction aux champs algébriques}
\item \hyperref[stacks-more-morphisms-section-phantom]{Compléments sur les morphismes de champs}
\item \hyperref[stacks-geometry-section-phantom]{La géométrie des champs}
\end{enumerate}
Thèmes de théorie des espaces de modules
\begin{enumerate}
\setcounter{enumi}{107}
\item \hyperref[moduli-section-phantom]{Champs de modules}
\item \hyperref[moduli-curves-section-phantom]{Espaces de modules de courbes}
\end{enumerate}
Divers
\begin{enumerate}
\setcounter{enumi}{109}
\item \hyperref[examples-section-phantom]{Exemples}
\item \hyperref[exercises-section-phantom]{Exercices}
\item \hyperref[guide-section-phantom]{Guide bibliographique}
\item \hyperref[desirables-section-phantom]{Desiderata}
\item \hyperref[coding-section-phantom]{Style de programmation}
\item \hyperref[obsolete-section-phantom]{Obsolète}
\item \hyperref[fdl-section-phantom]{Licence de documentation libre GNU}
\item \hyperref[index-section-phantom]{Index généré automatiquement}
\end{enumerate}
\end{multicols}


\bibliography{my}
\bibliographystyle{amsalpha}

\end{document}
